\documentclass[12pt]{book}
% Idioma y codificación
\usepackage[spanish, es-tabla]{babel}       %es-tabla para que se titule "Tabla"
\usepackage[utf8]{inputenc}

% Márgenes
\usepackage[a4paper,top=3cm,bottom=2.5cm,left=3cm,right=3cm]{geometry}

% Comentarios de bloque
\usepackage{verbatim}

% Paquetes de links
\usepackage[hidelinks]{hyperref}    % Permite enlaces
\usepackage{url}                    % redirecciona a la web

% Más opciones para enumeraciones
\usepackage{enumitem}

% Personalizar la portada
\usepackage{titling}


% Paquetes de tablas
\usepackage{multirow}


%------------------------------------------------------------------------

%Paquetes de figuras
\usepackage{caption}
\usepackage{subcaption} % Figuras al lado de otras
\usepackage{float}      % Poner figuras en el sitio indicado H.


% Paquetes de imágenes
\usepackage{graphicx}       % Paquete para añadir imágenes
\usepackage{transparent}    % Para manejar la opacidad de las figuras

% Paquete para usar colores
\usepackage[dvipsnames]{xcolor}
\usepackage{pagecolor}      % Para cambiar el color de la página

% Habilita tamaños de fuente mayores
\usepackage{fix-cm}


%------------------------------------------------------------------------

% Paquetes de matemáticas
\usepackage{mathtools, amsfonts, amssymb, mathrsfs}
\usepackage[makeroom]{cancel}     % Simplificar tachando
\usepackage{polynom}    % Divisiones y Ruffini
\usepackage{units} % Para poner fracciones diagonales con \nicefrac

\usepackage{pgfplots}   %Representar funciones
\pgfplotsset{compat=1.18}  % Versión 1.18

\usepackage{tikz-cd}    % Para usar diagramas de composiciones
\usetikzlibrary{calc}   % Para usar cálculo de coordenadas en tikz

%Definición de teoremas, etc.
\usepackage{amsthm}
%\swapnumbers   % Intercambia la posición del texto y de la numeración

\theoremstyle{plain}

\makeatletter
\@ifclassloaded{article}{
  \newtheorem{teo}{Teorema}[section]
}{
  \newtheorem{teo}{Teorema}[chapter]  % Se resetea en cada chapter
}
\makeatother

\newtheorem{coro}{Corolario}[teo]           % Se resetea en cada teorema
\newtheorem{prop}[teo]{Proposición}         % Usa el mismo contador que teorema
\newtheorem{lema}[teo]{Lema}                % Usa el mismo contador que teorema

\theoremstyle{remark}
\newtheorem*{observacion}{Observación}

\theoremstyle{definition}

\makeatletter
\@ifclassloaded{article}{
  \newtheorem{definicion}{Definición} [section]     % Se resetea en cada chapter
}{
  \newtheorem{definicion}{Definición} [chapter]     % Se resetea en cada chapter
}
\makeatother

\newtheorem*{notacion}{Notación}
\newtheorem*{ejemplo}{Ejemplo}
\newtheorem*{ejercicio*}{Ejercicio}             % No numerado
\newtheorem{ejercicio}{Ejercicio} [section]     % Se resetea en cada section


% Modificar el formato de la numeración del teorema "ejercicio"
\renewcommand{\theejercicio}{%
  \ifnum\value{section}=0 % Si no se ha iniciado ninguna sección
    \arabic{ejercicio}% Solo mostrar el número de ejercicio
  \else
    \thesection.\arabic{ejercicio}% Mostrar número de sección y número de ejercicio
  \fi
}


% \renewcommand\qedsymbol{$\blacksquare$}         % Cambiar símbolo QED
%------------------------------------------------------------------------

% Paquetes para encabezados
\usepackage{fancyhdr}
\pagestyle{fancy}
\fancyhf{}

\newcommand{\helv}{ % Modificación tamaño de letra
\fontfamily{}\fontsize{12}{12}\selectfont}
\setlength{\headheight}{15pt} % Amplía el tamaño del índice


%\usepackage{lastpage}   % Referenciar última pag   \pageref{LastPage}
\fancyfoot[C]{\thepage}

%------------------------------------------------------------------------

% Conseguir que no ponga "Capítulo 1". Sino solo "1."
\makeatletter
\@ifclassloaded{book}{
  \renewcommand{\chaptermark}[1]{\markboth{\thechapter.\ #1}{}} % En el encabezado
    
  \renewcommand{\@makechapterhead}[1]{%
  \vspace*{50\p@}%
  {\parindent \z@ \raggedright \normalfont
    \ifnum \c@secnumdepth >\m@ne
      \huge\bfseries \thechapter.\hspace{1em}\ignorespaces
    \fi
    \interlinepenalty\@M
    \Huge \bfseries #1\par\nobreak
    \vskip 40\p@
  }}
}
\makeatother

%------------------------------------------------------------------------
% Paquetes de cógido
\usepackage{minted}
\renewcommand\listingscaption{Código fuente}

\usepackage{fancyvrb}
% Personaliza el tamaño de los números de línea
\renewcommand{\theFancyVerbLine}{\small\arabic{FancyVerbLine}}

% Estilo para C++
\newminted{cpp}{
    frame=lines,
    framesep=2mm,
    baselinestretch=1.2,
    linenos,
    escapeinside=||
}



\usepackage{listings} % Para incluir código desde un archivo

\renewcommand\lstlistingname{Código Fuente}
\renewcommand\lstlistlistingname{Índice de Códigos Fuente}

% Definir colores
\definecolor{vscodepurple}{rgb}{0.5,0,0.5}
\definecolor{vscodeblue}{rgb}{0,0,0.8}
\definecolor{vscodegreen}{rgb}{0,0.5,0}
\definecolor{vscodegray}{rgb}{0.5,0.5,0.5}
\definecolor{vscodebackground}{rgb}{0.97,0.97,0.97}
\definecolor{vscodelightgray}{rgb}{0.9,0.9,0.9}

% Configuración para el estilo de C similar a VSCode
\lstdefinestyle{vscode_C}{
  backgroundcolor=\color{vscodebackground},
  commentstyle=\color{vscodegreen},
  keywordstyle=\color{vscodeblue},
  numberstyle=\tiny\color{vscodegray},
  stringstyle=\color{vscodepurple},
  basicstyle=\scriptsize\ttfamily,
  breakatwhitespace=false,
  breaklines=true,
  captionpos=b,
  keepspaces=true,
  numbers=left,
  numbersep=5pt,
  showspaces=false,
  showstringspaces=false,
  showtabs=false,
  tabsize=2,
  frame=tb,
  framerule=0pt,
  aboveskip=10pt,
  belowskip=10pt,
  xleftmargin=10pt,
  xrightmargin=10pt,
  framexleftmargin=10pt,
  framexrightmargin=10pt,
  framesep=0pt,
  rulecolor=\color{vscodelightgray},
  backgroundcolor=\color{vscodebackground},
}

%------------------------------------------------------------------------

% Comandos definidos
\newcommand{\bb}[1]{\mathbb{#1}}
\newcommand{\cc}[1]{\mathcal{#1}}

% I prefer the slanted \leq
\let\oldleq\leq % save them in case they're every wanted
\let\oldgeq\geq
\renewcommand{\leq}{\leqslant}
\renewcommand{\geq}{\geqslant}

% Si y solo si
\newcommand{\sii}{\iff}

% Letras griegas
\newcommand{\eps}{\epsilon}
\newcommand{\veps}{\varepsilon}
\newcommand{\lm}{\lambda}

\newcommand{\ol}{\overline}
\newcommand{\ul}{\underline}
\newcommand{\wt}{\widetilde}
\newcommand{\wh}{\widehat}

\let\oldvec\vec
\renewcommand{\vec}{\overrightarrow}

% Derivadas parciales
\newcommand{\del}[2]{\frac{\partial #1}{\partial #2}}
\newcommand{\Del}[3]{\frac{\partial^{#1} #2}{\partial^{#1} #3}}
\newcommand{\deld}[2]{\dfrac{\partial #1}{\partial #2}}
\newcommand{\Deld}[3]{\dfrac{\partial^{#1} #2}{\partial^{#1} #3}}


\newcommand{\AstIg}{\stackrel{(\ast)}{=}}
\newcommand{\Hop}{\stackrel{L'H\hat{o}pital}{=}}

\newcommand{\red}[1]{{\color{red}#1}} % Para integrales, destacar los cambios.

% Método de integración
\newcommand{\MetInt}[2]{
    \left[\begin{array}{c}
        #1 \\ #2
    \end{array}\right]
}

% Declarar aplicaciones
% 1. Nombre aplicación
% 2. Dominio
% 3. Codominio
% 4. Variable
% 5. Imagen de la variable
\newcommand{\Func}[5]{
    \begin{equation*}
        \begin{array}{rrll}
            #1:& #2 & \longrightarrow & #3\\
               & #4 & \longmapsto & #5
        \end{array}
    \end{equation*}
}

%------------------------------------------------------------------------

% Para poder incluir árboles
\usepackage{forest}
\usepackage{booktabs}

% Para poder hacer graficas en polares (espiral de arquimedes)
\usepackage{pgfplots}
\usepgfplotslibrary{polar}

% // TODO:
\newcommand{\showimages}{} % Descomentar esta línea para mostrar las imágenes

\begin{document}

    % 1. Foto de fondo
    % 2. Título
    % 3. Encabezado Izquierdo
    % 4. Color de fondo
    % 5. Coord x del titulo
    % 6. Coord y del titulo
    % 7. Fecha
    % 8. Autor

    
    % 1. Foto de fondo
% 2. Título
% 3. Encabezado Izquierdo
% 4. Color de fondo
% 5. Coord x del titulo
% 6. Coord y del titulo
% 7. Fecha

\newcommand{\portada}[7]{

    \portadaBase{#1}{#2}{#3}{#4}{#5}{#6}{#7}
    \portadaBook{#1}{#2}{#3}{#4}{#5}{#6}{#7}
}

\newcommand{\portadaExamen}[7]{

    \portadaBase{#1}{#2}{#3}{#4}{#5}{#6}{#7}
    \portadaArticle{#1}{#2}{#3}{#4}{#5}{#6}{#7}
}




\newcommand{\portadaBase}[7]{

    % Tiene la portada principal y la licencia Creative Commons
    
    % 1. Foto de fondo
    % 2. Título
    % 3. Encabezado Izquierdo
    % 4. Color de fondo
    % 5. Coord x del titulo
    % 6. Coord y del titulo
    % 7. Fecha
    
    
    \thispagestyle{empty}               % Sin encabezado ni pie de página
    \newgeometry{margin=0cm}        % Márgenes nulos para la primera página
    
    
    % Encabezado
    \fancyhead[L]{\helv #3}
    \fancyhead[R]{\helv \nouppercase{\leftmark}}
    
    
    \pagecolor{#4}        % Color de fondo para la portada
    
    \begin{figure}[p]
        \centering
        \transparent{0.3}           % Opacidad del 30% para la imagen
        
        \includegraphics[width=\paperwidth, keepaspectratio]{assets/#1}
    
        \begin{tikzpicture}[remember picture, overlay]
            \node[anchor=north west, text=white, opacity=1, font=\fontsize{60}{90}\selectfont\bfseries\sffamily, align=left] at (#5, #6) {#2};
            
            \node[anchor=south east, text=white, opacity=1, font=\fontsize{12}{18}\selectfont\sffamily, align=right] at (9.7, 3) {\textbf{\href{https://losdeldgiim.github.io/}{Los Del DGIIM}}};
            
            \node[anchor=south east, text=white, opacity=1, font=\fontsize{12}{15}\selectfont\sffamily, align=right] at (9.7, 1.8) {Doble Grado en Ingeniería Informática y Matemáticas\\Universidad de Granada};
        \end{tikzpicture}
    \end{figure}
    
    
    \restoregeometry        % Restaurar márgenes normales para las páginas subsiguientes
    \pagecolor{white}       % Restaurar el color de página
    
    
    \newpage
    \thispagestyle{empty}               % Sin encabezado ni pie de página
    \begin{tikzpicture}[remember picture, overlay]
        \node[anchor=south west, inner sep=3cm] at (current page.south west) {
            \begin{minipage}{0.5\paperwidth}
                \href{https://creativecommons.org/licenses/by-nc-nd/4.0/}{
                    \includegraphics[height=2cm]{assets/Licencia.png}
                }\vspace{1cm}\\
                Esta obra está bajo una
                \href{https://creativecommons.org/licenses/by-nc-nd/4.0/}{
                    Licencia Creative Commons Atribución-NoComercial-SinDerivadas 4.0 Internacional (CC BY-NC-ND 4.0).
                }\\
    
                Eres libre de compartir y redistribuir el contenido de esta obra en cualquier medio o formato, siempre y cuando des el crédito adecuado a los autores originales y no persigas fines comerciales. 
            \end{minipage}
        };
    \end{tikzpicture}
    
    
    
    % 1. Foto de fondo
    % 2. Título
    % 3. Encabezado Izquierdo
    % 4. Color de fondo
    % 5. Coord x del titulo
    % 6. Coord y del titulo
    % 7. Fecha


}


\newcommand{\portadaBook}[7]{

    % 1. Foto de fondo
    % 2. Título
    % 3. Encabezado Izquierdo
    % 4. Color de fondo
    % 5. Coord x del titulo
    % 6. Coord y del titulo
    % 7. Fecha

    % Personaliza el formato del título
    \pretitle{\begin{center}\bfseries\fontsize{42}{56}\selectfont}
    \posttitle{\par\end{center}\vspace{2em}}
    
    % Personaliza el formato del autor
    \preauthor{\begin{center}\Large}
    \postauthor{\par\end{center}\vfill}
    
    % Personaliza el formato de la fecha
    \predate{\begin{center}\huge}
    \postdate{\par\end{center}\vspace{2em}}
    
    \title{#2}
    \author{\href{https://losdeldgiim.github.io/}{Los Del DGIIM}}
    \date{Granada, #7}
    \maketitle
    
    \tableofcontents
}




\newcommand{\portadaArticle}[7]{

    % 1. Foto de fondo
    % 2. Título
    % 3. Encabezado Izquierdo
    % 4. Color de fondo
    % 5. Coord x del titulo
    % 6. Coord y del titulo
    % 7. Fecha

    % Personaliza el formato del título
    \pretitle{\begin{center}\bfseries\fontsize{42}{56}\selectfont}
    \posttitle{\par\end{center}\vspace{2em}}
    
    % Personaliza el formato del autor
    \preauthor{\begin{center}\Large}
    \postauthor{\par\end{center}\vspace{3em}}
    
    % Personaliza el formato de la fecha
    \predate{\begin{center}\huge}
    \postdate{\par\end{center}\vspace{5em}}
    
    \title{#2}
    \author{\href{https://losdeldgiim.github.io/}{Los Del DGIIM}}
    \date{Granada, #7}
    \thispagestyle{empty}               % Sin encabezado ni pie de página
    \maketitle
    \vfill
}
    \portada{ffccA4.jpg}{Ecuaciones\\Diferenciales I}{Ecuaciones Diferenciales I}{MidnightBlue}{-8}{28}{2024-2025}{José Juan Urrutia Milán}

    \newpage
    La parte de teoría del presente documento (es decir, excluyendo las relaciones de problemas) está hecha en función de los apuntes que se han tomado en clase.
    No obstante, recomendamos seguir de igual forma los apuntes del profesor de la asignatura, Rafael Ortega, disponibles en su sitio web personal
    \href{https://www.ugr.es/~rortega/}{https://www.ugr.es/~rortega/}. Estos apuntes no son por tanto una completa sustitución de dichos apuntes, sino tan solo un complemento.
    

    %%\setcounter{chapter}{-1}
\chapter*{Introducción}

La teoría de Ecuaciones Diferenciales es la teoría matemática relacionada con el movimiento. Esta trata de resolver ecuaciones cuyas soluciones son funciones. Podríamos pensar en llamar a este área ``Ecuaciones Funcionales'', pero gracias a que en física $F = m \cdot a$, podemos encontrar de forma natural y útil ecuaciones funcionales donde la información que aparezca sobre la función a buscar está relacionada con las relaciones que guardan las derivadas de dicha función.

\begin{ejemplo}
    Como ejemplo de ecuación diferencial que surge de forma natural podemos pensar en el movimiento de un péndulo:

    Para describir un péndulo, nos es suficiente con tres variables independientes, que podemos ver en la siguiente ilustración:
    \begin{itemize}
        \item La longitud del hilo del péndulo, a la que llamamos $l$.
        \item La aceleración gravitatoria del planeta en el que nos encontremos, a la que llamamos $g$.
        \item Y el ángulo que guarda el péndulo sobre la vertical, al que llamamos $\theta$.
    \end{itemize}

    % Péndulo

\ifdefined\showimages
    \begin{figure}[H]
        \centering
        \begin{tikzpicture}
            % Parámetros
            \def\L{3}        % Longitud del péndulo
            \def\LL{2.7}     % Longitud de la vertical
            \def\ang{30}     % Ángulo del péndulo (en grados)
    
            % Dibujar la línea discontinua vertical
            \draw[dashed] (0,0) -- (0,-\LL);
    
            % Dibujar el péndulo
            \draw[thick] (0,0) -- ({\L*sin(\ang)}, {-\L*cos(\ang)});
    
            % Dibujar el círculo al final del péndulo
            \filldraw[thick] ({\L*sin(\ang)}, {-\L*cos(\ang)}) circle (0.2);
    
            % Indicar el ángulo
            \draw[thick, -stealth] (0,-1) arc[start angle=270, end angle=300, radius=1];
            \node at (0.35, -1.4) {$\theta$};
            \node at (1, -1) {$l$};
            \node at (3.3, -1) {$g$};

            % Flecha hacia abajo para la gravedad
            \draw[thick, -stealth] (3, -0.5) -- (3, -1.5);
    
            % Dibujar el punto de suspensión
            \filldraw (0,0) circle (0.05);
        \end{tikzpicture}
    \end{figure}
\fi
    Como nuestro objetivo es describir el movimiento que describe un péndulo, tenemos que introducir una variable más, el tiempo ($t$), y ver ahora la variable $\theta$ como una variable dependiente en función de $t$:
    \begin{equation*}
        \theta = \theta(t)
    \end{equation*}
    La física nos dice que si $\theta(t)$ es una función que nos describe el movimiento de un péndulo en función del tiempo, entonces debe cumplir la siguiente ecuación\footnote{De dónde sale dicha ecuación no nos es relevante.}:
    \begin{equation}\label{eq:pendulo}
        \theta''(t) + \frac{g}{l} \sen \theta(t) = 0
    \end{equation}

    A partir de esta ecuación, nos preguntamos por las funciones $\theta$ que cumplan dicha ecuación, siendo esta la primera ecuación diferencial que trataremos de resolver.

    A simple vista, podemos decir acertadamente que dos soluciones para dicha ecuación son:
    \begin{equation*}
        \theta(t) = 0 \qquad \qquad  \theta(t) = \pi
    \end{equation*}
    Pensando que en ambas soluciones el péndulo se encuentra en un estado estático, en la primera este se encuentra quieto debajo y en la segunda, quieto arriba. De forma intuitiva podemos pensar que el primero es un equilibrio estable y el segundo un equilibrio inestable.

    A partir de dichas soluciones, podemos adivinar que también serán soluciones de~(\ref{eq:pendulo}) cualquier función de la forma:
    \begin{equation*}
        \theta(t) = n\pi, \quad n \in \mathbb{Z}
    \end{equation*}
    De esta forma, hemos encontrado una \textbf{familia de soluciones}, es decir, tenemos una función que es solución de~(\ref{eq:pendulo}) para cada $n\in \mathbb{Z}$.

    Hemos encontrado infinitas soluciones para la ecuación~(\ref{eq:pendulo}). Podemos pensar que cada una de estas soluciones describe un péndulo distinto (esto es, cada una de las distintas formas de tirar el péndulo). En el mundo de las ecuaciones diferenciales es común encontrar muchas soluciones para una sola ecuación.

    En el caso de la ecuación~(\ref{eq:pendulo}), se ha demostrado que a parte de la familia de soluciones que hemos dado, no pueden encontrarse más soluciones con fórmula (aunque pueden aproximarse).
\end{ejemplo}~\\

El orden de una ecuación diferencial es la derivada de mayor grado qeu aparezca en la fórmula de la ecuación. En el caso de~(\ref{eq:pendulo}), esta era de orden 2. En esta asignatura nos centraremos en el estudio de las ecuaciones diferenciales de orden 1.

Una ecuación diferencial de primer orden genérica es de la forma:
\begin{equation}\label{eq_generica}
    \Phi(t,x(t), x'(t))=0
\end{equation}
Es decir, es una relación entre una variable independiente ($t$, que usualmente podremos entender como el tiempo), una variable dependiente o función ($x$, en función de $t$), cuya expresión estamos interesados en buscar; y su derivada.

La $\Phi$ que aparece en~(\ref{eq_generica}) será una función:
\Func{\Phi}{D\subseteq \mathbb{R}^3}{\mathbb{R}}{(t,x,y)}{\Phi(t,x,y)}
Donde trataremos de ver $x$ como variable independiente que tenemos que hacer dependiente de $t$: $x=x(t)$, siendo $y$ su derivada.

\begin{ejemplo}
    Dada la siguiente ecuación diferencial:
    \begin{equation*}
        {(x(t))}^{2} + {(x'(t))}^{2} = 1
    \end{equation*}
    La función $\Phi$ en cuestión es:
    \Func{\Phi}{D=\mathbb{R}^3}{\mathbb{R}}{(t,x,y)}{x^2+y^2-1}
    Soluciones de dicha ecuación son (a simple vista):

    \begin{equation*}
    \begin{array}{ll}
        x(t) = \sen t & x(t) = \cos t \\
        x(t) = 1 & x(t) = -1
    \end{array}
    \end{equation*}

    Además, podemos deducir que una familia de soluciones que engloba a las dos primeras es:
    \begin{equation*}
        x(t) = \sen(t+c), \quad c\in \mathbb{R}
    \end{equation*}

    % Grafica de familia de soluciones, 1 y -1
\ifdefined\showimages
    \begin{center}
    \begin{tikzpicture}
      \begin{axis}[
          axis lines = center,    % Dibujar los ejes
          xmin = -2*pi, xmax = 2*pi,   % Definir el rango en el eje x
          ymin = -1.5, ymax = 1.5,         % Definir el rango en el eje y
          xlabel = {},              % Etiqueta del eje x
          ylabel = {},                 % Etiqueta del eje y
          domain = -2*pi:2*pi,         % Dominio de las funciones
          samples = 200,               % Número de muestras para las funciones
          legend pos = north west,     % Posición de la leyenda
          %grid = both,                 % Mostrar la cuadrícula
          major grid style = {lightgray},   % Estilo de la cuadrícula principal
          minor grid style = {dotted}       % Estilo de la cuadrícula secundaria
      ]

      % Función constante 1
      \addplot[red, thick] {1};

      % Función constante -1
      \addplot[blue, thick] {-1};

      % Función sen(x)
      \addplot[green, thick] {sin(deg(x))};

      % Función cos(x)
      \addplot[purple, thick] {cos(deg(x))};

      % Función sen(x + pi)
      \addplot[purple, thick] {sin(deg(x + pi))};

      % Función sen(x + pi)
      \addplot[purple, thick] {sin(deg(x + 2 * pi))};

      \end{axis}
    \end{tikzpicture}
    \end{center}
\fi

    Graficando las soluciones, podemos además construir una nueva solución con una función a trozos:
    \begin{equation*}
        x(t) = \left\{\begin{array}{ll}
                \cos t & t \geq 0 \\
                1 & t<0
        \end{array}\right.
    \end{equation*}
    Derivable en $\mathbb{R}^*$ por el carácter local de la derivabilidad y en 0 por coincidir los dos límites laterales de las derivadas. Sin embargo, observamos que no es dos veces derivable.
\end{ejemplo}~\\

Las ecuaciones diferenciales de orden 1 sin restricción alguna son demasiado generales como para construir una teoría formal que centre su estudio en estas. Por tanto, estudiaremos aquellas que admitan escribirlas en \textbf{forma normal}, es decir, que si nos dan una ecuación en la forma~(\ref{eq_generica}), esta pueda escribirse como:
\begin{equation}\label{eq_generica_normal}
    x'(t) = f(t,x(t))
\end{equation}
Para una cierta
\Func{f}{C\subseteq \mathbb{R}^2}{\mathbb{R}}{(t,x)}{f(t,x)}

\begin{ejemplo}
    Dadas las dos siguientes ecuaciones diferenciales, discutir cuál de ellas es más simple y resolver ambas.
    \begin{itemize}
        \item $x'(t) = 7x(t)$
        \item $x'(t) = 7t$
    \end{itemize}
    La segunda ecuación es más sencilla, ya que se trata de un cálculo de una primitiva para $x'$: $x'(t) = h(t)$, a lo que ya estamos acostumbrados.

    Para dicha ecuación, tenemos:
    \begin{gather*}
        \Phi(t,x,y) = y-7t \\
        f(t,x) = 7t
    \end{gather*}

    con $D=\mathbb{R}^3$ y $C=\mathbb{R}^2$.

    Las soluciones de dicha ecuación diferencial son, por tanto:
    \begin{equation*}
        x(t) = \frac{7}{2} t^2 + c, \quad c\in \mathbb{R}
    \end{equation*}
    Para la primera ecuación, tenemos:
    \begin{gather*}
        \Phi(t,x,y) = y-7x \\
        f(t,x) = 7x
    \end{gather*}

    con $D=\mathbb{R}^3$ y $C=\mathbb{R}^2$.

    De forma simple vemos dos primeras soluciones:
    \begin{equation*}
        x(t) = 0 \qquad  x(t) = e^{7t}
    \end{equation*}

    Y podemos deducir además una familia de soluciones:
    \begin{equation*}
        x(t) = c \cdot e^{7t}, \quad c\in \mathbb{R}
    \end{equation*}

    De hecho, demostraremos más adelante que dicha ecuación no tiene más soluciones además de las de dicha familia.

\end{ejemplo}
   
\begin{ejemplo}
    Dadas las dos siguientes ecuaciones diferenciales, discutir cuál de ellas es más simple e intentar resolverlas.
    \begin{itemize}
        \item $x'(t) = \sen t$
        \item $x'(t) = \sen x(t)$
    \end{itemize}
    En este caso, es la primera la que es más simple, ya que se vuelve a tratar de un cálculo de primitiva.

    Tenemos:
    \begin{gather*}
        \Phi(t,x,y) = y-\sen t \\
        f(t,x) = \sen t
    \end{gather*}

    Siendo las soluciones de la ecuación diferencial:
    \begin{equation*}
        x(t) = -\cos t + c, \quad c \in \mathbb{R}
    \end{equation*}
    En el caso de la segunda, tenemos:
    \begin{gather*}
        \Phi(t,x,y) = y-\sen x \\
        f(t,x) = \sen x
    \end{gather*}
    Ecuación que todavía no sabemos resolver.
\end{ejemplo}

    %\chapter{Ensamblador}
\section{Historia}
La linea de procesadores Intel, comúnmente conocida como \textit{x86}
ha seguido un largo proceso de evolución, manteniendo compatibilidad
hasta 8086, introducido en 1978. La familia de procesadores x86 ha seguido una evolución hasta
el día de hoy:
\begin{itemize}
	\item 8086 (1978) $\rightarrow$ 29K transistores, 16 bits, 5 MHz. \quad [Intel 16-bit, IBM PC]
	\item 386 (1985) $\rightarrow$ 275K transistores, 32 bits, 16 MHz. \quad [Intel 32-bit]
	\item Pentium 4E (2004) $\rightarrow$ 125M transistores, 64 bits, 3.4 GHz. \quad [Intel 64-bit, Hyperthreading, Implementación de x86-64 de AMD]
	\item Core i7 Haswell (2013) $\rightarrow$ 1.4G transistores, 64 bits, 3.9 Ghz \quad[Multicore, AVX2]
	\item Core i9 Raptor Lake (2022) $ \rightarrow$ +30G transistores, 64 bits, 5.2 GHz \quad[DDR5, PCIe 5.0, AVX-512]
\end{itemize}
A lo largo de los años, muchas compañias han creado procesadores
que se puedan comparar con los de Intel. Entre ellas, AMD ha sido
la más destacada, cayendo ligeramente detrás de Intel en cuanto a rendimiento
y ofreciendo un precio más competitivo. Sin embargo, en los últimos años, AMD ha
sido capaz de ofrecer procesadores con un rendimiento similar o superior a los de Intel,
concretamente, en 2024, Intel ha sufrido una gran decadencia por problemas de estabilidad
en los ultimos lanzamientos.
Nosotros nos centraremos en la arquitectura x86-64, que es la que se utiliza en el libro de
clase, sin embargo, haremos pequeñas menciones a IA32, cuando sea conveniente.
\section{Lenguaje C, ensamblador y máquina}
Las computadoras modernas utilizan distintas capas de abstracción, ocultando la complejidad
de la implementación. Dos de estas capas son realmente importantes para nosotros,
el formato y comportamiento de las instrucciones de máquina, que viene determinado
por la ISA (Instruction Set Architecture), esta define el formato de las instrucciones,
y el resultado que cada una de estas tendrá sobre el estado de la máquina.
La mayoria de ISAs, incluyendo x86-64, describen el comportamiento de un programa de
forma secuencial, como si se ejecutase una instrucción tras otra, sin embargo, en realidad,
el hardware de la CPU es mucho mas complejo, ejecutando instrucciones de forma paralelizada
o en otro orden, sin embargo, esto se verá con mas detenimiento en la asignatura \textit{Arquitectura de Computadores}.
La otra capa de abstracción ocurre sobre las direcciones de memoria, se nos presentan
como \textit{virtual addresses}, proporcionando un modelo de memoria que aparenta ser
un gran array de bytes, sin embargo, en realidad, la memoria se maneja de otra forma,
como se ve en la asignatura \textit{Sistemas Operativos}.
Recalcamos aunque no vayan a ser estudiadas en profundidad, que existen otras
ISA, como ARM (utilizada en dispositivos móviles), RISC-V (Nueva ISA open-source),
o la misma IA32. Aunque no se mencione más, se recomienda al lector la profundización
sobre ARM, que actualmente está ampliando su uso en el mercado de portatiles. (Intel Ultra, M1, etc.)
Recordamos que podemos clasificar en general las distintas ISAs en dos tipos:
\begin{itemize}
	\item \textbf{CISC} (Complex Instruction Set Computer): Conjunto de instrucciones complejas.
	\item \textbf{RISC} (Reduced Instruction Set Computer): Conjunto de instrucciones reducidas.
\end{itemize}
Realmente la mayoría de ISAs actuales son híbridas, convirtiendo
las instrucciones de CISC en instrucciones de RISC, para mejorar el rendimiento.
\newline
\newline
El codigo x86-64 difiere en gran medida del código C, destacamos las siguientes
partes del estado del procesador que normalmente no se ven en C:
\begin{itemize}
	\item \textbf{PC (Program Counter)}: Registra la dirección de la siguiente instrucción a ejecutar.
	      Se denomina \textit{\%rip} en x86-64.
	\item \textbf{Archivo de registros}: Contiene 16 registros etiquetados con valores de 64 bits.
	      Pueden ser utilizados para almacenar valores temporales, punteros, etc.
	\item \textbf{Codigos de condición}: Contienen información sobre la instrucción mas reciente.
	      Se utilizan para realizar saltos condicionales.
	\item \textbf{Pila}: Se utiliza para almacenar enteros o flotantes. Es un array direccionable por bytes.
\end{itemize}
Ahora, vamos a ver cómo se relaciona el código C con el código ensamblador. Para ello,
habrá que tener en cuenta que el código generado por el compilador puede variar
según el compilador, la versión, las optimizaciones, etc.
\subsection{Proceso de compilación}
Partimos de un número arbitrario de ficheros \textit{.c}, que contienen código C.
Para compilarlos, se pueden utilizar distintos compiladores, nosotros usaremos
\textit{gcc}:
% codigo en bash:
\begin{minted}{bash}
    gcc -Og programa.c -o programa
\end{minted}
Cuando lo ejecutemos:
\begin{center}
	\begin{tikzpicture}[node distance=1.6cm]

		% Styles within tikzpicture
		\tikzstyle{process} = [rectangle, rounded corners, minimum width=3cm, minimum height=1cm, text centered, draw=black, fill=green!30]
		\tikzstyle{arrow} = [thick,->,>=stealth]

		% Nodes
		\node (source) [process, fill=green!60] {Source File};
		\node (preprocessor) [process, below of=source] {Preprocessor};
		\node (compiler) [process, below of=preprocessor] {Compiler};
		\node (assembler) [process, below of=compiler] {Assembler};
		\node (linker) [process, below of=assembler] {Linker};
		\node (executable) [process, fill=green!60, below of=linker] {Executable};

		% Arrows
        \draw [-Stealth] (source) -- (preprocessor);
        \draw [-Stealth](preprocessor) -- node[right]{Preprocessed File (.i file)} (compiler);
        \draw [-Stealth](compiler) -- node[right]{Assembly code (.s file)} (assembler);
        \draw [-Stealth](assembler) -- node[right]{Object code (.o file)} (linker);
        \draw [-Stealth](linker) -- (executable);
	\end{tikzpicture}
\end{center}
En este proceso, gcc primero preprocesa el código, generando un archivo \textit{.i}, que
posteriormente compila, generando un archivo \textit{.s}, que se ensambla para generar
\textit{.o}, finalmente, se enlazan los distintos archivos \textit{.o} para generar el ejecutable.
Vamos a explicar un poco más detenidamente cada uno de los pasos:
\begin{itemize}
	\item \textbf{Preprocesador}: C proporciona ciertas herramientas del lenguaje mediante un preprocesador,
	      que es el primer paso en la compilación. Los más utilizados son \textit{\#include}, para incluir el contenido de un archivo
	      durante la compilación, y \textit{\#define}, para reemplazar un token por una secuencia arbitraria de texto
	      o valores. Además, elimina los comentarios de un archivo fuente. En gcc, podemos generar el código preprocesado con el flag \textit{-E}.
	\item \textbf{Compilación}: El compilador traduce el codigo preprocesado y genera código ensamblador, con gcc, podemos
	      generar el código ensamblador con el flag \textit{-S}.
	\item \textbf{Ensamblado}: El ensamblador traduce el código ensamblador a código de máquina, generando un archivo \textit{.o}.
	      En gcc, podemos generar el código de máquina con el flag \textit{-c}.
	\item \textbf{Enlazado}: El enlazador combina los distintos archivos \textit{.o} generados en un solo ejecutable. También se encarga de combinar
	      las librerias que se hayan utilizado en el código. En gcc, podemos generar el ejecutable con el flag \textit{-o}.
\end{itemize}


\begin{ejemplo}
	Veamos un ejemplo de cómo se relaciona el código C con el código ensamblador (solo se muestra la parte interesante):
	\begin{center} % Centramos el contenido
		\begin{minipage}{0.45\textwidth}
			\begin{minted}[fontsize=\Large]{c}
long plus(long x, long y);

void sumStore(long x, long y,
        long *dest) {
    long t = plus(x, y);
    *dest = t;
}
                \end{minted}
		\end{minipage}
		\begin{minipage}{0.45\textwidth}
			\begin{minted}[fontsize=\Large]{asm}
sumstore:
    pushq %rbx
    movq  %rdx, %rbx
    call  plus
    movq  %raw, (%rbx)
    popq  %rbx
    ret
                \end{minted}
		\end{minipage}
	\end{center}
\end{ejemplo}
Si lo compilamos en nuestra propia máquina, encontraremos muchas directivas que comienzan por . (ignorar)
\subsection{Representación Datos C, IA32, x86-64}
A continuación, mostramos los distintos tamaños en Bytes para los tipos de dato mas comúnes en C. Así como
su representación en Assembly. Debido a sus origenes como una arquitectura de 16 bits que expandió a 32 bits,
Intel usa el termino \textit{word} para referirse a un dato de 16 bits, y \textit{doubleword} para referirse a
una palabra de 32 bits. En x86-64, se ha añadido el termino \textit{quadword} para referirse a una palabra de 64 bits.
\begin{center}
	\begin{tabular}{|c|c|c|c|c|}
		\hline
		\textbf{Tipo de dato en C} & \textbf{Assembly sufijo} & \textbf{Normal 32 bits} & \textbf{IA32} & \textbf{x86-64} \\
		\hline
		\text{unsigned}            &                          & 4                       & 4             & 4               \\
		\text{int}                 & l                        & 4                       & 4             & 4               \\
		\text{long}                & q                        & 4                       & 4             & 8               \\
		\text{char}                & b                        & 1                       & 1             & 1               \\
		\text{short}               & w                        & 2                       & 2             & 2               \\
		\text{float}               & s                        & 4                       & 4             & 4               \\
		\text{double}              & l                        & 8                       & 8             & 8               \\
		\text{long double}         &                          & 8                       & 10/12         & 16              \\
		\text{char *}              &                          & 4                       & 4             & 8               \\
		\hline
	\end{tabular}
\end{center}
Los números de punto flotante se representan en el formato IEEE 754, ya visto en TOC. Se incluye
un diagrama de la representación de un número de punto flotante en simple y doble precision:
\begin{center}
	\includesvg[width=0.8\textwidth]{Imagenes/IEEE}
\end{center}
También cabe mencionar que existe un tipo de dato cuyo tamaño es 80 bits (long double) sin embargo su uso está
completamente desaconsejado debido a su portabilidad entre sistemas.
\section{Registros, operandos y operaciones}
\subsection{Registros}
Un procesador x86-64 tiene 16 registros de propósito general, que almacenan valores de 64 bits.
Estos registros se usan para almacenar valores temporales, punteros, etc. Sus nombres comienzan
con $\%r$, sin embargo, siguen convenciones de nombres debido a su legado. El primer 8086 tenía
8 registros de 16 bits, cada uno tenía un uso específico, de ahí su nombre, con la extensión a
IA32, los registros se extendieron a 32 bits, y con x86-64, a 64 bits, incluyendo además 8
nuevos registros, cuyos nombres van desde $\%r8$ hasta $\%r15$.

\begin{center}
	\includesvg[width=0.8\textwidth]{Imagenes/registers}
\end{center}
\subsection{Operandos}
La mayoría de instrucciones tienen uno o más operandos que especifican el \textit{source} y el destino de la operación.
En x86-64, los operandos pueden ser de tres tipos, (1) Inmediatos, (2) Registros, (3) Memoria. Veamos la sintaxis
para cada uno de ellos junto con la operación que se realiza:
\begin{center}
	\begin{tabular}{|l|l|l|l|}
		\hline
		\textbf{Type} & \textbf{Form}               & \textbf{Operand value}                      & \textbf{Name}       \\
		\hline
		Immediate     & \$\textit{Imm}              & \textit{Imm}                                & Immediate           \\
		Register      & $r_a$                       & $R[r_a]$                                    & Register            \\
		Memory        & \textit{Imm}                & $M[\textit{Imm}]$                           & Absolute            \\
		Memory        & $(r_a)$                     & $M[R[r_a]]$                                 & Indirect            \\
		Memory        & $\textit{Imm}(r_b)$         & $M[\textit{Imm} + R[r_b]]$                  & Base + displacement \\
		Memory        & $(r_b, r_i)$                & $M[R[r_b] + R[r_i]]$                        & Indexed             \\
		Memory        & $\textit{Imm}(r_b, r_i)$    & $M[\textit{Imm} + R[r_b] + R[r_i]]$         & Indexed             \\
		Memory        & $(r_i, s)$                  & $M[R[r_i] \cdot s]$                         & Scaled indexed      \\
		Memory        & $\textit{Imm}(, r_i, s)$    & $M[\textit{Imm} + R[r_i] \cdot s]$          & Scaled indexed      \\
		Memory        & $(r_b, r_i, s)$             & $M[R[r_b] + R[r_i] \cdot s]$                & Scaled indexed      \\
		Memory        & $\textit{Imm}(r_b, r_i, s)$ & $M[\textit{Imm} + R[r_b] + R[r_i] \cdot s]$ & Scaled indexed      \\
		\hline
	\end{tabular}
\end{center}
Los tipos de datos inmediatos se usan para valores constantes. En AT\&T se representan con un signo de dolar, siguiendo la notación
usual de C, por ejemplo \textit{\$0xFF} o \textit{\$323}. Los registros se representan con un signo de porcentaje, por ejemplo, $\%rax$.
Como ya hemos visto, los registros son de 64 bits, sin embargo, se pueden realizar operaciones de 32 bits, 16 bits o 8 bits, más adelante veremos cómo.
Los operandos de memoria se representan con paréntesis, por ejemplo, \textit{(\%rax)}, \textit{0x8(\%rax)}, etc. Podemos interpretarlo como una
desreferenciación de un puntero, es decir, acceder al valor almacenado en la dirección de memoria que se indica.
\begin{ejemplo}
	Veamos un ejemplo de cómo se realiza un swap en x86-64:
	\begin{center} % Centramos el contenido
		\begin{minipage}{0.45\textwidth}
			\begin{minted}[fontsize=\Large]{c}
void swap(long *xp, long *yp) 
{
    long x = *xp;
    *xp = *yp;
    *yp = x;
}
                \end{minted}
		\end{minipage}
		\begin{minipage}{0.45\textwidth}
			\begin{minted}[fontsize=\Large]{asm}
swap:
    movq (%rdi), %rax
    movq (%rsi), %rdx
    movq %rax, (%rsi)
    movq %rdx, (%rdi)
    ret
                \end{minted}
		\end{minipage}
	\end{center}
	Podemos ver que inicialmente, la disposición era $\% rdi \rightarrow xp$, $\% rsi \rightarrow yp$.
	Posteriormente, se almacena el \textbf{contenido} de $\%rdi$ en el registro $\%rax$, y el contenido de $\%rsi$ en $\%rdx$.
	Finalmente, se almacena el contenido de $\%rax$ en la dirección de memoria que apunta $\%rsi$, y el contenido de $\%rdx$ en la dirección de memoria que apunta $\%rdi$.
	Veamos un ejemplo para entender mejor el funcionamiento de los registros, supongamos que tenemos los siguientes valores:
	\begin{itemize}
		\item \textbf{Registro }$\%$rdi: 0x120
		\item \textbf{Registro }$\%$rsi: 0x130
		\item \textbf{Memoria en 0x120: } 125
		\item \textbf{Memoria en 0x130: } 450
	\end{itemize}
	Veamos ahora paso por paso el contenido de cada registro( $\rightarrow$ representa el contenido de la dirección de memoria):
	\[
		\begin{array}{|c|c|c|c|c|c|}
			\hline
			\textbf{Registro} & \textbf{Valor} & \textbf{movq}           & \textbf{movq}           & \textbf{movq}                                 & \textbf{movq}                         \\
			                  &                & \textbf{(\%rdi), \%rax} & \textbf{(\%rsi), \%rdx} & \textbf{\%rax, (\%rsi)}                       & \textbf{\%rdx, (\%rdi)}               \\
			\hline
			\%rdi             & 0x120          & 0x120                   & 0x120                   & 0x120                         \rightarrow 125 & 0x120 \textcolor{red}{\rightarrow450} \\
			\%rsi             & 0x130          & 0x130                   & 0x130                   & 0x130\textcolor{red}{\rightarrow125}          & 0x130 \rightarrow 125                 \\
			\%rax             & -              & \textcolor{red}{125}    & 125                     & 125                                           & 125                                   \\
			\%rdx             & -              & -                       & \textcolor{red}{450}    & 450                                           & 450                                   \\
			\hline
		\end{array}
	\]
	Es posible que el lector no entienda el por qué de las cuatro instrucciones en vez de tres como en C, esto se debe a que, como hemos visto antes, no se pueden realizar
	operaciones de memoria a memoria, por lo que se necesita un registro para almacenar el valor de una de las direcciones de memoria.
\end{ejemplo}
\subsection{Operaciones}
Veamos algunas de las operaciones más comunes en x86-64, notemos que algunas de ellas además pueden tener
variantes indicando el tamaño de los operandos. Antes de ver las operaciones de forma general, recalcamos
la instrucción \textbf{\textit{leaq}}:
\subsubsection{leaq}
La instrucción (\textit{load effective address}) se utiliza como variante de la instrucción \textit{movq},
la instrucción lee de la memoria la dirección efectiva de un operando y la almacena en un registro, sin acceder a la memoria en sí.
Por ejemplo \textit{leaq 0x8(\%rax), \%rdx} almacena en $\%rdx$ la dirección de memoria $\%rax + 0x8$. Esta instrucción es a menudo utilizada
por los compiladores para optimizaciones.
\begin{ejemplo}
	Veamos una forma para multiplicar un numero por 12 usando leaq
	\begin{center} % Centramos el contenido
		\begin{minipage}{0.3\textwidth}
			\begin{minted}[fontsize=\large]{c}
void m12(long x) 
{
    return x*12;
}
                \end{minted}
		\end{minipage}
		\begin{minipage}{0.69\textwidth}
			\begin{minted}[fontsize=\large]{asm}
swap:
    leaq (%rdi, %rdi, 2), %rax # t = x + 2x
    salq 2, %rax             # return t<<2
                \end{minted}
		\end{minipage}
	\end{center}
\end{ejemplo}

Veamos algunas de las operaciones del tipo aritmetico-lógicas
\begin{center}
    \begin{tabular}{|c|c|c|}
    \hline
    \textbf{Instrucción} & \textbf{Efecto}                       & \textbf{Descripción}     \\ \hline
    \texttt{lea S, D}         & $D \leftarrow \&S$                   & Load effective address   \\ \hline
    \texttt{inc D}            & $D \leftarrow D + 1$                 & Increment                \\ \hline
    \texttt{dec D}            & $D \leftarrow D - 1$                 & Decrement                \\ \hline
    \texttt{neg D}            & $D \leftarrow -D$                    & Negate                   \\ \hline
    \texttt{not D}            & $D \leftarrow \sim D$                & Complement               \\ \hline
    \texttt{add S, D}         & $D \leftarrow D + S$                 & Add                      \\ \hline
    \texttt{sub S, D}         & $D \leftarrow D - S$                 & Subtract                 \\ \hline
    \texttt{imul S, D}        & $D \leftarrow D * S$                 & Multiply                 \\ \hline
    \texttt{xor S, D}         & $D \leftarrow D \oplus S$            & Exclusive-or             \\ \hline
    \texttt{or S, D}          & $D \leftarrow D \vee S$              & Or                       \\ \hline
    \texttt{and S, D}         & $D \leftarrow D \wedge S$            & And                      \\ \hline
    \texttt{sal k, D}         & $D \leftarrow D \ll k$               & Left shift               \\ \hline
    \texttt{shl k, D}         & $D \leftarrow D \ll k$               & Left shift (same as sal) \\ \hline
    \texttt{sar k, D}         & $D \leftarrow D \gg_A k$             & Arithmetic right shift   \\ \hline
    \texttt{shr k, D}         & $D \leftarrow D \gg_L k$             & Logical right shift      \\ \hline
    \end{tabular}
\end{center}

Destacamos una característica peculiar de las operaciones de shift, en el caso de que hagamos un shift a la derecha sobre un
registro de 32 bits \texttt{shr \$16 $\%$edx}, $\%$rdx se rellenará con 0 en los primeros 32 bits también, de esta forma si teniamos $\$-1$
por ejemplo y realizamos \texttt{shr \$16 $\%$edx}, nos quedaremos con 0x0000 0000 0000 FFFF en vez de 0xFFFF FFFF 0000 FFFF.
\subsubsection{Operaciones aritméticas especiales}
Multiplicar dos enteros con/sin signo de 64 bits puede tener un resultado que necesite 128 bits para representarlo
es para ello que se incluyen las siguientes instrucciones especiales:
\begin{center}
    \begin{tabular}{|c|c|c|}
    \hline
        \textbf{Instruction} & \textbf{Effect} & \textbf{Description} \\ \hline
    \texttt{imulq S} & \texttt{R[\%rdx]:R[\%rax] ← S × R[\%rax]} & Signed full multiply \\ \hline
    \texttt{mulq S} & \texttt{R[\%rdx]:R[\%rax] ← S × R[\%rax]} & Unsigned full multiply \\ \hline
    \texttt{cqto} & \texttt{R[\%rdx]:R[\%rax] ← SignExtend(R[\%rax])} & Convert to oct word \\ \hline
    \texttt{idivq S} & 
    \begin{tabular}[t]{@{}l@{}}
        \texttt{R[\%rdx] ← R[\%rdx]:R[\%rax] mod S;} \\ 
        \texttt{R[\%rax] ← R[\%rdx]:R[\%rax] ÷ S}
    \end{tabular} & Signed divide \\ \hline
    \texttt{divq S} & 
    \begin{tabular}[t]{@{}l@{}}
        \texttt{R[\%rdx] ← R[\%rdx]:R[\%rax] mod S;} \\ 
        \texttt{R[\%rax] ← R[\%rdx]:R[\%rax] ÷ S}
    \end{tabular} & Unsigned divide \\ \hline
    \end{tabular}
\end{center}
Además de esto, añadimos ciertas peculiaridades que nos pueden ser útiles en la práctica:

\begin{itemize}
    \item Las instrucciones \texttt{movzbl} y \texttt{movslq} (la \texttt{z} indica \textit{zero extension} y la \texttt{s} \textit{sign extension}) se utilizan para extender cantidades a registros. Lo remarcable es que, mientras que \texttt{movslq} mueve 32 bits a 64 bits con extensión de signo, \texttt{movzlq} no existe ya que al mover 32 bits a 64 bits (\texttt{movl}) ya se realiza una extensión de ceros.
    \item Cuando se tiene el código desensamblado de un programa, podemos ver la dirección en hexadecimal, junto con varios pares de números en hexadecimal y a la derecha las instrucciones. Los pares de números en hexadecimal representan los bytes de la instrucción. Por ejemplo, si tenemos la línea siguiente:
    \begin{center}
        400544: \quad \texttt{e8 08 00 00 00} \quad \texttt{callq 400551 <f>}
    \end{center}
    Podremos ver que \texttt{e8} corresponde con la instrucción \texttt{callq} y los siguientes 4 bytes son la dirección a la que se llama. En este caso, \texttt{400544} + 8 = \texttt{400551}, como se puede ver en la instrucción a la derecha.
\end{itemize}


\section{Control}
Hasta ahora hemos considerado como se comporta el código cuando se ejecuta de forma secuencial, sin embargo,
algunas estructuras de $C$ como los bucles, las estructuras de control, etc. requieren de una ejecución condicional, esto es,
la secuencia de instrucciones a ejecutar depende de una condición.
En x86-64, veremos dos formas de implementar estas estructuras de control más adelante, primero, veamos
los códigos de condición que se utilizan en x86-64:
\subsection{Codigos de condición}
Además de los registros de propósito general, la CPU mantiene un conjunto de bits de estado que se utilizan para
describir atributos de la última operación aritmético-lógica realizada.
Veamos los códigos de condición que vamos a usar:
\begin{itemize}
	\item \textbf{CF (Carry Flag)}: El flag de acarreo, se activa si hay un acarreo en la última operación. Se utiliza para detectar overflow en operaciones sin signo.
	\item \textbf{ZF (Zero Flag)}: El flag de cero, se activa si el resultado de la última operación es cero.
	\item \textbf{SF (Sign Flag)}: El flag de signo, se activa si el resultado de la última operación es negativo.
	\item \textbf{OF (Overflow Flag)}: El flag de desbordamiento, se activa si hay un desbordamiento en la última operación en complemento a 2.
\end{itemize}
Estos códigos de condición se ajustan implicitamente por las operaciones aritméticas, de esta forma, \textit{leaq} no afecta a los códigos de condición.
Además de ajustarse por las operaciones que ya conocemos, existen instrucciones que ajustan explicitamente los códigos de condición sin alterar los registros de propósito general.
\begin{itemize}
	\item El conjunto de instrucciones \textbf{\textit{Compare}} (cmp) ajustan los códigos de condición de acuerdo a las diferencias entre sus dos operandos. Debemos tener en cuenta que en $AT\&T$ la sintaxis
	      es al revés: \textit{cmpq S1, S2} compara $S1$ con $S2$ mediante $S2- S1$.
	\item El conjunto de instrucciones \textbf{\textit{Test}} (test) ajustan los códigos de condición de acuerdo a las operaciones lógicas entre sus dos operandos.
	      Se realiza la operación $S1 \& S2$ y se ajustan los códigos de condición. Es común que se utilicen ambos operandos iguales, por ejemplo, para
	      comprobar si un número es 0.
\end{itemize}
\subsubsection{Consultar códigos de condición}
En vez de consultarlos directamente, se utilizan mediante tres formas:
\begin{enumerate}
	\item Se establece un único byte a 0 o 1 en función de una combinación de los códigos de condición.
	\item Se utilizan para realizar saltos condicionales.
	\item Se utilizan para transferir datos condicionalmente.
\end{enumerate}
Para la primera forma, se utilizan las instrucciones \textit{setcc} siguientes:
\begin{center}
	\begin{tabular}{|c|c|}
		\hline
		\textbf{Instruction} & \textbf{Effect}                                                          \\ \hline
		\texttt{sete D}      & D $\leftarrow$ ZF (Equal / zero)                                         \\ \hline
		\texttt{setne D}     & D $\leftarrow$ $\sim$ ZF (Not equal / not zero)                          \\ \hline
		\texttt{sets D}      & D $\leftarrow$ SF (Negative)                                             \\ \hline
		\texttt{setns D}     & D $\leftarrow$ $\sim$ SF (Nonnegative)                                   \\ \hline
		\texttt{setg D}      & D $\leftarrow$ $\sim$ (SF $\oplus$ OF) \& $\sim$ ZF (Greater, signed >)  \\ \hline
		\texttt{setge D}     & D $\leftarrow$ $\sim$ (SF $\oplus$ OF) (Greater or equal, signed $\geq$) \\ \hline
		\texttt{setl D}      & D $\leftarrow$ SF $\oplus$ OF (Less, signed <)                           \\ \hline
		\texttt{setle D}     & D $\leftarrow$ (SF $\oplus$ OF) $\lor$ ZF (Less or equal, signed $\leq$) \\ \hline
		\texttt{seta D}      & D $\leftarrow$ $\sim$ CF \& $\sim$ ZF (Above, unsigned >)                \\ \hline
		\texttt{setae D}     & D $\leftarrow$ $\sim$ CF (Above or equal, unsigned $\geq$)               \\ \hline
		\texttt{setb D}      & D $\leftarrow$ CF (Below, unsigned <)                                    \\ \hline
		\texttt{setbe D}     & D $\leftarrow$ CF $\lor$ ZF (Below or equal, unsigned $\leq$)            \\ \hline
	\end{tabular}
\end{center}
\begin{ejemplo}
	Veamos un ejemplo de su uso:
	\begin{center}
		\begin{minipage}{0.39\textwidth}
			\begin{minted}[fontsize=\Large]{c}
    long gt(long a, long b)
    {
        return a > b;
    }
            \end{minted}
		\end{minipage}
		\begin{minipage}{0.6\textwidth}
			\begin{minted}[fontsize=\Large]{asm}
    gt:
        cmpq %rsi, %rdi
        setg %al
        movzbq %al, %rax
        ret

            \end{minted}
		\end{minipage}
	\end{center}
	Donde movzbq es una instrucción que significa \textit{ Move with Zero-extend Byte to Long}.
\end{ejemplo}

Veamos ahora cómo se realizan los saltos condicionales.
Estos provocan que la ejecución del programa pueda saltar a otra parte del código en función de los códigos de condición, es muy posible que
en algún momento nos encontremos una instrucción en C de tipo goto, su uso está desaconsejado, sin embargo, por su similitud con el ensamblador,
puede resultar útil en ciertas situaciones. Veamos una tabla con las instrucciones de salto condicional y después un ejemplo de uso.

\begin{center}
	\begin{tabular}{|c|c|c|}
		\hline
		\textbf{Instruction} & \textbf{Jump Condition}                 & \textbf{Description}         \\ \hline
		jmp                  & Label                                   & Direct jump                  \\ \hline
		jmp                  & *Operand                                & Indirect jump                \\ \hline
		je                   & ZF                                      & Equal / zero                 \\ \hline
		jne                  & $\sim$ZF                                & Not equal / not zero         \\ \hline
		js                   & SF                                      & Negative                     \\ \hline
		jns                  & $\sim$SF                                & Nonnegative                  \\ \hline
        jg                   & $\sim$(SF \^{} OF) $\&$ $\sim$ZF & Greater (signed >)           \\ \hline
        jge                  & $\sim$(SF \^{} OF)                  & Greater or equal (signed >=) \\ \hline
        jl                   & SF \^{} OF                          & Less (signed <)              \\ \hline
        jle                  & (SF \^{}  OF) $|$ ZF              & Less or equal (signed <=)    \\ \hline
        ja                   & $\sim$CF $\&$ $\sim$ZF               & Above (unsigned >)           \\ \hline
		jae                  & $\sim$CF                                & Above or equal (unsigned >=) \\ \hline
		jb                   & CF                                      & Below (unsigned <)           \\ \hline
		jbe                  & CF $|$ ZF                            & Below or equal (unsigned <=) \\ \hline
	\end{tabular}
\end{center}
\begin{ejemplo}
	Veamos ahora un ejemplo de cómo hacer por ejemplo un valor absoluto de una diferencia:
	\begin{center}
		\begin{minipage}{0.39\textwidth}
			\begin{minted}[fontsize=\Large]{c}
    long abs(long x, long y)
    {
        if (x > y)
            return x - y;
        else
            return y - x;
    }
            \end{minted}
		\end{minipage}
		\begin{minipage}{0.6\textwidth}
			\begin{minted}[fontsize=\Large]{asm}
        abs:
            cmpq %rsi, %rdi
            jle .L1
            movq    %rdi, %rax
            subq    %rsi, %rax
            ret
        .L1:
            movq %rsi, %rax
            subq %rdi, %rax
            ret
			\end{minted}
		\end{minipage}
	\end{center}
	Si tratamos de \href{https://godbolt.org/z/aj7f7os8d}{replicar} el ejemplo anterior, veremos que en general, el compilador no produce este código,
	esto se debe a las optimizaciones que el compilador aplica, evitando en la medida de lo posible los saltos condicionales, ya que muchas veces,
	es mejor calcular ambas ramas, a tener que predecir y calcular cual de las ramas hay que seguir.
	Para obtener este mismo código, deberemos especificar la opción de compilación \textit{-fno-if-conversion}.
\end{ejemplo}
Nos detenemos brevemente en las optimizaciones mencionadas, lo que hace el compilador en vez de saltar ramas es asignar valores a las variables
en función de un test: \texttt{let val = test ? x : y}. \textit{gcc} intenta utilizar este tipo de optimizaciones siempre que sea seguro.
De esta forma, el código anterior se traduciría a:
\begin{center}
	\begin{minted}[fontsize=\Large]{asm}
        abs:
            movq %rdi, %rax # x
            subq %rsi, %rax # x - y
            movq %rsi, %rdx 
            subq %rdi, %rdx # y - x
            cmpq %rsi, %rdi # x : y
            cmovle %rdx, %rax # si x <= y, rax = rdx
    \end{minted}
\end{center}
Esto se verá como una optimización en la asignatura de \textit{Arquitectura de Computadores}, sin embargo, adelantamos
que por ejemplo \texttt{ let val = test ? p : *q} no se puede optimizar de la misma forma, ya que pudiera tener efectos no deseables.
Notemos además que este tipo de optimizaciones solo tienen sentido cuando los cálculos son sencillos, en cuyo caso es preferible calcular las
dos ramas a predecir cual de ellas seguir. Es por este tipo de opciones que \textit{gcc} es posible que genere un código mas lento cuando se activan
las optimizaciones \textit{-O3} ya que si la rama era muy predecible, la predicción de saltos es mas rápida que calcular las dos ramas.
Tampoco se debe utilizar en casos donde una rama u otra tiene efectos colaterales, ya que eso pudiera cambiar el comportamiento del programa.
\subsection{Bucles}
La mayoría de lenguajes ofrecen varias formas de implementar bucles: \textit{ do-while, while, for, loop, etc.} En x86-64, no existen instrucciones
similares, en su lugar, se combinan diferentes tests y saltos condicionales para implementar los bucles.
\subsubsection{Bucles Do-While}
En C, un bucle \textit{do-while} se implementa de la siguiente forma:
\begin{center}
	\begin{minipage}{0.45\textwidth}
		\begin{minted}[fontsize=\Large]{c}
long pcount(long x)
{
    long result = 0;
    do
    {
        result += x & 0x1;
        x >>= 1;
    } while (x);
    return result;
}
        \end{minted}
		\descriptionlabel{Versión C}
	\end{minipage}
	\begin{minipage}{0.45\textwidth}
		\begin{minted}[fontsize=\Large]{c}
long pcount(long x)
{
    long result = 0;
    loop:
        result += x & 0x1;
        x >>= 1;
        if(x) goto loop;
    return result;
}
        \end{minted}
		\descriptionlabel{Versión goto}

	\end{minipage}
\end{center}
En este ejemplo se ha hecho un clásico programa popcount que cuenta el número de bits a 1 en un número.
El nombre \textit{popcount} viene de \textit{population count}, que es el nombre que se le da a esta operación.
Como ya sabemos, el efecto del loop es ejecutar el cuerpo del bucle, evaluar la condición y si el resultado de la condición es distinto de 0,
seguir ejecutando el bucle. Notemos que se ejecuta al menos una vez.
\subsubsection{Bucles While}
\begin{ejemplo}
	Veamos ahora cómo se implementa un bucle \textit{while}:
	\begin{center}
		\begin{minipage}{0.45\textwidth}
			\begin{minted}[fontsize=\Large]{c}
long pcount(long x)
{
    long result = 0;
    while(x) {
        result += x & 0x1;
        x >>= 1;
    }

    return result;
}
\end{minted}
			\descriptionlabel{Versión C}
		\end{minipage}
		\begin{minipage}{0.45\textwidth}
			\begin{minted}[fontsize=\Large]{c}

long pcount(long x)
{
        long result = 0;
        goto test;
    loop:
        result += x & 0x1;
        x >>= 1;
    test:
        if (x) goto loop;
        return result;
}
        \end{minted}
			\descriptionlabel{Versión goto}

		\end{minipage}
	\end{center}
\end{ejemplo}
En general, while evalúa la condición antes de ejecutar el bucle, por lo que si la condición no se cumple, el bucle no se ejecuta.
\subsubsection{Bucles For}
Sabemos ya que en general un bucle \textit{for} se compone de tres partes, la inicialización, la condición y la actualización.
No adjuntamos el código en C ya que es trivial sabiendo los anteriores, notemos que un bucle \textit{for} es equivalente a un bucle \textit{while} de la siguiente forma:
\textit{for (init; test; update) body} es equivalente a \textit{init; while(test) {body; update;}}
Por último, veamos las secuencias switch.
\subsection{Secuencias Switch}
Una secuencia switch proporciona una ramificación condicional múltiple basada en el valor de una expresión entera.
Son muy útiles cuando trabajamos con tests que pueden tener una gran cantidad de resultados, haciendo el código más legible
y eficiente. Este último punto puede sorprender, por ello, vamos a desarrollarlo.
\\ \\
Una \textit{tabla de salto} es una estructura de datos en forma de array donde $i$ es la dirección de un segmento de codigo que ejecuta la acción
que el programa debería seguir cuando el valor de la expresión es $i$. La diferencia entre multiples \textit{if-else} y \textit{switch} radica en
que el tiempo que se tarda en ejecutar un \textit{switch} es constante, mientras que en un \textit{if-else} es lineal en función del número de casos.

\begin{center}
	\begin{minipage}{0.45\textwidth}
		\centering
		\textbf{Código en C}

		\begin{minted}[fontsize=\large]{c}

void switch_eg(long x) {
    switch (x) {
        case 1:
            printf("Code 1\n");
            break;
        case 2:
            printf("Code 2\n");
            break;
        case 3:
            printf("Code 3\n");
            break;
        default:
            printf("Code 4\n");
    }
}

    \end{minted}

	\end{minipage}
	\hfill
	\begin{minipage}{0.45\textwidth}
		\centering
		\textbf{Assembly Code}

		\begin{minted}[fontsize=\large]{asm}
cmpq $3, %rdi
ja .L4 # Si x > 3, Default
jmp *.L1(,%rdi,8) # Dirección L1 + 8x 
.L1:
.L2:
.L3:
.L4 # Default
        \end{minted}

		\vspace{1em}
		\textbf{Branch Table (Jump Table)}

		\begin{tabular}{|c|c|}
			\hline
			\textbf{Índice} & \textbf{Dirección de Salto} \\
			\hline
			1               & Code Block 1                \\
			\hline
			2               & Code Block 2                \\
			\hline
			3               & Code Block 3                \\
			\hline
			Default         & Code Block 4                \\
			\hline
		\end{tabular}
	\end{minipage}
\end{center}
Como podemos ver, la tabla de saltos comienza en este caso en .L1, para acceder
a el resto de bloques de código, se suma 8 veces el valor de x, ya que cada dirección
ocupa 8 bytes. Si x es mayor que 3, se salta a la dirección .L4, que es el bloque por defecto.
Si por ejemplo en vez de hacer break, se quisiera hacer un fall through, obtendríamos algo
del siguiente estilo:
\begin{center}
	\begin{minipage}{0.45\textwidth}
		\begin{minted}[fontsize=\large]{c}
long w = 1;
switch(x) {
    ...
  case 2:
    w = y/z;
    /* Fall through */ 
case 3:
    w += z;
    break;
    ...
}
\end{minted}
	\end{minipage}
	\begin{minipage}{0.45\textwidth}
		\begin{equation*}
			\%rdi \rightarrow x \quad \%rsi \rightarrow y \quad \%rdx \rightarrow z
		\end{equation*}
		\begin{minted}[fontsize=\large]{asm}
.L5: # Case 2
    movq %rsi, %rax
    cqto # Convert Quad To Oct
    idivq %rdx 
    jmp .L6
.L9: # Case 3
    movl $1, %eax
.L6: # Merge
    addq %rcx, %rax # w += z
    ret 

\end{minted}
	\end{minipage}
\end{center}

\section{Procedimientos}
Los procedimientos son una abstracción clave en el desarrollo de software, permiten dividir un programa en pequeñas partes
donde cada una de ellas proporciona una funcionalidad concreta al proporcionarle una serie de argumentos.
Pueden ser llamados desde cualquier parte del programa, lo que permite reutilizar el código y hacerlo más modular. Los procedimientos vienen
de distintas formas en cada lenguaje: funciones, métodos, subrutinas, lambdas, etc, pero todos ellos comparten ciertos elementos comunes.
\\ \\
Existen muchos atributos que hay que manejar cuando se proporciona soporte en lenguaje ensamblador para procedimientos, para verlos, supongamos un proceso
$P$ que llama a un procesor $Q$, que se ejecuta y tras esto vuelve a $P$, estas acciones incluyen:
\begin{itemize}
	\item \textbf{Pasar Control}: El PC tiene que fijarse en la dirección de comienzo de $Q$, y tras esto, cuando $Q$ termina, el PC tiene que volver a la dirección de la instrucción que sigue a la llamada.
	\item \textbf{Parámetros}: $P$ tiene que pasar los parámetros a $Q$, y $Q$ tiene que devolver un valor a $P$.
	\item \textbf{Reserva y Limpieza de Memoria}: $Q$ puede necesitar reservar memoria para sus propias variables, y tras la ejecución, limpiarla.
\end{itemize}
La implementación de estos procedimientos en x86-64 involucra una combinación de instrucciones especiales y convenciones de cómo usar los recursos de la máquina.
Detrás de esto hay un gran esfuerzo por minimizar el overhead que involucra una llamada, esto provoca que al final, solo se ejecuta lo mínimo indispensable mencionado
antes. Veamos ahora diferentes mecanismos paso por paso, comenzando por la pila:
\subsection{Stack (Pila)}
Es una región de memoría gestionada con aspectos de LIFO, es decir, el último elemento en entrar es el primero en salir, crece hacia direcciones inferiores, en x86-64
el registro $\%rsp$ contiene la dirección mas baja de la pila. Los datos, pueden ser accedidos y extraidos de la pila usando las instrucciones \textit{pushq} y \textit{popq}.
Por ejemplo para reservar espacio para datos sin valor inicial, se puede decrementar la pila la cantidad necesaria, análogamente, se puede liberar el espacio incrementando el stack pointer.
En x86-64, cuando un procedimiento requiere más espacio de lo que puede guardar en registros, reserva espacio en el stack. A esta región se la llama \textit{procedure space} para las variables locales.
En interés del espacio y la eficiencia, los procedimientos de x86-64 reservan solo el espacio necesario.

\subsection{Pasando el Control}
Pasar el control de una función $P$ a otra $Q$ requiere únicamente de poner el PC en la dirección de inicio de $Q$,
sin embargo, cuando $Q$ termina, el PC tiene que tener alguna manera de saber dónde debe reanudar la ejecución de $P$, esta información se guarda en las máquinas
x86-64 usando la instrucción call $Q$. Esta instrucción añade una dirección $A$ al stack y pone el PC en el comienzo de $Q$. La dirección $A$ se denomina \textit{return address} y se calcula como
la instrucción que sigue inmediatamente a call. La instrucción contraria ret, elimina una dirección $A$ de la pila y pone el PC en $A$.
(Notemos que estas instrucciones nos las encontraremos como callq y retq, donde el sufijo $q$ simplemente enfatiza que son versiones de x86-64 no de IA32).
Veamos un esquema con la estructura general del stack:


\begin{center}
	\includesvg[width=0.6\textwidth]{Imagenes/stack}
	\label {fig:stack}
\end{center}
\subsection{Pasando los datos}
Además de pasar el control a un procedimiento cuando es invocado, y de vuelta cuando termina, las llamadas a procedimientos
pueden involucrar pasar datos como argumentos, además un procedimiento puede devolver también un valor. En x86-64, la mayoria de
este flujo de datos es llevada a cabo mediante registros, ya hemos visto numerosos ejemplos en los que $\%rdi$ y $\%rsi$ se pasan como
parámetros y valores que se devuelven mediante $\%rax$. En x86-64, se pueden pasar hasta 6 argumentos mediante registros, que se usan en un orden
específico con el nombre del registro en función del tamaño del dato que contiene.
Cuando una función tiene más de 6 argumentos, los argumentos restantes se pasan en la pila tal que el argumento 7 se coloca abajo, el 8 encima, y así sucesivamente.
\subsection{Gestionando datos locales}
Hasta ahora, la mayoría de ejemplos que hemos visto no han requerido de variables locales, ya sea por optimización del compilador o por la sencillez del código,
sin embargo, sabemos que muchas veces es necesario almacenar variables locales.
Normalmente, en lenguajes basados en pila, o que soportan la recursión (e.g C, C++, Java, Rust, etc), las variables locales se almacenan en la pila.
Esta pila, se reserva en marcos (allocated in frames), que son bloques de memoria que contienen las variables locales de un procedimiento, estos marcos de memoria
se organizan tal que contienen información de retorno, almacenamiento local si es necesario y espacio temporal si es necesario, además hay un puntero de marco opcional:
$\%rbp$. Veamos un ejemplo de cadena de llamadas a procedimientos:
% \begin{center}
% 	\begin{tikzpicture}[>=Stealth, scale=1.5]
%
% 		\matrix (m1) [matrix of nodes, nodes={draw, minimum size=1cm}, column sep=0cm] {
% 			\node[fill=green!20] {\phantom{a}}; \\
% 			\node[fill=green!20] {\phantom{a}}; \\
% 			\node[fill=green!20] {\phantom{a}}; \\
% 			\node[fill=green!20] {\phantom{a}}; \\
% 		};
%
% 		\matrix (m2) [matrix of nodes, right=of m1, nodes={draw, minimum size=1cm}, column sep=0cm] {
% 			\node[fill=green!20] {\phantom{a}}; \\
% 			\node[fill=green!40] {fn A};                    \\
% 			\node[fill=green!20] {\phantom{a}}; \\
% 			\node[fill=green!20] {\phantom{a}}; \\
% 		};
%
% 		\draw[->] (m1) -- (m2);
%
% 		\matrix (m3) [matrix of nodes, right=of m2, nodes={draw, minimum size=1cm}, column sep=0cm] {
% 			\node[fill=green!20] {\phantom{a}}; \\
% 			\node[fill=green!40] {fn A};                    \\
% 			\node[fill=green!70] {fn B};                    \\
% 			\node[fill=green!0] {\phantom{a}};  \\
% 		};
%
% 		\draw[->] (m2) -- (m3);
% 		\matrix (m4) [matrix of nodes, right=of m3, nodes={draw, minimum size=1cm}, column sep=0cm] {
% 			\node[fill=green!20] {\phantom{a}}; \\
% 			\node[fill=green!40] {fn A};                    \\
% 			\node[fill=green!70] {fn B};                    \\
% 			\node[fill=green!70] {fn B};                    \\
% 		};
%
% 		\draw[->] (m3) -- (m4);
%
% 		\matrix (m5) [matrix of nodes, right=of m4, nodes={draw, minimum size=1cm}, column sep=0cm] {
% 			\node[fill=green!20] {\phantom{a}}; \\
% 			\node[fill=green!40] {fn A};                    \\
% 			\node[fill=green!70] {fn B};                    \\
% 			\node[fill=green!20] {\phantom{a}}; \\
% 		};
%
% 		\draw[->] (m4) -- (m5);
% 		\matrix (m6) [matrix of nodes, right=of m5, nodes={draw, minimum size=1cm}, column sep=0cm] {
% 			\node[fill=green!20] {\phantom{a}}; \\
% 			\node[fill=green!40] {fn A};                    \\
% 			\node[fill=green!20] {\phantom{a}}; \\
% 			\node[fill=green!20] {\phantom{a}}; \\
% 		};
%
% 		% Flecha entre la cuarta y quinta tabla
% 		\draw[->] (m5) -- (m6);
% 	\end{tikzpicture}
% \end{center}
Donde en cada uno de los pasos, $\%rsp$ apunta en la función más abajo existente y $\%rbp$ apunta al marco de la función inmediata anterior.
Los registros de programa actuan como un único recurso compartido por todos los procedimientos, aunque un solo procedimiento puede estar activo en un momento dado,
tenemos que asegurarnos de que cuando un procedimiento (caller) llama otro procedimiento (callee), el callee no sobreescriba los registros del caller. Veamos un ejemplo de esto:
\begin{center}
	\begin{minipage}{0.45\textwidth}
		\begin{minted}[fontsize=\large]{asm}
funcA:
    movq $15213, %rdx
    call funcB
    addq %rdx, %rax 
    ret
\end{minted}
	\end{minipage}
	\begin{minipage}{0.45\textwidth}
		\begin{minted}[fontsize=\large]{asm}
funcB: 
    subq $18213, %rdx
    ret 
    \end{minted}

	\end{minipage}
\end{center}
Como podemos ver, pudiera pasar que $\%rdx$ fuese sobreescrito en la función $funcB$, para evitar esto,
existen convenciones para coordinar el uso de registros.
Los registros $\%rbx, \%rbp, \%r12-\%r15$ se denominan \textit{callee-saved} ya que el callee debe guardar su valor si los modifica.
Para este fin, el callee puede o bien no modificar los registros, o guardando su valor en la pila y restaurándolo antes de salir, a este espacio se le llama
\textit{Saved Registers}, podemos verlo en el diagrama \ref{fig:stack}.
El resto de registros, excepto $\%rsp$, se clasifican como \textit{Caller-saved}, lo que
significa que pueden ser modificados por cualquier función, el nombre viene de un procedimiento $P$ que llama a un $Q$, como
$Q$ puede modificar libremente estos registros, $P$ debe guardar su valor si los necesita.
\subsection{Recursividad}
Las convenciones que hemos visto hasta ahora son suficientes para permitir a los procedimientos de x86-64 llamarse a sí mismos.
Cada llamada a procedimiento tiene su propio espacio en el stack, luego las variables locales de las múltiples llamadas no interfieren entre ellas.
Veamos un ejemplo del código C y ensamblador generado para una función recursiva:
\begin{center}
	\begin{minipage}{0.45\textwidth}
		\begin{minted}[fontsize=\large]{c}
long pcount_r(unsigned long x) {
    if (x == 0) {
        return 0;
    } else {
        return (x & 1) 
            + pcount_r(x >> 1);
    }
}
\end{minted}
	\end{minipage}
	\begin{minipage}{0.45\textwidth}
		\begin{minted}[fontsize=\large]{asm}
pcount_r:
    movl $0, %eax
    testq %rdi, %rdi
    je .L2
    pushq %rbx 
    movq %rdi, %rbx
    andl $1, %ebx
    shrq %rdi
    call pcount_r
    addq %rbx, %rax
    popq %rbx
  .L2
    rep; ret # Para branch prediction 
  \end{minted}
	\end{minipage}
\end{center}
Veamos más detenidamente lo que está pasando:
%Error 3
% \begin{itemize}
% 	\item Las primeras tres instrucciones sirven para saber cuando acaba la recursión, si x es 0, se devuelve 0.
% 	\item El pushq está ya que se va a modificar $\%rbx$ y se necesita restaurar su valor al salir de la función.
% 	\item Desde movq hasta shrq se calcula x $\&$ 1 y x >> 1.
% 	\item x >> 1 se pasa como argumento a la función $pcount_r$ (llamada recursiva)
% 	\item Se suma el resultado de la llamada recursiva con el valor de x \& 1.
% \end{itemize}
\section{Datos}
% Marker [DE AQUI HACIA ABAJO TODO CORRECTO]
\subsection{Arrays}
Los arrays en C se usan para almacenar colecciones de datos escalares en variables más grandes, la implementación particular
de C de los arrays es muy directa, una característica importante es que en C podemos manipular punteros a elementos de un array
y realizar operaciones con ellos. Los compiladores son extremadamente buenos simplificando el cálculo de direcciones 
usado cuando accedemos a un array, es por ello que la correspondencia entre C y ensamblador puede llegar a ser dificil de descifrar.
\subsubsection{Principios Básicos}
Sea $T$ un tipo de dato cualquiera, y $N$ una constante, consideramos la declaración $T A[N]$, entonces $A$ es un array de $N$ elementos.
La declaración reserva una región contigua de memoria de $T \cdot N$ bytes e introduce un identificador $A$ que se usa como puntero
al comienzo de esta región de memoria. Los elementos de este array se acceden mediante el uso de un índice entre 0 y N-1.
\begin{ejemplo}
    Consideramos \texttt{int arr[5]}, veamos los siguientes ejemplos:
    \begin{center}
        \includesvg[width=0.5\textwidth]{Imagenes/array}
        \end{center}
    \begin{itemize}
        \item \texttt{arr[4]} (int) se refiere al quinto elemento del array (valor 3).
        \item \texttt{arr} (int *) comienzo del array (posición arr)
        \item \texttt{arr+1} (int *) comienzo del array + 1 (posición arr + 4) 
        \item \texttt{\&arr[2]} (int *) dirección de val[2] (posición arr + 8)
        \item \texttt{arr[5]} (int) fuera de los límites del array, \href{https://en.cppreference.com/w/c/language/behavior}{undefined behavior.}
        \item \texttt{*(arr+1)} (int) valor de la dirección val + 1 (valor 5)
    \end{itemize}
\end{ejemplo}
Destacamos que declaramos arrays inmediatamente después de otro, no está garantizado que estos se almacenen en memoria de forma contigua. Así como hemos podido ver en el ejemplo anterior,
que C permite operaciones con los punteros, así, si $p$ es un puntero a un array de tipo $T$ y la dirección de memoria de $p$ es $x_p$, entonces la expresión $p + i$ tiene el valor \texttt{$x_p$ +  $T \cdot i$ = $i \cdot$ sizeof(T)}.
Recordamos que los operadores $\& $ y $*$ permiten la generación y dereferenciación de punteros respectivamente.
\begin{ejemplo}
    Ejemplo de acceso a elementos de array:
    \begin{center}
        \begin{minipage}{0.45\textwidth}
            \begin{minted}[fontsize=\large]{c}
int get_num(int *arr, int i) 
{
    return arr[i];
}
        \end{minted}
        \end{minipage}
        \begin{minipage}{0.45\textwidth}
        \begin{minted}[fontsize=\large]{asm}
    get_num:
        movslq  %esi, %rsi
        movl (%rdi, %rsi, 4), %eax
        ret
            \end{minted}
        \end{minipage}
    \end{center}
    Recordamos que \texttt{movslq} es una instrucción que significa \textit{Move with Sign-extend Long to Quad}, se utiliza ya que para hacer 
    operaciones con la dirección de memoria, es necesario que los registros sean de 64 bits.
\end{ejemplo}

\subsubsection{Arrays Anidados}
Los principios de los arrays multidimensionales son los mismos que los de los arrays unidimensionales, de esta manera, la declaración
\texttt{T A[R][C]} declara un array de $R$ elementos, cada uno de los cuales es un array de $C$ elementos. En general, un array bidimensional de elementos de tipo T. De otra forma,
podemos decir que tiene $R$ filas y $C$ columnas. Notemos que si el tamaño de $T$ es $K$ bytes, entonces el tamaño de $A$ es $ C \cdot R \cdot K$ bytes.
La disposición de los elementos en memoria es por filas, \textit{(row-major-order)}:
\begin{center}
\includesvg[width=0.8\textwidth]{Imagenes/array2d.svg}
\end{center}
De esta forma, si tenemos un array \texttt{T A[R][C]} el elemento 
\texttt{T[i][j]} está en la posición de memoria \texttt{$\&$A[i][j] = $x_A$ + L(C $\cdot i + j$)}.
\begin{ejemplo}
    Veamos un ejemplo en el que el flag --no-pie debe ser considerado
    \begin{center}
        \begin{minipage}{0.45\textwidth}
            \begin{minted}[fontsize=\large]{c}
    int *get_row(int index) 
    {
        return arr[index];
    }
            \end{minted}
        \end{minipage}
        \begin{minipage}{0.45\textwidth}
            \begin{minted}[fontsize=\large]{c}
#define ROWS 4
// Zip_dig -> Array de 5 elementos
zip_dig arr[ROWS] = {
    {1, 2, 7, 4, 5},
    {6, 7, 2, 9, 1},
    {1, 9, 3, 1, 5},
    {1, 7, 8, 9, 2}};
            \end{minted}
        \end{minipage}
        \vspace{0.5cm}
        \vspace{0.5cm}
        \begin{minipage}{0.45\textwidth}
            \begin{minted}[fontsize=\large]{asm}
    get_num:
        leaq (%rdi, %rdi, 4), %rdx
        leaq arr(,%rdx, 4), %rax
        ret
            \end{minted}
        \end{minipage}
        \begin{minipage}{0.45\textwidth}
            \begin{minted}[fontsize=\large]{asm}
    get_num:
        movslq %edi, %rdi
        leaq arr(%rdi ,%rdi, 4), %rdx
        leaq arr(%rip), %rax
        leaq (%rax, %rdx, 4), %rax
        ret
        \end{minted}
        \end{minipage}
    \end{center}
    A la izquierda, como se usa --no-pie que desactiva el \textit{Position Independent Code}, la dirección de memoria es conocida en tiempo de compilación, por lo que
    se puede acceder directamente con $ 5 \cdot \text{index}$ por ser 5 elementos por fila y $4 \cdot \text{index}$ por ser un tipo de dato entero de 4 bytes.
    En el segundo caso, el direccionamiento es relativo a contador de programa $\%rip$. Notemos que --no-pie resulta más eficiente sin embargo es más vulnerable a ataques de buffer overflow.
\end{ejemplo}
\subsubsection{Array Multinivel}
De manera similar a los arrays bidimensionales, tenemos los arrays multinivel, que tienen la peculiaridad de que 
se tiene un array de punteros a arrays, de esta forma, para acceder a un dato son necesarios dos accesos a memoria, uno para acceder al puntero y otro para acceder al dato.
\begin{ejemplo}
    Un ejemplo de array multinivel en C: 
    \begin{center}
        \begin{minted}{c}
zip_dig cmu = {0, 5, 2, 7, 3};
zip_dig mit = {1, 2, 3, 4, 5};
zip_dig ucb = {9, 3, 7, 5, 1};

# define UCOUNT 3
int *univ[UCOUNT] = {cmu, mit, ucb};
        \end{minted}
    \end{center}
\end{ejemplo}
\subsection{Estructuras}
C permite dos formas de definir tipos de datos compuestos: \textit{struct} que agrupan múltiples datos en un 
solo objeto y \textit{union} que permiten que un objeto sea referenciado usando diferentes tipos de datos. Veamos en detalle las estructuras.
\\ \\ 
El struct en C crea un tipo de dato que agrupa múltiples variables que pueden ser de diferentes tipos en un solo objeto,
los distintos elementos de la estructura se referencian mediante nombres, con la sintaxis \texttt{struct.field, struct->field}.
El struct se parece a los arrays en el sentido de que se garantiza que los datos de un mismo struct están contiguos en memoria.
Es el compilador el que mantiene información sobre cada estructura para indicar el offset de cada campo. El tamaño de los 
structs es el tamaño de todos sus campos sumados, teniendo en cuenta que es posible que el compilador reserve espacio adicional tras un 
tipo de dato para alinear la memoria.
\begin{ejemplo}
    Veamos un ejemplo de struct en C y cómo se distribuye su memoria:
    \begin{center}
        \begin{minipage}{0.45\textwidth}
            \begin{minted}[fontsize=\large]{c}
        struct rec {
            int a[4];
            size_t i;
            struct rec *next;
        }
            \end{minted}
        \end{minipage}
        \begin{minipage}{0.48\textwidth}
            \includesvg[width=0.8\textwidth]{Imagenes/struct}
        \end{minipage}
    \end{center}
        Notemos que si en el struct el array de enteros fuese de tamaño 3, el tamaño del struct se mantendría en 32
        B por el alineamiento de la memoria. El puntero $r$ indica el primer byte de la estructura, se accede
        mediante el offset, que como hemos dicho lo determina el compilador en tiempo de compilación.
        \begin{center}
        \begin{minipage}{0.48\textwidth}
            \begin{minted}[fontsize=\large]{c}
    void set_i(struct rec *r, size_t val)
    {
        r->i = val;
    }
            \end{minted}
        \end{minipage}
        \begin{minipage}{0.45\textwidth}
            \begin{minted}[fontsize=\large]{asm}
        set_i:
            movq %rsi, (%rdi, 16)
            ret
            \end{minted} 
        \end{minipage}
    \end{center}
\end{ejemplo}
Hemos introducido ya el alineamiento de los datos, detallamos los principios que rigen su funcionamiento:
\begin{itemize}
    \item Si el tipo de dato primitivo requiere $K$ bytes, la dirección debe ser múltiplo de $K$.
    \item El alineamiento es requisito en algunas maquinas, en x86-64 es recomendado por motivos de eficiencia (Es ineficiente cargar o almacenar datos que cruzan fronteras de quad word y
        la memoria virtual es delicada cuando un dato se extiende a 2 páginas).
    \item En la mayoría de lenguajes compilados (C, C++, Rust, etc) es el compilador el que inserta huecos en la estructura para alinear los campos.
    \item El alineamiento se hace:
        \begin{itemize}
            \item Para el struct, se toma el máximo valor miembro para alinear, esto se llama alineamiento externo.
            \item Para los datos primitivos, se toma el tamaño del dato, es decir se añade padding antes del dato miembro si fuese necesario para que se encuentre en una posición múltiplo de su tamaño.
        \end{itemize}
\end{itemize}
El compilador no reordena los datos con el fin de ahorrar espacio, sin embargo
el programador debería considerar siempre cómo ordenar los datos para ocupar el menor espacio posible, esto es 
especialmente importante ya que a diferencia de otras optimizaciones que el compilador hace por nosotros, esta no se puede hacer ya que rompería la coherencia y predictibilidad del código.
\subsection{Uniones}
Las uniones proporcionan una manera de circunvalar el sistema de tipos de C, permitiendo a un solo objeto ser referenciado por múltiples tipos. La sintaxis de la unión es idéntica a la de un struct, pero 
sus semánticas son diferentes. En una unión, todos los campos comparten la misma dirección de memoria, lo que significa que si se modifica un campo, se modifica el resto, para hacer esto, el compilador reserva
tamaño de acuerdo al elemento más grande.
\begin{ejemplo}
    Veamos un ejemplo de una unión en comparación con un struct:
    \begin{center}
        \begin{minipage}{0.45\textwidth}
            \begin{minted}[fontsize=\large]{c}
    union U1 {
        char c;
        int i[2];
        double v;
    } *up;
    \end{minted}
    \end{minipage}
    \begin{minipage}{0.45\textwidth}
        \begin{minted}[fontsize=\large]{c}
    struct S1 {
        char c;
        int i[2];
        double v;
    } *sp;
        \end{minted}
    \end{minipage}
\end{center} 
Si compilamos ambos códigos en una máquina Linux en x86-64, el offset y tamaño total de las estructuras sería el siguiente:
\begin{center}
    \begin{tabular}{ccccc}
        Tipo & c & i & v & Total Size \\ 
        \hline 
        U1 & 0 & 0 & 0 & 8 \\ 
        S1 & 0 & 4 & 16 & 24 \\ 
    \end{tabular}
\end{center}
Este gran poder conlleva una gran responsabilidad, ya que el programador debe ser consciente de que 
si usa datos como \texttt{float} e \texttt{int} en la misma unión, 
e intenta interpretar los bits del entero como si fuese un float,
el resultado no será el esperado, ya que no tendremos el mismo efecto que el cast (float) sobre el entero.
\end{ejemplo}
Recordamos que las palabras se almacenan en memoria como un número N de bytes, en esta organización, el byte mas 
significativo puede ser el que está en la dirección más baja (Big Endian) o en la dirección más alta (Little Endian).
Esta diferencia puede causar problemas al intercambiar datos binarios entre máquinas, también mencionamos Middle endian, que son  
aquellas arquitecturas que se pueden configurar de ambas maneras, aparece en ARM, sin embargo, la mayoría de las máquinas modernas son Little Endian.
\\ \\ 
Para finalizar este capítulo, presentamos un resumen de los tipos compuestos en C:
\begin{itemize}
    \item \textbf{Arrays}
        \begin{itemize}
            \item Reservan memoria contigua para almacenar elementos
            \item Se usa aritmética de indexación para localizar elementos individuales
            \item Puntero al primer elemento
            \item No comprueban límites
        \end{itemize}
\item \textbf{Structs}
    \begin{itemize}
        \item Reserva de una sola región de memoria, los campos en el mismo orden declarado 
        \item Se accede usando desplazamientos determinados por el compilador
        \item Alineamiento interno y externo mediante relleno
    \end{itemize}
\item \textbf{Unions}
    \begin{itemize}
        \item Declaraciones superpuestas
        \item Proveen de una manera para saltarse el sistema de tipos de C
    \end{itemize}
\end{itemize}

    \chapter{Sincronización en memoria compartida. Monitores}
Suponiendo que existe una memoria común para los distintos procesos que ejecutan un programa concurrente, este Capítulo trata sobre la sincronización de los mismos usando para ello instrucciones que usan dicha memoria compartida.\\

Estudiaremos en profundidad el problema de la exclusión mutua, que ya obtuvo una solución en la asignatura de Arquitectura de Computadores, usando para ello cerrojos hardware o instrucciones máquina que nos proveían de funcionalidades deseadas a la hora de implementar una exclusión mutua.\\

Posteriormente, continuaremos con el problema de la sincronización de procesos en un programa concurrente, usando para ello conceptos con un mayor nivel de abstracción, los cuales nos permitirán resolver problemas más complejos de forma sencilla. Nos centraremos en el uso de los monitores, construcciones de alto nivel que nos ofrecen mayor libertad que los semáforos.\\

\section{Problema de la exclusión mutua}
En esta Sección tratamos de resolver el problema de la exclusión mutua mediante soluciones software, de forma que la solución no dependa del repertorio de instrucciones de una máquina, sino que sea una solución portable a cualquier dispositivo, de forma que podamos asegurar sobre los procesos de nuestros programas concurrentes todas las propiedades deseadas.\\

Consideraremos solo soluciones al problema en el que el acceso a la sección crítica se resuelva mediante instrucciones básicas de lectura y escritura sobre una o varias variables compartidas en memoria.\\

Como mecanismo para realizar la espera de los procesos en el acceso a la sección crítica usaremos la \textit{espera ocupada}, es decir, meteremos a los procesos que no deben entrar a la sección crítica todavía en un bucle que realice iteraciones ``vacías'' (sin ninguna utilidad) con la finalidad a que esperen a que el proceso que se encuentre en la sección crítica abandone la misma y deje pasar al siguiente.

Hemos de comentar que la espera ocupada no es la mejor solución de espera para los procesos, ya que introduce un uso innecesario de los procesadores con el fin de que ciertos procesos esperen. Puede considerarse una solución aceptable cuando el sistema no disponga de muchos procesos, pero en otro caso podríamos considerar otro tipo de esperas, como que el propio Sistema Operativo suspenda a los procesos.

\subsection{Condiciones de Dijkstra}
Dijkstra enunció que para obtener una solución parcialmente correcta al problema de la exclusión mutua, debían cumplirse 4 condiciones:
\begin{enumerate}
    \item \textit{No hacer ninguna suposición acerca de las instrucciones o número de procesos soportados por el multiprocesador.} Esto es, solo podremos hacer uso de operaciones que entendemos como básicas, tales como leer o escribir en una variable compartida para resolver el problema.

        Dichas instrucciones se ejecutarán de forma atómica, de forma que si dos procesos distintos intentan acceder a la vez a una misma posición de memoria, será el controlador de memoria quien determine de forma arbitraria qué proceso accederá antes y qué proceso después, de forma que el acceso a memoria se lleve a cabo secuencialmente, pero no de una forma predecible.
    \item \textit{No hacer ninguna suposición acerca de la velocidad de ejecución de los procesos}, salvo que esta no es cero, para que se cumpla la hipótesis de Progreso Finito.
    \item \textit{Cuando un proceso se encuentra ejecutando código fuera de la sección crítica, no puede impedir que otros entren a la misma.}
    \item \textit{La sección crítica será alcanzada finalmente por alguno de los procesos que quieran entrar.} Esta condición asegura la propiedad de \textit{alcanzabilidad}, que excluye la posibilidad de que los procesos lleguen a una situación de interbloqueo.

        Esta propiedad no asegura que todos los procesos entren alguna vez a la sección crítica, y mucho menos que lo hagan de forma equitativa.
\end{enumerate}

\subsection{Método de refinamiento sucesivo}
Dijkstra propuso a su vez una forma de obtener una solución al problema de la exclusión mutua para dos procesos, basada en 4 pasos, modificaciones o etapas a partir de un esquema inicial para obtener la solución de forma razonada, que terminará en una quinta etapa, denominada \textit{algoritmo de Dekker}.

\subsubsection{Primera etapa}
Inicialmente, se presupone que los procesos alternarán su entrada en la sección crítica según indique el valor de una variable compartida llamada \verb|turno|. Dicha variable contendrá el identificador del proceso que en cada momento puede entrar a la sección crítica.\\

En un escenario con dos procesos que se disponen a ejecutar una sección crítica, la primera etapa consta del código de la Figura~\ref{fig:cod_primera_etapa}.

\begin{figure}
    \centering
    \setlength{\columnsep}{1cm}
    \begin{multicols}{2}
\begin{minted}{pascal}
var turno : integer;

Process P1();
begin
   while true do
   begin
      { Acceso a la SC }
      while turno <> 1 do
      begin
         null;
      end do
      { Sección crítica }
      turno := 2;
   end do
end
\end{minted}
\begin{minted}{pascal}


Process P2();
begin
   while true do
   begin
      { Acceso a la SC }
      while turno <> 2 do
      begin
         null;
      end do
      { Sección crítica }
      turno := 1;
   end do
end
\end{minted}
\end{multicols}
\caption{Algoritmo para la primera etapa de refinamiento sucesivo.}
\label{fig:cod_primera_etapa}
\end{figure}


Esta solución garantiza el acceso en exclusión mutua de los procesos a la sección crítica de forma independiente a la velocidad de ejecución de los mismos, por lo que la solución es segura. Sin embargo, no cumple la tercera condición de Dijkstra, ya que la solución obliga a la alternancia entre los procesos en la entrada a la sección crítica. 

De esta forma, si el proceso $P1$ entra a la sección crítica mientras que $P2$ se encuentra haciendo otras operaciones independientes, el proceso $P1$ no podrá volver a entrar a la sección crítica hasta que no lo haga $P2$.

\subsubsection{Segunda etapa}
La alternancia que obtuvimos en la etapa anterior y que nos impedía cumplir con todas las condiciones de Dijkstra se debía a que para decidir qué proceso entraba en la sección crítica era necesario almacenar información global del estado del programa.

Para evitar esto, la idea ahora es asociar a cada proceso una variable que contenga su información de estado, variable que llamaremos \verb|clave|, la cual indicará de forma binaria si el proceso se encuentra o no en la sección crítica en dicho instante de ejecución del algoritmo, mediante dos estados:
\begin{itemize}
    \item Estado pasivo, el proceso no intenta acceder a la sección crítica, representado con un 1.
    \item El proceso intenta acceder a la sección crítica, representado con un 0.
\end{itemize}
El código quedaría como el de la Figura~\ref{fig:cod_segunda_etapa}.\\
\begin{figure}
    \centering
\setlength{\columnsep}{1cm}
\begin{multicols}{2}
    \begin{minted}{pascal}
        var c1, c2 : integer;
    
        Process P1();
        begin
           c1 := 1;
    
           while true do
           begin
              { Acceso a la SC }
              { Si P2 entró }
              while c2 = 0 do   
              begin
                 null;
              end do
    
              { Entra a la sección crítica }
              c1 := 0;
              { Sección crítica }
              c1 := 1;
           end do
        end
    \end{minted}
    \begin{minted}{pascal}

    
        Process P2();
        begin
           c2 := 1;
    
           while true do
           begin
              { Acceso a la SC }
              { Si P1 entró }
              while c1 = 0 do
              begin
                 null;
              end do
    
              { Entra a la sección crítica }
              c2 := 0;
              { Sección crítica }
              c2 := 1;
           end do
        end
    \end{minted}
\end{multicols}
\caption{Algoritmo para la segunda etapa de refinamiento sucesivo.}
\label{fig:cod_segunda_etapa}
\end{figure}
Como podemos ver, el protocolo de entrada consiste en leer el valor de la clave del otro proceso con la finalidad de consultar si dicho proceso ha entrado o no en la sección crítica, y esperar mientras este se encuentre dentro de la sección crítica.\\

Sin embargo, en este caso la solución no es segura, ya que si $P1$ y $P2$ se ejecutan a la misma velocidad, entonces ambos entrarían a la vez a la sección crítica, debido a que los dos verían que el estado del otro es pasivo, con lo que ninguno entraría en el bucle de espera ocupada. Notemos que esto sucede porque cambiamos el estado de un proceso a 0 justo antes de entrar a la sección crítica, por lo que es ya tarde para impedir la entrada a otro proceso.

Como la bondad de la solución depende de la velocidad de ejecución relativa entre los procesos, se dice que es inaceptable por incumplir la segunda condición de Dijkstra.

\subsubsection{Tercera etapa}
Para esta etapa planteamos una sencilla modificación sobre la etapa anterior, que consiste en cambiar el valor de la variable clave a 0 antes de consultar el valor de la variable clave del otro proceso. De esta forma, para que un proceso pueda entrar a la sección crítica, debe primero cambiar su estado a 0, con el fin de recuperar la condición de seguridad de que solo un proceso pueda entrar a la vez a la sección crítica. El código resultante es el de la Figura~\ref{fig:cod_tercera_etapa}.\\
\begin{figure}
\setlength{\columnsep}{1cm}
\begin{multicols}{2}
    \begin{minted}{pascal}
        var c1, c2 : integer;
    
        Process P1();
        begin
           c1 := 1;
    
           while true do
           begin
              { Acceso a la SC }
              c1 := 0;
              { Si P2 entró }
              while c2 = 0 do
              begin
                 null;
              end do
              { Sección crítica }
              c1 := 1;
           end do
        end
    \end{minted}
    \begin{minted}{pascal}

    
        Process P2();
        begin
           c2 := 1;
    
           while true do
           begin
              { Acceso a la SC }
              c2 := 0;
              { Si P1 entró }
              while c1 = 0 do
              begin
                 null;
              end do
              { Sección crítica }
              c2 := 1;
           end do
        end
    \end{minted}
\end{multicols}
\caption{Algoritmo para la tercera etapa de refinamiento sucesivo.}
\label{fig:cod_tercera_etapa}
\end{figure}
Sin embargo, si ambos procesos tienen la misma velocidad, puede suceder que ambos cambien el valor de su clave a 0 al mismo tiempo, con lo que se de una situación de interbloqueo, que incumpliría la cuarta condición de Dijkstra, al no poder alcanzar nunca la sección crítica (además de la segunda, ya que la bondad de la solución depende de la velocidad de ejecución de los procesos).

\subsubsection{Cuarta etapa}
Lo que causó el problema en la tercera etapa fue que puede suceder que un proceso cambie el valor de su clave a la vez que el otro de forma concurrente, sin que este se de cuenta de que el otro lo ha hecho a la vez. La solución que se propone en esta etapa es permitir a un proceso volver a cambiar el valor de su clave a 1 si después de asignar su clave a 0, comprueba que el otro proceso también cambió su clave al mismo valor. De esta forma, planteamos el código de la Figura~\ref{fig:cod_cuarta_etapa}.\\
\begin{figure}
\setlength{\columnsep}{1cm}
\begin{multicols}{2}
    \begin{minted}{pascal}
    var c1, c2 : integer;

    Process P1();
    begin
       c1 := 1;

       while true do
       begin
          { Acceso a la SC }
          c1 := 0;
          { Si P2 entró }
          while c2 = 0 do
          begin
             c1 := 1;
             while c2 = 0 do
             begin
                null;
             end do
             c1 := 0;
          end do
          { Sección crítica }
          c1 := 1;
       end do
    end
\end{minted}
\begin{minted}{pascal}


    Process P2();
    begin
       c2 := 1;

       while true do
       begin
          { Acceso a la SC }
          c2 := 0;
          { Si P1 entró }
          while c1 = 0 do
          begin
             c2 := 1;
             while c1 = 0 do
             begin
                null;
             end do
             c2 := 0;
          end do
          { Sección crítica }
          c2 := 1;
       end do
    end
\end{minted}
\end{multicols}
\caption{Algoritmo para la cuarta etapa de refinamiento sucesivo.}
\label{fig:cod_cuarta_etapa}
\end{figure}
Sin embargo, si ambos procesos se ejecutasen a la misma velocidad, se podría seguir produciendo un interbloqueo entre ambos procesos, aunque esta situación ahora sea más improbable. La solución no sería válida por incumplir tanto la segunda como la cuarta condición de Dijkstra.\\

La conclusión a la que llegamos tras todas estas etapas es que las variables \verb|c1| y \verb|c2| nos son útiles para coordinar la entrada a la sección crítica, pero no son suficientes para dar una solución correcta al problema que tratamos de resolver.\\

\subsection{Algoritmo de Dekker}
Se podría considerar como una quinta etapa del método de refinamiento sucesivo, pero esta vez obteniendo una solución válida del problema. El algoritmo de Dekker junta las ideas presentes en la primera y cuarta etapa de refinamiento de Dijkstra:
\begin{itemize}
    \item La primera etapa producía una solución segura, pero obligaba a la alternancia en el acceso de los procesos a la sección crítica.
    \item Por otra parte, la cuarta etapa no cuenta con dicha alternancia en el acceso, pero puede llevar a un interbloqueo de los procesos del programa concurrente.
\end{itemize}
Para resolver el problema, se considera el código de la cuarta etapa de Dijkstra y se le añade un orden establecido en la entrada mediante una variable \verb|turno|, para desempatar la situación en la que los dos procesos quieran entrar exactamente al mismo tiempo en la sección crítica.\\

De esta forma, un proceso que quiera entrar en la sección crítica asignará primero su clave a 0, y si el otro proceso también tiene su clave a 0, lo primero que hará es comprobar de quién es el turno y si no dispone del mismo, cambiará su clave a 1, pasando a esperar y dejando al otro proceso continuar con la ejecución de la sección crítica. Vemos el código en la Figura~\ref{fig:cod_dekker}.

\begin{figure}[H]
\setlength{\columnsep}{1cm}
\begin{multicols}{2}
    \begin{minted}{pascal}
    var c1, c2, turno : integer;

    Process P1();
    begin
       while true do
       begin
          { Acceso a la SC }
          c1 := 0;

          while c2 = 0 do
          begin
             if turno = 2 then
             begin

                c1 := 1;
                while turno = 2 do
                begin
                   null;
                end do
                c1 := 0;

             end
          end do

          { Sección crítica }
          turno := 2;
          c1 := 1;
       end do
    end
\end{minted}
\begin{minted}{pascal}


    Process P2();
    begin
       while true do
       begin
          { Acceso a la SC }
          c2 := 0;

          while c1 = 0 do
          begin
             if turno = 1 then
             begin

                c2 := 1;
                while turno = 1 do
                begin
                   null;
                end do
                c2 := 0;

             end
          end do

          { Sección crítica }
          turno := 1;
          c2 := 1;
       end do
    end
\end{minted}
\end{multicols}
\caption{Algorimo de Dekker.}
\label{fig:cod_dekker}
\end{figure}

\subsubsection{Propiedades de corrección}
En esta sección, demostraremos que se cumple siempre el acceso en exclusión mutua a la sección crítica por parte de los procesos que intervienen en los programas concurrentes, así como la propiedad de alcanzabilidad de la sección crítica y vivacidad de los procesos que tratan de hacer uso de la misma:

\begin{description}
    \item [Exclusión mutua.]~\\
    El proceso $P_i$ (con $i = 1$ ó $i=2$) entrará en la sección crítica solo si el otro proceso, $P_j$ mantiene su clave $cj$ a 1. Dado que la clave de un proceso solo la puede modificar el propio proceso y que el proceso $P_i$ comprueba la clave $cj$ solo después de asignar su propia clave $ci$ a 0, si el proceso $P_i$ entra en sección crítica, se ha de cumplir la condición $ci = 0 \land cj = 1$. Notemos que esta situación es incompatible con la condición de que el proceso $P_j$ entre en la sección crítica: $cj = 0 \land ci = 1$.
    \item [Alcanzabilidad de la sección crítica.]~\\
    Para demostrar la alcanzabilidad de la sección crítica, distinguimos casos:
    \begin{itemize}
        \item Si suponemos que el proceso $P_i$ intenta entrar solo en la sección crítica, entonces el otro proceso $P_j$ se mantendrá en estado pasivo, con lo que el valor de su clave $cj$ será 1. De esta forma, el proceso $P_i$ puede entrar a la sección crítica.
        \item Sin embargo, si tanto $P_i$ como $P_j$ intentan entrar a la vez a la sección crítica y suponemos que $turno = i$, entonces:
            \begin{itemize}
                \item Si $P_j$ encuentra la clave $ci$ a 1, entonces $P_j$ entrará en la sección crítica.
                \item Si $P_j$ encuentra la clave $ci$ a 0, como $turno = i$, entonces $P_j$ entrará en el segundo bucle interno para realizar la espera ocupada, poniendo antes su clave $cj$ a 1, que permitirá pasar al proceso $P_i$.
                \item Si $P_i$ encuentra la clave $cj$ a 1, entonces $P_i$ entrará en la sección crítica.
                \item Si $P_i$ encuetran la clave $cj$ a 0, se mantendrá realizando iteraciones en el bucle de espera ocupada más externo con $ci$ a 0, hasta que lea el valor de $cj$ a 1, que sucederá por el punto superior, con lo que $P_i$ entrará en la sección crítica.
            \end{itemize}
    \end{itemize}

    
    \item [Vivacidad.]~\\
    Dependiendo del hardware de control de acceso a memoria, el algoritmo de Dekker puede llegar a provocar la inanición de uno de los dos procesos:

    Supongamos que tenemos a los procesos $P1$ y $P2$ ejecutando su código, queriendo acceder continuamente a la sección crítica. Supongamos además que el proceso $P2$ se ejecuta a una velocidad bastante lenta en comparación al proceso $P1$. Nos encontramos en el caso en el que ambos procesos cambiaron sus claves al mismo tiempo y el turno inicial era 1, con lo que $P1$ pasó a ejecutar el código de la sección crítica y $P2$ se quedó esperando en el bucle más interno, con el valor de su clave \verb|c2| a 1.\\

    Debemos recordar que anteriormente mencionamos que el acceso al módulo de memoria no se hace de forma paralela, sino que se hace de forma secuencial, de forma que si dos procesos intentan acceder a la vez a una misma posición de memoria es el controlador de memoria quien determina el acceso a un proceso de forma arbitraria.\\

    Supongamos pues que $P1$ termina de ejecutar el código de la sección crítica, con lo que cambia el valor de la variable \verb|turno| a 2, \verb|c1| a 1 y cambia también \verb|c1| a 0. Posteriormente, como $P1$ cambió \verb|turno| a 2, el proceso $P2$ sale del bucle más interno, con lo que se dispone a cambiar el valor de su clave a 0.

    Sin embargo, en este momento sucede que tanto $P1$ como $P2$ intentan acceder a la vez al valor de \verb|c2|, $P1$ para leer (en la condición del \verb|while| exterior) y $P2$ para escribir. Si en dicho momento el controlador de memoria da prioridad a las lecturas, $P1$ volvería a introducirse en la sección crítica.

    Inmediatamente, $P2$ se dispondría a cambiar el valor de su clave a 0, pero como mencionamos anteriormente, $P2$ es muy lento, con lo que resulta que le da tiempo a $P1$ a ejecutar la sección crítica y volver a la lectura de \verb|c2| en el bucle más externo a la vez que $P2$, con lo que el controlador de memoria puede volver a darle prioridad.\\

    Si este escenario sucede de forma indefinida, tenemos una falta de vivacidad en el proceso $P2$, ya que mientras $P1$ esté en funcionamiento, no podrá avanzar en su ejecución.

    \item [Equidad del protocolo.]~\\
    Como hemos comentado en el apartado de vivacidad, la equidad del algoritmo de Dekker dependerá de la equidad del hardware de la máquina en el que ejecutemos el programa concurrente. Si existen peticiones de acceso simultáneo a una misma dirección de memoria compartida por dos procesos, uno para lectura y otro para escritura de forma que el hardware da prioridad a las lecturas, no se puede demostrar que el algoritmo de Dekker sea equitativo, pudiendo llegar al escenario de inanición de uno de los procesos como ya se ha descrito anteriormente.
\end{description}

\subsection{Algoritmo de Dijkstra}
Una vez visto el algoritmo de Dekker, primera solución aceptable al problema de la exclusión mutua (aunque no cumpla con las propiedades deseables de vivacidad y equidad), nos encontramos con que no es generalizable para $n$ procesos, con lo que mostramos a continuación el algoritmo de Dijkstra, que resuelve el problema de la exclusión mutua para $n$ procesos utilizando un array de $n$ posiciones en las que cada proceso almacena su estado, destacando tres posibles estados:
\begin{itemize}
    \item Proceso pasivo, el proceso no intenta acceder al protocolo de entrada.
    \item Solicitando, el proceso intenta acceder al protocolo de entrada.
    \item En SC, el proceso está dentro de la sección crítica.
\end{itemize}
El algoritmo de Dijkstra podemos verlo en la Figura~\ref{fig:cod_dijkstra}.
\begin{figure}[H]
\begin{minted}{pascal}
    var c : array[0..n-1] of (pasivo, solicitando, en_SC);
        turno : 0..n-1;

    Process Pi();
    begin
       while true do
       begin
          { Acceso a la sección crítica }
          repeat
             c[i] := solicitando;

             { A }
             while turno <> i do
             begin
                if c[turno] = pasivo then turno := i;
             end do

             c[i] := en_SC;

             { B }
             j := 0;
             while (j<n) and (j=i or c[j] <> en_SC) do
             begin
                j := j+1;
             end do
          until j >= n

          { Sección crítica }
          c[i] := pasivo;
       end do
    end
\end{minted}
\caption{Algoritmo de Dijkstra, código para el $i$-ésimo proceso.}
\label{fig:cod_dijkstra}
\end{figure}

De esta forma, lo primero que hace un proceso $P_i$ al llegar al protocolo de entrada a la sección crítica es cambiar su estado a \verb|solicitando|. Posteriormente:
\begin{itemize}
    \item Si es su turno, no realizará ninguna iteración del bucle \verb|A|, con lo que pasará inmediatamente a la ejecución del bucle \verb|B|, tras cambiar su clave a \verb|en_SC|.
    \item Si no es su turno:
        \begin{itemize}
            \item Si el proceso que tiene el turno está en estado pasivo, entonces el proceso $P_i$ pone el turno a $i$, siendo ahora su turno, con lo que pasa del bucle \verb|A|.
            \item Si el proceso que tiene el turno no está en estado pasivo, entonces el proceso $P_i$ esperará hasta que este lo esté, es decir, hasta que el proceso $P_{\text{turno}}$ abandone la sección crítica.
        \end{itemize}
\end{itemize}
Una vez superado el bucle \verb|A|, entonces el estado del proceso $P_i$ pasará a \verb|en_SC|, y solo tendrá que superar el bucle \verb|B| para poder entrar a la sección crítica.\\

Puede sueceder que dos (o más, en cuyo caso es una explicación análoga) procesos, $P_i$ y $P_j$ (con $i,j\notin \{turno\}$) lleguen al bucle \verb|A| de forma que (suponiendo que \verb|c[turno] = pasivo|) ambos lo ejecuten a la misma velocidad, con lo que ambos pasen dicho bucle quedando la variable \verb|turno| asignada al identificador de cualquiera de los dos procesos (a $i$ o a $j$). Posteriormente, ambos cambiarán su estado a \verb|en_SC|. En dicho caso, contamos con el bucle \verb|B| para cumplir con la condición de seguridad de tener un único proceso ejecutando la sección crítica a la vez.\\

Lo que hace el bucle \verb|B| es comprobar todos los estados de los procesos antes de dejar pasar al proceso $P_i$ entrar a la sección crítica, de forma que si todos los demás procesos tienen un estado distinto de \verb|en_SC|, entonces el contador \verb|j| llegará hasta \verb|n|, con lo que $P_i$ saldrá del bucle \verb|B| así como del \verb|repeat| de entrada a la sección crítica.\\

En el caso en el que dos (o más) procesos estén con estado \verb|en_SC|, entonces la variable \verb|j| de cada proceso no llegará a aumentar hasta \verb|n| (lo que le permitiría salir del bucle \verb|repeat| y entrar en sección crítica), con lo que los procesos no podrán salir del \verb|repeat|, teniendo que volver a pasar por el bucle \verb|A|, donde el turno estará ya fijado en un proceso con un estado distinto de pasivo (el último que actualizó el valor de dicha variable).

\subsubsection{Propiedades de corrección}
\noindent
Pasamos ahora a demostrar las propiedades de corrección del algoritmo de Dijsktra:

\begin{description}
    \item [Exclusión mutua.]~\\
        La propiedad de exclusión mutua está garantizada gracias al bucle \verb|B|:\\
        El proceso $P_i$ (con $i \in \{0,\ldots,n-1\}$) entrará en la sección crítica solo si el resto de procesos $P_j$ con $j\neq i$ tienen su estado distinto de \verb|en_SC|. En dicho caso, $P_i$ (y solo él) pasará a ejecutar la sección crítica, poniendo su estado a \verb|en_SC|. 

        Supuesto que ahora otro proceso $P_j$ intenta acceder a la sección crítica, solo podrá hacerlo si el resto de procesos tienen su estado distinto de \verb|en_SC|, algo que no puede suceder hasta que el proceso $P_i$ salga de la sección crítica.
    \item [Alcanzabilidad de la sección crítica.]~\\
        Supuesto que disponemos de $m$ procesos $P_1, P_2, \ldots, P_m$ de forma que tratan de acceder a la sección crítica al mismo tiempo, entonces el valor de la variable compartida \verb|turno| será un cierto $k \in \{1,\ldots,m\}$, correspondiente al identificador del último proceso que cambió su valor. 

        Los procesos que superen el bucle \verb|A| junto con $P_k$ no podrán superar el bucle \verb|B| con un valor de \verb|j| igual o superior a \verb|n|, por lo que todos los procesos (junto con $P_k$) volverán al bucle \verb|A|.

        En dicho instante, el turno estará fijado a $k$ y \verb|c[k]| tendrá un valor distinto de pasivo, con lo que el proceso $P_k$ será ahora el único que consiga pasar el bucle \verb|A|, con lo que completará el bucle \verb|B| con un valor de \verb|j| igual a \verb|n|, lo que le permitirá acceder a la sección crítica.

    \item [Vivacidad.]~\\
        El algoritmo de Dijkstra también puede llevar a la inanición de uno de los procesos del programa concurrente, al igual que el algoritmo de Dekker.

        Para ver dicha situación, supongamos que tenemos tres procesos, $P_0$, $P_1$ y $P_2$, los tres intentando acceder a la sección crítica. Supongamos que inicialmente tenemos $turno = 0$ y que la traza de ejecución obtenida es aquella que hace que $P_0$ salte el bucle \verb|A|, antes de que $P_1$ y $P_2$ cambien su estado a \verb|solicitando|, con lo que ambos quedarán bloqueados en dicho bucle.

        Esto permite a $P_0$ entrar en la sección crítica. Cuando este proceso salga y cambie su estado a \verb|pasivo|, provocará que los dos procesos anteriores pasen el bucle \verb|A|, en un orden que puede ser primero $P_2$ (cambia $turno = 2$) y luego $P_1$ (cambia $turno = 1$), de forma que ambos lleguen al bucle \verb|B|, que obligue a ambos a pasar nuevamente por el bucle \verb|A|, donde el proceso $P_2$ quedará bloqueado y será $P_1$ quien consiga llegar a la sección crítica.

        Suponiendo que ahora $P_0$ vuelve a intentar acceder a la sección crítica y que después $P_1$ cambia su estado a \verb|pasivo| (sale de la sección crítica), puede suceder la situación anterior pero en este caso con $P_0$, es decir, que sea $P_2$ quien salga antes del bucle \verb|A| y $P_0$ después, con lo que el turno queda asignado finalmente a 0, lo que permitirá a $P_0$ volver a entrar en la sección crítica.

        Esta situación se puede repetir tantas veces como queramos, de forma que $P_2$ nunca llegue a ejecutar la sección crítica, por lo que no es posible demostrar la vivacidad del algoritmo de Dijkstra, ya que este puede llevar a la inanición de uno de los procesos que ejecutan el protocolo de entrada a la sección crítica.
    \item [Equidad.]~\\
        Como no podemos garantizar la vivacidad de la solución, mucho menos podremos garantizar la equidad de los procesos en el acceso a la sección crítica.
\end{description}

\subsection{Algoritmo de Knuth}
Con la finalidad de no caer en los errores de los algoritmos de Dekker y Dijkstra (los cuales pueden llevar a un estado en el que un proceso sufra inanición), Donald Knuth añadió una quinta condición a las condiciones de Dijkstra para la corrección parcial de una solución al problema de la exclusión mutua:
\begin{enumerate}[label=5.]
    \item \textit{Se tiene que poder demostrar que todos los procesos pueden sufrir un retraso máximo en el acceso a la sección crítica que sea calculable.}

        Con el fin de garantizar que ningún proceso sufrirá de inanición, ya que en algún momento conseguirá entrar a la sección crítica, lo que nos da la propiedad de equidad de la solución.
\end{enumerate}

Donald Knuth se dio cuenta de que la propiedad de vivacidad no podía ser demostrada debido a que la variable compartida \verb|turno| utilizada en los algoritmos de Dekker y Dijkstra sufría condiciones de carrera. El algoritmo propuesto por Knuth hace uso de una variable local \verb|j| que copia el valor de la variable compartida \verb|turno|, con el fin de evitar dichas condiciones de carrera. Por lo demás, el algoritmo es similar al de Dijkstra, tal y como podemos ver en la Figura~\ref{fig:cod_knuth}.

\begin{figure}[H]
\begin{minted}{pascal}
    var c : array[0..n-1] of (pasivo, solicitando, en_SC);
        turno : 0..n-1;

    Process Pi();
    var j : 0..n-1;
    begin
       while true do
       begin
          repeat
             c[i] := solicitando;
             j := turno;

             { A }
             while j <> i do
             begin
                if c[j] <> pasivo then j := turno;
                else j := (j-1) mod n;
             end do

             c[i] := en_SC;

             { B }
             k := 0;
             while (k<n) and (k=i or c[k] <> en_SC) do
             begin
                k := k+1;
             end do
          until k >= n;

          turno := i;
          { Sección crítica }
          turno := (i-1) mod n;
          c[i] := pasivo;
       end do
    end
\end{minted}
\caption{Algoritmo de Knuth.}
\label{fig:cod_knuth}
\end{figure}
De esta forma, suponiendo que somos el proceso $P_i$ y que nos disponemos a ejecutar la sección crítica:\\
Si tenemos que \verb|turno = i|, entonces $P_i$ se salta el bucle \verb|A|. Si no (\verb|turno != i|), entonces:
\begin{itemize}
    \item Si el proceso que tiene el turno no está en estado pasivo, entonces el proceso se bloqueará realizando una espera activa, copiando el valor de \verb|turno| en \verb|j|. De esta forma, cuando el proceso que tenía el turno salga de la sección crítica, cambiará el valor de \verb|turno| y el proceso cambiará su copia en \verb|j|, con lo que saldrá de la espera ocupada e intentará volver a realizar el protocolo de entrada para la sección crítica.
    \item Si el proceso que tiene el turno está en estado pasivo, entonces la copia local del valor de \verb|turno|, es decir, la variable \verb|j| ciclará su valor, de forma que volvamos a ejecutar el protocolo de entrada, destacando tres posibilidades:
        \begin{itemize}
            \item Si al ciclar el valor de \verb|j| ahora tenemos que \verb|j = i|, entonces $P_i$ pasa el bucle \verb|A|.
            \item Si al ciclar el valor de \verb|j| se detecta que $P_j$ no está en estado pasivo, entonces $P_i$ realizará una espera activa, de forma que actualice \verb|j| a \verb|turno|, luego ciclará el valor de \verb|j| y detectará que $P_j$ no estará en estado pasivo.
            \item Si al ciclar el valor de \verb|j| se detecta que $P_j$ está en estado pasivo, se cicla el valor de \verb|j| otra vez, volviendo a destacar tres posibilidades.
        \end{itemize}
\end{itemize}
De esta forma, Knuth controla el número de veces que los procesos pueden adelantarse los unos a los otros dentro del bucle \verb|A|.\\

Por el resto, es idéntico al funcionamiento del algoritmo de Dijkstra, salvo que solo modifica el valor de la variable \verb|turno| una vez el proceso se encuentra dentro de la sección crítica, evitando condiciones de carrera.

\subsubsection{Propiedades de corrección}
\noindent
Veamos las propiedades de corrección del algoritmo de Knuth:
\begin{description}
    \item [Exclusión mutua.]~\\
        Está garantizada gracias al bucle \verb|B|, tal y como razonamos en el algoritmo de Dijkstra.
    \item [Alcanzabilidad de la sección crítica.]~\\
        Supuesto que tenemos $m$ procesos, $P_1$, $P_2$, \ldots, $P_m$ que tratan de acceder a la sección crítica a la vez, y suponiendo que \verb|turno = k|:
        \begin{itemize}
            \item Si $k \in \{1,\ldots,m\}$, entonces, $P_k$ superará el bucle \verb|A| y tratará de entrar a la sección crítica, distinguiendo dos posibilidades:
                \begin{itemize}
                    \item Si el resto de procesos vieron que \verb|c[k]| era distinto de pasivo, los $m-1$ procesos restantes quedaron bloqueados en el bucle \verb|A|, con lo que solo el proceso $P_k$ tendrá su clave a \verb|en_SC|, con lo que conseguirá superar el bucle \verb|B|, entrando a la sección crítica.
                    \item Si hay algún proceso (o varios) que vieron \verb|c[k]| a pasivo (esto es, antes de que el proceso $k$ cambiase el valor de su clave), estos irán ciclando el valor de \verb|j|, de forma que puede haber algún conjunto de procesos que supere el bucle \verb|A| junto con $P_k$. En dicho caso, tras superar el bucle \verb|B|, ninguno de esos obtendrá \verb|k = n|, con lo que han de volver a iterar en el bucle de \verb|repeat|, realizando otra vez el bucle \verb|A|. En este caso, estamos en las condiciones del punto superior, con lo que $P_k$ accede a la sección crítica.
                \end{itemize}
            \item Si $k\notin \{1,\ldots,m\}$, entonces tenemos que \verb|c[k] = pasivo|. En dicho caso, todos los procesos irán ciclando su valor de \verb|j| decrementándolo, de forma que si algún proceso $P_i$ consigue tener $i = j$, entonces superará el bucle \verb|A| y destacamos dos posibilidades:
                \begin{itemize}
                    \item Si solo un proceso $P_i$ consigue tener $i = j$, entonces conseguirá superar el bucle \verb|B| con $k = n$, con lo que entrará a ejecutar la sección crítica.
                    \item Si hay dos o más procesos que consiguieron tener $i = j$ (puede suceder, si los procesos con $i$ más lejana a \verb|turno| ejecutaron el bucle \verb|A| de forma más rápida), al superar el bucle \verb|B|, ninguno conseguirá tener $k = n$, con lo que todos volverán a ejecutar el bucle \verb|A|. En este momento, todos estos procesos que tenían otro proceso más cercano a \verb|turno| quedarán bloqueados en el bucle \verb|A|, ya que observarán que \verb|c[j] != pasivo| en el proceso con el índice $j$ más cercano a \verb|turno|, por lo que solo el proceso $P_i$ con índice más cercano a \verb|turno| conseguirá superar el bucle \verb|A|, quien superará el bucle \verb|B| con un valor $k = n$, lo que le permitirá ejecutar la sección crítica.
                \end{itemize}
        \end{itemize}
    \item [Vivacidad.]~\\
        Para demostrar la vivacidad del algoritmo de Knuth lo que tenemos que hacer es demostra la ausencia de inanición, es decir, que no podemos tener un proceso que nunca alcance la sección crítica.

        Por reducción al absurdo, supongamos pues, que en un escenario con $n$ procesos tenemos un proceso $P_{victima}$ (con $victima \in \{1,\ldots,n\}$) que nunca puede alcanzar la sección crítica, debido a que existe un proceso $P_i$ (con $i\neq victima$) que adelanta a $P_{victima}$ de forma indefinida en el acceso a la sección crítica.

        Entonces, ha de existir un proceso $P_k$ con el índice $k$ posterior a $i$ pero anterior a $victima$ en el protocolo de entrada a la sección crítica\footnote{Notemos que en este punto no estamos suponiendo de forma implícita que $i>k>victima$, ya que hablamos en sentido rotatorio, con lo que podemos tener $victima > i$ pero en dicho caso tendremos $k>victima$. Simplemente indicamos que el turno pasa antes a $k$ por $i$ que a $victima$.}, ya que si no (esto es, $P_{victima}$ es el siguiente en turno rotatorio tras $P_i$) tras salir $P_i$ de la seccion crítica ejecutaría la sentencia \verb|turno := (i-1) mod n|, lo que provocaría que el siguiente proceso a ejecutar la sección crítica fuera $P_{victima}$, pero habíamos supuesto que $P_{victima}$ era adelantado de forma indefinida por el proceso $P_i$.

        De esta forma, tenemos que también $P_k$ adelanta de forma indefinida al proceso $P_{victima}$ (porque $P_i$ lo adelanta de forma indefinida y el turno tras $P_i$ pasa a un proceso posterior, de forma que $P_k$ se posiciona antes de que el turno llegue a $P_{victima}$), con lo que se encuentra en la misma situación del proceso $P_i$ anterior, pudiendo encontrar un proceso $P_h$ con índice posterior a $k$ pero anterior a $victima$ que también adelante a $P_{victima}$ de forma indefinida.

        Como podemos observar, podemos repetir este razonamiento tantas veces como queramos, pudiendo obtener una cantidad de procesos superior a $n$ que ejecuta el algoritmo de Knuth. Sin embargo, como la cantidad de procesos que ejecuta el algoritmo no puede ser superior a $n$, llegamos a contradicción con que $P_{victima}$ era un proceso que nunca puede entrar a la sección crítica, lo que demuestra la propieda de vivacidad del algoritmo de Knuth.
    \item [Equidad.]~\\
        Para demostrar la propiedad de equidad del algoritmo de Knuth lo que haremos será acotar el número máximo de adelantamientos por otros procesos que puede sufrir un proceso que quiera entrar a la sección crítica. Para ello, sea $f:\mathbb{N}\rightarrow\mathbb{N}$ una función que contabiliza el número máximo de adelantamientos, $f(n)$ para un escenario con $n$ procesos, entonces:
        \begin{itemize}
            \item $f(0) = 0$, ya que en un escenario con ningún proceso, ningún proceso puede ser adelantado.
            \item $f(1) = 0$, ya que en un escenario con un proceso, este no puede ser adelantado por ningún otro.
            \item $f(2) = 1$, ya que si tenemos dos procesos, puede suceder que nuestro proceso sea adelantado por el otro proceso.
        \end{itemize}
        Para calcular $f(3)$ mostraremos el siguiente razonamiento, que nos permitirá generalizarlo, para conseguir obtener $f(n)$ a partir de $f(n-1)$, obteniendo una recurrencia y pudiendo calcular el valor de $f(n)$ para cualquier $n\in \mathbb{N}$.\\

        Para ello, hemos de buscar cuál es la peor situación que se pueda dar en un escenario con 3 procesos de forma que un proceso obtenga el mayor número de adelantamientos en su entrada a la sección crítica. No demostraremos que se trata de la peor situación, pero puede intuirse que no hay peor situación que la que estamos a punto de describir:

        Supongamos que tenemos 3 procesos: $P_0$, $P_1$, $P_2$ de forma que inicialmente $turno = 1$ y $P_2$ será nuestro proceso víctima (el que vamos a intentar que sea adelantado el máximo número de veces por $P_0$ y por $P_1$). Para ello, supongamos que los tres intentan acceder de forma simultánea al protocolo de entrada de la sección crítica, obteniendo la siguiente traza de ejecución:
        \begin{itemize}
            \item El proceso $P_0$ es el primero que cambia su clave \verb|c[0]| a \verb|solicitando|, de forma que cuando llega al bucle \verb|A|, observa que \verb|c[1]| está pasivo, decrementa el valor de \verb|j| (ahora, $j=0$), con lo que el proceso $P_0$ pasa del bucle \verb|A|.
            \item Ahora, el proceso $P_1$ cambia su clave a \verb|solicitando| y al ser $turno = 1$, pasa del bucle \verb|A|.
            \item En este momento, $P_2$ cambia su clave a \verb|soclitando| y como \newline \verb|c[1] = solicitando|, el proceso se queda bloqueado en el bucle \verb|A|.
            \item Supongamos que $P_0$ es muy rápido y le da tiempo a asignarse el estado \verb|en_SC|, de forma que antes de que $P_1$ cambie su clave a \verb|en_SC|, le da tiempo a comprobar que no hay ningún otro proceso con clave \verb|en_SC|, con lo que entra en la sección crítica. Ejecuta dicho código muy rápido, cambiando $turno$ a 2 y regresando al bucle \verb|A| para volver a entrar en la sección crítica.
            \item Si en este momento el proceso $P_1$ cambia su clave a \verb|en_SC| y le da tiempo a completar el bucle \verb|B| antes de que el proceso $P_2$ llegue a salir del bucle \verb|A|, el proceso $P_1$ conseguirá salir del bucle \verb|B| con $k = n$, lo que le permita ejecutar la sección crítica.
            \item Cuando el proceso $P_1$ abandone la seccion crítica, dejará $turno$ a 0, con lo que $P_2$ seguirá bloqueado en el bucle \verb|A| y será ahora $P_0$ quien vuelva a ejecutar la sección crítica.
            \item Finalmente $P_0$ sale de la sección crítica y cambia $turno$ a 2, lo que ahora sí permitirá a $P_2$ entrar a la sección crítica.
        \end{itemize}
        De esta forma, el proceso $P_2$ ha sido adelantado 3 veces, obteniendo finalmente que los identificadores de los procesos que han ejecutado la sección crítica han sido:
        \begin{equation*}
            0\ 1\ 0\ 2
        \end{equation*}
        Concluimos que $f(3) = 3$.\\

        Si ahora disponemos que $n$ procesos: $P_0$, $P_1$, \ldots, $P_{n-1}$ de forma que inicialmente $turno = n-2$ y $P_{n-1}$ será nuestro proceso víctima, podemos obtener la siguiente traza de ejecución, de forma análoga al caso para 3 procesos:
        \begin{itemize}
            \item Inicialmente, todas las claves están en \verb|pasivo|, ya que los $n$ procesos accederán a la vez al protocolo de entrada de la sección crítica.
            \item El proceso $P_{n-3}$ verá a todos pasivos, con lo que decrementará su índice en sentido rotatorio hasta llegar a tener $j = n-2$, con lo que pasará del bucle \verb|A|.
            \item $P_{n-4}$ tendrá un comportamiento similar a $P_{n-3}$, y así con todos los procesos hasta $P_0$, quedando $P_{n-1}$ y $P_{n-2}$ que todavía están en estado pasivo.
            \item Si ahora $P_{n-2}$ accede antes al bucle \verb|A|, como tiene el turno, saltará el bucle sin problema y será el proceso $P_{n-1}$ quien quedará bloqueado en el mismo.
            \item Buscamos el peor caso, luego el momento en el que $turno = n-1$ y $P_{n-1}$ pueda ejecutar la sección crítica queremos alejarlo tanto como podamos. Para ello, buscamos el peor momento para tener $turno = n-2$, con lo que el turno ha de ciclar por todos los procesos anteriores para luego llegar a $n-1$.

                De esta forma, suponemos que el acceso a la sección crítica ha sido aquél que nos daba el mayor número de adelantamientos para el caso $n-1$ (si pensamos que $n=4$, el acceso que nos daba el mayor número de adelantamientos para $n-1=3$ era $P_0\ P_1\ P_0$). Después de este acceso, el proceso $P_{n-2}$ podrá entrar a la sección crítica, con lo que el proceso $P_{n-1}$ lleva ya $f(n-1)+1$ adelantamientos (los adelantamientos realizados al proceso $P_{n-2}$ y el propio adelantamiento de $P_{n-2}$).
            \item Una vez salga el proceso $P_{n-2}$ de la sección crítica, asignará $turno$ a $n-3$, y queremos buscar el peor caso de ejecución para el proceso $P_{n-1}$. Para ello, observemos que si tenemos $turno = n-3$, nos encontramos en las hipótesis de poder encontrar nuevamente el peor caso de ejecución para que entre de nuevo $P_{n-2}$, pero ahora este no conseguirá entrar, ya que después del último proceso que adelante a $P_{n-2}$ (que será $P_0$), este cambiará el turno a $n-1$ (en vez del caso $n-2$), con lo que después de tener otra vez los adelantamientos correspondientes al caso $n-1$, el proceso $P_{n-1}$ conseguirá entrar a la sección crítica. 

                De esta forma, a la cantidad anterior de $f(n-1)+1$ adelantamientos, hemos de sumar nuevamente los adelantamientos que se realizan al proceso $P_{n-2}$ en el caso para $n-1$ procesos, es decir, $f(n-1)$. En este punto, $P_{n-1}$ será capaz de entrar en la sección crítica.
        \end{itemize}
        Por tanto, concluimos que el número máximo de adelantamientos en el caso $n$ es:
        \begin{equation*}
            f(n) = f(n-1) + 1 + f(n-1) = 2f(n-1) + 1 \qquad \forall n\geq 2
        \end{equation*}
        Recurrencia que podemos resolver\footnote{Tal y como aprendimos en asignaturas anteriores}, obteniendo que:
        \begin{equation*}
            f(n) = 2^{n-1} -1 \qquad \forall n\geq 1
        \end{equation*}
        Con lo que el número máximo de adelantamientos que puede sufrir un proceso en el acceso a la sección crítica está acotado, lo que demuestra la equidad del algoritmo.
\end{description}

\begin{ejemplo}
    Observemos que lo que estamos haciendo en la última demostración de la propiedad de equidad es encontrar las peores trazas de ejecución de forma que un proceso sufra el máximo retardo posible en la entrada a la sección crítica. Hemos podido encontrar un patrón repetitivo en estas trazas, el cual está de forma implícita en la demostración y que ahora mostramos de forma explícita (suponiendo que en el caso $n$, es el proceso $P_{n-1}$ quien sufre el número de adelantamientos $f(n)$):
    \begin{description}
        \item [Para $n=0$:] No hay adelantamientos posibles.
        \item [Para $n=1$:] No hay adelantamientos posibles.
        \item [Para $n=2$:] $P_0$.
        \item [Para $n=3$:] $P_0\ P_1\ P_0$.
        \item [Para $n=4$:] $P_0\ P_1\ P_0\ P_2\ P_0\ P_1\ P_0$.
        \item [Para $n=5$:] $P_0\ P_1\ P_0\ P_2\ P_0\ P_1\ P_0\ P_3\ P_0\ P_1\ P_0\ P_2\ P_0\ P_1\ P_0$.
        \item [Para $n$:] (los que adelantan a $P_{n-2}$) $P_{n-2}$ (los que adelantan a $P_{n-2}$)
    \end{description}
\end{ejemplo}

\subsection{Algoritmo de Peterson para 2 procesos}
Peterson propuso una forma simple de resolver el problema de la exclusión mutua para dos procesos, de forma que a partir de entonces se consideró la solución canónica a dicho problema. Esta es generalizable a $n$ procesos.\\

La idea es la que se aplica en una pescadería: los clientes que quieren comprar pescado (en nuestro caso, los dos procesos que quieren entrar a la sección crítica), van asignándose sus turnos en orden de llegada, de forma que el último que se asigna el turno es quien debe esperar a que el otro pase primero (a la sección crítica). Una vez que este primero haya terminado, avisa al siguiente de que ya puede avanzar.

De esta forma, el algoritmo cuenta con las variables compartidas:
\begin{minted}{pascal}
    var c : array[1..2] of boolean;
        turno : 1..2;
\end{minted}
Donde el array de claves \verb|c| se inicializa con todos sus valores a \verb|false|, de forma que \verb|c[i]| indica si el $i$-ésimo proceso está o no solicitando acceder a la sección crítica. La variable \verb|turno| sirve para resolver conflictos en el caso de que ambos procesos intenten acceder a la vez a la sección crítica.\\

Cuando un proceso quiere acceder a la sección crítica, lo que hace primero es cambiar su clave a ``solicitando sección crítica'' (representado con \verb|true|). Posteriormente, se asigna el turno a dicho proceso. En el caso en el que ambos intenten acceder a la vez a la sección crítica, uno de ellos habrá sido necesariamente el último en asignarse el turno, y será el proceso que espere a que el otro ejecute la sección crítica.

\begin{figure}[H]
    \centering
    \setlength{\columnsep}{1cm}
    \begin{multicols}{2}
        \begin{minted}{pascal}
            Process P1();
            begin
               while true do
               begin
                  { Acceso a la SC }
                  c[1] := true;
                  turno := 1;

                  while (c[2] and turno = 1) do
                  begin
                     null;
                  end do

                  { Sección crítica }
                  c[1] := false;
               end do
            end
        \end{minted}
        \begin{minted}{pascal}
            Process P2();
            begin
               while true do
               begin
                  { Acceso a la SC }
                  c[2] := true;
                  turno := 2;

                  while (c[1] and turno = 2) do
                  begin
                     null;
                  end do

                  { Sección crítica }
                  c[2] := false;
               end do
            end
        \end{minted}
    \end{multicols}
    \caption{Algoritmo de Peterson para 2 procesos.}
\end{figure}

\subsubsection{Propiedades de corrección}
A continuación, mostramos que la solución de Peterson no solo es una solución para el problema de la exclusión mutua, sino que además garantiza la equidad en el acceso a la misma para los procesos que intervienen en un programa concurrente que usa dicho algoritmo.
\begin{description}
    \item [Exclusión mutua.]~\\
        Suponiendo que nos encontramos en el proceso $P_i$ (con $i=1$ ó $i=2$), sea $j$ el otro proceso ($j =1$ si $i=2$ ó $j=2$ si $i=1$), supongamos que el proceso $P_i$ consigue entrar a la sección crítica. Esto sucede solo si $c[j]$ es falso o si \verb|turno| es $j$. Es decir, si $P_j$ no está intentando acceder a la sección crítica o si $P_j$ fue el último proceso en asignarse el turno. 

        De esta forma, si el proceso $P_j$ intenta acceder a la vez a la sección crítica (entonces, \verb|c[j] = true|, con lo que \verb|turno = j|), quedará bloqueado, por ser \verb|turno = j| y \verb|c[i] = 1|, con lo que solo un proceso conseguirá al final entrar a la sección crítica.
    \item [Alcanzabilidad de la sección crítica.]~\\
        Supusto que dos procesos intenta acceder a la sección crítica, ambos cambiarán su valor de \verb|c| a \verb|true|, y cambiarán el valor de \verb|turno| a su identificador, quedando uno de ellos bloqueado (el último que modificó dicha variable) y el otro proceso consigue pasar el bucle, con lo que alcanza la sección crítica.

        Si solo un proceso $P_i$ intenta acceder a la sección crítica, si notamos por $j$ al identificador del otro proceso, entonces $P_i$ observará \verb|c[j] = false|, con lo que $P_i$ entrará a la sección crítica.
    \item [Equidad.]~\\
        El algoritmo de Peterson no solo garantiza la condición de vivacidad para los procesos del programa concurrente sino que garantiza la equidad en su acceso a la sección crítica:

        Supongamos que un proceso $P_i$ está bloqueado esperando a entrar a la sección crítica y un proceso $P_j$ está continuamente entrando y saliendo de la misma monopolizando su uso. Sin embargo, cuando $P_j$ salga de la sección crítica por primera vez, cambiará el valor de \verb|turno| a $j$, con lo que $P_j$ quedará bloqueado en el acceso a la sección crítica y $P_i$ conseguirá entrar a la misma, contradicción con que un proceso pueda quedar bloqueado.
\end{description}

\subsection{Algoritmo de Peterson para $n$ procesos}
Lo que haremos para obtener el algoritmo de Peterson que resuelve el problema de la exclusión mutua para cualquier número de procesos $n$ será generalizar el caso de dos procesos, suponiendo que ahora tenemos $n$ procesos que quieren entrar a la sección crítica.

\subsubsection{Idea}
La idea para pasar de 2 procesos a $n$ es repetir $n-1$ veces el proceso de decisión que hicimos anteriormente con 2 procesos. 

Es decir, primero llegarán $n$ procesos al protocolo de acceso a la sección crítica y cada uno de estos irá asignando la variable compartida \verb|turno| a su identificador de proceso, de forma que habrá un proceso de todos esos $n$ que será el último en haber cambiado el valor de la variable compartida \verb|turno|. Este proceso será el que se quedará esperando en un bucle y dejaremos a los $n-1$ procesos restantes pasar a realizar este procedimiento otra vez, obteniendo $n-2$ procesos. Repetiremos este procedimiento un total de $n-1$ veces, para obtener un único proceso de todos los que intentaron acceder a la vez a la sección crítica.\\

A este procedimiento que repetimos para eliminar un proceso del total de procesos le llemaremos \textit{etapa}, de forma que dispondremos de $n-1$ etapas para, en el caso de que lleguen a la vez los $n$ procesos al protocolo de acceso a la sección crítica, dejar pasar solo a uno de ellos a la sección crítica, como resultado de pasar los $n$ procesos por las $n-1$ etapas distintas. Una vez que un proceso haya completado las $n-1$ etapas, se encontrará en la ``etapa $n$'', que es la propia sección crítica.

\subsubsection{Implementación}
Para implementar un número de etapas que coincida con el número de procesos que en un futuro ejecutarán nuestro programa concurrente, nos vemos obligados a crear un bucle que itere $n-1$ veces, de forma que cada iteración del bucle será una nueva etapa. Así mismo, contaremos con una variable $j$, que será la que nos indique en qué etapa nos encontramos.

Además, necesitaremos de las siguientes variables compartidas:
\begin{minted}{pascal}
    var c : array[0..n-1] of -1..n-2;
        turno : array[0..n-2] of 0..n-1;
\end{minted}
donde:
\begin{itemize}
    \item \verb|c| es un array que contiene las claves de los $n$ procesos que ejecutarán nuestro programa concurrente, de forma que \verb|c[i]| será la clave del proceso $i$-ésimo, la cual podrá tomar los valores:
        \begin{itemize}
            \item $-1$, que indica que el proceso está en estado pasivo (no está intentando acceder a la sección crítica).
            \item $k$ con $k \in \{0,\ldots,n-1\}$, que indica en qué etapa se encuentra el proceso $i$-ésimo, entendiendo que la primera etapa es la $0$ y la última la $n-2$; siendo la etapa $n-1$ la correspondiente a la sección crítica.
        \end{itemize}
    \item \verb|turno| es un array de forma que \verb|turno[j]| indica el valor del turno para la etapa $j$-ésima, cuyo valor es cualquiera de los identificadores de los procesos que ejecutan nuestro programa (como tendremos $n$ procesos, estos están identificados desde el $0$ hasta el $n-1$).
\end{itemize}
Procedemos ahora a mostrar el código del algoritmo de Peterson para $n$ procesos, el cual pasaremos a explicar inmediatamente:
\begin{minted}{pascal}
    Process Pi();
    begin
       while true do
       begin
          { - Acceso a la SC - }
          for j=0 to n-2 do
          begin
             c[i] := j;
             turno[j] := i;

             { Existe k <> i con c[k] >= j AND turno[j] = i }
             for k=0 to n-1 do
             begin
                if (k = i) then continue;
                while (c[k] >= j and turno[j] = i ) do
                   null;
                end do
             end
          end

          c[i] := n-1;     { Proceso en etapa n-1 (SC) }
          { - Sección Crítica - }
          c[i] := -1;
       end do
    end
\end{minted}
De esta forma, cuando un proceso llega al protocolo de entrada de la sección crítica, lo que hace es entrar en la primera etapa (\verb|j=0|). En esta, actualiza el valor de su clave a \verb|j| y actualiza la variable \verb|turno| a su identificador.

Posteriormente, ejecuta el \verb|for| indicado en el código, el cual es equivalente a:
\begin{minted}{pascal}
    while(Exists k <> i : c[k] >= j AND turno[j] = i) do 
       null;
    end do
\end{minted}
Es decir, si \verb|turno[j] = i| (lo que significa que el proceso $i$ fue el último en cambiar el valor de la variable \verb|turno|) y si hay otro proceso $k$ de forma que \verb|c[k] >= j|, es decir, que existe un proceso en una etapa superior, entonces el proceso $i$ se quedará esperando en la etapa $j$, hasta que:
\begin{itemize}
    \item Llegue un nuevo proceso a la etapa $j$, con lo que el proceso $i$ ya no habrá sido el último en asignar el turno (\verb|turno[j] != i|).
    \item No haya ningún proceso en ninguna etapa siguiente (es decir, que el proceso que antes había ha terminado ya la ejecución de la sección crítica (ya que es la única forma de avanzar), con lo que no existe un $k$ de forma que \verb|c[k] >= j|).
\end{itemize}

\begin{ejemplo}
    Para motivar el buen funcionamiento del algoritmo, el cual demostraremos después, mostramos ya un ejemplo de ejecución del mismo, donde suponemos que tenemos $n=4$ procesos: $P_0$, $P_1$, $P_2$ y $P_3$ que quieren acceder a la sección crítica.

    Suponiendo que observamos la traza de ejecución:
    \begin{minted}{pascal}
        turno[0] = 2; turno[0] = 0; turno[0] = 3; turno[0] = 1;
    \end{minted}
    Entonces, obtendríamos el siguiente comportamiento de los procesos (siendo $x$, $y$ y $z$ valores entre 0 y 3, cuyo valor no nos es relevante):
    \begin{table}[H]
    \centering
    \begin{tabular}{|c|c|c|c|c|c|c|c|}
        \hline
        $t$ & \verb|c[0]| & \verb|c[1]| & \verb|c[2]| & \verb|c[3]| & \verb|turno[0]| & \verb|turno[1]| & \verb|turno[2]| \\
        \hline
        0 & -1 & -1 & -1 & -1 & $x$ & $y$ & $z$ \\
        \hline
        1 & 0 & 0 & 0 & 0 & $2$ & $y$ & $z$ \\
        \hline
        2 & 0 & 0 & 1 & 0 & $0$ & $2$ & $z$ \\
        \hline
        3 & 1 & 0 & 2 & 0 & $3$ & $0$ & $2$ \\
        \hline
        4 & 2 & 0 & 3 & 1 & $1$ & $3$ & $0$ \\
        \hline
        5 & 3 & 1 & -1 & 2 & $1$ & $1$ & $3$ \\
        \hline
        6 & -1 & 2 & -1 & 3 & $1$ & $1$ & $1$ \\
        \hline
        7 & -1 & 3 & -1 & -1 & $1$ & $1$ & $1$ \\
        \hline
        7 & -1 & -1 & -1 & -1 & $1$ & $1$ & $1$ \\
        \hline
    \end{tabular}
    \caption{Evoluión de los estados de los procesos.}
    \end{table}
    Donde hemos supuesto que todos los procesos llevan una velocidad de ejecución similar, de forma que en cada unidad de tiempo, un proceso es capaz de cambiar su clave, asignar su turno y comprobar la condición del bucle \verb|while|. De esta forma (cada índice indica el valor de $t$):
    \begin{enumerate}[label=(\arabic*), start=0]
        \item Ningún proceso accede al protocolo de entrada a la sección crítica.
        \item Todos los procesos acceden a la vez al protocolo de entrada, a la etapa 0, de forma que $P_2$ es el primero en asignar el valor de \verb|turno[0]|.
        \item Como no hay ningún proceso en la etapa 1 (no existe $k$ de forma que \verb|c[k] >= j|), entonces el proceso $P_2$ pasa a la siguiente etapa, la 1. En este mismo instante, $P_0$ cambia el valor de \verb|turno[0]|.
        \item Ahora, $P_3$ cambia el valor de \verb|turno[0]|, lo que provoca que $P_2$ no haya sido el último proceso en cambiar el turno en la etapa 0, con lo que avanza a la etapa 1, cambiando su valor, lo que hace que $P_2$ ya no haya sido el último, avanza a la etapa 2.
        \item $P_2$ ve que no hay ningún proceso en la etapa 3 (la sección crítica), con lo que avanza a dicha etapa. En este momento, $P_1$ cambia el valor de \verb|turno[0]|, con lo que $P_3$ avanza a la siguiente etapa, lo que provoca que $P_0$ también avance de etapa.
        \item $P_2$ abandona la sección crítica, pasando ahora a estar pasivo, lo que hace que $P_0$ pueda entrar en la sección crítica.
    \end{enumerate}
    Notemos que lo que hemos hecho ha sido elegir una traza muy concreta para mostrar este ejemplo de ejecución, en el que hemos intentado dar ejemplos de las dos razones por las que un proceso puede avanzar a la siguiente etapa:
    \begin{itemize}
        \item Por no haber ningún proceso en ninguna etapa siguiente.
        \item Por no ser el último proceso en haber cambiado la variable \verb|turno| en la etapa actual.
    \end{itemize}
\end{ejemplo}

\subsubsection{Propiedades de corrección}
A continuación, mostraremos las propiedades de corrección del algoritmo de Peterson para $n$ procesos. Para ello, será necesario antes introducir una definición y varios lemas:

\begin{definicion}[Precedencia]
    Dados dos procesos, $P_i$ y $P_k$, decimos que $P_i$ precede a $P_k$ si $c[i] > c[k]$. Es decir, si $P_i$ se encuentra en un estado posterior al estado de $P_k$.
\end{definicion}

Además, será común nombrar ``condición de espera de la etapa $j$'', con lo que nos estaremos refiriendo a la condición: 
\begin{center}
\verb|Exists k <> i : c[k] >= j AND turno[j] = i|.
\end{center}

\begin{lema}\label{lema:1}
    Un proceso que precede a todos los demás puede avanzar al menos una etapa.
    \begin{proof}
        Sea $P_i$ un proceso que precede a todos los demás, entonces tendremos que $c[i] = j > c[k]$ $\forall k \neq i$. De esta forma, no existirá ningún proceso en una etapa superior, con lo que la condición de espera en la etapa $j$ será falsa, permitiendo a dicho proceso avanzar a la etapa $j+1$.\\

        Sin embargo, como la evaluación de la condición de espera en una etapa no se realiza de forma atómica (ya que debemos implementar dicha condición mediante un bucle \verb|for| que itere sobre $k$), puede ocurrir que uno o varios procesos alcancen también la etapa $j$ en la que se encuentra $P_i$ mientras comprueba la condición de espera de la etapa $j$. Sin embargo, tan pronto como esto suceda, algún proceso cambiará el valor de la variable \verb|turno[j]|, lo que hará que el proceso $P_i$ evalúe la condición de espera de la etapa $j$ como falsa, lo que le permita avanzar a la etapa $j+1$.
    \end{proof}
\end{lema}

\begin{lema}\label{lema:2}
    Cuando un proceso pasa de la etapa $j$ a la $j+1$, se ha de verificar alguna de las dos siguientes condiciones:
    \begin{enumerate}[label=\alph*)]
        \item Que dicho proceso preceda a todos los demás.
        \item Que dicho proceso no estuviese solo en la etapa $j$.
    \end{enumerate}
    \begin{proof}
        Supuesto que el proceso $P_i$ está a punto de pasar de la etapa $j$ a la $j+1$, entonces la condición de espera de la etapa $j$ ha de ser falsa para $P_i$, con lo que:
        \begin{itemize}
            \item No existe $k \neq i$ que verifique $c[k]\geq j$, es decir, que $c[k] < j = c[i]$ para todo $k \neq i$, con lo que $P_i$ precede a todos los demás procesos.
            \item $turno[j] \neq i$, con lo que existirá otro proceso $P_k$ de forma que haya sido el último en modificar el valor de $turno[j]$, con lo que dicho proceso $P_k$ ha de estar necesariamente en la etapa $j$, al igual que $P_i$.
        \end{itemize}
    \end{proof}
\end{lema}

\begin{lema}\label{lema:3}
    Si tenemos al menos dos procesos en la etapa $j$, entonces ha de haber, al menos, un proceso en cada etapa anterior.
    \begin{proof}
        Distinguimos dos casos:
        \begin{description}
            \item [$j=1$.] Supuesto que tenemos dos procesos en la etapa 1, $P_i$ y $P_k$, entonces alguno de ellos habrá sido el último en pasar de la etapa 0 a la 1. No perdemos generalidad suponiendo que dicho proceso es $P_i$. En dicho momento, y por el Lema~\ref{lema:2}, se cumplirá que o bien $P_i$ precede a todos los demás o bien que no estaba solo en la etapa 0.

                Sin embargo, como $P_k$ pasó antes a la etapa 1 que $P_i$, es imposible que $P_i$ preceda a todos los demás al pasar a la etapa 1, por no preceder a $P_k$. Por ello, $P_i$ no estaba solo cunado pasó de la etapa 0 a la 1, con lo que habrá otro proceso, $P_l$, en la etapa 0.

                De esta forma, tenemos demostrado el lema para el caso $j=1$.
            \item [Supuesto que el lema es cierto para $j-1$, veámosolo para $j$.] Por un razonamiento análogo al anterior, supuesto que tenemos dos procesos, $P_i$ y $P_k$ en la etapa $j$, y supuesto que $P_i$ fue el último de ellos en avanzar desde la etapa $j-1$ a la $j$, por el Lema~\ref{lema:2}, dicho proceso no estaba solo en la etapa $j-1$, con lo que hay al menos un proceso en cada etapa, desde la 0 hasta la $j-1$.
        \end{description}
    \end{proof}
\end{lema}

\begin{lema}\label{lema:4}
    El número máximo de procesos que puede haber en la etapa $j$ es de $n-j$ procesos.
    \begin{proof}
        Supuesto que tenemos al menos dos procesos en la etapa $j$, entonces por el Lema~\ref{lema:3} las etapas $0$, $1$, \ldots, $j-1$ (son $j$ etapas) han de contener al menos un proceso, con lo que ninguno de estos procesos puede estar en la etapa $j$, con lo que el número máximo de procesos en la etapa $j$ es $n-j$.
    \end{proof}
\end{lema}

Estamos ya listos para probar todas las propiedades de corrección del algoritmo de Peterson:
\begin{description}
    \item [Exclusión mutua.]\ 
        \begin{itemize}
            \item Supuesto que tenemos un proceso en la etapa $n-2$ y uno en la $n-1$ (en la sección crítica), entonces el proceso de la etapa $n-2$ no puede avanzar, ya que este proceso no precede a todos los demás y está solo en dicha etapa, con lo que de avanzar, no se cumpliría el Lema~\ref{lema:2}.

                Sin embargo, cuando el proceso de la etapa $n-1$ abandone dicha etapa, entonces el proceso de la etapa $n-2$ se encontraría en las hipótesis del Lema~\ref{lema:1}, con lo que podría avanzar a dicha etapa.
            \item Supuesto que tenemos dos procesos en la etapa $n-2$, en cuanto uno avance, el otro no podrá hacerlo, tal y como acabamos de razonar.
            \item Gracias al Lema~\ref{lema:2}, sabemos que el número máximo de procesos en la etapa $n-2$ es de 2, por lo que la distinción de casos está ya completa.
        \end{itemize}
    \item [Alcanzabilidad de la sección crítica.]~\\
        Por reducción al absurdo, supongamos que hay una traza de ejecución del programa concurrente en la que la sección crítica es inalcanzable, es decir, que ningún proceso puede llegar a la sección crítica. En dicho caso, suponemos que tenemos un proceso en una etapa $j$ y distingamos casos:
        \begin{itemize}
            \item Si dicho proceso precede a todos los demás, entonces por el Lema~\ref{lema:1} este puede avanzar.
            \item Si no, hay al menos un proceso que le precede, es decir, que está en una etapa superior a $j$:
                \begin{itemize}
                    \item Si dicho proceso o cualquier otro precede a todos los demás, por el Lema~\ref{lema:1}, este puede avanzar.
                    \item Si no hay ningún proceso que preceda a todos los demás, entonces hay al menos dos procesos en la etapa más avanzada, con lo que alguno fue el último en asignar el valor de la variable \verb|turno| de dicha etapa, por lo que para el resto de procesos la condición de espera de dicha etapa será falsa, permitiéndoles avanzar.
                \end{itemize}
        \end{itemize}
        En cualquier caso, habrá algún proceso que consiga avanzar en cualquier etapa, con lo que dicha situación es imposible.
    \item [Equidad]~\\
        No solo demostraremos la vivacidad del algoritmo de Peterson, sino también la equidad de los procesos en el acceso a la sección crítica. Para ello, supongamos el caso más desfavorable en la entrada a la sección crítica y calculamos el número máximo de turnos que el proceso ha de esperar para poder entrar en la sección crítica:

        El caso más desfavorable posible se da cuando los $n$ procesos intentan acceder a la vez a la sección crítica, de forma que el proceso $P_i$ es el último de los $n$ procesos en modificar el valor de la variable \verb|turno[0]|, con lo que se bloqueará en dicha etapa y dejará pasar al resto de los $n-1$ procesos. Estos tardarán 1 un turno en propagarse por cada etapa y como tenemos $n-2$ etapas restantes (más el turno de la etapa 0), llevamos ya $n-1$ turnos de espera para $P_i$. Este permanecerá bloqueado en esta etapa hasta que:
        \begin{itemize}
            \item Preceda a todos los demás procesos, con lo que todos los $n-1$ procesos anteriores hayan ejecutado ya la sección crítica.

                En este caso, $P_i$ podrá entrar a la sección crítica, sin esperar un turno más.
            \item Venga un nuevo proceso a la etapa 0, es decir, haya un proceso $P_k$ que ya haya ejecutado la sección crítica y entre de nuevo en el protocolo de entrada, lo que permitirá a $P_i$ avanzar, al menos en una etapa.

                En este caso, $P_i$ tendrá que esperar tantos turnos como antes para llegar a la sección crítica menos uno, ya que ha avanzado de etapa.
        \end{itemize}
        En cualquier caso, para llegar a la sección crítica, $P_i$ debe esperar a que el proceso que se bloqueó en la etapa 1 llegue a la sección crítica.

        De esta forma, si $r(n)$ es el número de turnos de espera para $n$ procesos, entonces tenemos que en el peor caso, el proceso $P_i$ deberá dejar pasar a $n-1$ procesos (con lo que pasarán $n-1$ turnos), así como esperar los turnos correspondientes a que el proceso que se bloquea en la etapa 1 entre a la sección crítica, que resulta en un problema similar pero en este caso de $r(n-1)$ turnos de espera. De esta forma:
        \begin{equation*}
            r(n) = n-1+r(n-1)
        \end{equation*}
        Por tanto, el número de turnos máximos de espera para entrar en la sección crítica según el algoritmo de Peterson está limitado por el valor
        \begin{equation*}
            r(n) = \dfrac{n(n-1)}{2}
        \end{equation*}
\end{description}

Recordando el algoritmo de Knuth, este era un algoritmo más complicado que también resolvía el problema de la exclusión mutua. Sin embargo, pese a la sencillez del algoritmo de Peterson, es más eficiente en cuanto al acceso a la sección crítica, ya que el orden del tiempo de espera en el algoritmo de Knuth era del orden $O(2^n)$, siendo este $O(n^2)$ en el algoritmo de Peterson.\\

Por tanto, decimos que el algoritmo de Peterson es la solución canónica al problema de la exclusión mutua, al ser el algoritmo más eficiente de los dos estudiados en esta Sección que cumplen con la quinta propiedad de corrección (la de demostrar un retraso máximo en el acceso de los procesos a la sección crítica).

% // TODO: problema de la exclusion mutua en capitulo 3 en distribuidos.

\section{Definición de un monitor}
El concepto de semáforo se desarrolló previamente en el Seminario 1 de prácticas\footnote{Por lo que el lector debería estar familiarizado con ellos.}. Los semáforos presentan dos grandes limitaciones:
\begin{enumerate}
    \item Están basados en variables compartidas del programa, por lo que no fomentan la modularidad de los programas, impidiendo su reutilización.
    \item Las operaciones de los semáforos (\verb|sem_wait| y \verb|sem_signal|) se encuentran dispersas a lo largo del código del programa concurrente. Además, estas instrucciones no solo afectan al bloque de código en el que se encuentran, sino a cualquier otro bloque que use el mismo semáforo.
\end{enumerate}
En definitiva, los semáforos no son un buen mecanismo de programación concurrente, y además la verificación de programas que usan semáforos es muy complicada. 

Era necesario encontrar un nuevo mecanismo de programación concurrente que permitiera la encapsulación de la información y de la sincronización entre procesos, así como programar las operaciones de sincronización (como \verb|wait| o \verb|signal|) dentro de bloques o procedimientos que se ejecuten con instrucciones atómicas, para que las instrucciones de sincronización no se encuentren desperdigadas por el programa.
Fue Charles Antony Richard Hoare quien inventó los monitores, concepto en el que ahondaremos a lo largo de este Capítulo.\\

La idea básica de monitor es un módulo que contiene un conjunto de variables a las que llamaremos \textit{variables permanentes}\footnote{A pesar de su nombre, no serán constantes, sino que podremos modificar su valor.}, de forma que dichas variables solo podrán ser alteradas dentro de los procedimientos del módulo monitor. Garantizaremos que la ejecución de cada uno de esos procedimientos se ejecute la mayor parte del tiempo como una única instrucción atómica, salvo que se produzca algo por lo que interrumpir la ejecución del procedimiento.\\

Podemos pensar en un monitor como en un tipo de dato abstracto que define tipos y variables permanentes propias del monitor, así como un conjunto de procedimientos dentro de dicho módulo. No debemos pensar en los monitores como en una clase, ya que no pueden hacer lo mismo que ellas (no se pueden instanciar y tampoco existe polimorfismo o ligadura dinámica).

\subsubsection{Ventajas}
A continuación, los programas concurrentes estarán formados tanto por procesos que se ejecutarán de forma concurrente, como por monitores, los cuales velarán por la sincronización y acceso a variables compartidas de dichos procesos, de forma que no se produzcan condiciones de carrera o comportamientos indeseados. Podremos modelar tantas relaciones de interacción entre los procesos de un programa concurrente como queramos. De esta forma, el uso de los monitores o de procedimientos asociados a monitores no restringen las posibilidades del modelado de un sistema concurrente. 

Los procesos de un programa concurrente no tendrán que llamar a operaciones de sincronización, sino que llamarán a procedimientos del monitor, los cuales realizarán la funcionalidad deseada sobre las variables compartidas garantizando la sincronización entre los procesos.\\

Además, los monitores nos permiten una alta reusabilidad de código, ya que podremos reutilizar un monitor ya creado para resolver problemas similares. Sin embargo, la reutilización de código no es similar a la usada en programación orientada a objetos mediante instancias de una misma clase, sino que se hará por copias parametrizables: tendremos una definición de un monitor basada en parámetros, y cuando necesitemos usar un monitor, crearemos una copia de dicha definición parametrizándola (pasándole los parámetros que necesitemos para resolver nuestro problema). De esta forma, no es reutilización por instanciación, sino por \textit{parametrización}.\\

Los procesos que usemos en los progrmas concurrentes no verán el acceso a las variables compartidas, sino que será realizado por los procedimientos del monitor, garantizando que se hacen como deben hacerse, evitando condiciones de carrera. De esta forma, los monitores garantizan la ocultación de las variables compartidas, haciéndolas transparentes a los procesos del sistema concurrente.\\

Finalmente, existen unos axiomas que nos permiten verificar los programas concurrentes que usen monitores de forma sencilla. Dichas demostraciones estarán basadas en el uso de los invariables globales. Ahondaremos en la verificación de programas concurrentes que utilicen monitores más adelante.

\subsection{Concepto de monitor}\label{sec:concepto_de_monitor}
A modo de resumen para comenzar a definir lo que es un monitor, podemos decir que:
\begin{itemize}
    \item Es un módulo con un conjunto de variables permanentes que solo pueden ser modificadas por los procedimientos del monitor.
    \item Cada uno de los procedimientos\footnote{Podemos pensar en ellos como en los ``métodos'' de una clase, haciendo hincapié en que los monitores \textbf{no son} clases.} de un monitor se ejecutan en exclusión mutua (garantizando el acceso a las variables compartidas sin condiciones de carrera). Sin embargo, estos no tienen por qué ejecutarse completamente, sino que pueden interrumpirse y en algún momento futuro seguir ejecutándose en exclusión mutua.
    \item La ejecución de los procedimientos de un monitor modifican el estado interno del mismo (esto es, el conjunto de las variables permanentes asociadas al monitor).
    \item El estado inicial del monitor (de sus variables permanentes) se establece mediante la ejecución de un procedimiento especial, al que llamaremos \textit{código de inicialización}. Este se ejecuta tras la declaración de una variable de tipo monitor y da valores iniciales a las variables permanentes.
\end{itemize}
De esta forma, un monitor puede visualizarse de forma intuitiva en la Tabla~\ref{esq:monitor_1}, como un conjunto que englobla:
\begin{itemize}
    \item Un conjunto de variables, llamadas \textit{variables permanentes}, que no son accesibles desde fuera del monitor.
    \item Un conjunto de procedimientos que el monitor proporciona como servicio a los procesos de un programa concurrente (para por ejemplo, acceder a las variables permanentes que serán las variables que compartan dichos procesos), llamados \textit{procedimientos exportados} o \textit{exportables}.
    \item Un procedimiento especial llamado \textit{código de inicialización}, que permite inicializar las variable permanentes.
\end{itemize}

\begin{table}[H]
\centering
\begin{tabular}{|l|}
\hline
Variables \\
permanentes \\
\hline
Procedimientos \\
exportados \\
\hline
Código de \\ 
inicialización \\
\hline
\end{tabular}
\caption{Esquema de un monitor.}
\label{esq:monitor_1}
\end{table}

\begin{ejemplo}
    Aunque todavía no entendemos muy bien qué es un monitor, daremos a continuación un ejemplo de uso del mismo para ilustrar la definición que queremos dar de monitor, pese a que algunas cosas del ejemplo no podamos entenderlas todavía y deberemos dejarlas para más adelante\footnote{Como el tipo de dato \texttt{cond}.}.\\

    En este ejemplo, queremos solventar un problema mediante el paradigma productor/consumidor. Tendremos dos procesos, un productor y un consumidor, de forma que el productor escribirá en un buffer (o vector) que usaremos como cola cíclica (esto es, que si nos pasamos de la posición final, volvemos al inicio y con planificación FIFO), mientras que el consumidor irá leyendo los datos de dicho buffer. Siendo \verb|Buf| una variable de tipo monitor que luego definiremos en este ejemplo, el código del productor y del consumidor será el siguiente (pensando en que tenemos que usar procedimientos del monitor para el acceso a las variables compartidas, en este caso el buffer):
    \begin{minted}[escapeinside=\#\#]{pascal}
        Proceso Prod1::
          var d : tipo_dato;

          while true do begin
            d = producir();
            Buf.insertar(d); {mete d en el buffer}
          end do
    \end{minted}
    \begin{minted}[escapeinside=\#\#]{pascal}
        Proceso Cons1::
          var x : tipo_dato;

          while true do begin
            Buf.retirar(x); {retira del buffer en x}
            consumir(x);
          end do
    \end{minted}
    El código del monitor será el siguiente en pseudo-pascal (hemos omitido el código de inicialización):
    \begin{minted}[escapeinside=\#\#]{pascal}
        Monitor Buf
          var
            -elementos_ocupados : int;
            -frente, atras: 0..N-1;
            -no_vacio, no_lleno : cond;

          +insertar(d : tipo_dato);
          +retirar(var x : tipo_dato);
    \end{minted}
    Donde vemos 5 variables permanentes: \verb|elementos_ocupados|, que mide la cantidad de posiciones ocupadas del buffer, \verb|frente|, que marca la casilla en la que el productor insertará el próximo dato (por tanto, ha de estar siempre vacía), \verb|atras|, que marca la casilla de la que leerá el consumidor, \verb|no_vacio| y \verb|no_lleno|, variables de tipo \verb|cond|, las cuales aprenderemos lo que hacen más adelante.
    
    Contamos además con dos procedimientos: \verb|insertar|, que inserta un dato en el buffer en caso de que haya hueco (si no hay hueco, se bloquea hasta que el consumidor lea un dato y deje un hueco libre):
    \begin{minted}[escapeinside=\#\#]{pascal}
        procedure insertar(d : tipo_dato) begin
          if((frente + 1) mod N = atras) then no_lleno.wait();
          introducir(buf, frente, d);  {inserta d en la posicion frente en el buffer}
          elementos_ocupados += 1;
          frente = (frente + 1) mod N;
          no_vacio.signal();
        end
    \end{minted}
    Y con el procedimiento \verb|retirar|, que retira un dato del buffer y lo devuelve como resultado del procedimiento, siempre que esto sea posible (es decir, si no hay ningún dato que leer en el buffer, se bloquea esperando a que el productor ponga algún dato):
    \begin{minted}[escapeinside=\#\#]{pascal}
        procedure retirar(var x : tipo_dato) begin
          if(frente = atras) then no_vacio.wait();
          eliminar(buf, atras, x);  {inserta buf[atras] en x y lo borra del buffer}
          elementos_ocupados -= 1;
          atras = atras mod N + 1;
          no_lleno.signal();
        end
    \end{minted}
    Como hemos ya comentado mientras mostrábamos los pseudocódigos del ejemplo, hay que establecer condiciones que identifiquen las dos condiciones inseguras del ejemplo: que el buffer esté lleno o que el buffer esté vacío:
    \begin{itemize}
        \item Si \verb|frente = atras|, entonces el último dato que se ha de consumir está en una casilla vacía en la que el productor escribirá. Se trata de la situación en la que el buffer está vacío. Debemos por tanto, evitar que el consumidor lea un dato del buffer.
        \item Si \verb|(frente + 1) mod N = atras|, entonces el siguiente dato a introducir en el buffer está justo delante del dato a consumir. Se trata de la situación en la que el buffer está lleno. Debemos por tanto, evitar que el productor introduzca un dato en el buffer\footnote{Definimos anteriormente que \texttt{frente} siempre apunta a una casilla vacía, por lo que como máximo el buffer tendrá ocupados $N-1$ elementos.}.
    \end{itemize}
    Los procesos del programa llaman a los procedimientos del monitor, y no tienen acceso directo al buffer, por lo que no pueden saber cuándo este está lleno o vacío. De esta forma, lo que sucederá es que los procedimientos internos del monitor realizarán una sincronización interna mediante el uso de llamadas bloqueantes:
    \begin{itemize}
        \item Si el buffer está lleno y el productor se dispone a escribir un dato, quedará el proceso bloqueado hasta que un consumidor lea un dato. Este señalará (\verb|signal|) al productor, desbloqueándolo.
        \item Si el buffer está vacío y el consumidor se disopne a leer un dato, quedará bloqueado el proceso que ejecute el procedimiento del monitor. Cuando el productor escriba un dato, enviará una señal al consumidor, desbloqueándolo.
    \end{itemize}
    Esta funcionalidad se consigue mediante las variables de tipo \verb|cond|. Se verán a continuación, pero para entenderlas por ahora digamos que necesitamos tener una variable de tipo \verb|cond| por cada razón por la que queremos bloquear un proceso\footnote{En el caso de productor/consumidor, queremos bloquear un proceso si sucede alguno de los dos puntos superiores, condiciones inseguras, luego nos harán falta dos variables de tipo \texttt{cond}. En otros problemas, el número de variables de tipo \texttt{cond} podría ser otro.}.\\

    El código de los procedimientos es ejecutado por los propios procesos que ejecutan cada proceso (productor o consumidor, en este caso) del programa concurrente. Por tanto, si el productor ejecuta un procedimiento del monitor con un \verb|wait|, dicho proceso se bloqueará y no podrá ejecutar código hasta desbloquearse.\\

    Para que el código que hemos visto funcione adecuadamente, nos falta introducir un último concepto en los monitores, y es que mientras se ejecuta un procedimiento de un monitor, no se puede ejecutar ningún otro, sino que han de ejecutarse en \textbf{exclusión mutua}.
\end{ejemplo}

\subsection{Características de programación con monitores}
Una vez ilustrado el uso de la herramienta que estamos construyendo en este Capítulo mediante el ejemplo anterior, vamos ahora a introducir la noción de que solo puede ejecutarse a la vez un único procedimiento de un monitor.\\

Como ya hemos visto, los procedimientos de los monitores no tienen por qué ejecutarse de principio a fin, sino que un proceso puede comenzar a ejecutar un procedimiento, bloquearse (dejando por tanto libre al monitor) y que otro proceso comience a ejecutar un procedimiento de dicho monitor, sucediéndose un entrelazamiento de las trazas de ejecución de los procedimientos.\\

Cuando un proceso se encuentra ejecutando un procedimiento del monitor, decimos que el monitor está \emph{ocupado}. En caso contrario, diremos que este está \emph{libre}. Notemos que si un proceso se bloquea mientras ejecuta un procedimiento del monitor, el monitor tiene que quedar libre, ya que si no no habría forma de volver a despertar a dicho proceso (tenemos que ejecutar un \verb|signal| sobre la misma variable \verb|cond| que bloqueó al proceso\footnote{Se explicará más adelante.}). La situación de bloquear a un proceso y dejar que entre otro al monitor es delicada y debe hacerse con cuidado, para garantizar que solo haya un único proceso ejecutando un procedimiento del monitor al mismo tiempo.\\

Los monitores son objetos \emph{pasivos}. Esto es, no tienen una hebra dentro que ejecute su código, sino que simplemente proporciona código (sus procedimientos) a otros procesos para que sean ellos quien ejecuten el código del monitor.\\

Para implementar una librería con monitores en un lenguaje de programación base, este debe tener la propiedad de ser \emph{reentrante}.
\begin{definicion}
    Un lenguaje de programación tiene la propiedad de ser reentrante si, siempre que tengamos un proceso ejecutando una función y este se bloquea, sea capaz de conservar la siguiente instrucción a ejecutar y el valor de sus variables locales tras desbloquearse. Es decir, el proceso no debe enterarse localmente\footnote{Las variables locales a la función deben mantenerse, pero puede haber variables globales que sí hayan cambiado.} de que nada haya cambiado mientras estaba bloqueado.
\end{definicion}
Notemos que debemos tener esta propiedad en el lenguaje de programación con el que trabajemos para poder hacer uso de funciones bloqueantes (como \verb|wait|) dentro de los procedimientos de un monitor, algo básico en el funcionamiento de este. Afortunadamente, actualmente todos los lenguajes de programación que encontramos en el mercado son reentrantes.\\

\subsubsection{Copias paramétricas de un monitor}
El siguiente ejemplo nos ilustra cómo podemos crear nuevos monitores a partir de uno ya creado, fijando parámetros que use el código de inicialización.

\begin{ejemplo}
    Aunque los monitores están pensados para programas concurrentes (ya que no tiene sentido su uso en programas secuenciales), usaremos en este ejemplo un monitor en un programa secuencial, ya que solo nos interesa la forma en la que los monitores inicializan sus variables permanentes\footnote{Además, no hemos terminado de desarrollar cómo es que solo puede ejecutarse a la vez un único procedimiento del monitor, por lo que no entendemos hasta ahora cómo es que sirven para sincronizar programas concurrentes.}.\\

    Tenemos un programa en el que necesitamos dos variables, las cuales queremos consultar e incrementar mediante un incremento previamente fijado que no cambiará. Para ello, creamos un monitor de acceso a una variable, con parámetros de entrada, para luego poder crear dos copias parametrizadas del mismo. El código del monitor será algo parecido a:
    \begin{minted}[escapeinside=\#\#]{pascal}
        class monitor VariableProtegida(inicio, incremento : integer);
          var x, inc : integer;

          procedure incremento();
          begin
            x = x + inc;
          end

          procedure valor(var v : integer);
          begin
            v = x;
          end
          
          begin
            x = inicio; inc = incremento;
          end
    \end{minted}
    De esta forma, podemos usar dos copias del monitor de la forma:
    \begin{minted}[escapeinside=\#\#]{pascal}
        var mv1 : VariableProtegida(0,1);   {empieza en 0 e incrementa en 1}
            mv2 : VariableProtegida(10,4);  {empieza en 10 e incrementa en 4}
            a, b : integer;
        begin
          mv1.incremento();  {+=1}
          mv1.valor(a);      {a=1}
          mv2.incremento();  {+=4}
          mv2.valor(b);      {b=14}
        end
    \end{minted}
\end{ejemplo}

\subsection{Exclusión mutua en los procedimientos de un monitor}
Si tenemos varios procesos del programa concurrente que quieren hacer uso de procedimientos del monitor a la vez, solo podremos dejar pasar un proceso al monitor (suponiendo que este se encuentre libre). Para los otros procesos, almacenaremos su llamada al procedimiento. 

Para ello, todos los monitores tienen implementada una cola con planificación FIFO, llamada \emph{cola de entrada al monitor}. Si tenemos dos procesos que quieren acceder a un procedimiento de un monitor libre, solo podrá hacerlo un proceso. La llamada al procedimiento del monitor del otro proceso quedará almacenada en la cola de entrada al monitor, y este pasará a ejecutar el procedimiento deseado una vez el proceso anterior haya dejado libre el monitor. 
\begin{observacion}
    En esta asignatura, supondremos que la cola de entrada al monitor es suficientemente larga como para albergar a todos los procesos que necesiten esperar a que el monitor quede libre.
\end{observacion}

Podemos representar la vida de un proceso de un programa concurrente que hace uso de monitores para sincronizar a sus procesos con el siguiente diagrama:
\begin{figure}[H]
\centering
\begin{tikzpicture}[node/.style={rectangle, draw, minimum size=1cm, align=center},
                    edge/.style={-stealth}]
    \node[node] (A) {Procedimiento\\terminado};
    \node[node, left =of A] (B) {\begin{tabular}{c}Ejecutando código\\del proceso\\(fuera del monitor)\end{tabular}};
    \node[node, left =of B] (C) {\begin{tabular}{c}Llamando a\\procedimiento\\del monitor\end{tabular}};
    \node[node, below =of C] (D) {\begin{tabular}{c}Esperando en\\cola de entrada\\al monitor\end{tabular}};
    \node[node, right =of D] (E) {\begin{tabular}{c}Monitor\\libre\end{tabular}};
    \node[node, right =of E, xshift=.6cm] (F) {\begin{tabular}{c}Ejecutando procedimiento\\(dentro del monitor)\end{tabular}};
    
    \draw[edge] (A) -- (B);
    \draw[edge] (B) -- (C);
    \draw[edge] (C) -- (D);
    \draw[edge] (D) -- (E);
    \draw[edge] (E) -- (F);
    \draw[edge] (F) -- (A);
\end{tikzpicture}
\caption{Vida de un proceso en un programa concurrente con monitores.}
\end{figure}

De esta forma, podemos ahora reescribir la descripción gráfica de monitor que hicimos en la tabla~\ref{esq:monitor_1}, incluyendo ahora la cola de entrada al monitor, tal y como vemos en la tabla~\ref{esq:monitor_2}
\begin{table}[H]
\centering
\begin{tabular}{|l|}
\hline
Cola del\\
monitor \\
\hline
Variables \\
permanentes \\
\hline
Procedimientos \\
exportados \\
\hline
Código de \\ 
inicialización \\
\hline
\end{tabular}
\caption{Esquema de un monitor incluyendo la cola de entrada.}
\label{esq:monitor_2}
\end{table}

\subsection{Operaciones de sincronización}
Las operaciones de sincronización entre los procesos de un programa concurrente se programan, como ya hemos visto, dentro de los procedimientos del monitor. Son instrucciones que permiten detener la ejecución de un procedimiento de un monitor y bloquear en una cola al proceso que ha hecho la llamada del procedimiento del monitor. Tenemos para realizar esta acción dos operaciones principales: \verb|wait| y \verb|signal|.\\

Sin embargo, las operaciones \verb|wait| y \verb|signal| que manejamos en monitores no se parecen a las que usábamos en los semáforos:
\begin{itemize}
    \item En los semáforos, la ejecución de \verb|wait| ofrecía la posibilidad de bloquear al proceso, ya que no lo hacía si el entero de dentro del semáforo era mayor estricto que $0$. Por contra, en monitores la llamada \verb|wait| siempre será bloqueante.
    \item Las operaciones \verb|wait| y \verb|signal| eran relativas a un semáforo: hacía falta usar un semáforo por cada razón que tuviéramos dentro de un programa concurrente para bloquear a uno o varios procesos (en el caso del productor/consumidor, usar dos semáforos). Sin embargo, con un solo monitor podremos bloquear procesos por tantas razones como queramos, usando un nuevo tipo de dato.
\end{itemize}

\subsubsection{Tipo de dato \texttt{cond}}
En los monitores, para poder usar las operaciones de \verb|wait| y \verb|signal|, será necesario utilizar una variable de tipo de dato condición, o \verb|cond|. 

Las variables tipo \verb|cond| se encuentran junto con las variables permanentes de un monitor. Estas no se inicializan a ningún valor.\\

En nuestro monitor, tendremos varias razones por las que queramos bloquear a los procesos concurrentes de nuestro programa por alguna determinada razón, hasta que se cumpla una condición determinada, a la que llamaremos \textit{condición de sincronización}. Por ejemplo, en el problema del productor/consumidor:
\begin{itemize}
    \item Queremos bloquear a cualquier productor que intente escribir si la estructura de datos intermedia que usamos está llena. Desbloquearemos a un proceso productor cuando se vacíe un hueco en dicha estructura.
    \item Además, queremos bloquear a cualquier consumidor que intente leer de la estructura de datos intermedia cuando esta esté vacía. Desbloquearemos a un consumidor cuando algún productor haya escrito algún dato.
\end{itemize}
Por cada razón o condición distinta por la que queramos bloquear a los procesos de un programa concurrente en relación a una misma variable compartida (para evitar estados inseguros), crearemos una variable de tipo \verb|cond|. Es decir, una variable por cada una de las razones por las que queramos que esperen los procesos. En el ejemplo del productor/consumidor, son necesarias únicamente dos variables de tipo \verb|cond|.\\

Las variables de tipo \verb|cond| admiten 4 métodos (aunque solo recomendamos usar los dos primeros):
\begin{description}
    \item [wait] Bloquea al proceso que ejecuta este método. Dicho proceso pasa a una cola asociada a la variable condición correspondiente con planificación FIFO\@.
    \item [signal] En caso de haber algún proceso bloqueado en la cola asociada a la variable condición correspondiente, lo desbloquea\footnote{Hemos de tener cuidado con esto, se explicará más adelante.}. Si esta cola está vacía, es equivalente a una operación nula\footnote{Esto es, equivalente a la instrucción \texttt{;}.}.
    \item [queue] Devuelve un booleano que indica (\verb|true|) si la cola asociada a la variable condición contiene al menos un proceso bloqueado.
    \item [signal\_all] Desbloquea de una sola vez a todos los procesos bloqueados en la cola asociada a la variable condición. El orden de dicha cola no se mantiene para realizar la petición de acceso al monitor, por lo que se produce competencia entre los procesos para entrar al monitor, incumpliendo la propiedad de equidad entre procesos. Depende de la semántica de las señales del lenguaje\footnote{Se explicará más adelante qué es esto.}. Se recomienda \textbf{no usarla}.
\end{description}
De esta forma, la representación gráfica final de un monitor es la que se muestra en la tabla~\ref{esq:monitor_3}:
\begin{table}[H]
\centering
\begin{tabular}{|l|}
\hline
Cola del\\
monitor \\
\hline
Variables \\
permanentes \\
\hline
Variables \\
condición y\\
colas de procesos\\
bloqueados\\
\hline
Procedimientos \\
exportados \\
\hline
Código de \\ 
inicialización \\
\hline
\end{tabular}
\caption{Esquema de un monitor incluyendo las variables condición.}
\label{esq:monitor_3}
\end{table}

\subsubsection{Semánticas de señales}
Como hemos comentado ya, los monitores solo permiten que un único proceso se encuentre ejecutando un procedimiento del mismo. En este caso, decíamos que el monitor está ocupado. En caso de que un proceso que estaba ejecutando el procedimiento ejecute un \verb|wait| (o salga del procedimiento), hay que dejar el monitor libre para dejar pasar a otro. Se trata de un momento muy delicado, ya que se pueden producir condiciones de carrera entre los procesos que quieran conseguir el monitor. Esta situación la hemos solucionado ya con la cola de entrada al monitor, ya que con la planificación FIFO, solo podrá entrar un único proceso al monitor.\\

Si ahora el proceso nuevo que ejecuta el monitor ejecuta un \verb|signal| sobre una variable condición (recordemos que tenemos al menos un proceso bloqueado), desbloqueará al primer proceso de su cola asociada, que estaba ejecutando código del monitor, por lo que ahora tenemos dos procesos ejecutando código del monitor: el proceso que señala y el señalado (el recién desbloqueado). Esta condición no puede darse, ya que los procedimientos de un monitor deben ejecutarse en exclusión mutua. Una solución a este problema es que el procedimiento que señala se bloquee en la cola de entrada al monitor\footnote{Veremos más soluciones.}, dejando paso al recién desbloqueado.

Esta solución plantea un problema, y es que si el proceso señalador se bloquea, puede que en algún momento deje al monitor libre, por lo que se meta un proceso de la cola de entrada al monitor, planteando nuevamente la situación en la que tenemos dos procesos en el monitor (el desbloqueado y el primero que estaba esperando entrar al monitor). Deberá haber un mecanismo que elija quién de los dos acaba finalmente con el monitor. En caso de que sea el proceso que estaba esperando en la cola de entrada, diremos que se produce un \emph{robo de señal}, donde el proceso recién desbloqueado debe irse al final de la cola de entrada al monitor.\\

Para solucionar este segundo problema, algunos lenguajes implementan una \emph{semántica desplazante}\footnote{Se trabajará más adelante sobre las semánticas de las señales.} en las señales: el proceso que ejecuta el \verb|signal| le pasa el monitor al proceso que recibió la señal (el primero en la cola de bloqueados de la variable condición correspondiente), sin liberar en ningún momento el monitor, de forma que el proceso señalado tiene prioridad. Se dice que la señal usada con la operación \verb|signal| tiene \emph{semántica desplazante}.\\

Cabe destacar que \textbf{no todos los lenguajes con monitores tienen señales con semántica desplazante}, por lo que en dichos lenguajes pueden sucederse robos de señales. En esta sección y en la siguiente, supondremos que estaremos trabajando siempre con señales \verb|signal| con semántica desplazante, de forma que el proceso señalador se bloqueará tras ejecutar un \verb|signal| y podemos pensar que es enviado a la cola de entrada al monitor. En una futura sección veremos los tipos de semánticas de señales que podemos encontrarnos (cada uno hará que el comportamiento de \verb|signal| sea distinto).\\

\noindent
Como comentario final a la descripción de un monitor y para motivar la siguiente sección:
\begin{itemize}
    \item Se presupone que el programador de monitores es un programador experto, de forma que el compilador en ningún momento se dedicará a comprobar si hemos programado de forma correcta un monitor o un procedimiento de él, más allá de la sintaxis del código.
    \item No deben programarse operaciones \verb|wait| indebidas ni omitirse operaciones \verb|signal| necesarias. Para comprobar esto, usaremos nuestro sistema de verificación formal.
\end{itemize}

\section{Verificación de programas con monitores}
En la verificación de los programas concurrentes que hemos manejado hasta ahora, hemos primero demostrado la corrección secuencial de cada proceso que forma parte de un programa secuencial, para luego demostrar la no interferencia entre los mismos. 

Sin embargo, ahora que introducimos los monitores, esto no podrá ser nunca más así, ya que un programador nunca puede conocer a priori la traza que genera un proceso que forma parte de un programa concurrente con monitores, ya que al ejecutar procedimientos de monitores, estos pueden quedar bloqueados y se ejecutarían en medio instrucciones de otros procesos que podrían alterar las variables compartidas del programa, falseando alguna precondición o poscondición del proceso bloqueado, por lo que tras desbloquearse, no podemos esperar nada de dicho proceso.\\

Es por tanto que ahora la estrategia a seguir en las demostraciones es mediante un Invariante de Monitor.

\subsection{Invariante de monitor}
\begin{definicion}[Invariante de Monitor]
    Un Invariante de Monitor (IM) es una relación entre las variables permanentes de un monitor que debe ser cierta en cualquier estado del programa concurrente, excepto cuando un proceso esté ejecutando código de un procedimiento del monitor.
\end{definicion}
De esta forma, un IM puede no ser cierto durante la ejecución de un procedimiento por parte de un proceso, pero este ha de cumplirse antes y después de la ejecución de dicho procedimiento.\\

Si conseguimos probar la existencia de un IM en un programa concurrente, entonces bastará con probar cada una de las secciones de código secuenciales entre llamadas a procedimientos del monitor. Para probar finalmente la corrección de los procesos, usaremos que los IM se mantienen antes y después de las llamadas a procedimientos, para conseguir probar finalmente la corrección de cada uno de los procesos. Si nuestro IM estaba relacionado con la solución al problema, como el acceso a variables compartidas estará controlado por los monitores, al final del programa todos los IMs demostrados se seguirán compliendo, por lo que tendremos probada la corrección de nuestro programa concurrente.

Es decir, primero demostraremos que por cada monitor que usamos se verifica un IM, y luego pasaremos a probar la corrección de cada proceso que interviene en el programa concurrente, usando para ello dichos IMs. Finamente, tendremos probado el programa concurrente.

\subsubsection{Esquema de demostración}
Suponiendo que hemos encontrado una relación matemática entre las variables permanentes de un monitor y queremos probar que se trata de un IM\footnote{A continuación, llamaremos a dicha condición \texttt{IM}, pese a no haber demostrado aún que se trate de verdad de un IM.}, lo primero será probar que $IM$ se cumple en el estado inicial del monitor, esto es, justo después de la inicialización de las variables permanentes, por lo que tendremos que probar que se verifica el \textbf{triple de inicialización de variables}:
\begin{equation*}
    \{V\}\ \text{código de inicialización}\ \{IM\}
\end{equation*}
Posteriormente, deberemos probar que $IM$ se mantiene antes y después de la llamada a cada procedimiento. Es decir, notando por $IN$ a las precondiciones que tenemos antes de la ejecución de un procedimiento y por $OUT$ a las poscondiciones que deseamos tener tras dicho procedimiento, debemos demostrar los \textbf{triples de procedimientos del monitor}, es decir, demostrar un triple
\begin{equation*}
    \{IM \land IN\}\ \text{procedimiento}\ \{IM \land OUT\}
\end{equation*}
por cada procedimiento que tenga nuestro monitor.

\subsubsection{Cuidado con las intereferencias}
Terminaremos de ver esto más adelante, pero es necesario darnos cuenta de un detalle, y es que si un procedimiento modifica el valor de alguna variable compartida que se usa en otro proceso, debemos demostrar la no interferencia entre dichas instrucciones. Ilustramos esto con el siguiente ejemplo.
\begin{ejemplo}
    Si tenemos un monitor llamado \verb|Buf| con un procedimientos \verb|retirar(x)|, de forma que modifica el valor del parámetro que le pasamos, ante el siguiente código (si \verb|x| es una variable compartida):
    \begin{minted}[escapeinside=\#\#]{pascal}
        cobegin y = x; #||# Buf.retirar(x); coend
    \end{minted}
    Tenemos que probar que al cambiar el valor de \verb|x| con el procedimiento \verb|retirar|, no hay interferencia con la instrucción de la izquierda. Es decir, tenemos que probar:
    \begin{align*}
        &NI(pre(y=x), Buf.retirar(x)) \\
        &NI(pos(y=x), Buf.retirar(x))
    \end{align*}
    Sin embargo, en caso de ejecutar el siguiente código:
    \begin{minted}[escapeinside=\#\#]{pascal}
        z = x;
        cobegin y=z; #||# Buf.retirar(x); coend
    \end{minted}
    No tendríamos que hacerlo, ya que el uso de variables disjuntas nos garantiza la no interferencia entre dichas instrucciones.
\end{ejemplo}

\subsection{Axiomas para operaciones de sincronización desplazantes}\label{sec:ax_desplazantes}
Sabemos ya demostrar toda la corrección de un programa secuencial que usa monitores, salvo por un detalle, y es que no sabemos nada sobre cómo demostrar los triples:
\begin{align*}
    \{P\}\ &c.wait();\ \{Q\} \\
    \{P\}\ &c.signal();\ \{Q\}
\end{align*}
para cualesquiera asertos $P$ y $Q$.

En esta subsección, trataremos de dar axiomas para la comprobación de dichos triples, razonándolos de forma intuitiva y mediante el uso de Invariantes de Monitores.\\

\subsubsection{Axioma de operación \texttt{wait}}
Comenzaremos primero con el triple $\{P\}\ c.wait();\ \{Q\}$. Para necesitar ejecutar una instrucción \verb|wait| en un procedimiento de un monitor, lo que sucede es que estamos cerca de un estado inseguro del programa (intuitivamente, que $IM$ está a punto de incumplirse), pero no llegamos a él, porque para ello ejecutamos esta operación, para impedir que el proceso ejecute una instrucción que falsee el $IM$. Por tanto, el proceso se bloquea, dejando libre el monitor, por lo que entra otro proceso a ejecutar otro procedimiento. 

Solo podremos desbloquear al proceso cuando nos alejemos de dicho estado inseguro, por lo que además de cumplirse el $IM$, deberá cumplirse una condición un tanto más estricta que el $IM$ (que nos indique que estamos lejos de aquel estado inseguro por el cual se bloqueó el proceso). Dicha condición recibe el nombre de \textit{condición de sincronización}, y la notaremos por $C$\footnote{Notemos que según hemos definido $C$, ha de verificarse que $C\rightarrow IM$.}. 

Resumiendo:
\begin{itemize}
    \item Antes de ejecutar la operación \verb|wait|, hemos de estar en un estado seguro del programa, por lo que ha de cumplirse el $IM$.
    \item Tras ejecutar la operación \verb|wait| (es decir, después de que el proceso haya sido desbloqueado), ha de cumplirse la condición de sincronización $C$.
\end{itemize}
Teniendo en cuenta que además se puede cumplir un invariante local al que llamamos $L$ (esto es, relaciones entre variables locales del procedimiento del monitor) antes y después\footnote{Gracias a que estamos en lenguajes reentrantes.} de dicha instrucción \verb|wait|.

De esta forma, acabamos de razonar de forma intuitiva el \textbf{Axioma de la operación wait}:
\begin{equation*}
    \{IM \land L\}\ c.wait();\ \{C \land L\}
\end{equation*}

\subsubsection{Axioma de operación \texttt{signal}}
Si nos disponemos a ejecutar una instrucción \verb|signal| en nuestro código, es porque el estado del programa se ha alejado de la condición insegura de la que hablábamos en la subsección anterior, que falsearía el valor de verdad de $IM$. Por tanto, el programa ha llegado a un punto en el que se cumple la condición de sincronización $C$, y ya puede desbloquear al proceso que anteriormente bloqueó. Tras su desbloqueo, este proceso podría ejecutar una instrucción que volviera a acercarnos a un estado inseguro, pero sin llegar a él (ya que $C$ era suficientemente restrictiva), por lo que como poscondición de la instrucción \verb|signal| no podremos garantizar $C$, sino solo podremos asegurar que se sigue cumpliendo $IM$.

Añadiendo la posibilidad de tener un invariante local $L$ y que si la cola de la variable condición está vacía, la operación \verb|signal| es una instrucción nula, llegamos al \textbf{Axioma de la operación signal}:
\begin{equation*}
    \{\lnot vacio(c) \land C \land L\}\ c.signal();\ \{IM \land L\}
\end{equation*}

o equivalentemente:
\begin{equation*}
    \{c.queue() \land C \land L\}\ c.signal();\ \{IM \land L\}
\end{equation*}
En caso de cumplirse que $c.queue() = false$, entonces negaría la precondición del triple, haciéndolo la regla cierta por un razonamiento por vacuidad.\\

\begin{observacion}
    Notemos que el axioma de la operación \verb|signal| funciona porque hemos suspuesto que \textbf{tenemos semántica desplazante}, ya que después de ejecutar \verb|signal| desbloqueamos al proceso que teníamos bloqueado, cediéndole el monitor, por lo que dicho proceso seguirá procesando su procedimiento. Cuando el proceso señalador vuelva al monitor, solo podremos garantizar que se cumple $IM$, ya que tanto el proceso señalado como cualquier otro que se introduzca en el monitor después del señalado (veremos más adelante si esto es posible), pueden cambiar la condición $C$, por lo que solo podemos esperar $IM$.
\end{observacion}~\\

Una vez vistos ya todos los axiomas sobre verificación de operaciones de sincronización de monitores, estamos listos para desmotrar la corrección de un $IM$. Lo haremos en el siguiente ejemplo.

\begin{ejemplo}
    En este ejemplo, queremos programar un monitor que simule el funcionamiento de un semáforo. Para ello, se nos ha ocurrido el siguiente código:
    \begin{minted}[escapeinside=\#\#]{pascal}
        Monitor Semaforo;
          var s : integer;
              c : cond;

          procedure P;
          begin
            if s=0 then
              c.wait;
            else
              null;
            end if
            s = s - 1;
          end

          procedure V;
          begin
            s = s + 1;
            c.signal;
          end

          begin   {código de inicialización}
            s = 0;
          end
    \end{minted}
    Donde hemos llamado \verb|P| a la función \verb|sem_wait| del semáforo y por \verb|V| a la función \verb|sem_signal|.\\

    Procedemos a realizar la demostración de que existe un Invariante de Monitor que se mantiene tras la inicialización de las variables permanentes de nuestro monitor y antes y después de cada procedimiento, con la finalidad de poder usar dicho IM en las demostraciones de cualquier programa concurrente que use el semáforo que acabamos de implementar mediante un monitor.
    \begin{proof}
        Tratamos de demostrar que este monitor tiene como IM el aserto
        \begin{equation*}
            IM \equiv \{s \geq 0\}
        \end{equation*}

        \begin{enumerate}
            \item Primero, tenemos que demostrar el triple de inicialización de variables:
                \begin{equation*}
                    \{V\}\ s=0;\ \{s\geq 0\}
                \end{equation*}
                Como el triple $\{V\}\ s=0;\ \{s=0\}$ es cierto por el axioma de asignación y tenemos que $\{s=0\}\rightarrow\{s\geq 0\}$, usando la primera regla de la consecuencia tenemos demostrado el triple.
            \item Posteriormente, demostraremos el triple de procedimiento del monitor para el procedimiento \verb|P|: $\{IM\}\ P\ \{IM\}$. Para ello, primero tendremos que probar el triple
                \begin{equation*}
                    \{IM\}\ \texttt{if\ }s=0\texttt{\ then\ }c.wait;\texttt{\ else\ }null;\texttt{\ end if}\ \{s>0\}
                \end{equation*}
                Luego usaremos la regla del \verb|if|, por lo que será suficiente con probar los triples:
                \begin{gather*}
                    \{IM \land s=0\}\ c.wait;\ \{s>0\} \\
                    \{IM \land s\neq0\}\ null;\ \{s>0\}
                \end{gather*}
                \begin{enumerate}
                    \item Comenzamos por el segundo, por ser más sencillo. Tenemos:
                        \begin{equation*}
                            {\{IM \land s\neq0\} \equiv \{s\geq 0 \land s\neq0\}\equiv \{s>0\}}
                        \end{equation*}
                        Por tanto, basta probar el triple ${\{s>0\}\ null;\ \{s>0\}}$, que es cierto por el axioma de la sentencia nula.
                    \item Para el primer triple, buscamos aplicar el axioma de la operación \verb|wait|, por lo que tenemos que buscar la condición de sincronización. Para ello, buscamos la precondición del \verb|signal| asociado a la misma variable condición, que se encuentra en el procedimiento \verb|V|. Para hallar la precondición de la instrucción \verb|c.signal|, tendremos que demostrar alguna instrucción de dicho procedimiento, con el fin de hallar la precondición.

                        Sobre el código de \verb|V|, vemos que antes de \verb|c.signal| se ejecuta una primera instrucción \verb|s=s+1;|. Suponemos que \verb|V| tiene como precondición $IM$, por lo que buscamos una poscondición para \verb|s=s+1;|:
                        \begin{equation*}
                            \{IM\}\equiv \{s\geq 0\}\ s=s+1;
                        \end{equation*}
                        Puede comprobarse con el axioma de asignación que la poscondición buscada es $\{s>0\}$. Por tanto, esta será la condición de sincronización de la variable condición \verb|c|:
                        \begin{equation*}
                            C \equiv \{s>0\}
                        \end{equation*}
                        Por tanto, por el axioma de la operación \verb|wait| usado con $L=\{V\}$, tenemos que el siguiente triple es cierto:
                        \begin{equation*}
                            \{IM\}\ c.wait;\ \{s>0\}
                        \end{equation*}
                        Como $\{IM \land s=0\}\to \{IM\}$, por la segunda regla de la consecuencia, tenemos demostrado el triple que buscábamos.
                \end{enumerate}
                Una vez demostrados los dos triples, tenemos probado el triple del \verb|if|, por lo que solo faltará probar el triple 
                \begin{equation*}
                    \{s>0\}\ s=s-1;\ \{IM\}
                \end{equation*}
                Para tener probado el triple del procedimiento \verb|P|.

                Como $\{IM\} \equiv \{s\geq 0\}$, basta aplicar el axioma de asignación, para obtener $\{s>0\}\ s=s-1;\ \{s\geq 0\}$.

                Aplicando finalmente la regla de composición sobre el triple del \verb|if| y este último triple, tenemos ya probado $\{IM\}\ P\ \{IM\}$.

            \item Finalmete, hemos de probar el triple $\{IM\}\ V\ \{IM\}$ para garantizar al fin que $IM$ es un IM\@. Para ello, hemos de probar el triple
                \begin{equation*}
                    \{IM\}\ s=s+1; c.signal;\ \{IM\}
                \end{equation*}
                Basta con probar los triples
                \begin{gather*}
                    \{IM\}\ s=s+1;\ \{s>0\} \\
                    \{s>0\}\ c.signal;\ \{IM\}
                \end{gather*}
                y aplicar la regla de composición. El primer triple ya lo demostramos en la demostración del triple del procedimiento \verb|P|, luego bastará probar el segundo, el cual es cierto gracias al axioma de la operación \verb|signal|.

                Acabamos de probar que $\{IM\}\ V\ \{IM\}$, que era el último procedimiento del monitor, luego $IM$ es un IM\@.
        \end{enumerate}
    \end{proof}
\end{ejemplo}

\begin{ejercicio*}
    Se pide demostrar que el siguiente monitor funciona como un semáforo de Habermann. Un semáforo de Habermann se trata de un semáforo normal que lleva la cuenta de:
    \begin{itemize}
        \item El número de llamadas realizadas a \verb|signal|, \verb|nv|.
        \item El número de llamadas realizadas a \verb|wait|, \verb|na|.
        \item El número de llamadas completadas a \verb|wait|, \verb|np|.
    \end{itemize}
    En este caso, llamaremos \verb|P| al procedimiento que simule la operación \verb|wait| y \verb|V| al procedimiento que simule la operación \verb|signal|.
    
    \begin{minted}[escapeinside=\#\#]{pascal}
        Monitor Semaforo;
          var na, np, nv : int;
              c : cond;

          procedure P;
          begin
            na = na + 1;
            if(na > nv) then c.wait();
            np = np + 1;
          end

          procedure V;
          begin
            nv = nv + 1;
            if(na > np) then c.signal();
          end

          begin
            na = 0; np = 0; nv = 0;
          end
    \end{minted}
    Buscamos un IM para preceder a la demostración del mismo. Notemos que las variables permanentes han de cumplir:
    \begin{itemize}
        \item $np \leq na$, ya que para completar una llamada \verb|P| hay que realizar una llamada.
        \item $np \leq nv$, ya que, como inicialmente $np = na = nv = 0$, para poder completar una llamada a \verb|P|, hay que previamente haber hecho una llamada a \verb|V|.
        \item $np \geq \min(na,nv)$, para no detener innecesariamente a los preocesos, cumpliendo la hipótesis de progreso finito.
    \end{itemize}
    De estas tres propiedades, deducimos que el invariante a usar es:
    \begin{equation*}
        \{IM\} \equiv \{np = \min(na,nv)\}
    \end{equation*}
    Pasemos ahora a la demostración del monitor:
    \begin{enumerate}
        \item En primer lugar, probaremos el triple de inicialización de variables:
            \begin{gather*}
                \{V\} \\
                na = 0;\ np=0;\ nv=0; \\
                \{na = 0 \land np = 0 \land nv = 0\} \rightarrow \{IM\}
            \end{gather*}
        \item Posteriormente, probaremos el triple del procedimiento \verb|P|:
            \begin{equation*}
                \{IM \land L\}\ P\ \{C \land L\}
            \end{equation*}

            para ello:
            \begin{gather*}
                \{IM\}\equiv \{np=\min(na,nv)\} \\
                na = na + 1;\\
                \{np = \min(na-1,nv)\} \\
                \texttt{if\ }(na > nv) \texttt{\ then\ } \\
                \{na > nv \land np = \min(na-1,nv)\} \rightarrow \{na-1\geq nv \land np = \min(na-1,nv)\} \rightarrow \\ 
                \rightarrow\{na > nv \land np = nv\} \rightarrow \{np = \min(na,nv)\} \\
                c.wait();
            \end{gather*}
            Donde en la precondición de la operación \verb|c.wait();| hemos necesitado comprobar que $IM$ se seguía compliendo.

            Y para poder seguir, hemos de buscar la precondición de la operación \verb|c.signal();| asociada a la misma variable condición, con lo que comenzamos a demostrar el procedimiento \verb|V|:
            \begin{gather*}
                \{IM\} \equiv \{np = \min(na,nv)\} \\
                nv = nv + 1; \\
                \{np = \min(na,nv-1)\} \\
                \texttt{if\ } (na > np) \texttt{\ then\ } \\
                \{na > np \land np = \min(na,nv-1)\} \rightarrow \{na > np \land np = nv-1\}
            \end{gather*}
            Luego tenemos ya la poscondición de la operación \verb|c.wait();|, con lo que podemos volver por donde íbamos:
            \begin{gather*}
                \{IM\} \equiv \{np = \min(na,nv)\} \\
                na = na + 1;\\
                \{np = \min(na-1,nv)\} \\
                \texttt{if\ }(na > nv) \texttt{\ then\ } \\
                \{na > nv \land np = \min(na-1,nv)\} \rightarrow \{na-1\geq nv \land np = \min(na-1,nv)\} \rightarrow \\ 
                \rightarrow\{np = nv\} \rightarrow \{np = \min(na,nv)\} \\
                c.wait(); \\
                \{na > np \land np = nv - 1\} \\
                \texttt{else do} \\
                \{na \leq nv \land np = \min(na-1,nv)\} \rightarrow \{na-1 < nv \land np = \min(na-1,nv)\} \rightarrow \\
                \rightarrow \{na-1 < nv \land np = na-1\} \rightarrow \{np < nv \land np = na-1\} \\
                \texttt{null;} \\
                \{np < nv \land np = na -1\} \\
                \texttt{endif}
            \end{gather*}
            Como las poscondiciones de cada bloque del \verb|if| son distintas, debemos relajarlas para buscar dos poscondiciones iguales, para poder aplicar la regla del \verb|if|. Notemos que:
            \begin{gather*}
                \{na > np \land np = nv-1\} \rightarrow \{na > nv-1 \land np = nv-1\} \rightarrow \\ \rightarrow \{na \geq nv \land np = nv - 1\} \\
                \{np < nv \land np = na-1\} \rightarrow \{na-1 < nv \land np = na-1\} \rightarrow\\\rightarrow \{na \leq nv \land np = na-1\}
            \end{gather*}
            Podemos por tanto, relajar ambas poscondiciones a la poscondición
            \begin{equation*}
                \{np+1 = \min(na,nv)\}
            \end{equation*}
            Con lo que ahora sí podemos aplicar la regla del \verb|if|, con lo que podemos finalizar la demostración del triple del procedimiento \verb|P|:
            \begin{gather*}
                \{IM\} \equiv \{np = \min(na,nv)\} \\
                na = na + 1;\\
                \{np = \min(na-1,nv)\} \\
                \texttt{if\ }(na > nv) \texttt{\ then\ } \\
                \{np = \min(na,nv)\} \\
                c.wait(); \\
                \{np+1 = \min(na,nv)\} \\
                \texttt{else do} \\
                \{np+1 = \min(na,nv)\} \\
                \texttt{null;} \\
                \{np+1 = \min(na,nv)\} \\
                \texttt{endif} \\
                \{np+1 = \min(na,nv)\} \\
                np = np + 1; \\
                \{np = \min(na,nv)\} \equiv \{IM\}
            \end{gather*}
        \item Para probar ahora el triple del procedimiento \verb|V|:
            \begin{gather*}
                \{IM\} \equiv \{np = \min(na,nv)\} \\
                nv = nv + 1; \\
                \{np = \min(na,nv-1)\} \\
                \texttt{if\ } (na > np) \texttt{\ then\ } \\
                \{na > np \land np = \min(na,nv-1)\} \rightarrow \{na > np \land np = nv-1\}\\
                c.signal(); \\
                \{np = \min(na,nv)\} \\
                \texttt{else do} \\
                \{na \leq np \land np = \min(na,nv-1)\}\rightarrow\{np = na\} \rightarrow \{np = \min(na,nv)\} \\
                \texttt{null;} \\
                \{np = \min(na,nv)\} \\
                \texttt{endif} \\
                \{np = \min(na,nv)\} \equiv \{IM\}
            \end{gather*}
    \end{enumerate}
    Con lo que tenemos provado que $IM$ es un IM\@.
\end{ejercicio*}

\subsection{Regla de concurrencia para programas con monitores}
Consideramos un programa concurrente en el que tenemos $n$ procesos ejecutándose que podemos representar como triples de Hoare ciertos $\{P_i\}\ S_i\ \{Q_i\}$ con \newline $i \in \{1, \ldots, n\}$ de forma que ninguna variable en $P_i$ o en $Q_i$ es modificada por ningún $S_j$ con $i\neq j$. Si en dicho código tenemos $m$ monitores de forma que para cada uno hemos conseguido probar un IM, $IM_k$ con $1 \leq k \leq m$, entonces podemos aplicar la \textbf{regla de concurrencia para programas con monitores}:
\begin{equation*}
    \dfrac{\{P_i\}\ S_i\ \{Q_i\} \quad 1 \leq i \leq n}{\begin{array}{c}
            \{MI_1 \land \ldots \land MI_m \land P_1 \land \ldots\land P_n\} \\
            cobegin\ S_1\ ||\ S_2\ ||\ \ldots\ ||\ S_n\ coend \\
            \{MI_1 \land \ldots \land MI_m \land Q_1 \land  \ldots\land Q_n\} 
        \end{array}}
\end{equation*}
Obteniendo así la verificación de nuestro programa concurrente.

\section{Patrones de uso de un monitor}
Al programar programas concurrentes que nos resuelvan un problema, encontramos muchas veces ciertos pequeños problemas a resolver que se repiten a menudo. En esta sección, destacamos tres de estos problemas, describiéndolos, planteando una solución mediante un monitor a utilizar en ellos y demostrando que el monitor realiza el funcionamiento esperado.

\subsection{Espera única}
Puede suceder que en un programa concurrente queramos que un proceso espere a que otro proceso haya realizado una cierta acción para seguir con su ejecución. Llamaremos al proceso que tiene que esperar al otro \textit{consumidor} y a dicho otro \textit{productor}.\\

El problema se resuelve de forma muy sencilla, con una variable compartida que indique si el productor ha realizado ya su acción por la que el consumidor debe esperar o si no lo ha hecho todavía. Cuando el consumidor se acerque a la zona en la que debe esperar al productor, consultará la variable compartida y en caso de que esta indique que no se ha realizado la acción, bloquearemos al proceso. Si se ha realizado la acción, no haremos nada.

Además, cuando el productor haya realizado la acción que ha de realizar, cambiaremos el valor de dicha variable compartida, indicando que ya se ha realizado la acción. En caso de que el consumidor se haya bloqueado antes de realizar la acción, lo desbloquearemos.\\

Observemos que estamos haciendo uso de una variable compartida, que es modificada por un proceso. Debemos por tanto acceder a ella en exclusión mutua. Esto es garantizado por el uso del monitor.

\subsubsection{Monitor a usar}
El monitor que usaremos tendrá dos procedimentos exportables, uno llamado \verb|esperar| que será el que use el proceso consumidor, y otro llamado \verb|notificar|, que será el que use el proceso productor para avisar al consumidor de que ya ha realizado la acción.

De esta forma, podemos ver el código monitor en pseudo código:
\begin{minted}[escapeinside=\#\#]{pascal}
    monitor EU;
    var terminado : boolean;  {si terminado = true, se ha realizado la acción}

        {variable auxiliar para la demostracion}
        autorizado : boolean; {si autorizado = true, consumidor puede ejecutarse}

        c : cond;

       procedure esperar(); begin
          if (not terminado) then
             c.wait();
          autorizado = true;
       end

       procedure notificar(); begin
          terminado = true; 
          c.signal();
       end

       begin {codigo de inicializacion}
          terminado = false; autorizado = false;
       end
\end{minted}
% // TODO: Capel incorpora un procedimiento terminar que hace lo mismo que el codigo de inicializacion, no se por qué.

Para su demostración, usaremos el invariante
\begin{equation*}
    \{IM\} \equiv \{terminado = false \Longrightarrow autorizado = false\}
\end{equation*}
La demostración se deja como ejercicio al lector.
% // TODO: Hacer la demostración y quitar lo de que se deja como ejercicio

\subsection{Exclusión mutua}
Es muy habitual que en programas concurrentes tengamos una o varias regiones de código que queramos que se ejecuten en exclusión mutua, llamadas secciones críticas. Es decir, mientras que un proceso se encuentre ejecutando una sección crítica, ningún otro proceso podrá estar ejecutando a la vez la misma sección crítica\footnote{Podemos tener dos secciones críticas con distintos códigos pero que sean referidas al acceso para la misma variable compartida. En dicho caso, pensamos que las secciones críticas son iguales, ya que solo puede haber un proceso que ejecute una u otra a la vez.}.\\

Usualmente, querremos tener exclusiones mutuas cuando varios procesos de un programa hagan uso de un recurso compartido, tal como una variable compartida, una salida a un fichero o imprimir información en un entorno gráfico.

\subsubsection{Monitor a usar}
El monitor que resuelve el problema de la exclusión mutua tendrá dos procedimientos exportables, \verb|entrar|, que se ejecutará antes de cualquier sección crítica, y \verb|salir|, que se ejecutará al final de cada sección crítica.

De esta forma, tenemos el monitor:
\begin{minted}[escapeinside=\#\#]{pascal}
    monitor EM;
    var ocupada : boolean;  {ocuapda = true si hay un proceso en seccion critica}
        cola : cond;

       procedure entrar(); begin
          if ocupada then
             cola.wait();
          ocupada = true;
       end

       procedure salir(); begin
          ocupada = false;
          cola.signal();
       end

       begin
          ocupada = false;
       end
\end{minted}
Para demostrarlo, primero definiremos la variable $num_{sc}$ como el número de procesos que se encuentran ejecutando la sección crítica. Una vez definido, el invariante a usar será
\begin{equation*}
    \{IM\} \equiv \{(ocupada = false \Longleftrightarrow num_{sc} = 0) \land 0\leq num_{sc} \leq 1\}
\end{equation*}
La demostración se deja como ejercicio al lector.
% // TODO: Hacer la demostración y quitar lo de que se deja como ejercicio

\subsection{Productores/Consumidores}
Volvemos otra vez al paradigma del productor/consumidor, el cual hemos explicado varias veces ya en este documento. Ahora, admitiremos la existencia de varios procesos productores que querrán generar datos y escribirlos en una variable compartida, así como varios consumidores, que querrán leer dicha variable compartida y realizar los cálculos pertinentes.

\subsubsection{Monitor a usar}
Usaremos un monitor con dos procedimientos: \verb|escribir|, que permitirá escribir en la variable compartida el valor indicado, y \verb|leer|, que permitirá leer el valor de la variable compartida. La ventaja de usar un monitor es que los códigos de sincronización solo los tenemos que realizar en el monitor, dejando limpios los códigos de los procesos.

Como tenemos dos condiciones de sincronización, bloquear a productores que quieran escribir en una variable que no se ha leído, o bloquear a consumidores que quieran leer de una variable cuyo valor ya ha sido leído, necesitaremos dos variables de tipo \verb|cond|:
\begin{minted}[escapeinside=\#\#]{pascal}
    monitor PC;
    var valor : integer;  {variable compartida a usar}
        pendiente : boolean;  {si pendiente = true, valor escrito y no leido}
        cola_prod, cola_cons : cond;

       procedure escribir(v : integer); begin
          if pendiente then
             cola_prod.wait();
          valor = v;
          pendiente = true;
          cola_cons.signal();
       end

       procedure leer() : integer; begin
          if (not pendiente) then
             cola_cons.wait();
          result = valor;
          pendiente = false;
          cola_prod.signal();
       end

       begin
          pendiente = false;
       end
\end{minted}
Para la demostración, debemos definir primero:
\begin{align*}
    E &= \text{\ número\ de\ llamadas\ a\ escribir\ \textbf{completadas}.} \\
    L &= \text{\ número\ de\ llamadas\ a\ leer\ \textbf{completadas}.} \\
\end{align*}
De esta forma, podemos definir el invariante
\begin{gather*}
    \{IM\} \equiv 
    \left\{E-L = \left\{\begin{array}{rl}
                0 & \text{\ si\ } pendiente = false \\
                1 & \text{\ si\ } pendiente = true \\
    \end{array}\right\}\right\} \equiv  \\
    \equiv \{(pendiente = false \land E-L=0) \lor (pendiente = true \land E-L=1)\}
\end{gather*}
La demostración se deja como ejercicio para el lector.
% // TODO: Hacer la demostración y quitar lo de que se deja como ejercicio

Asímismo, notemos que la demostración es similar a la del monitor para exclusión mutua, teniendo en cuenta que:
\begin{align*}
    E - L &= num_{sc} \\
    pendiente &= \lnot libre
\end{align*}

\begin{ejercicio*}
    Si ahora los productores no escriben sobre una misma variable compartida sino sobre un buffer (por ejemplo, un array con planificación FIFO), plantear un monitor que solucione el problema de los productores/consumidores, así como desmostrar que dicho monitor funciona correctamente.\newline
    (\textbf{Pista:} para la demostración, sustituir en el IM ``1'' por el tamaño del buffer)
\end{ejercicio*}

\section{Semánticas de señales}
Como comentamos ya al inicio del Capítulo, las operaciones \verb|signal| de las variables condición son operaciones delicadas, ya que son ejecutadas por un proceso que se encuentra ejecutando algún procedimiento del monitor y que lo que hacen es desbloquear a algún proceso se que se encontraba ejecutando código del monitor, por lo que (si no hacemos nada), tendremos dos procesos distintos ejecutando a la vez procedimientos (podría ser el mismo procedimiento) de un mismo monitor, algo que no puede darse.\\

Para solucionar el problema que nos plantea la operación \verb|signal|, plantearemos distintas \textit{semánticas de señales} \verb|signal|. Esto es, plantearemos varios paradigmas en los que la señal \verb|signal| tendrá un significado u otro, de forma que su finalidad sea siempre \textit{sacar al primero proceso de la cola de la variable condición} en cuestión y ponerlo en otro sitio que permita que dicho proceso entre al monitor en algún futuro próximo, garantizando la propiedad de vivacidad de que dicho proceso en algún momento volverá a entrar al monitor.\\

Además, clasificaremos los distintos tipos de semánticas de señales que veremos en \textbf{no desplazantes} y \textbf{desplazantes}.

\begin{definicion}
    Diremos que \textbf{una semántica de señal es desplazante} si, siempre que un proceso ejecute una operación \verb|signal| sobre una variable condición \verb|c|, en dicho instante cederá el monitor al primer proceso de la cola de bloqueados de la variable condición \verb|c| sin que el monitor quede libre en ningún momento\footnote{Evitando que entre al monitor cualquier otro proceos de la cola de entrada al monitor.}. Además, debe garantizarse la propiedad de vivacidad de que el proceso señalador volverá a entrar al monitor en algún momento.
\end{definicion}

Veremos ahora todas las posibles semánticas de señales que podemos encontrarnos, entendiendo que el proceso \textit{señalador} es aquel que ejecuta la operación \verb|signal| sobre una variable condición; y que el proceso \textit{señalado} es aquel que estaba primero en la cola de bloqueados de la misma variable condición.

\subsection{Señalar y Continuar (SC)}
El proceso señalador no se bloquea tras ejecutar \verb|signal|, sino que sigue con la ejecución del procedimiento en cuestión. El proceso señalado se bloquea hasta que se pueda adquirir de nuevo el monitor.\\

Como podemos ver, se trata de una semántica no desplazante, ya que el proceso señalador no se bloquea tras ejecutar la instrucción \verb|signal|. En relación al proceso señalado, se produce una \textit{competición} contra el resto de procesos que esperan en la cola de entrada al monitor. Esta ``competición'' que hemos mencionado depende de la implementación que se haga, destacando dos posibilidades:
\begin{enumerate}
    \item El proceso señalado pasa al final de la cola de entrada al monitor.
    \item Después de que el proceso señalador deje libre el monitor (ya sea porque termina el procedimiento o se bloquea debido a una instrucción \verb|wait|), se desbloquea al proceso señalado y al primero de la cola de entrada al monitor, se sucede una condición de carrera entre ambos y el vencedor (el que primero llega al monitor), acaba ejecutándolo. El perdedor acaba al final de la cola de entrada al monitor.

        Como podemos entender, la primera opción es mejor, ya que en la segunda puede suceder que un proceso pierda varias veces la competición por el monitor, incumpliendo la condición de equidad entre procesos; además de que no sabemos qué sucede tras desbloquear a dicho proceso, debido a que la competición es transparente para el programador.
\end{enumerate}
De esta forma, podemos representar la vida de un proceso que forma parte de un programa concurrente que usa monitores con semántica de señales de señalar y continuar con el diagrama de la Figura~\ref{fig:sc}.

% // TODO: Ver si fijar o no estos 4 grafos
\begin{figure}%[H]
\centering
\begin{tikzpicture}[node/.style={rectangle, draw, minimum size=1cm, align=center},
                    edge/.style={-stealth}]
    \node[node] (A) {Procedimiento\\terminado};
    \node[node, left =of A] (B) {\begin{tabular}{c}Ejecutando código\\del proceso\\(fuera del monitor)\end{tabular}};
    \node[node, left =of B] (C) {\begin{tabular}{c}Llamando a\\procedimiento\\del monitor\end{tabular}};
    \node[node, below =of C] (D) {\begin{tabular}{c}Esperando en\\cola de entrada\\al monitor\end{tabular}};
    \node[node, right =of D] (E) {\begin{tabular}{c}Monitor\\libre\end{tabular}};
    \node[node, right =of E, xshift=.6cm] (F) {\begin{tabular}{c}Ejecutando procedimiento\\(dentro del monitor)\end{tabular}};

    \node[node, below right =of F] (G) {\texttt{c.signal()}};
    \node[node, below =of F] (H) {\texttt{c.wait()}};
    \node[node, left =of H] (I) {\begin{tabular}{c}Bloqueado en una\\cola condición $c$\end{tabular}};
    \node[node, left =of I, xshift=-1.3cm] (J) {Señalado};
    
    \draw[edge] (A) -- (B);
    \draw[edge] (B) -- (C);
    \draw[edge] (C) -- (D);
    \draw[edge] (D) -- (E);
    \draw[edge] (E) -- (F);
    \draw[edge] (F) -- (A);

    % \draw[edge] (F) to[out=0, in=100] (G);
    % \draw[edge] (G) to[out=180, in=-50] (F);

    \draw[edge] (F) to[out=0, in=90] (G);
    \draw[edge] (G) to[out=180, in=-90] ([xshift=1.5cm]F.south);

    \draw[edge] (F) -- (H);
    \draw[edge] (H) -- (I);
    \draw[edge] (I) -- (J);
    \draw[edge] (J) -- (D);
\end{tikzpicture}
\caption{Vida de un proceso en un programa con monitores con semántica SC.}
\label{fig:sc}
\end{figure}

\subsubsection{A tener en cuenta para bloquear}
Como se trata de una semántica no desplazante, debemos tener cuidado con un detalle:

En semántica desplazante, bloqueamos a un proceso porque se acerca a un estado inseguro del programa. Cuando nos alejemos de dicho estado, realizaremos la instrucción \verb|signal|, por lo que el proceso que bloqueamos se podrá ejecutar al habernos alejado del estado inseguro del programa.\\

Por otra parte, ahora también desbloquearemos al proceso por alejarnos de dicho estado inseguro, pero el proceso señalado no se ejecutará tras esto, por lo que puede que se ejecuten procesos en medio que vuelvan a acercarse a este estado inseguro, luego tendremos que volver a comprobar si estamos cerca o no de dicho estado inseguro.\\

\noindent
Por tanto, en semánticas se señales no desplazantes, en vez de bloquear a procesos con:
\begin{minted}[escapeinside=\#\#]{pascal}
    if (cerca de estado inseguro) then
       c.wait();
\end{minted}
Tendremos que plantear el código
\begin{minted}[escapeinside=\#\#]{pascal}
    while (cerca de estado inseguro) do begin
       c.wait();
    end
\end{minted}
Ya que cuando el proceso señalado vuelva no sabremos si la condición por la que lo bloqueamos es cierta o no. Puede suceder que bloqueemos a un proceso por acercarse a un estado inseguro, que lo desbloqueemos cuando el estado del programa se aleje de dicho estado, que se ejecuten procesos que se acerquen a dicho estado y que cuando el proceso que fue señalado entre al monitor, se encuentre cerca del mismo estado inseguro por el que tuvo que bloquearse, teniendo que hacerlo nuevamente.

\subsection{Señalar y Salir (SS)}
El proceso señalador \textbf{finaliza} la ejecución de su procedimiento tras la ejecución de una instrucción \verb|signal|. El proceso señalado se desbloquea posteriormente y en todo este tiempo el monitor no queda libre en ningún momento. Se trata, por tanto, de una semántica desplazante, ya que el proceso señalador cede el monitor al proceso señalado.\\

Notemos que en señalar y salir, el proceso señalador \textbf{finaliza} la ejecución de su procedimiento. Es decir, cualquier instrucción que haya tras una operación \verb|signal| no se ejecutará nunca. Obliga por tanto a una disciplina de programación en la que si queremos realizar un \verb|signal| dentro de un procedimiento, esta instrucción debe ser la última dentro del procedimiento, para poder ejecutar antes todas las instrucciones que deseemos. 

Si recordamos el ejemplo de la Sección~\ref{sec:concepto_de_monitor}, observamos que en este las instrucciones \verb|signal| se encuentran al final de los procedimientos de forma natural. Por tanto, la semántica SS funciona bien para este monitor.\\

Al igual que hicimos para SC, podemos representar la vida de un proceso que forma parte de un programa concurrente que usa monitores con semántica de señales de señalar y salir con el diagrama de la Figura~\ref{fig:ss}.
\begin{figure}%[H]
\centering
\begin{tikzpicture}[node/.style={rectangle, draw, minimum size=1cm, align=center},
                    edge/.style={-stealth}]
    \node[node] (A) {Procedimiento\\terminado};
    \node[node, left =of A] (B) {\begin{tabular}{c}Ejecutando código\\del proceso\\(fuera del monitor)\end{tabular}};
    \node[node, left =of B] (C) {\begin{tabular}{c}Llamando a\\procedimiento\\del monitor\end{tabular}};
    \node[node, below =of C, yshift=-1.5cm] (D) {\begin{tabular}{c}Esperando en\\cola de entrada\\al monitor\end{tabular}};
    \node[node, right =of D] (E) {\begin{tabular}{c}Monitor\\libre\end{tabular}};
    \node[node, right =of E, xshift=.6cm] (F) {\begin{tabular}{c}Ejecutando procedimiento\\(dentro del monitor)\end{tabular}};

    \node[node, below =of B] (G) {\texttt{c.signal()}};
    \node[node, below right =of F] (H) {Señalado};
    \node[node, below left =of F] (I) {\texttt{c.wait()}};
    \node[node, below =of F, yshift=.1cm] (J) {\begin{tabular}{c}Bloqueado en una\\cola condición $c$\end{tabular}};

    \draw[edge] (A) -- (B);
    \draw[edge] (B) -- (C);
    \draw[edge] (C) -- (D);
    \draw[edge] (D) -- (E);
    \draw[edge] (E) -- (F);
    \draw[edge] (F) -- (A);

    \draw[edge] (F) -- (G);
    \draw[edge] (G) -- (B);

    \draw[edge] (F) -- (I);
    \draw[edge] (I) -- (J);

    \draw[edge] (H) -- (F);
    \draw[edge] (J) -- (H);
\end{tikzpicture}
\caption{Vida de un proceso en un programa con monitores con semántica SS.}
\label{fig:ss}
\end{figure}

\subsection{Señalar y Esperar (SE)}
El proceso señalador se bloquea al final de la cola de entrada al monitor y cede el monitor al proceso señalado. El monitor no queda libre en ningún momento. Se trata de otro tipo de señal con semántica desplazante.\\

Puede considerarse una semántica \textit{injusta}, ya que cada vez que un proceso ejecuta la operación \verb|signal|, debe irse al final de la cola de entrada al monitor, teniendo que esperar entre todos los procesos nuevamente (el proceso probablemente tuvo ya que esperar para entrar al monitor para ejecutar el procedimiento que contenía la operación \verb|signal|).

Podemos representar la vida de un proceso que forma parte de un programa concurrente que usa monitores con semántica de señales de señalar y esperar con el diagrama de la Figura~\ref{fig:se}.

\begin{figure}%[H]
\centering
\begin{tikzpicture}[node/.style={rectangle, draw, minimum size=1cm, align=center},
                    edge/.style={-stealth}]
    \node[node] (A) {Procedimiento\\terminado};
    \node[node, left =of A] (B) {\begin{tabular}{c}Ejecutando código\\del proceso\\(fuera del monitor)\end{tabular}};
    \node[node, left =of B] (C) {\begin{tabular}{c}Llamando a\\procedimiento\\del monitor\end{tabular}};
    \node[node, below =of C] (D) {\begin{tabular}{c}Esperando en\\cola de entrada\\al monitor\end{tabular}};
    \node[node, right =of D] (E) {\begin{tabular}{c}Monitor\\libre\end{tabular}};
    \node[node, right =of E, xshift=.6cm] (F) {\begin{tabular}{c}Ejecutando procedimiento\\(dentro del monitor)\end{tabular}};

    \node[node, below =of F, yshift=-1.5cm] (J) {\begin{tabular}{c}Bloqueado en una\\cola condición $c$\end{tabular}};
    \node[node, right =of J] (H) {Señalado};
    \node[node, left =of J] (I) {\texttt{c.wait()}};

    \node[node, below =of E, yshift=.3cm] (G) {\texttt{c.signal()}};

    \draw[edge] (A) -- (B);
    \draw[edge] (B) -- (C);
    \draw[edge] (C) -- (D);
    \draw[edge] (D) -- (E);
    \draw[edge] (E) -- (F);
    \draw[edge] (F) -- (A);

    \draw[edge] ([xshift=-2cm]F.south) to[out=-90, in=0] (G);
    \draw[edge] (G) to[out=180, in=-90] (D);

    \draw[edge] (F) -- (I);
    \draw[edge] (I) -- (J);

    \draw[edge] (H) -- (F);
    \draw[edge] (J) -- (H);
\end{tikzpicture}
\caption{Vida de un proceso en un programa con monitores con semántica SE.}
\label{fig:se}
\end{figure}

\subsection{Señalar y Espera Urgente (SU)}
El proceso señalador se bloquea en una nueva cola del monitor llamada \textit{cola de procesos urgentes}, cediendo el monitor al proceso señalado. El monitor no queda libre en ningún momento, por lo que se trata de una señal con semántica desplazante.\\

La nueva cola de procesos urgentes actúa como una nueva cola de entrada al monitor, pero con mayor preferencia que la cola de entrada que ya teníamos. De esta forma, es similar a la semántica de señalar y esperar pero sin ser tan \textit{injusta}, ya que da preferencia a los procesos que se bloquearon tras ejecutar un \verb|signal| y deja en segunda posición a los proceos que esperan para ejecutar procedimientos del monitor.\\

Esta semántica modifica la última representación gráfica que teníamos de un monitor en la Tabla~\ref{esq:monitor_3}, dejándola finalmente en la que observamos en la Tabla~\ref{esq:monitor_4}.

\begin{table}[H]
\centering
\begin{tabular}{|l|}
\hline
Cola del\\
monitor \\
\hline
Cola de\\
urgentes \\
\hline
Variables \\
permanentes \\
\hline
Variables \\
condición y\\
colas de procesos\\
bloqueados\\
\hline
Procedimientos \\
exportados \\
\hline
Código de \\ 
inicialización \\
\hline
\end{tabular}
\caption{Esquema de un monitor incluyendo la cola de urgentes.}
\label{esq:monitor_4}
\end{table}

Podemos representar la vida de un proceso que forma parte de un programa concurrente que usa monitores con semántica de señalar y espera urgente con el diagrama de la Figura~\ref{fig:su}.\\

Finalmente, recordemos que con semántica SU, si la cola de una variable condición está vacía y realizamos un \verb|signal| sobre la misma, será equivalente a una operación nula, con lo que en este caso el proceso señalador no se irá a la cola de bloqueados.

\begin{figure}%[H]
\centering
\begin{tikzpicture}[node/.style={rectangle, draw, minimum size=1cm, align=center},
                    edge/.style={-stealth}]
    \node[node] (A) {Procedimiento\\terminado};
    \node[node, left =of A] (B) {\begin{tabular}{c}Ejecutando código\\del proceso\\(fuera del monitor)\end{tabular}};
    \node[node, left =of B] (C) {\begin{tabular}{c}Llamando a\\procedimiento\\del monitor\end{tabular}};
    \node[node, below =of C] (D) {\begin{tabular}{c}Esperando en\\cola de entrada\\al monitor\end{tabular}};
    \node[node, right =of D] (E) {\begin{tabular}{c}Monitor\\libre\end{tabular}};
    \node[node, right =of E, xshift=.6cm] (F) {\begin{tabular}{c}Ejecutando procedimiento\\(dentro del monitor)\end{tabular}};

    \node[node, below =of D] (G) {\texttt{c.wait()}};
    \node[node, below =of G] (H) {\begin{tabular}{c}Bloqueado en una\\cola condición $c$\end{tabular}};
    \node[node, right =of H] (I) {Señalado};

    \node[node, below =of F, xshift=4cm] (J) {\begin{tabular}{c}Monitor\\libre\end{tabular}};
    \node[node, below =of J] (K) {\begin{tabular}{c}Esperando en\\cola de urgentes\end{tabular}};
    \node[node, left =of K] (L) {\texttt{c.signal()}};

    \draw[edge] (A) -- (B);
    \draw[edge] (B) -- (C);
    \draw[edge] (C) -- (D);
    \draw[edge] (D) -- (E);
    \draw[edge] (E) -- (F);
    \draw[edge] (F) -- (A);

    \draw[edge] ([xshift=-2cm]F.south) to[out=-90, in=0] (G);
    \draw[edge] (G) -- (H);
    \draw[edge] (H) -- (I);
    \draw[edge] (I) to[out=0, in=-90] ([xshift=-1cm]F.south);

    % \draw[edge] (F) to[out=0, in=90] (J);
    % \draw[edge] (J) -- (K);
    % \draw[edge] (K) -- (L);
    % \draw[edge] (L) -- (F);

    \draw[edge] (J) to[out=90, in=0] (F);
    \draw[edge] (K) -- (J);
    \draw[edge] (L) -- (K);
    \draw[edge] (F) -- (L);
\end{tikzpicture}
\caption{Vida de un proceso en un programa con monitores con semántica SU.}
\label{fig:su}
\end{figure}

\subsection{Comparativa}\label{cap:comp}
Antes de comparar las semánticas de señales que ya hemos descrito, destacamos un último tipo de semántica de señales, las \textit{señales automáticas} (SA). En estas, el programador programa las operaciones \verb|wait| y es el compilador quien analiza los códigos programados y decide cuándo desbloquear a los procesos, por lo que se trata de unas señales implícitas (ya que en ningún moemnto se especifica cuándo hacer una instrucción \verb|signal|). 

A pesar de que este tipo de semántica aparece en la bibliografía, no aparece en ningún lenguaje de programación conocido, por lo que no tendrá relevancia en la asignatura ni en la comparativa que ahora realizaremos.\\

En primer lugar, hemos de destacar que cualquiera de las primeras 4 semánticas vistas es capaz de resolver cualquier problema, por lo que nos guiaremos por la sencillez de uso de cada una y la eficiencia que cada una nos aporta a los programas concurrentes para elegir una u otra semántica.
\begin{itemize}
    \item En cuato a sencillez de uso, SC, SE y SU presentan la misma facilidad, mientras que SS nos condiciona a un paradigma de programación en el que las instrucciones \verb|signal| sean las últimas que se ejecuten en los procedimientos de un monitor que contenga alguna operación \verb|signal|.
    \item En cuanto a eficiencia:
        \begin{itemize}
            \item Las semánticas SE y SU resultan ineficientes cuando las instrucciones \verb|signal| se encuentran al final de los procedimientos (ya que habrá procesos que ejecutan un \verb|signal|, se bloquearán, esperarán en la cola correspondiente de entrada al monitor, y cuando por fin tengan acceso a él no realizarán nada, sino que simplemente terminarán la ejecución del procedimiento), por lo que se saturará la cola de entrada al monitor y se realizarán cambios de contexto de formas innecesarias.
            \item La semántica SC es poco eficiente, al tener que obligarnos a usar un bucle de comprobación para comprobar si estamos cerca de una situación insegura cada vez que entramos al monitor, como consecuencia de no tener semántica desplazante, pudiendo producirse robos de señal.
        \end{itemize}
        Por tanto, si no nos resulta incómodo tener que colocar todas las instrucciones \verb|signal| como última instrucción de los procedimientos en los que aparecen, SS será probablemente la semántica de señales más eficientes.
\end{itemize}

Para concluir la comparativa de las semánticas de señales, veremos el siguiente ejemplo, que muestra que a veces la semántica elegida condiciona la forma en la que programamos los monitores.

\begin{ejemplo}
    Queremos programar en un programa concurrente una \textit{barrera parcial}.

    A partir del concepto de barrera que ya manejamos en la asignatura de Arquitectura de Computadores\footnote{Consultar los apuntes si no se conoce el concepto.}, ahora no queremos tener a los procesos esperando a que todos los procesos que intervienen en el programa lleguen a un punto en específico, sino solo queremos que lo hagan $n$ procesos.

    De esta forma, cuando el primer proceso llegue al punto de la barrera, este se bloqueará. Estos les sucederá a los primeros $n-1$ procesos en llegar a la barrera. Cuando el proceso $n$-ésimo llegue a esta, se abrirá la barrera, dejando pasar a estos $n$ primeros procesos que llegaron a esta. La barrera no debe dejar pasar al proceso número $n+1$, sino que este deberá esperar al proceso número $2n$ para voler a abrir la barrera.\\

    Una vez descrito el problema a resolver, planteamos el siguiente monitor como primera solución:
    \begin{minted}[escapeinside=\#\#]{pascal}
        Monitor BP;
           var cola : cond;
               contador : integer;

           procedure barrera(); begin
              contador = contador + 1;
              if (contador < n) then
                 cola.wait();
              else begin
                 for i=1 to n-1 do
                    cola.signal();
                 end
                 contador = 0;
              end
              {Funcionalidad tras pasar la barrera parcial}
           end

           begin
              contador = 0;
           end
    \end{minted}
Veamos ahora el comportamiento del monitor, según la semántica de señales que hayamos escogido:
\begin{description}
    \item [Señalar y Continuar.] Cuando llega el $n$-ésimo proceso de un grupo, este desbloquea a los últimos $n-1$ procesos bloqueados, que pasan a competir por el monitor. Puede suceder que algún proceso del siguiente grupo adelante a uno de este grupo, por lo que no sería una solución válida.
    \item [Señalar y Salir.] El $n$-ésimo proceso solo podría despertar al primer proceso que llegó a la barrera (ya que tras ejecutar \verb|signal|, este terminaría la ejecución de su procedimiento), sin reestablecer el contador a 0, por lo que cualquier proceso que pase por la barrera despertaría al siguiente proceso bloqueado, de forma que los últimos $n$ procesos que ejecuten el protocolo \verb|barrera| quedarán bloqueados de forma indefinida.
    \item [Señalar y Esperar.] El $n$-ésimo proceso solo podría despertar al primer proceso que llegó a la barrera, volviendo el proceso señalador a la cola de entrada al monitor. El $n+1$-ésimo proceso solo podría despertar al segundo proceso, y se iría a la cola de entrada al monitor. Cuando el $2n$-ésimo proceso entre al monitor, este ya no ejecutará \verb|signal| (ya que como la cola está vacía, es equivalente a una instrucción nula), poniendo por fin el contador a 0 y comenzando otra vez con dicho comportamiento.

        No es una solución válida.
    \item [Señalar y Espera Urgente.] El $n$-ésimo proceso desbloquea al primer proceso que llegó a la barrera y a continuación se bloquearía en la cola de urgentes, cediendo el monitor al primer proceso, que finalizaría la ejecución del procedimiento y volvería a entrar el $n$-ésimo proceso, por tener la cola de urgentes mayor prioridad que la cola de entrada al monitor. El proceso se repetiría hasta que el $n$-ésimo proceso finalie la ejecución del procedimiento, restaurando la barrera parcial para los siguientes $n$ procesos.

        Sería un funcionamiento correcto de la barrera, pero hace que el $n$-ésimo proceso realice $n-1$ cambios de contexto innecesarios.
\end{description}
Mostramos ahora una versión alternativa para el monitor que resuelve el problema de la barrera parcial:
\begin{minted}[escapeinside=\#\#]{pascal}
    Monitor BP;
       var cola : cond;
           contador : integer;

       procedure barrera(); begin
          contador = contador + 1;
          if (contador < n) then
             cola.wait();
          contador = contador - 1;
          {Funcionalidad tras pasar la barrera parcial}
          if (contador > 0) then
             cola.signal();
       end

       begin
          contador = 0;
       end
\end{minted}
Vemos su comportamiento:
\begin{itemize}
    \item No funciona con la semántica SC\@.
    \item Funciona con cualquier semántica desplazante, ya que tras llegar a $n$ procesos en la barrera, el $n$-ésimo desbloquea al primero, que desbloquea al segundo, que desbloquea al tercero, \ldots, hasta llegar al $n-1$, que realiza su funcionalidad y viceversa hasta el $n$-ésimo proceso, que deja la barrera en el estado inicial para los siguientes grupos de procesos.
\end{itemize}
\end{ejemplo}~\\

En general, tenemos que tener cuidado con la semántica de señales que estemos empleando, sobre todo si las instrucciones \verb|signal| tienen instrucciones detrás. Además, la semántica SC suele complicar como norma general los diseños de los monitores.

\begin{ejercicio*}
    Determinar para qué tipo de semánticas (SC o desplazantes) funcionan los siguientes monitores que tratan de simular el comportamiento de un semáforo:
    \setlength{\columnsep}{1cm}
    \begin{multicols}{2}
        \begin{minted}[escapeinside=\#\#]{pascal}
            Monitor semaforo_FIFO1;
            var c : cond;
                s : int;

            procedure P;
            begin
                if (s = 0) then
                    c.wait();
                s := s - 1;
            end;

            procedure V;
            begin
                s := s + 1;
                c.signal();
            end

            begin
                s := 0;
            end
        \end{minted}
        \begin{minted}[escapeinside=\#\#]{pascal}
            Monitor semaforo_FIFO2;
            var c : cond;
                s : int;

            procedure P;
            begin
                while (s = 0) do
                    c.wait();
                end do;
                s := s - 1;
            end;

            procedure V;
            begin
                c.signal();
                s := s + 1;
            end

            begin
                s := 0;
            end
        \end{minted}
    \end{multicols}
    En primer lugar, notamos que el monitor de la izquierda realiza una instrucción \verb|if| antes de llamar a la operación \verb|wait|, luego solo sirve para semánticas desplazantes.

    Posteriormente, observamo que el monitor de la derecha realiza la operación \verb|wait| dentro de un bucle, por lo que este sirve para semánticas SC\@. Sin embargo, no sirve para semánticas desplazantes, ya que la operación \verb|signal| se realiza antes de incrementar la variable \verb|s|, por lo que nunca se realizaría dicho incremento, al bloquear al proceso señalador.
\end{ejercicio*}

\subsection{Axiomas para operaciones de sincronización no desplazantes}
Como comentamos ya en la Sección~\ref{sec:ax_desplazantes}, los axiomas que sabemos para demostrar las operaciones de sincronización en monitores solo aplican cuando trabajamos con señales con semánticas desplazantes. Por tanto, si nos encontramos trabajando con semánticas de señales no desplazantes como SC, necesitamos unas nuevas reglas a aplicar.

\begin{description}
    \item [Axioma de la operación \texttt{wait}.]~\\
        Sea $c$ una variable de tipo condición, $L$ un invariante local del procedimiento en el que nos encontremos y $IM$ el invariante del monitor que usamos para demostrar la corrección del mismo, se verifica que
        \begin{equation*}
            \{IM \land L\}\ c.wait();\ \{IM \land L\}
        \end{equation*}
        Es decir, antes y después de la operación \verb|wait| solo podemos asegurar que se cumple el IM (ya que podrían producirse robos de señal), por lo que para asegurarnos de que estamos lejos de un estado inseguro del programa, debemos hacer uso de un bucle \verb|while|, tal y como expusimos anteriormente:
        \begin{minted}[escapeinside=\#\#]{pascal}
            while (cerca de estado inseguro) do begin
               c.wait();
            end
        \end{minted}
    \item [Axioma de la operación \texttt{signal}.]~\\
        Como estamos ante una semántica no desplazante, el proceso señalador seguirá ejecutando el procedimiento que estaba ejecutando cuando ejecutó la operación \verb|signal|, por lo que nada habrá cambiado:
        \begin{equation*}
            \{P\}\ c.signal();\ \{P\}
        \end{equation*}
        Por tanto, la operación \verb|signal| tiene el mismo comportamiento que la instrucción nula en semánticas no desplazantes, ya que podemos poner como pre y poscondición el mismo aserto $P$ y el triple será cierto independientemente del aserto escogido.
\end{description}

\subsection{Intercambio de señales en programas que usan monitores}
En el Capítulo~\ref{cap:comp}, mencionamos que todas las semánticas de señales eran equivalentes para resolver cualquier problema. Esto se debe a que si se cumplen unas determinadas condiciones, entonces todas las señales (las automáticas, las desplazantes y las no desplazantes) son equivalentes. Esto significa que si estamos trabajando con una semántica, podemos reemplazar las señales de dicha semántica por cierto código que asemeja el comportamiento de otra semántica de señales.\\

Las condiciones que deben cumplirse dentro de un monitor para que podamos simular una semántica mediante otra son:
\begin{enumerate}
    \item Se tiene que exigir como poscondición de la operación \verb|wait| el invariante del monitor y nada más estricto (esto es, no podremos usar directamente la condición de sincronización tras la operación \verb|wait|, sino solamente el invariante del monitor\footnote{Notemos que esto no inclumple el axioma de la operación \texttt{wait}, ya que $C\rightarrow IM$.}).
    \item Con señales continuas, cuando realizamos un \verb|signal| no hay desplazamiento, por lo que para que se puedan asemejar a las señales desplazantes, hemos de obligar a que la llamada a \verb|signal| se haga siempre antes de suspender el proceso y salir del monitor.

        Esto es, que la operación \verb|signal| sea la última instrucción en el procedimiento de un monitor o que tras dicha operación haya una instrucción \verb|wait|.
    \item Finalmente, para que sean equivalentes los códigos, no podemos usar la operación \verb|signal_all|.
\end{enumerate}

\subsection{Señales \texttt{wait} con prioridad}
Como hemos visto hasta ahora, tras ejecutar una operación \verb|wait| sobre una variable condición \verb|c|, pasamos a la cola de bloqueados de la misma, con planificación FIFO\@. Sin embargo, en algunos lenguajes de programación aparece la operación \verb|wait| con prioridad sobre las variables de tipo condición, con el fin de solucionar situaciones en las que no queremos que tras una operación \verb|signal| se desbloquee el primero de la cola, sino aquel que tuviera más prioridad y lleve más tiempo en la misma.\\

En el caso descrito, la operación \verb|wait| aceptará un parámetro entero no negativo llamado \textit{prioridad}, de forma que a menor sea el valor de dicho parámetro, mayor prioridad tendrá en la cola de bloqueados. En este caso, la cola FIFO de la variable condición se sustituye por una cola con prioridad.\\

Mostramos a continuación un ejemplo para motivar por qué algunos lenguajes implementan esta operación \verb|wait|.
\begin{ejemplo}
    Queremos programar una alarma para procesos de tal manera que el proceso se bloquee al fijar la alarma y que este se desbloquee al cumplirse la hora previamente fijada. Para ello, implementaremos la alarma haciendo uso de un monitor con dos procedimientos:
    \begin{itemize}
        \item \verb|tick|, cableado a la interrupción de reloj del sistema operativo, con el fin de que el reloj sepa el instante en el que se encuentra en cada momento.
        \item \verb|despiertame|, que servirá para que cada proceso fije el momento en el que quiere ser despertado.
    \end{itemize}
    Para simplificar el ejemplo, usaremos como unidad de medida un ``tick''. Planteamos el siguiente monitor como solución al problema, usando señales \verb|wait| con prioridad:
    \begin{minted}[escapeinside=\#\#]{pascal}
        monitor despertador;
        var ahora : Long_integer;
            despertar : cond;  {variable condición prioritaria}

        procedure despiertame(n : integer);
        var alarma : Long_integer;
        begin
            alarma := ahora + n;    {hora para despertar}
            while ahora < alarma do 
                despertar.wait(alarma);
            end do;
            despertar.signal();
        end

        procedure tick();   {cableada a INT CLK}
        begin
            ahora := ahora + 1;
            despertar.signal();
        end

        begin
            ahora := 0;
        end
    \end{minted}
    De esta forma, lo que estamos haciendo es ordenar la cola de bloqueados de la variable condición \verb|despertar| poniendo al inicio de la misma el primer proceso a desbloquear.

    Por ejemplo, ante el uso del monitor por los procesos que observamos en la Tabla~\ref{tab:ejm_uso_monitor},
    \begin{table}
    \centering
    \begin{tabular}{|c|c|c|c|}
        \hline
        Proceso & Instante & Llamada & Cola de variable condición \\
        \hline
        P1 & 0 & \verb|despiertame(10)| & (P1) \\
        \hline
        P2 & 1 & \verb|despiertame(3)| & (P2 | P1) \\
        \hline
        P3 & 2 & \verb|despiertame(5)| & (P2 | P3 | P1) \\
        \hline
        P4 & 3 & \verb|despiertame(1)| & (P2 | P4 | P3 | P1) \\
        \hline
    \end{tabular}
    \caption{Ejemplo de uso del monitor}
    \label{tab:ejm_uso_monitor}
    \end{table}
    Cuando la interrupción de reloj del sistema llame al procedimiento \verb|tick|, se desbloqueará al primer proceso de la cola de la variable condición, P2, que saldrá del bucle \verb|while| y desbloqueará al siguiente proceso de la cola, P4, que hará lo mismo, desbloqueando a P3, y como para P3 \verb|ahora < alarma|, volverá a bloquearse, quedando la cola de la variable condición como (P3 | P1).\\

    Notemos que con el uso de la operación \verb|wait| con prioridad minimizamos el número de procesos a desbloquear en cada nuevo instante, ya que tenemos la cola de bloqueados ordenada de forma que al inicio están los procesos que antes se desbloquearán, con lo que si encontramos el primer proceso que no se desbloquea, hemos terminado en dicho instante de despertar a los procesos.\\

    Si tratamos de programar este mismo problema sin operaciones \verb|wait| con prioridad, sino con la operación \verb|wait| que venimos usando a lo largo de este Capítulo, el procedimiento \verb|despiertame| quedaría como:
    \begin{minted}[escapeinside=\#\#]{pascal}
        procedure despiertame(n : integer);
        var alarma : Long_integer;
        begin
            alarma := ahora + n;    {hora para despertar}
            while ahora < alarma do 
                despertar.signal();
                despertar.wait();
            end do;
            despertar.signal();
        end
    \end{minted}
    Lo que sucede es que tenemos que estar continuamente desbloqueando a los procesos para comprobar si deben o no despertarse, al no tener ningún orden en dicha cola.
\end{ejemplo}

\section{Implementación de los monitores}
Podemos implementar todas las funcionalidades de un monitor usando semáforos. En esta sección, explicaremos cómo implementar un monitor con semántica de señales SU, que nos garantice el comportamiento de un proceso descrito en la Figura~\ref{fig:su}.\\

Para conseguir simular el comportamiento de un monitor, hemos de conseguir:
\begin{itemize}
    \item Tener una cola de entrada al monitor, que implementaremos usando un semáforo llamado \verb|mutex|.
    \item Tener una cola de procesos urgentes, que implementaremos con un semáforo llamado \verb|next|.
    \item Contabilizar el número de procesos bloqueados en la cola de urgentes\footnote{Con la finalidad de saber si dicha cola está o no vacía}, con una variable entera \verb|next_count|.
    \item Implementación de las variables condición.

        Para ello, por cada variable condición que queramos tener en un monitor, crearemos un semáforo nuevo y una variable entera que controle el número de procesos que bloquea dicho semáforo. Además, crearemos unas nuevas funciones \verb|wait| y \verb|signal| (llamadas \verb|x_wait| y \verb|x_signal|) que simulen el comportamiento de las operaciones \verb|wait| y \verb|signal| de las variables compartidas.
\end{itemize}

Inicializaremos las variables que nos permiten controlar el monitor de la forma:
\begin{minted}[escapeinside=\#\#]{pascal}
    mutex := 1;
    next := 0;
    next_count := 0;
\end{minted}

\subsubsection{Procedimientos del monitor}
Cada vez que nos dispongamos a crear un procedimiento nuevo para el monitor, deberemos ejecutar cierto código antes y después del mismo, con la finalidad de garantizar la exclusión mutua dentro del monitor. Para ello, crearemos dos funciones, \verb|entrada| y \verb|salida|, las cuales deberemos invocar antes y después del cuerpo del procedimiento a programar:
\begin{minted}{pascal}
    procedure P();
    begin
        entrada();
        {cuerpo del procedimiento}
        salida();
    end
\end{minted}
De esta forma, el código de entrada al monitor sería:
\setlength{\columnsep}{1cm}
\begin{minted}{pascal}
    procedure entrada();
    begin
        sem_wait(mutex);    {garantizar exclusión mutua}
    end
\end{minted}    
Y el de salida:
\begin{minted}{pascal}
    procedure salida();
    begin
        if(next_count <> 0)    {hay procesos en cola de urgentes}
            sem_signal(next);
        else
            sem_signal(mutex);   {liberamos el monitor}
        end if
    end
\end{minted}

\subsubsection{Variables condición}
Por cada variable condición a usar dentro del monitor necesitamos tener un par semáforo, entero. Mostramos ahora cómo podemos implementar las operaciones \verb|wait| y \verb|signal| de las variables condición usando semáforos.

Para ello, crearemos una función por cada operación a simular, la cual recibirá dos parámetros: el semáforo que simula la variable condición (\verb|x_sem|) y la cantidad de procesos que dicho semáforo tiene bloqueados (\verb|x_count|):
\begin{minted}{pascal}
    procedure x_wait(x_sem : semaphore, x_count : integer);
    begin
        x_count := x_count + 1;  {un bloqueado más}

        if(next_count <> 0) then  {hay procesos en cola de urgentes}
            sem_signal(next);
        else
            sem_signal(mutex);  {deja libre el monitor}
        end if

        sem_wait(x_sem);
        x_count := x_count - 1;
    end
\end{minted}
\begin{minted}{pascal}
    procedure x_signal(x_sem : semaphore, x_count : integer);
    begin
        if(x_count <> 0) then  {si no hay bloqueado no hace nada}
            next_count := next_count + 1;
            sem_signal(x_sem);
            sem_wait(next);
            next_count := next_count - 1;
        end if
    end
\end{minted}
Y con todas estas funciones tenemos \textbf{casi} implementados los monitores. Hay que tener en cuenta que los semáforos no tienen una política de planificación fija, sino que es el sistema operativo quien planifica los procesos desbloqueados por el semáforo.

Es necesario por tanto, usar semáforos con colas FIFO para tener totalmente implementados los monitores.


    \fancyhead[R]{\helv \nouppercase{\rightmark}}
    \chapter{Relaciones de Problemas}
    %\section{Introducción}

\begin{ejercicio}
    Considerar el siguiente fragmento de programa para 2 procesos \verb|P1| y \verb|P2|: Los dos procesos
    pueden ejecutarse a cualquier velocidad. ¿Cuáles son los posibles valores resultantes para la
    variable \verb|x|? Suponer que \verb|x| debe ser cargada en un registro para incrementarse y que cada
    proceso usa un registro diferente para realizar el incremento.
    \setlength{\columnsep}{2cm} % Ajusta el espacio entre columnas
    \begin{multicols}{2}
        \begin{minted}{pascal}
        {variables compartidas}
        var x : integer := 0 ;
        Process P1;
        var i: integer;
        begin
          begin
            for i:= 1 to 2 do begin
              x:= x + 1;
            end
          end
        end
        \end{minted}
        
        \begin{minted}{pascal}
            

        Process P2;
        var j: integer;
        begin
          begin
            for j:= 1 to 2 do begin
              x:= x + 1;
            end
          end
        end
        \end{minted}
    \end{multicols}

Observando el código, cada proceso hace 2 lecturas y dos escrituras (incrementos) en \verb|x|.
\begin{itemize}
    \item Como cada proceso aumenta dos veces el valor de \verb|x|, el valor de \verb|x| ha de ser, como mínimo, $2$.
    \item Como en total se hacen 4 incrementos, el valor de \verb|x| ha de ser $4$ como máximo.
\end{itemize}

Notando por $l_{ij}$ a la $j$-ésima lectura del proceso $i$ y por $e_{ij}$ a la $j$-ésima escritura del proceso $i$, ambas referidas a la variable \verb|x|, podemos obtener cualquiera de las siguientes trazas de ejecución:

\begin{table}[H]
\centering
\begin{tabular}{|c|c|c||c|c|c||c|c|c||c|c|c|}
    \hline
    \verb|P1| & \verb|P2| & \verb|x| & \verb|P1| & \verb|P2| & \verb|x| & \verb|P1| & \verb|P2| & \verb|x| & \verb|P1| & \verb|P2| & \verb|x| \\
    \hline
    \hline
    $l_{11}$ & - & 0 & $l_{11}$ & - & 0 & $l_{11}$ & - & 0 & $l_{11}$ & - & 0 \\
    \hline
    $e_{11}$ & - & 1 & - & $l_{21}$ & 0 & $e_{11}$ & - & 1 & - & $l_{21}$ & 0 \\
    \hline
    - & $l_{21}$ & 1 & $e_{11}$ & - & 1 & - & $l_{21}$ & 1 & $e_{11}$ & - & 1 \\
    \hline
    - & $e_{21}$ & 2 & - & $e_{21}$ & 1 & - & $e_{21}$ & 2 & - & $e_{21}$ & 1 \\
    \hline
    $l_{12}$ & - & 2 & $l_{12}$ & - & 1 & $l_{12}$ & - & 2 & $l_{12}$ & - & 1 \\
    \hline
    $e_{12}$ & - & 3 & $e_{12}$ & - & 2 & - & $l_{22}$ & 2 & - & $l_{22}$ & 1 \\
    \hline
    - & $l_{22}$ & 3 & - & $l_{22}$ & 2 & $e_{12}$ & - & 3 & $e_{12}$ & - & 2 \\
    \hline
    - & $e_{22}$ & 4 & - & $e_{22}$ & 3 & - & $e_{22}$ & 3 & - & $e_{22}$ & 2  \\
    \hline
\end{tabular}
\end{table}

Luego los posibles valores resultantes para \verb|x| son: 2, 3 y 4.

\end{ejercicio}


\begin{ejercicio}
    ¿Cómo se podría hacer la copia del fichero \verb|f| en otro \verb|g|, de forma concurrente, utilizando la
    instrucción concurrente \verb|cobegin-coend|? Para ello, suponer que:
    \begin{enumerate}
        \item Los archivos son una secuencia de items de un tipo arbitrario \verb|T|, y se encuentran ya abiertos
        para lectura (\verb|f|) y escritura (\verb|g|). Para leer un ítem de \verb|f| se usa la llamada a función \verb|leer(f)| y
        para saber si se han leído todos los ítems de \verb|f|, se puede usar la llamada \verb|fin(f)| que devuelve
        verdadero si ha habido al menos un intento de leer cuando ya no quedan datos. Para
        escribir un dato \verb|x| en \verb|g| se puede usar la llamada a procedimiento \verb|escribir(g,x)|.

        \item El orden de los items escritos en \verb|g| debe coincidir con el de \verb|f|.
        \item Dos accesos a dos archivos distintos pueden solaparse en el tiempo.
    \end{enumerate}~\\

\noindent
La copia del fichero \verb|f| en el fichero \verb|g| se podría realizar siguiendo el paradigma productor/consumidor que hemos visto en teoría en el Tema 1, mediante el uso de dos procesos:
\begin{itemize}
    \item Uno que lea un ítem del fichero \verb|f| y lo escriba en una variable compartida.
    \item Otro que lea dicha variable compartida y escriba el ítem en el fichero \verb|g|.
\end{itemize}
En dicho código, debemos garantizar que:
\begin{itemize}
    \item El consumidor no lea la variable antes de que el productor escriba en ella.
    \item En la segunda escritura del productor, debemos esperar a que antes la haya leído el consumidor.
    \item En la segunda lectura del consumidor, debemos esperar a que antes haya modificado la variable el productor.
\end{itemize}
Siguiendo estos pasos, obtendríamos un código como el siguiente:
\begin{minted}{pascal}
    process CopiaFicheros ;
    var ant, sig : T ;
    begin
      sig = leer(f) ;
      while not fin(f) do begin
        ant = sig ;
        cobegin
          escribir(g, anterior) ;
          sig = leer(f) ;
        coend
      end
    end
\end{minted}


\end{ejercicio}

\begin{ejercicio}\label{ej:3}
    Construir, utilizando las instrucciones concurrentes \verb|cobegin-coend| y \verb|fork-join|, programas concurrentes que se correspondan con los grafos de precedencia que se muestran en la figura~\ref{fig:grafoEj3}.
    \begin{figure}
        \centering
        \begin{subfigure}{0.3\textwidth}
            \centering
            \resizebox{\linewidth}{!}{
                \begin{tikzpicture}[
                    node/.style={circle, draw, minimum size=0.5cm},
                    edge/.style={-stealth}
                    ]
        
                    % Nodos
                    \node[node] (P0) {P0};
                    \node[node, below left=of P0] (P1) {P1};
                    \node[node, below right=of P0] (P2) {P2};
                    \node[node, below=of P1] (P3) {P3};
                    \node[node, below left=of P3] (P4) {P4};
                    \node[node, below right=of P3] (P5) {P5};
                    \node[node, below=of P5] (P6) {P6};
        
                    % Aristas
                    \draw[edge] (P0) -- (P1);
                    \draw[edge] (P0) -- (P2);
                    \draw[edge] (P1) -- (P3);
                    \draw[edge] (P3) -- (P4);
                    \draw[edge] (P3) -- (P5);
                    \draw[edge, bend right] (P4) to (P6);
                    \draw[edge] (P5) -- (P6);
                    \draw[edge, bend left] (P2) to (P6);
                
                \end{tikzpicture}
            }
            \caption{DAG del apartado~\ref{ej:3.1}.}
            \label{fig:grafoEj3.1}
            
        \end{subfigure}
        \begin{subfigure}{0.3\textwidth}
            \centering
            \resizebox{\linewidth}{!}{
                \begin{tikzpicture}[
                    node/.style={circle, draw, minimum size=1cm},
                    edge/.style={-stealth}
                    ]
        
                    % Nodos
                    \node[node] (P0) {P0};
                    \node[node, below left=of P0] (P1) {P1};
                    \node[node, below right=of P0] (P2) {P2};
                    \node[node, below=of P1] (P3) {P3};
                    \node[node, below left=of P1] (P4) {P4};
                    \node[node, below right=of P1] (P5) {P5};
                    \node[node, below=of P5] (P6) {P6};
        
                    % Aristas
                    \draw[edge] (P0) -- (P1);
                    \draw[edge] (P0) -- (P2);
                    \draw[edge] (P1) -- (P3);
                    \draw[edge] (P1) -- (P4);
                    \draw[edge] (P1) -- (P5);
                    \draw[edge, bend right] (P4) to (P6);
                    \draw[edge] (P5) -- (P6);
                    \draw[edge, bend left] (P2) to (P6);
                    \draw[edge, bend right] (P3) to (P6);
                
                \end{tikzpicture}
            }
            \caption{DAG del apartado~\ref{ej:3.2}.}
            \label{fig:grafoEj3.2}
            
        \end{subfigure}
        \begin{subfigure}{0.3\textwidth}
            \centering
            \resizebox{\linewidth}{!}{
                \begin{tikzpicture}[
                    node/.style={circle, draw, minimum size=1cm},
                    edge/.style={-stealth}
                    ]
        
                    % Nodos
                    \node[node] (P0) {P0};
                    \node[node, below left=of P0] (P1) {P1};
                    \node[node, below right=of P0] (P2) {P2};
                    \node[node, below=of P1] (P3) {P3};
                    \node[node, below left=of P3] (P4) {P4};
                    \node[node, below right=of P3] (P5) {P5};
                    \node[node, below=of P5] (P6) {P6};
        
                    % Aristas
                    \draw[edge] (P0) -- (P1);
                    \draw[edge] (P0) -- (P2);
                    \draw[edge] (P1) -- (P3);
                    \draw[edge] (P3) -- (P4);
                    \draw[edge] (P3) -- (P5);
                    \draw[edge, bend right] (P4) to (P6);
                    \draw[edge] (P5) -- (P6);
                    \draw[edge] (P2) to (P5);
                
                \end{tikzpicture}
            }
            \caption{DAG del apartado~\ref{ej:3.3}.}
            \label{fig:grafoEj3.3}
            
        \end{subfigure}
        \caption{Grafos de precedencia del ejercicio~\ref{ej:3}.}
        \label{fig:grafoEj3}
    \end{figure}


    \begin{enumerate}
        \item \label{ej:3.1}
         Grafo de precedencia de la figura~\ref{fig:grafoEj3.1}:

         \begin{figure}[H]
             \centering
             \begin{subfigure}[b]{0.45\textwidth}
                \centering
                 \begin{minted}{pascal}
                     begin
                       P0;
                       fork P2; P1;
                       P3;
                       fork P5; P4;
                       join P2; join P5;
                       P6;
                     end
                 \end{minted}
             \end{subfigure}\hfill
             \begin{subfigure}[b]{0.45\textwidth}
                \centering
                 \begin{minted}{pascal}
                     begin
                       P0;
                       cobegin
                         P2;
                         begin
                           P1;
                           P3;
                           cobegin P4; P5; coend
                         end
                       coend
                       P6;
                     end
                 \end{minted}
             \end{subfigure}
         \end{figure}
                 
        \item \label{ej:3.2}
        Grafo de precedencia de la figura~\ref{fig:grafoEj3.2}:
        \begin{figure}[H]
            \centering
            \begin{subfigure}[b]{0.45\textwidth}
                \centering
                \begin{minted}{pascal}
                begin
                  P0;
                  fork P2; P1;
                  fork P5; fork P3; P4;
                  join P2; join P5; join P3;
                  P6;
                end
                \end{minted}
            \end{subfigure}\hfill
            \begin{subfigure}[b]{0.45\textwidth}
                \centering
                \begin{minted}{pascal}
                begin
                  P0;
                  cobegin
                    P2;
                    begin
                      P1;
                      cobegin P4; P3; P5; coend
                    end
                  coend
                  P6;
                end
                \end{minted}
            \end{subfigure}
        \end{figure}
        
         \item \label{ej:3.3}
         Grafo de precedencia de la figura~\ref{fig:grafoEj3.3}:
         \begin{minted}{pascal}
         begin
           P0;
           fork P2; P1;
           P3;
           fork P4; join P2; P5;
           join P4;
           P6;
         end
         \end{minted}

         Sin embargo, no podemos hacer al 100\% el DAG de la figura~\ref{fig:grafoEj3.3}, ya que tras \verb|P3| debemos crear una estructura \verb|cobegin-coend|. Sin embargo, este debe esperar a \verb|P2|, por lo que la estructura \verb|cobegin-coend| tendrá que esperar a \verb|P2|, pero es que \verb|P4| no necesita que \verb|P2| termine. 

         Por tanto, no se puede programar con creación de hebras de forma estructurada. Sin embargo, podemos ofrecer dos soluciones, cada una que impone algo que el grafo no nos dice:
         \begin{enumerate}
             \item Si obligamos a que \verb|P4| también espere a \verb|P2|, obtendríamos el código:
                 \begin{minted}{pascal}
                     begin
                       P0;
                       cobegin
                         P2;
                         begin
                           P1; P3;
                         end
                       coend
                       cobegin P4; P5; coend
                       P6;
                     end
                 \end{minted}
                 
             \item Si ahora queremos ejecutar de forma concurrente el flujo que tiene a \verb|P1|, \verb|P3| y \verb|P4| con el flujo que tiene a \verb|P2|, entonces obligamos a que \verb|P5| espere a \verb|P4| (que no nos lo especifica el DAG, pero lo necesitamos para poder programarlo de forma estructurada):
                 \begin{minted}{pascal}
                     begin
                       P0;
                       cobegin
                         begin P1; P3; P4; end
                         P2;
                       coend
                       P5;
                       P6;
                     end
                 \end{minted}
                 
         \end{enumerate}
    \end{enumerate}

\end{ejercicio}



\begin{ejercicio} \label{ej:4}
    Dados los siguientes fragmentos de programas concurrentes, obtener sus grafos de precedencia asociados:
    \begin{figure}[H]
        \centering
        \begin{subfigure}[b]{0.45\textwidth}
            \centering
            \begin{minted}{pascal}
                begin
                    P0 ;
                    cobegin
                        P1 ;
                        P2 ;
                        cobegin
                            P3 ; P4 ; P5 ; P6 ;
                        coend ;
                        P7 ;
                    coend
                    P8 ;
                end
            \end{minted}
            \caption{Programa 1.}
            \label{code:prog1_Ej4}
        \end{subfigure}\hfill
        \begin{subfigure}[b]{0.45\textwidth}
            \centering
            \begin{minted}{pascal}
                begin
                    P0 ;
                    cobegin
                        begin
                            cobegin
                                P1 ; P2 ;
                            coend
                            P5 ;
                        end
                        begin
                            cobegin
                                P3 ; P4 ;
                            coend
                            P6 ;
                        end
                    coend
                    P7 ;
                end
            \end{minted}
            \caption{Programa 2.}
            \label{code:prog2_Ej4}
        \end{subfigure}
        \caption{Programas concurrentes del ejercicio~\ref{ej:4}.}
    \end{figure}
    
    \begin{enumerate}
        \item Programa de la figura~\ref{code:prog1_Ej4}.

            \begin{figure}[H]
            \centering
            \begin{tikzpicture}[
                    node/.style={circle, draw, minimum size=1cm},
                    edge/.style={-stealth}
                    ]
                \node[node] (P0) {P0};
                \node[node, below=of P0] (P4) {P4};
                \node[node, left=of P4] (P3) {P3};
                \node[node, left=of P3] (P2) {P2};
                \node[node, left=of P2] (P1) {P1};
                \node[node, right=of P4] (P5) {P5};
                \node[node, right=of P5] (P6) {P6};
                \node[node, right=of P6] (P7) {P7};
                \node[node, below=of P4] (P8) {P8};

                \draw[edge] (P0) -- (P1);
                \draw[edge] (P0) -- (P2);
                \draw[edge] (P0) -- (P3);
                \draw[edge] (P0) -- (P4);
                \draw[edge] (P0) -- (P5);
                \draw[edge] (P0) -- (P6);
                \draw[edge] (P0) -- (P7);

                \draw[edge] (P1) -- (P8);
                \draw[edge] (P2) -- (P8);
                \draw[edge] (P3) -- (P8);
                \draw[edge] (P4) -- (P8);
                \draw[edge] (P5) -- (P8);
                \draw[edge] (P6) -- (P8);
                \draw[edge] (P7) -- (P8);
            \end{tikzpicture}
            \caption{DAG para la figura~\ref{code:prog1_Ej4}.}
            \end{figure}
        \item Programa de la figura~\ref{code:prog2_Ej4}.
            \begin{figure}[H]
            \centering
            \begin{tikzpicture}[node/.style={circle, draw, minimum size=1cm},
                                edge/.style={-stealth}]
                \node[node] (P0) {P0};
                \node[node, below left=of P0] (P2) {P2};
                \node[node, left=of P2] (P1) {P1};
                \node[node, below right=of P0] (P3) {P3};
                \node[node, right=of P3] (P4) {P4};
                \node[node, below right=of P1] (P5) {P5};
                \node[node, below left=of P4] (P6) {P6};

                \node[below=of P0] (A) {};
                \node[below=of A] (B) {};
                \node[node, below=of B] (P7) {P7};
                
                \draw[edge] (P0) -- (P1);
                \draw[edge] (P0) -- (P2);
                \draw[edge] (P0) -- (P3);
                \draw[edge] (P0) -- (P4);
                \draw[edge] (P1) -- (P5);
                \draw[edge] (P2) -- (P5);
                \draw[edge] (P3) -- (P6);
                \draw[edge] (P4) -- (P6);
                \draw[edge] (P5) -- (P7);
                \draw[edge] (P6) -- (P7);
            \end{tikzpicture}
            \caption{DAG para la figura~\ref{code:prog2_Ej4}.}
            \end{figure}
    \end{enumerate}
\end{ejercicio}

\begin{ejercicio} \label{ej:5}
    Suponer un sistema de tiempo real que dispone de un captador de impulsos conectado a un
    contador de energía eléctrica. La función del sistema consiste en contar el número de impulsos
    producidos en 1 hora (cada Kwh consumido se cuenta como un impulso) e imprimir este número
    en un dispositivo de salida. Para ello se dispone de un programa concurrente con 2 procesos: un
    proceso acumulador (lleva la cuenta de los impulsos recibidos) y un proceso escritor (escribe
    en la impresora). En la variable común a los 2 procesos \verb|n| se lleva la cuenta de los impulsos. El
    proceso acumulador puede invocar un procedimiento \verb|Espera_impulso| para esperar a que llegue
    un impulso, y el proceso escritor puede llamar a \verb|Espera_fin_hora| para esperar a que termine
    una hora. El código de los procesos de este programa podría ser el descrito en el Código Fuente~\ref{code:ej5}.
    \begin{observacion}
        En el programa se usan sentencias de acceso a la variable \verb|n| encerradas entre los símbolos \verb|<| y
        \verb|>|. Esto significa que cada una de esas sentencias se ejecuta en exclusión mutua entre los dos
        procesos, es decir, esas sentencias se ejecutan de principio a fin sin entremezclarse entre ellas.
        Supongamos que en un instante dado el acumulador está esperando un impulso, el escritor está
        esperando el fin de una hora, y la variable \verb|n| vale \verb|k|. Después se produce de forma simultánea
        un nuevo impulso y el fin del periodo de una hora.
    \end{observacion}

    Obtener las posibles secuencias de interfolicación de las instrucciones (1),(2), y (3) a partir de
    dicho instante, e indicar cuales de ellas son correctas y cuales incorrectas (las incorrectas son
    aquellas en las cuales el impulso no se contabiliza).
    \begin{listing}
        \begin{minted}{pascal}
            { variable compartida: }
            var n : integer; { contabiliza impulsos }
            begin
            while true do begin
                Espera_impulso();
                < n := n+1 > ; { (1) }
                end
            end
            process Escritor ;
            begin
            while true do begin
                Espera_fin_hora();
                write( n ) ; { (2) }
                < n := 0 > ; { (3) }
                end
            end
        \end{minted}
        \caption{Código acumulador-escritor del ejercicio~\ref{ej:5}.}
        \label{code:ej5}
    \end{listing}

    En primer lugar, notemos que la intrucción $2$ siempre se ejecutará
    antes que la instrucción $3$, ya que están ambas en el mismo proceso (el escritor).
    Por tanto, las secuencia de instrucciones que se pueden dar son las intercalaciones de la instrucción $1$ con las otras dos instrucciones; es decir: $A)~1-2-3$, $B)~2-1-3$ y $C)~2-3-1$. El análisis de cada una de ellas está en la Tabla~\ref{tab:ej5}.
    \begin{table}
        \centering
        \begin{tabular}{|lcc|lcc|lcc|}
            \hline
            \multicolumn{3}{|c|}{Opción $A$} & \multicolumn{3}{|c|}{Opción $B$} & \multicolumn{3}{|c|}{Opción $C$} \\ \hline
            Operación & \verb|n| & Salida & Operación & \verb|n| & Salida & Operación & \verb|n| & Salida\\
            \hline
            $-$ & $k$ & $-$ & $-$ & $k$ & $-$ & $-$ & $k$ & $-$ \\
            \verb|1|$\to$ \verb|n:=n+1| & $k+1$ & $-$ & \verb|2|$\to$ \verb|write(n)| & $k$ & $k$ & \verb|2|$\to$ \verb|write(n)| & $k$ & $k$ \\
            \verb|2|$\to$ \verb|write(n)| & $k+1$ & $k+1$ & \verb|1|$\to$ \verb|n:=n+1| & $k+1$ & $k$ & \verb|3|$\to$ \verb|n:=0| & $0$ & $k$ \\
            \verb|3|$\to$ \verb|n:=0| & $0$ & $k+1$ & \verb|3|$\to$ \verb|n:=0| & $0$ & $k$ & \verb|1|$\to$ \verb|n:=n+1| & $1$ & $k$ \\
            \hline
        \end{tabular}
        \caption{Tabla de opciones del ejercicio~\ref{ej:5}.}
        \label{tab:ej5}
    \end{table}
    Notemos que tanto la opción $A$ como la $C$ son válidas, ya que en ambas se contabiliza el impulso. No obstante, difieren entre sí, puesto que en la opción $A$ se contabiliza el impulso en la hora que termina (imprimiéndose entonces), mientras que en la opción $C$ se contabiliza el impulso en la hora siguiente (no imprimiéndose entonces en esta salida). No obstante, la opción $B$ es incorrecta, ya que no se contabiliza el impulso en la hora que termina ni en la siguiente.
\end{ejercicio}



\begin{ejercicio} \label{ej:6}
    Supongamos un programa concurrente en el cual hay, en memoria compartida dos vectores \verb|a| y
    \verb|b| de enteros y con tamaño par, declarados como sigue:
    \begin{minted}{pascal}
        var a,b : array[1..2*n] of integer ; { n es una constante predefinida }
    \end{minted}
    Queremos escribir un programa para obtener en \verb|b| una copia ordenada del contenido de \verb|a| (nos
    da igual el estado en que queda \verb|a| después de obtener \verb|b|). Para ello disponemos de la función
    \verb|Sort| que ordena un tramo de \verb|a| (entre las entradas \verb|s| y \verb|t|, ambas incluidas). También disponemos
    la función \verb|Copiar|, que copia un tramo de \verb|a| (desde \verb|s| hasta \verb|t|) en \verb|b| (a partir de \verb|o|). Estas funciones
    se muestran en el Código Fuente~\ref{code:ej_6SortCopiar}.
    \begin{listing}
        \begin{minted}{pascal}
            procedure Sort( s,t : integer );
                var i, j : integer ;
                begin
                    for i := s to t do
                    for j:= s+1 to t do
                        if a[i] < a[j] then
                            swap( a[i], b[j] ) ;
                end

            procedure Copiar( o,s,t : integer );
                var d : integer ;
                begin
                    for d := 0 to t-s do
                        b[o+d] := a[s+d] ;
                end
        \end{minted}
        \caption{Procedimientos \mintinline{pascal}{Sort} y \mintinline{pascal}{Copiar} del ejercicio~\ref{ej:6}.}
        \label{code:ej_6SortCopiar}
    \end{listing}

    El programa para ordenar se puede implementar de dos formas:
    \begin{enumerate}
        \item Ordenar todo el vector \verb|a|, de forma secuencial con la función \verb|Sort|, y después copiar cada
        entrada de \verb|a| en \verb|b|, con la función \verb|Copiar|.
        \item Ordenar las dos mitades de \verb|a| de forma concurrente, y después mezclar dichas dos mitades
        en un segundo vector \verb|b| (para mezclar usamos un procedimiento \verb|Merge|).
    \end{enumerate}
    En el Código Fuente~\ref{code:ej6_2versiones} se muestra el código de ambas versiones.
    \begin{listing}
        \begin{minted}{pascal}
            procedure Secuencial() ;
                var i : integer ;
                begin
                    Sort( 1, 2*n ); { ordena a }
                    Copiar( 1, 2*n ); { copia a en b }
                end

            procedure Concurrente() ;
                begin
                    cobegin
                        Sort( 1, n );
                        Sort( n+1, 2*n );
                    coend
                    Merge( 1, n+1, 2*n );
                end
        \end{minted}
        \caption{Procedimientos \mintinline{pascal}{Secuencial} y \mintinline{pascal}{Concurrente} del ejercicio~\ref{ej:6}.}
        \label{code:ej6_2versiones}
    \end{listing}

    El código de la función \verb|Merge|, disponible en el Código Fuente~\ref{code:ej6_Merge}, se encarga de ir leyendo las dos mitades de \verb|a|, en cada paso, seleccionar el menor elemento de los dos siguientes por leer (uno en cada mitad), y escribir dicho menor elemento en la siguiente mitad del vector mezclado \verb|b|.
    \begin{listing}
        \begin{minted}{pascal}
            procedure Merge( inferior, medio, superior: integer ) ;
                { siguiente posicion a escribir en b }
                var escribir : integer := 1 ;
                { siguiente pos. a leer en primera mitad de a }
                var leer1 : integer := inferior ;
                { siguiente pos. a leer en segunda mitad de a }
                var leer2 : integer := medio ;
                begin
                    { mientras no haya terminado con alguna mitad }
                    while leer1 < medio and leer2 <= superior do begin
                        if a[leer1] < a[leer2] then begin { minimo en la primera mitad }
                            b[escribir] := a[leer1] ;
                            leer1 := leer1 + 1 ;
                        end else begin { minimo en la segunda mitad }
                            b[escribir] := a[leer2] ;
                            leer2 := leer2 + 1 ;
                        end
                        escribir := escribir+1 ;
                    end
                    { se ha terminado de copiar una de las mitades,
                    copiar lo que quede de la otra }
                    if leer2 > superior then
                        { copiar primera } Copiar( escribir, leer1, medio-1 );
                    else Copiar( escribir, leer2, superior ); { copiar segunda }
                end
        \end{minted}
        \caption{Procedimiento \mintinline{pascal}{Merge} del ejercicio~\ref{ej:6}.}
        \label{code:ej6_Merge}
    \end{listing}

    Llamaremos $T_s(k)$ al tiempo que tarda el procedimiento \verb|Sort| cuando actúa sobre un segmento del vector con $k$ entradas. Suponemos que el tiempo que (en media) tarda cada iteración del bucle interno que hay en \verb|Sort| es la unidad (por definición). Es evidente que ese bucle tiene $\dfrac{k(k-1)}{2}$ iteraciones, luego:
    \[
        T_s(k) = \dfrac{k(k-1)}{2} = \dfrac{1}{2}\cdot k^2 - \dfrac{1}{2}\cdot k
    \]

    El tiempo que tarda la versión secuencial sobre $2n$ elementos (llamaremos $S$ a dicho tiempo) será evidentemente $T_s(2n)$, luego:
    \[
        S = T_s(2n) = \dfrac{1}{2}\cdot (2n)^2 - \dfrac{1}{2}\cdot 2n = 2n^2 - n
    \]

    Con estas definiciones, calcular el tiempo que tardará la versión paralela, en los dos siguientes casos. Para esto, hay que suponer que cuando el procedimiento \verb|Merge| actúa sobre un vector con $p$ entradas, tarda $p$ unidades de tiempo en ello, lo cual es razonable teniendo en cuenta que en esas circunstancias \verb|Merge| copia $p$ valores desde \verb|a| hacia \verb|b|. Si llamamos a este tiempo $T_m(p)$, podemos escribir $T_m(p) = p$. Escribe también una comparación cualitativa de los tres tiempos ($S$, $P_1$ y $P_2$). 
    \begin{enumerate}
        \item Las dos instancias concurrentes de \verb|Sort| se ejecutan en el mismo procesador (llamamos $P_1$ al tiempo que tarda).
        
        En este caso, tenemos que hay ganancia de tiempo, ya que las dos instancias de \verb|Sort| no pueden ejecutarse simultáneamente. Por tanto, tenemos que:
        \begin{equation*}
            P_1 = 2\cdot T_s(n) + T_m(2n) = 2\cdot \left( \dfrac{1}{2}\cdot n^2 - \dfrac{1}{2}\cdot n \right) + 2n = n^2 - n + 2n = n^2 + n
        \end{equation*}
        \item Cada instancia de \verb|Sort| se ejecuta en un procesador distinto (lo llamamos $P_2$).
        
        Al poder ejecutarse ahora de forma simultánea, tenemos que:
        \begin{equation*}
            P_2 = \max\{ T_s(n), T_s(n) \} + T_m(2n) = \left( \dfrac{1}{2}\cdot n^2 - \dfrac{1}{2}\cdot n \right) + 2n = \dfrac{1}{2}\cdot n^2 + \dfrac{3}{2}\cdot n
        \end{equation*}
    \end{enumerate}

    Como podemos ver, en los tres casos la eficiencia es del orden cuadrático. No obstante,
    el coeficiente de $n^2$ es distinto en cada caso, siendo el mayor en la versión secuencial. Deducimos por tanto que las versiones concurrentes son más eficientes que la secuencial, y estas mejoras son significativas para valores de $n$ grandes.    

\end{ejercicio}

\begin{ejercicio} \label{ej:7}
    Supongamos que tenemos un programa con tres matrices (\verb|a|, \verb|b| y \verb|c|) de valores flotantes declaradas
    como variables globales. La multiplicación secuencial de \verb|a| y \verb|b| (almacenando el resultado en \verb|c|)
    se puede hacer mediante un procedimiento \verb|MultiplicacionSec| declarado como aparece aquí:
    \begin{minted}{pascal}
        var a, b, c : array[1..3,1..3] of real ;
        procedure MultiplicacionSec()
            var i,j,k : integer ;
            begin
                for i := 1 to 3 do
                    for j := 1 to 3 do begin
                        c[i,j] := 0 ;
                        for k := 1 to 3 do
                            c[i,j] := c[i,j] + a[i,k]*b[k,j] ;
                    end
            end
    \end{minted}
    Escribir un programa con el mismo fin, pero que use 3 procesos concurrentes. Suponer que
    los elementos de las matrices \verb|a| y \verb|b| se pueden leer simultáneamente, así como que elementos distintos de \verb|c| pueden escribirse simultáneamente.\\

    Esta solución puede llevarse a cabo de dos formas (tal y como vimos en la asigantura de AC). Crearemos tres procesos, y cada uno puede calcular una fila de la matriz \verb|c| (Código Fuente~\ref{code:ej7_fila}) o bien cada uno puede calcular una columna de la matriz \verb|c| (Código Fuente~\ref{code:ej7_columna}). No obstante, en la asignatura de AC se vio que, por norma general, es más eficiente calcular por filas que por columnas, ya que en el primer caso se accede a la matriz \verb|a| de forma secuencial, aprovechando así la localidad espacial y provocando un número menor de fallos de caché.
    \begin{figure}
        \centering
        \begin{subfigure}{0.45\textwidth}
            \centering
            \begin{minted}{pascal}
                var a, b, c : array[1..3,1..3] of real ;
                process CalcularFila[ i : 1..3] ;
                    var j, k : integer ;
                    begin
                        for j := 1 to 3 do begin
                            c[i,j] := 0 ;
                            for k := 1 to 3 do
                                c[i,j] := c[i,j] + a[i,k]*b[k,j] ;
                        end
                    end
                end
            \end{minted}
            \caption{Procesos concurrentes para calcular por filas.}
            \label{code:ej7_fila}
        \end{subfigure}\hfill
        \begin{subfigure}{0.45\textwidth}
            \centering
            \begin{minted}{pascal}
                var a, b, c : array[1..3,1..3] of real ;
                process CalcularColumna[j : 1..3] ;
                    var i, k : integer ;
                    begin
                        for i := 1 to 3 do begin
                            c[i,j] := 0 ;
                            for k := 1 to 3 do
                                c[i,j] := c[i,j] + a[i,k]*b[k,j] ;
                        end
                    end
                end
            \end{minted}
            \caption{Procesos concurrentes para calcular por columnas.}
            \label{code:ej7_columna}
        \end{subfigure}
        \caption{Códigos para el ejercicio~\ref{ej:7}.}
    \end{figure}
    % // TODO: Se lanzan solos?
\end{ejercicio}

\begin{ejercicio}\label{ej:8}
    Un trozo de programa ejecuta nueve rutinas o actividades (\verb|P1|, \verb|P2|, \ldots, \verb|P9|), repetidas veces,
    de forma concurrentemente con \verb|cobegin-coend| (ver trozo de código de la figura~\ref{code:ej8_enunciado}), pero que requieren
    sincronizarse según determinado grafo (ver la figura~\ref{fig:ej8_grafo}).
    \begin{figure}
        \centering
        \begin{subfigure}{0.45\textwidth}
            \centering
            \begin{minted}{pascal}
                while true do
                cobegin
                    P1 ; P2 ; P3 ;
                    P4 ; P5 ; P6 ;
                    P7 ; P8 ; P9 ;
                coend
            \end{minted}
            \caption{Código del ejercicio~\ref{ej:8}.}
            \label{code:ej8_enunciado}
        \end{subfigure} \hfill
        \begin{subfigure}{0.45\textwidth}
            \centering
            \begin{tikzpicture}[
                node/.style={circle, draw, minimum size=0.5cm},
                edge/.style={-Stealth}
                ]
    
                % Nodos
                \node[node] (P1) {P1};
                \node[node, right =of P1] (P2) {P2};
                \node[node, right =of P2] (P3) {P3};
                \node[node, below =of P1] (P4) {P4};
                \node[node, right =of P4] (P5) {P5};
                \node[node, right =of P5] (P6) {P6};
                \node[node, below =of P4] (P7) {P7};
                \node[node, right =of P7] (P8) {P8};
                \node[node, right =of P8] (P9) {P9};

                % Aristas
                \draw[edge] (P1) -- (P4);
                \draw[edge] (P1) -- (P2);
                \draw[edge] (P2) -- (P3);
                \draw[edge] (P2) -- (P5);
                \draw[edge] (P3) -- (P6);
                \draw[edge] (P4) -- (P5);
                \draw[edge] (P4) -- (P7);
                \draw[edge] (P5) -- (P6);
                \draw[edge] (P5) -- (P8);
                \draw[edge] (P6) -- (P9);
                \draw[edge] (P7) -- (P8);
                \draw[edge] (P8) -- (P9);
            \end{tikzpicture}
            \caption{DAG del ejercicio~\ref{ej:8}.}
            \label{fig:ej8_grafo}
        \end{subfigure}
        \caption{Figuras del ejercicio~\ref{ej:8}.}
    \end{figure}

    Supón que queremos realizar la sincronización indicada en el grafo, usando para ello llamadas
    desde cada rutina a dos procedimientos (\verb|EsperarPor| y \verb|Acabar|). Se dan los siguientes hechos:
    \begin{itemize}
        \item El procedimiento \verb|EsperarPor(i)| es llamado por una rutina cualquiera (la número $k$) para esperar a que termine la rutina número $i$, usando espera ocupada. Por tanto, se usa por la rutina $k$ al inicio para esperar la terminación de las otras rutinas que corresponda según el grafo.
        \item El procedimiento \verb|Acabar(i)| es llamado por la rutina número $i$, al final de la misma, para indicar que dicha rutina ya ha finalizado.
        \item Ambos procedimientos pueden acceder a variables globales en memoria compartida.
        \item Las rutinas se sincronizan única y exclusivamente mediante llamadas a estos procedimientos, siendo la implementación de los mismos completamente transparente para las rutinas.
    \end{itemize}
    Escribe una implementación de \verb|EsperarPor| y \verb|Acabar| (junto con la declaración e inicialización de las variables compartidas necesarias) que cumpla con los requisitos dados.\\

    Usaremos para ello un vector de variables booleanas \verb|finalizado|, donde \verb|finalizado[i]| indica si la rutina $i$ ha finalizado o no. Inicialmente estará inicializado a \verb|false|, puesto que ningún proceso ha finalizado. Es decir:
    \begin{minted}{pascal}
        var finalizado : array[1..9] of boolean := (false,...,false) ;
    \end{minted}

    El procedimiento \verb|EsperarPor| se implementa de la siguiente forma:
    \begin{minted}{pascal}
        procedure EsperarPor( i : integer ) ;
        begin
            while not finalizado[i] do
            begin
                { no hacer nada }
            end;
        end.
    \end{minted}
    % // TODO: ¿Cuando es procedure y cuándo es process?

    Respecto a la función \verb|Acabar|, podríamos pensar que tan solo se necesita cambiar el valor de la variable \verb|finalizado[i]| a \verb|true|. No obstante, hemos de tener en cuenta que este programa de ejecuta de forma repetida, por lo que si no se reinicia el vector \verb|finalizado| al final de cada ejecución, en la siguiente ejecución ya no habrá la sincronización necesaria. Por tanto, tras el fin del proceso \verb|P9|, se reinicia el vector \verb|finalizado| a \verb|false|.
    \begin{minted}{pascal}
        procedure Acabar( i : integer ) ;
        begin
            finalizado[i] := true ;
            if i >= 9 then
                finalizado := (false,...,false) ;
        end
    \end{minted}
    
\end{ejercicio}


\begin{ejercicio}
    En el ejercicio~\ref{ej:8} los procesos \verb|P1|, \verb|P2|, \ldots, \verb|P9| se ponen en marcha usando \verb|cobegin-coend|.
    Escribe un programa equivalente, que ponga en marcha todos los procesos, pero que use declaración
    estática de procesos, usando un vector de procesos \verb|P|, con índices desde 1 hasta 9, ambos incluidos. El proceso \verb|P[n]| contiene una secuencia de instrucciones desconocida, que llamamos \verb|S_n|, y además debe incluir las llamadas necesarias a \verb|Acabar| y \verb|EsperarPor| (con la misma implementación que antes) para lograr la sincronización adecuada. Se incluye aquí una plantilla:
    \begin{minted}{pascal}
        Process P[ n : 1..9 ]
        begin
            ..... { esperar (si es necesario) a los procesos que corresponda }
            S_n ; { sentencias especificas de este proceso (desconocidas) }
            ..... { senalar que hemos terminado }
        end
    \end{minted}

    En primer lugar, hemos de especificar, en función de $n$, a qué procesos ha de esperar cada uno, lo cual se hará mediante una matriz compartida. Tenemos que:
    \begin{minted}{pascal}
        var espera : array[1..9,1..2] of integer := (
            (-1, -1),   { P1 }
            (1, -1),    { P2 }
            (2, -1),    { P3 }
            (1, -1),    { P4 }
            (2, 4),     { P5 }
            (3, 5),     { P6 }
            (4, -1),    { P7 }
            (5, 7),     { P8 }
            (6, 8)      { P9 }
        ) ;
        Process P[ n : 1..9 ]
        begin
            for i := 1 to 2 do
                if espera[n,i] <> -1 then { != -1 }
                    EsperarPor( espera[n,i] ) ;
            
            S_n ;
            Acabar( n ) ;
        end
    \end{minted}
    % // TODO: Sintaxis en Pascal? Relevante?
\end{ejercicio}

\begin{ejercicio}
    Para los siguientes fragmentos de código, obtener la \emph{poscondición} adecuada para convertirlo en un triple demostrable con la Lógica de Programas:
    \begin{enumerate}
        \item $\{i < 10\} \quad i = 2 \ast i + 1 \quad \{ \}$\\
            Obtenemos la poscondición de este triple razonando matemáticamente:
            \begin{align*}
                i &< 10 \\
                2\ast i &< 20 \\
                i' = 2\ast i + 1 &< 21
            \end{align*}
            donde hemos notado por $i'$ al nuevo valor que adopta la variable $i$.\\

            Por tanto, la poscondición del triple es: $i < 21$.

            Pasamos ahora a demostrar el triple siguiente:
            \begin{equation*}
               \{i < 10\} \quad i = 2 \ast i + 1 \quad \{i<21\}
            \end{equation*}
            
            Usando el axioma de asignación, tenemos que:
            \begin{equation*}
                \{i<21\}^i_{2\ast i + 1}\quad i=2\ast i + 1\quad \{i<21\}
            \end{equation*}

            No obstante, de la definición de Sustitución Textual, tenemos que:
            \begin{equation*}
                \{i<21\}^i_{2\ast i + 1} \equiv \{2\ast i + 1 < 21\} \equiv \{i < 10\} 
            \end{equation*}
            
            Uniendo ambas ecuaciones, obtenemos que el triple es cierto.
            \begin{equation*}
                \{i<10\}\quad i=2\ast i + 1\quad \{i<21\} \\
            \end{equation*}

        \item $\{i > 0\} \quad i = i - 1; \quad \{ \}$\\
            Obtenemos la poscondición de este triple razonando matemáticamente:
            \begin{align*}
                i &> 0 \\
                i' = i - 1 &> -1
            \end{align*}
            donde hemos notado por $i'$ al nuevo valor que adopta la variable $i$.\\

            Por tanto, la poscondición del triple es: $i > -1$.

            Pasamos ahora a demostrar el triple siguiente:
            \begin{equation*}
                \{i > 0\} \quad i = i - 1; \quad \{i > -1\}
            \end{equation*}

            Usando el axioma de asignación, tenemos que:
            \begin{equation*}
                \{i > -1\}^i_{i-1}\quad i=i-1\quad \{i>-1\}
            \end{equation*}

            No obstante, de la definición de Sustitución Textual, tenemos que:
            \begin{equation*}
                \{i > -1\}^i_{i-1} \equiv \{i-1 > -1\} \equiv \{i > 0\}
            \end{equation*}

            Uniendo ambas ecuaciones, obtenemos que el triple es cierto.
            \begin{equation*}
                \{i > 0\}\quad i = i - 1\quad \{i > -1\}
            \end{equation*}

        \item $\{i > j\} \quad i = i + 1;~j = j + 1 \quad \{ \}$\\
            De forma matemática y notando por $i'$ y $j'$ a las modificaciones de $i$ y $j$, respectivamente:
            \begin{align*}
                i &> j \\
                i' = i+1 &> j+1=j'\\
                i' &> j'
            \end{align*}

            Por tanto, la poscondición del triple es: $i > j$. Pasamos a demostrar el triple:
            \begin{equation*}
                \{i > j\} \quad i = i + 1;~j = j + 1 \quad \{i>j\}
            \end{equation*}

            Usando la regla de la composición, tenemos que:
            \begin{equation*}
                \dfrac{\{P\}S_1\{Q\}, \{Q\}S_2\{R\}}{\{P\}S_1;S_2\{R\}} \\
            \end{equation*}

            Identificando $Q$ con $i>j+1$, tenemos que bastará con probar los triples:
            \begin{enumerate}
                \item $\{i>j\}\ i=i+1\ \{i>j+1\}$
                
                Mediante el axioma de asignación, tenemos que:
                \begin{equation*}
                    \{i>j\} \equiv \{i+1>j+1\} \equiv \{i>j+1\}^i_{i+1}\ i=i+1\ \{i>j+1\}
                \end{equation*}

                \item $\{i>j+1\}\ j=j+1\ \{i>j\}$
                
                Mediante el axioma de asignación, tenemos que:
                \begin{equation*}
                    \{i>j+1\} \equiv \{i>j\}^j_{j+1}\ j=j+1\ \{i>j\}
                \end{equation*}
            \end{enumerate}

            Como ambos son ciertos, el triple que queríamos demostrar también lo es gracias a la regla de composición.
        \item $\{\text{falso}\} \quad a = a + 7; \quad \{ \}$\\
            En este caso, partimos de un estado del programa inalcanzable, por lo que en la poscondición podemos poner cualquier estado del programa, es decir, $\{\text{verdad}\}$.
        \item $\{\text{verdad}\} \quad i = 3;~j = 2 \ast i \quad \{ \}$\\
            Como partimos de cualquier estado del programa y sólo se realizan asignaciones, es fácil intuir cuál será la poscondición:
            \begin{gather*}
                i = 3 \\
                j = 2\ast i = 2\ast 3 = 6
            \end{gather*}
            Pasamos a demostrar el triple
            \begin{equation*}
                \{\text{verdad}\} \quad i = 3;~j = 2 \ast i \quad \{i=3\ \land\ j=6\}
            \end{equation*}
            Usando la regla de composición, nos será suficiente probar los triples:
            \begin{equation*}
                \{\text{verdad}\}\ i=3\ \{i=3\} \qquad \{i=3\}\ j=2\ast i\ \{i=3\ \land\ j=6\}
            \end{equation*}
            \begin{enumerate}
                \item Para el primer triple, usamos el axioma de asignación:
                \begin{equation*}
                    \{\text{verdad}\} \equiv \{3=3\} \equiv \{i=3\}^i_3\ i=3\ \{i=3\}
                \end{equation*}
                \item Para el segundo, volvemos a usar el axioma de asignación:
                \begin{equation*}
                    \{i=3\ \land\ j=6\}^j_{2\ast i}\ j=2\ast i\ \{i=3\ \land\ j=6\}
                \end{equation*}

                No obstante, de la definición de Sustitución Textual, tenemos que:
                \begin{multline*}
                    \{i=3\ \land\ j=6\}^j_{2\ast i} \equiv \{i=3\ \land\ 2\ast i=6\} \equiv \{i=3\ \land 2\ast 3=6\} \equiv \\ \equiv \{i=3\ \land\ 6=6\} \equiv \{i=3\}
                \end{multline*}

                Uniendo ambas ecuaciones, obtenemos que el triple es cierto.
            \end{enumerate}
            Ambos triples son ciertos, luego por la regla de la composición tenemos demostrado nuestro triple.
        \item $\{\text{verdad}\} \quad c = a + b;~c = \nicefrac{c}{2} \quad \{ \}$\\
            Notando por $c'$ al nuevo valor de $c$, tenemos que:
            \begin{align*}
                c &= a + b \\
                c' = \nicefrac{c}{2} &= \frac{a+b}{2}
            \end{align*}
            Tratamos por tanto de probar el siguiente triple
            \begin{equation*}
                \{\text{verdad}\} \quad c = a + b;~c = \nicefrac{c}{2} \quad \left\{c=\frac{a+b}{2}\right\}
            \end{equation*}
            Usando la regla de la composición, basta con probar los triples:
            \begin{equation*}
                \{\text{verdad}\}\ c=a+b;\ \{c=a+b\} \qquad \{c=a+b\}\ c=\nicefrac{c}{2};\ \left\{c=\frac{a+b}{2}\right\}
            \end{equation*}
            \begin{enumerate}
                \item Para el primero, usamos el axioma de asignación:
                    \begin{equation*}
                        \{\text{verdad}\}\equiv\{a+b=a+b\}\equiv \{c=a+b\}^c_{a+b}\ c=a+b\ \{c=a+b\}
                    \end{equation*}
                \item Para la segunda, también usamos el axioma de asignación:
                    \begin{equation*}
                        \{c=a+b\}\equiv \left\{\frac{c}{2} = \frac{a+b}{2}\right\} \equiv \left\{c=\frac{a+b}{2}\right\}^c_{\nicefrac{c}{2}}\ c=\nicefrac{c}{2}\ \left\{c=\frac{a+b}{2}\right\}
                    \end{equation*}
            \end{enumerate}
            Usando la relga de composición, tenemos demostrado nuestro triple.
    \end{enumerate}
\end{ejercicio}

\begin{ejercicio}
    ¿Cuáles de los siguientes triples no son demostrables con la Lógica de Programas?
    (Considerando que $i,x,a \in \mathbb{Z}$)
    \begin{enumerate}
        \item $\{i > 0\} \quad i = i - 1; \quad \{i \geq 0\}$\\
            El siguiente triple sabemos que es cierto:
            \begin{equation*}
                \{i>0\}\ i=i-1\ \{i>-1\} 
            \end{equation*}
            $\{i>-1\}\rightarrow \{i\geq 0\}$, luego es cierto por la primera regla de la consecuencia.
        \item $\{x \geq 7\} \quad x = x + 3; \quad \{x \geq 9\}$\\
            El siguiente triple sabemos que es cierto:    
            \begin{equation*}
                \{x\geq 7\}\ x=x+3\ \{x\geq 10\}
            \end{equation*}
            $\{x\geq 10\}\rightarrow\{x\geq 9\}$, luego es cierto por la primera regla de la consecuencia.
        \item $\{i < 9\} \quad i = 2 \ast i + 1; \quad \{ i \leq 20\}$\\
            El siguiente triple sabemos que es cierto:
            \begin{equation*}
                \{i<9\}\ i=2\ast i+1\ \{i<19\}
            \end{equation*}
            $\{i<19\}\rightarrow\{i\leq 20\}$, luego es cierto por la primera regla de la consecuencia.
        \item $\{a > 0\} \quad a = a - 7; \quad \{a > -6\}$\\
            \begin{equation*}
                \{a>0\}\ a=a-7\ \{a>-7\}
            \end{equation*}
            Pero $\{a>-7\}\not\rightarrow\{a>-6\}$, luego este triple no es demostrable.
    \end{enumerate}
\end{ejercicio}

\begin{ejercicio}
    Si el triple $\{P\} C \{Q\}$ es demostrable, indicar por qué los siguientes triples también lo son (o no se pueden demostrar y por qué):
    \begin{enumerate}
        \item $\{P\} C \{Q \lor P\}$\\
           Es demostrable, ya que $\{Q\}\rightarrow\{Q \lor P\}$ y por la primera regla de la consecuencia, tomando $R=Q \lor P$:
            \begin{equation*}
                \dfrac{\{P\}C\{Q\}, \{Q\}\rightarrow\{R\}}{\{P\}C\{R\}}
            \end{equation*}
            Tenemos que se debilita la poscondición.

        \item $\{P \land D\} C \{Q\}$\\
            Es demostrable, ya que $\{P \land D\}\rightarrow\{P\}$ y por la segunda regla de la consecuencia, tomando $R = P \land D$:
            \begin{equation*}
                \dfrac{\{P\}\rightarrow\{R\},\{R\}C\{Q\}}{\{P\}C\{Q\}}
            \end{equation*}
            Tenemos que se fortalece la precondición.
        \item $\{P \lor D\} C \{Q\}$\\
            No es demostrable, porque se debilita la precondición.

        \item $\{P\} C \{Q \lor D\}$\\
            Al igual que hemos hecho en el apartado 1, es demostrable ya que ${\{Q\}\rightarrow\{Q \lor D\}}$ y usando la primera regla de la consecuencia.
        \item $\{P\} C \{Q \land P\}$\\
            No podemos demostrarlo, ya que se fortalece la poscondición.
    \end{enumerate}
\end{ejercicio}

\begin{ejercicio}
    Si el triple $\{P\} C \{Q\}$ es demostrable, ¿cuál de los siguientes triples no se puede demostrar?
    \begin{enumerate}
        \item $\{P \land D\} C \{Q\}$\\
            Sabemos que $\{P \land D\}\rightarrow\{P\}$, luego puede demostrarse por la segunda regla de la consecuencia (se fortalece la precondición).
        \item $\{P \lor D\} C \{Q\}$\\
            No puede demostrarse, porque se debilita la precondición.
        \item $\{P\} C \{Q \lor D\}$\\
            Puede demostrarse mediante la primera regla de la consecuencia, ya que se tiene que ${\{Q\}\rightarrow\{Q \lor D\}}$.

        \item $\{P\} C \{Q \lor P\}$\\
        Puede demostrarse mediante la primera regla de la consecuencia, ya que se tiene que ${\{Q\}\rightarrow\{Q \lor P\}}$.
    \end{enumerate}
\end{ejercicio}

\begin{ejercicio}
    Dado el siguiente programa, obtener:
    \begin{minted}{pascal}
        int x = 5, y = 2;
        cobegin
            < x = x + y >;
            < y = x * y >;
        coend
    \end{minted}
    \begin{enumerate}
        \item Valores finales de $x$ e $y$.
        Tenemos dos posibles trazas de ejecución:
        \begin{enumerate}
            \item Primero se ejecuta la primera instrucción, por lo que obtendríamos $x = 7$ y $y = 14$.
            \item Primero se ejecuta la segunda instrucción, por lo que obtendríamos $x = 15$ y $y = 10$.
        \end{enumerate}
        \item Valores finales de $x$ e $y$ si quitamos los símbolos \verb|< >| de instrucción atómica.
        
        Encontramos cada uno de los dos estados anteriores, además de $x = 7$ y $y= 10$.
    \end{enumerate}
\end{ejercicio}

\begin{ejercicio}
    Comprobar si la demostración del siguiente triple interfiere con los teoremas siguientes:
    \[
        \{x \geq 2\} \quad < x = x - 2 > \quad \{x \geq 0\}
    \]
    Es decir, queremos comprobar si $R\equiv <x = x-2>$ con $pre(R) = \{x\geq 2\}$ interfiere con los triples siguientes:
    \begin{enumerate}
        \item $\{x \geq 0\} \quad < x = x + 3 > \quad \{x \geq 3\}$\\
            Comprobamos en primer lugar su interferencia con la precondición:
            \begin{equation*}
                \{x\geq 0 \land x \geq 2\}\ <x=x-2>\ \{x\geq 0\}
            \end{equation*}
            Este triple es correcto por la segunda regla de la consecuencia, luego no interfiere con la precondición.\\

            Comprobemos ahora su interferencia con la poscondición:
            \begin{equation*}
                \{x\geq 3 \land x \geq 2\}\ <x=x-2>\ \{x\geq 1\}
            \end{equation*}
            En este caso, $\{x\geq 1\}\not\rightarrow\{x\geq 3\}$, luego este triple no es demostrable y $R$ interfiere con la poscondición del triple en cuestión.

        \item $\{x \geq 0\} \quad < x = x + 3 > \quad \{x \geq 0\}$\\
            Comprobamos en primer lugar su interferencia con la precondición:
            \begin{equation*}
                \{x\geq 0 \land x \geq 2\}\ <x=x-2>\ \{x\geq 0\}
            \end{equation*}
            Este triple es correcto por la segunda regla de la consecuencia, luego no interfiere con la precondición.
            Además, como la precondición y la poscondición son iguales, $R$ tampoco interfiere con la poscondición, luego no interfiere con este triple.
        \item $\{x \geq 7\} \quad < x = x + 3 > \quad \{x \geq 10\}$\\
            Comprobamos en primer lugar su interferencia con la precondición:
            \begin{equation*}
                \{x\geq 7 \land x \geq 2\}\ <x=x-2>\ \{x\geq 5\}
            \end{equation*}
            No obstante, como $\{x\geq 5\}\not\rightarrow\{x\geq 7\}$, $R$ interfiere con la precondición de este triple.
        \item $\{y \geq 0\} \quad < y = y + 3 > \quad \{y \geq 3\}$\\
            $R$ no interfiere con este triple, ya que son variables disjuntas.
        \item $\{x \text{ es impar}\} \quad < y = x + 1 > \quad \{y \text{ es par}\}$\\
            Comprobamos en primer lugar su interferencia con la precondición:
            \begin{equation*}
                \{x \text{ es impar } \land x\geq 2\}\ <x=x-2>\ \{x \text{ es impar } \land x \geq 0\}
            \end{equation*}
            Por la $2^a$ regla de la consecuencia, como $\{x \text{ es impar } \land x \geq 0\} \rightarrow \{x \text{ es impar}\}$, tenemos que es correcto y $R$ no interfiere con la precondición.

            Comprobamos ahora su interferencia con la poscondición:
            \begin{equation*}
                \{y \text{ es par } \land x \geq 0\}\ <x=x-2>\ \{y \text{ es par } \land x \geq -2\}
            \end{equation*}
            Por la $2^a$ regla de la consecuencia, como $\{y \text{ es par } \land x \geq -2\} \rightarrow \{y \text{ es par}\}$, tenemos que es correcto y $R$ no interfiere con la poscondición. Por tanto, $R$ no interfiere con este triple.
            
    \end{enumerate}
\end{ejercicio}

\begin{ejercicio}
    Dado el siguiente triple:
    \begin{gather*}
        \{x = 0\} \\
        \text{cobegin} \\
        <x = x + a> || <x = x + b> || <x = x + c> \\
        \text{coend} \\
        \{x = a + b + c\}
    \end{gather*}
    
    Demostrarlo utilizando la lógica de asertos para cada una de las tres instrucciones atómicas y después que se llega a la poscondición final $x = a + b + c$ utilizando para ello la regla \emph{de la composición concurrente} de instrucciones atómicas.\\

Inicialmente, demostraremos los $3$ siguientes triples, uno por cada instrucción atómica. Hemos de notar que, en cada uno de ellos, como no sabemos en qué orden se ejecutan, tenemos que incluir en las precondiciones y las poscondiciones todas las posibilidades.
\begin{enumerate}
    \item El correspondiente a la primera instrucción atómica:
    \begin{gather*}
        \{x=0 \lor x=b \lor x=c \lor x=b+c\}\ <x=x+a>\ \\ \{x=a \lor x=a+b \lor x=a+c \lor x=a+b+c\}
    \end{gather*}

    Mediante el axioma de asignación, tenemos que:
    \begin{gather*}
        \{x=a \lor x=a+b \lor x=a+c \lor x=a+b+c\}_{x+a}^a <x=x+a>\\ \{x=a \lor x=a+b \lor x=a+c \lor x=a+b+c\}
    \end{gather*}

    No obstante, de la definición de Sustitución Textual, tenemos:
    \begin{gather*}
        \{x=a \lor x=a+b \lor x=a+c \lor x=a+b+c\}_{x+a}^a \equiv \\
        \equiv \{x+a=a \lor x+a=a+b \lor x+a=a+c \lor x+a=a+b+c\} \equiv \\
        \equiv \{x=0 \lor x=b \lor x=c \lor x=b+c\}
    \end{gather*}

    Por tanto, el triple en cuestión es cierto.

    \item El correspondiente a la segunda instrucción atómica:
    \begin{gather*}
        \{x=0 \lor x=a \lor x=c \lor x=a+c\}\ <x=x+b>\ \\ \{x=b \lor x=a+b \lor x=b+c \lor x=a+b+c\}
    \end{gather*}

    Es cierto, y su demostración es análoga al primer caso.
    \item El correspondiente a la tercera instrucción atómica:
    \begin{gather*}
        \{x=0 \lor x=b \lor x=a \lor x=a+b\}\ <x=x+c>\ \\ \{x=c \lor x=a+c \lor x=b+c \lor x=a+b+c\}
    \end{gather*}
    Es cierto, y su demostración es análoga al primer caso.
\end{enumerate}

Seguidamente, tenemos que ver que dichos $3$ triples están libres de interferencias. Para ello, hemos de probar $12$ triples, ya que hay $3$ instrucciones atómicas, cada una de ellas con $2$ asertos, por lo que por cada instrucción hemos de comprobar $4$ asertos:
\begin{align*}
    &NI(x=0 \lor x=a \lor x=c \lor x=a+c, <x=x+a>)\\
    &NI(x=b \lor x=a+b \lor x=b+c \lor x=a+b+c, <x=x+a>)\\
    &NI(x=0 \lor x=b \lor x=a \lor x=a+b, <x=x+a>)\\
    &NI(x=c \lor x=a+c \lor x=b+c \lor x=a+b+c, <x=x+a>)\\ \\
    &NI(x=0 \lor x=b \lor x=c \lor x=b+c, <x=x+b>)\\
    &NI(x=a \lor x=a+b \lor x=a+c \lor x=a+b+c, <x=x+b>)\\
    &NI(x=0 \lor x=b \lor x=a \lor x=a+b, <x=x+b>)\\
    &NI(x=c \lor x=a+c \lor x=b+c \lor x=a+b+c, <x=x+b>)\\ \\
    &NI(x=0 \lor x=b \lor x=c \lor x=b+c, <x=x+c>)\\
    &NI(x=a \lor x=a+b \lor x=a+c \lor x=a+b+c, <x=x+c>)\\
    &NI(x=0 \lor x=a \lor x=c \lor x=a+c, <x=x+c>)\\
    &NI(x=b \lor x=a+b \lor x=b+c \lor x=a+b+c, <x=x+c>)\\
\end{align*}

Demostremos ahora el primero, ya que el resto son idénticos.
\begin{align*}
    NI&(x=0 \lor x=a \lor x=c \lor x=a+c, <x=x+a>)\equiv \\
    &\equiv \{(x=0 \lor x=a \lor x=c \lor x=a+c) \land (x=0 \lor x=b \lor x=c \lor x=b+c)\}\\&\qquad <x=x+a> \{x=0 \lor x=a \lor x=c \lor x=a+c\}\equiv \\
    &\equiv \{x=0 \lor x=c\} <x=x+a> \{x=0 \lor x=a \lor x=c \lor x=a+c\}
\end{align*}
Este triple efectivamente es cierto, lo cual se puede demostrar empleando en primer lugar el Axioma de Sustitución y, posteriormente, la primera regla de la consecuencia.\\

Por tanto, y tras aplicar la Regla de la Composición Concurrente, tenemos de forma directa que:
    \begin{gather*}
        \{x = 0\} \\
        \text{cobegin} \\
        <x = x + a> || <x = x + b> || <x = x + c> \\
        \text{coend} \\
        \{x = a + b + c\}
    \end{gather*}
\end{ejercicio}

\begin{ejercicio}
    El siguiente triple:
    \begin{gather*}
        \{x=0 \land y = 0 \land z = 0\} \\
        <x = z + a> || <y = x+b> \\
        \{(x = a) \land (y = b \lor y = a +b) \land z = 0\}
    \end{gather*}
    \begin{enumerate}[label=(\alph*)]
        \item Es indemostrable salvo que se cumpla siempre que $a=0$.
        \item El triple anterior es demostrable para cualquier valor de las variables $a$ o $b$.
        \item Es indemostrable salvo que se cumpla siempre que $b = 0$.
        \item Es indemostrable salvo que se cumpla siempre que $a = 0 \land b = 0$.
    \end{enumerate}

    Veamos si podemos demostrarlo. Para ello, notamos cada instrucción atómica de la siguiente forma:
    \begin{align*}
        S_1 &= <x = z + a> \\
        S_2 &= <y = x + b>
    \end{align*}

    Veamos cuál ha de ser la precondición de cada instrucción atómica:
    \begin{align*}
        P_1 &= x=0 \land (y=0 \lor y=b) \land z=0 \\
        P_2 &= (x=0\lor x=a) \land y=0 \land z=0
    \end{align*}

    Veamos cuál ha de ser la poscondición de cada instrucción atómica:
    \begin{align*}
        Q_1 &= x=a\land (y=0 \lor y=b)\land z=0 \\
        Q_2 &= [(x=0\land y=b) \lor (x=a\land y=a+b)] \land z=0
    \end{align*}

    Vemos por tanto que ambos triples son ciertos:
    \begin{enumerate}
        \item $\{P_1\} S_1 \{Q_1\}$
        
        De la definición de Sustitución Textual, tenemos que:
        \begin{multline*}
            \{x=a \land (y=0 \lor y=b) \land z = 0\}_{z+a}^x \equiv \{z+a=a \land (y=0 \lor y=b) \land z = 0\} \equiv\\\equiv \{(y=0 \lor y=b) \land z=0\}
        \end{multline*}

        Por tanto, del Axioma de Asignación, tenemos que:
        \begin{equation*}
            \{(y=0 \lor y=b) \land z = 0\} <x = z + a> \{x=a \land (y=0 \lor y=b) \land z = 0\}
        \end{equation*}

        Finalmente, usando la segunda regla de la consecuencia, como se tiene que $\{P_1\}\rightarrow\{(y=0 \lor y=b) \land z = 0\}$, tenemos que el triple es cierto.
        
        \item $\{P_2\} S_2 \{Q_2\}$
        
        De la definición de Sustitución Textual, tenemos que:
        \begin{multline*}
            \{[(x=0\land y=b) \lor (x=a\land y=a+b)] \land z=0\}_{x+b}^y \equiv \\ \equiv \{[(x=0\land x+b=b) \lor (x=a\land x+b=a+b)] \land z = 0\} \equiv \\ \equiv \{(x=0\lor x=a) \land z = 0\}
        \end{multline*}

        Por tanto, del Axioma de Asignación, tenemos que:
        \begin{multline*}
            \{(x=0\lor x=a) \land z = 0\} <y = x + b> \\ \{[(x=0\land y=b) \lor (x=a\land y=a+b)] \land z=0\}
        \end{multline*}

        Por tanto, usando la segunda regla de la consecuencia, como se tiene que $\{P_2\}\rightarrow\{(x=0\lor x=a) \land z = 0\}$, tenemos que el triple es cierto.
    \end{enumerate}

    Veamos ahora que no interfieren entre sí. Como tenemos dos instrucciones atómicas, cada una con dos asertos, hemos de comprobar $4$ asertos:
    \begin{align*}
        NI(P_2, S_1) & \equiv \{P_1\land P_2\} S_1 \{P_2\} \\
        NI(Q_2, S_1) & \equiv \{P_1\land Q_2\} S_1 \{Q_2\} \\
        NI(P_1, S_2) & \equiv \{P_2\land P_1\} S_2 \{P_1\} \\
        NI(Q_1, S_2) & \equiv \{P_2\land Q_1\} S_2 \{Q_1\}
    \end{align*}

    Veamos por tanto en qué queda cada uno de ellos:
    \begin{align*}
        NI(P_2, S_1) & \equiv \{x=0 \land y=0 \land z=0\} <x=z+a> \{(x=0\lor x=a) \land y=0 \land z=0\} \\
        NI(Q_2, S_1) & \equiv \{x=0 \land y=b \land z=0\} <x=z+a>\\&\hspace{5cm} \{[(x=0\land y=b) \lor (x=a\land y=a+b)] \land z=0\} \\
        NI(P_1, S_2) & \equiv \{x=0 \land y=0 \land z=0\} <y=x+b> \{x=0 \land (y=0\lor y=b) \land z=0\} \\
        NI(Q_1, S_2) & \equiv \{x=a \land y=0 \land z=0\} <y=x+b> \{x=a \land (y=0\lor y=b) \land z=0\}
    \end{align*}

    Intentemos demostrar el segundo. Usando la definición de Sustitución Textual, tenemos que:
    \begin{multline*}
        \{[(x=0\land y=b) \lor (x=a\land y=a+b)] \land z=0\}_{z+a}^x \equiv \\ \equiv \{[(z+a=0\land y=b) \lor (z+a=a\land y=a+b)] \land z=0\} \equiv \\ \equiv \{y=a+b \land z=0\}
    \end{multline*}
    Por tanto, del Axioma de Asignación, tenemos que:
    \begin{multline*}
        \{y=a+b \land z=0\} <x=z+a> \{[(x=0\land y=b) \lor (x=a\land y=a+b)] \land z=0\}
    \end{multline*}
    No obstante, de forma general, tenemos que $\{x=0\land y=b\land z=0\}\not\rightarrow\{y=a+b \land z=0\}$, por lo que $NI(Q_2, S_1)$ no es demostrable. No obstante, si $a=0$, entonces sí que sería demostrable. Supongamos por tanto a partir de ahora que \ul{$a=0$}.

    Intentemos ahora demostrar el cuarto. Usando la definición de Sustitución Textual, tenemos que:
    \begin{multline*}
        \{x=a \land (y=0 \lor y=b) \land z = 0\}_{x+b}^y \equiv \\ \equiv \{x=a \land (x+b=0 \lor x+b=b) \land z = 0\} \equiv \\ \equiv \{x=a \land (x=-b\lor x=0) \land z = 0\} \equiv 
        \{x=0 \land z = 0\}
    \end{multline*}
    donde en la última igualdad hemos usado que \ul{$a=0$}. Por tanto, del Axioma de Asignación, tenemos que:
    \begin{multline*}
        \{x=0 \land z = 0\} <y=x+b> \{x=a \land (y=0 \lor y=b) \land z = 0\}
    \end{multline*}
    Además, como \ul{$a=0$}, tenemos que $\{x=0\land z=0\}\rightarrow\{x=a \land y=0 \land z=0\}$, tenemos que es cierto. Por tanto, $NI(P_1, S_2)$ es demostrable.

    Los otros dos triples son análogamente ciertos, por lo que podemos aplicar la regla de la composición concurrente y llegar al siguiente triple:
    \begin{gather*}
        \{x=0 \land y = 0 \land z = 0\} \\
        <x = z + a> || <y = x+b> \\
        \{x = a \land  y = a +b \land z = 0\}
    \end{gather*}
    Este es el triple que queríamos demostrar, suponiendo que \ul{$a=0$}. Por tanto, la respuesta correcta es la \textbf{a)}, ya que no es demostrable salvo que se cumpla siempre que $a=0$.

    % // TODO: JJ Revisar. Me ha parecido bastante difícil
\end{ejercicio}

\begin{ejercicio}
    Suponer que $\{suma>1\}\ suma=suma+4\ \{suma>5\}$ es demostrable, entonces: ¿cuál de los siguientes triples es también demostrable? (indicar por qué)
    \begin{enumerate}
        \item $\{suma>2\}\ suma=suma+4\ \{suma>5\}$.\\
            Es demostrable, ya que $\{suma>2\}\rightarrow\{suma>1\}$ y podemos aplicar la segunda regla de la consecuencia.
        \item $\{suma\geq 1\}\ suma=suma+4\ \{suma > 5\}$.\\
            No es demostarble, ya que debilita la precondición.
        \item $\{suma>0\}\ suma=suma+4\ \{suma > 5\}$.\\
            No es demostarble, ya que debilita la precondición.
        \item $\{suma>1\}\ suma=suma+4\ \{suma > 6\}$.\\
            No es demostrable, ya que fortalece la poscondición.
    \end{enumerate}
\end{ejercicio}

\begin{ejercicio}
    Suponer que $\{x<y\}\ C_1\ \{u<v\}$ es demostrable, entonces: ¿cuáles de los siguientes triples son también demostrables? (indicar por qué)
    \begin{enumerate}
        \item $\{x\leq y\}\ C_1\ \{u<v\}$.\\
            No es demostrable, ya que debilita la precondición.
        \item $\{x\leq y-2\}\ C_1\ \{u<v\}$.\\
            Es demostrable, ya que $\{x\leq y-2\}\rightarrow\{x+2\leq y\}\rightarrow\{x<y\}$, y mediante la segunda regla de la consecuencia se tiene que es cierto.
        \item $\{x\leq y\}\ C_1\ \{u\leq v\}$.\\
            El triple $\{x<y\}\ C_1\ \{u\leq v\}$ sí que es demostrable ya que relaja la poscondición, pero el triple que se nos dice no es demostrable, ya que también relaja la precondición.
            Como además no tenemos relación entre $x$ y $u$ ni entre $y$ y $v$, no podemos inferir nada.
        \item $\{x<y\}\ C_1\ \{u<v-2\}$.\\
            No es demostrable, ya que fortalecemos la poscondición.
    \end{enumerate}
\end{ejercicio}

\begin{ejercicio}
    Seleccionar el valor correcto de las 2 variables ($x$ e $y$) después de ejecutarse el siguiente programa concurrente:
    \begin{minted}{pascal}
        int x=5, y=2;
        cobegin <x=x+y>; <y=x*y>; <x=x-y>; coend;
    \end{minted}
    \begin{enumerate}[label=(\alph*)]
        \item $x=7$ y $y=14$.
        \item $x=5$ y $y=10$.
        \item $x=-7$ y $y=14$.
        \item $x=-3$ y $y=10$.
    \end{enumerate}

    Numeramos las instrucciones atómicas de la siguiente forma:
    \begin{enumerate}
        \item \verb|<x=x+y>|
        \item \verb|<y=x*y>|
        \item \verb|<x=x-y>|
    \end{enumerate}

    Veamos ahora, en función del orden de ejecución, cuál sería el valor de las variables $x$ e $y$:
    \begin{itemize}
        \item \ul{1, 2, 3}: $x=-7$ y $y=14$.
        \item \ul{1, 3, 2}: $x=5$ y $y=10$.
        \item \ul{3, 1, 2}: $x=5$ y $y=10$.
        \item \ul{2, 1, 3}: $x=5$ y $y=10$.
        \item \ul{2, 3, 1}: $x=5$ y $y=10$.
        \item \ul{3, 2, 1}: $x=9$ y $y=6$.
    \end{itemize}

    Por tanto, las respuestas $b$ y $c$ son correctas.
    % // TODO: Lo ha hecho Arturo. Revisar JJ
\end{ejercicio}

\begin{ejercicio}
    El siguiente código concurrente no puede ser demostrado directamente con la lógica de aserciones (pre y poscondiciones). Elegir la respuesta que explica correctamente la razón de que ocurra esto.
    \begin{minted}{pascal}
        {x=0} cobegin <x=x+a>; <x=x+a> coend; {x=2*a}
        (a es un valor entero positivo)
    \end{minted}
    \begin{enumerate}[label=(\alph*)]
        \item Porque la poscondición que se propone $\{x=2\ast a\}$ es falsa.
        \item Porque falta incluir la posibilidad de que el valor final de $x$ sea también ${\{x=a\}}$.
        \item Porque al aplicar directamente la regla de inferencia de la \textit{composición concurrente} utilizo unas condiciones (pre y post-condiciones) demasiado débiles.
        \item Porque tengo que incluir en los asertos el valor del contador de programa de cada procesador.
    \end{enumerate}

    Notamos cada instrucción atómica de la siguiente forma:
    \begin{equation*}
        S_1 = S_2 = <x=x+a>
    \end{equation*}

    Veamos cuál ha de ser la precondición de cada instrucción atómica:
    \begin{equation*}
        P_1 = P_2 = x=0 \lor x=a
    \end{equation*}

    Veamos cuál ha de ser la poscondición de cada instrucción atómica:
    \begin{equation*}
        Q_1 = Q_2 = x=a \lor x=2a
    \end{equation*}

    Veamos ahora que cada triple es cierto. Como son los mismos, hemos de demostrar:
    \begin{equation*}
        \{x=0 \lor x=a\} <x=x+a> \{x=a \lor x=2a\}
    \end{equation*}
    
    Este se demuestra de forma directa. Además, también hemos de demostrar que no interfieren entre sí. Tenemos que demostrar:
    \begin{align*}
        NI(P_1, S_2) &\equiv \{x=0 \lor x=a\} <x=x+a> \{x=0 \lor x=a\} \\
        NI(Q_1, S_2) &\equiv \{x=a\} <x=x+a> \{x=a \lor x=2a\}
    \end{align*}

    Estos también son ciertos, por lo que podemos aplicar la regla de la composición concurrente y llegar al siguiente triple:
    \begin{equation*}
        \{x=0 \lor x=a\} cobegin <x=x+a>; <x=x+a> coend; \{x=a\lor x=2a\}
    \end{equation*}

    Por la primera regla de la consecuencia, podemos debilitar la precondición, llegando al siguiente triple:
    \begin{equation*}
        \{x=0\} cobegin <x=x+a>; <x=x+a> coend; \{x=a\lor x=2a\}
    \end{equation*}
    
    No obstante, la poscondición no se puede debilitar, por lo que la respuesta correcta es la \textbf{b)}, ya que falta incluir en los asertos el valor final de $x$ sea también $\{x=a\}$.

    % // TODO: Preguntar a Capel
    % {x=0} cobegin <x=x+a>; <x=x+a> coend; {x=2*a}       no se puede demostrar
    % {x=0} cobegin <x=x+a>; <x=x+a> coend; {x=2*a  o   x=a}    es cierto
\end{ejercicio}

\begin{ejercicio}
   Estudiar cuáles son los valores finales de las variables \verb|x| e \verb|y| en el siguiente programa. Insertar los asertos adecuados entre llaves, antes y después de cada sentencia, para poder obtener una traza de demostración del programa, que incluya en su último aserto los valores finales de las variables. 
   \begin{minted}{c++}
       int x = c1;
       int y = c2;
       x = x + y;
       y = x * y;
       x = x - y;
   \end{minted}

   Tenemos que cada triple, por orden, es:
   \begin{gather*}
       \{x = c_1 \land y = c_2\}\\
       x = x + y \\
       \{x = c_1 + c_2 \land y =c_2\} \\
       y = x * y \\
       \{x = c_1 + c_2 \land y = (c_1 + c_2)\cdot c_2 \} \\
       x = x - y \\
       \{x = (c_1 + c_2) - (c_1 + c_2)\cdot c_2 = (c_1 + c_2)\cdot (1-c_2) \land y = (c_1 + c_2)\cdot c_2 \}
   \end{gather*}
\end{ejercicio}

\begin{ejercicio}
    Demostrar que el siguiente triple es cierto:
    \begin{gather*}
        \{x=0\}\\ 
        \texttt{cobegin} \\
        <x=x+1> || <x=x+2> || <x=x+4> \\
        \texttt{coend} \\
        \{x=7\}
    \end{gather*}

    Notamos cada instrucción atómica de la siguiente forma:
    \begin{align*}
        S_1 &= <x=x+1> \\
        S_2 &= <x=x+2> \\
        S_3 &= <x=x+4>
    \end{align*}

    Veamos cuál ha de ser la precondición de cada instrucción atómica:
    \begin{align*}
        P_1 = x=0 \lor x=2 \lor x=4 \lor x=6 \\
        P_2 = x=0 \lor x=1 \lor x=4 \lor x=5 \\
        P_3 = x=0 \lor x=1 \lor x=2 \lor x=3
    \end{align*}

    Veamos cuál ha de ser la poscondición de cada instrucción atómica:
    \begin{align*}
        Q_1 = x=1 \lor x=3 \lor x=5 \lor x=7 \\
        Q_2 = x=2 \lor x=3 \lor x=6 \lor x=7 \\
        Q_3 = x=4 \lor x=5 \lor x=6 \lor x=7
    \end{align*}

    Cada triple es directamente cierto por el Axioma de Asignación. Veamos ahora que no interfieren entre sí. Tenemos que demostrar:
    \begin{align*}
        NI(P_2, S_1) &\equiv \{x=0 \lor x=4\} <x=x+1> \{x=0\lor x=1\lor x=4\lor x=5\} \\
        NI(Q_2, S_1) &\equiv \{x=2 \lor x=6\} <x=x+1> \{x=2\lor x=3\lor x=6\lor x=7\} \\
        NI(P_3, S_1) &\equiv \{x=0 \lor x=2\} <x=x+1> \{x=0\lor x=1\lor x=2\lor x=3\} \\
        NI(Q_3, S_1) &\equiv \{x=4 \lor x=6\} <x=x+1> \{x=4\lor x=5\lor x=6\lor x=7\} \\ \\
        NI(P_1, S_2) &\equiv \{x=0 \lor x=4\} <x=x+2> \{x=0\lor x=2\lor x=4\lor x=6\} \\
        NI(Q_1, S_2) &\equiv \{x=1 \lor x=5\} <x=x+2> \{x=1\lor x=3\lor x=5\lor x=7\} \\
        NI(P_3, S_2) &\equiv \{x=0 \lor x=1\} <x=x+2> \{x=0\lor x=1\lor x=2\lor x=3\} \\
        NI(Q_3, S_2) &\equiv \{x=4 \lor x=5\} <x=x+2> \{x=4\lor x=5\lor x=6\lor x=7\} \\ \\
        NI(P_1, S_3) &\equiv \{x=0 \lor x=2\} <x=x+4> \{x=0\lor x=2\lor x=4\lor x=6\} \\
        NI(Q_1, S_3) &\equiv \{x=1 \lor x=3\} <x=x+4> \{x=1\lor x=3\lor x=5\lor x=7\} \\
        NI(P_2, S_3) &\equiv \{x=0 \lor x=1\} <x=x+4> \{x=0\lor x=1\lor x=4\lor x=5\} \\
        NI(Q_2, S_3) &\equiv \{x=2 \lor x=3\} <x=x+4> \{x=2\lor x=3\lor x=6\lor x=7\}
    \end{align*}

    Todos estos son ciertos, por lo que podemos aplicar la regla de la composición concurrente y llegar al siguiente triple:
    \begin{gather*}
        \{x=0\} \\
        \texttt{cobegin} \\
        <x=x+1> || <x=x+2> || <x=x+4> \\
        \texttt{coend} \\
        \{x=7\}
    \end{gather*}

    % // TODO: JJ Revisar
    % // Este ha sido fácil. Intuición? Cómo verlo directo sin escribir todo?
\end{ejercicio}

\begin{ejercicio}
    Dada la siguiente construcción de composición concurrente \verb|P|:
    \begin{gather*}
        \texttt{cobegin} \\
        <x=x-1>;<x=x+1>;||<y=y-1>;<y=y+1>; \\
        \texttt{coend} \\
    \end{gather*}
    demostrar que se cumple la invarianza de $\{x=y\}$, es decir, que $\{x=y\}\ P\ \{x=y\}$ es un triple cierto.\\

Para ello, comenzamos demostrando los siguientes triples (se ha de mantener la invarianza en el código secuencial)
\begin{gather*}
    \{x=y\}\ x=x-1;x=x+1;\ \{x=y\} \\
    \{x=y\}\ y=y-1;y=y+1;\ \{x=y\} \\
\end{gather*}
\begin{enumerate}
    \item Respecto al primero, tenemos que los siguientes triples son ciertos:
        \begin{gather*}
            \{x=y\}\ x=x-1\ \{x+1=y\} \\
            \{x+1=y\}\ x=x+1\ \{x=y\}
        \end{gather*}
        Usando la regla de la composición, se tiene.
    \item Respecto al segundo, tenemos que los siguientes triples son ciertos:
        \begin{gather*}
            \{y=x\}\ y=y-1\ \{y+1=x\} \\
            \{y+1=x\}\ y=y+1\ \{y=x\}
        \end{gather*}
        Usando la regla de la composición, se tiene.
\end{enumerate}
Ahora, los triples son libres de interferencia por tener variables disjuntas. Podemos aplicar por tanto la regla de la composición concurrente, llegando a lo que queríamos probar:
\begin{gather*}
    \{x=y\} \\
    \texttt{cobegin} \\
    <x=x-1>;<x=x+1>;||<y=y-1>;<y=y+1>; \\
    \texttt{coend} \\
    \{x=y\}
\end{gather*}
\end{ejercicio}

\begin{ejercicio}
    Usando la regla de la conjunción, demostrar que 
    \begin{equation*}
        \{i>2\}\ i=2\ast i\ \{i>4\}
    \end{equation*}
    % // TODO: JJ. Regla de la conjunción?. Se demuestra directamente con el axioma de asignación
\end{ejercicio}

\begin{ejercicio}
    Se dan los siguientes triples de Hoare:
    \begin{gather*}
        \{j>1\}\ i=i+2;\ j=j+3;\ \{j>4\} \\
        \{i>2\}\ i=i+2;\ j=j+3;\ \{i>4\}
    \end{gather*}
    Demostrar que estos triples implican que
    \begin{equation*}
        \{j>1 \land i>2\}\ i=i+2;\ j=j+3\ \{j>4 \land i>4\}
    \end{equation*}
    ¿Qué regla se debe utilizar para la demostración?

    % // TODO: JJ. Cómo??
\end{ejercicio}

\begin{ejercicio}
    Sean $A$ y $B$ los valores iniciales de $a$ y $b$ respectivamente. Escribir un fragmento de código que tenga $\{a=A+B \land b =A-B\}$ como poscondición y demostrar que el código es correcto.\\

    En este caso, nos piden un código $C$ que cumpla:
    \begin{equation*}
        \{a=A \land b=B\}\ C\ \{a=A+B \land b =A-B\}
    \end{equation*}

    Sea $C=<a=a+b; b=a-2b>$. Buscamnos entonces demostrar los siguientes triples:
    \begin{gather*}
        \{a=A \land b=B\}\ a=a+b\ \{a=A+B \land b=B\} \\
        \{a=A+B \land b=B\}\ b=a-2b\ \{a=A+B \land b=A-B\}
    \end{gather*}

    Demostramos cada uno por separado:
    \begin{enumerate}
        \item Usando el axioma de asignación:
            \begin{align*}
                \{a=A+B \land b=B\}_{a+b}^a
                &\equiv \{a+b=A+B \land b=B\}
                \equiv\\&\equiv \{a=A\land b=B\}\ a=a+b\ \{a=A+B \land b=B\}
            \end{align*}
        \item Usando el axioma de asignación:
            \begin{align*}
                \{a=A+B &\land b=A-B\}_{a-2b}^b
                \equiv \{a=A+B \land a-2b=A-B\}
                \equiv\\& \equiv \{a=A+B \land A+B-2b=A-B\}
                \equiv\\& \equiv \{a=A+B \land b=B\}\ b=a-2b\ \{a=A+B \land b=A-B\}
            \end{align*}
    \end{enumerate}
    Usando la regla de la composición, tenemos que el código es correcto.

\end{ejercicio}

\begin{ejercicio}
    Demostrar que la siguiente sentencia tiene la poscondición ${\{x\geq 0, x^2\geq a^2\}}$.
    \verb|if a > 0 then x = a else x = -a|. Es decir, probar el triple:
    \begin{equation*}
        \{V\}\ \texttt{if a >\ 0 then x = a else x = -a}\ \{x\geq 0, x^2\geq a^2\}
    \end{equation*}~\\
    Para ello, tenemos que usar la regla del \verb|if|:
    \begin{equation*}
        \dfrac{\{P\land B\}S_1\{Q\}, \{P\land \lnot B\}S_2\{Q\}}{\{P\}\ \texttt{if B then}\ S_1\ \texttt{else}\ S_2\ \{Q\}}
    \end{equation*}
    Luego bastará con probar los triples
    \begin{gather*}
        \{a>0\} \equiv \{V \land a > 0\}\ x=a\ \{x\geq 0 \land x^2\geq a^2\} \\
        \{a\leq 0\} \equiv \{V \land a \leq 0\}\ x=-a\ \{x\geq 0 \land x^2\geq a^2\}
    \end{gather*}
    \begin{enumerate}
        \item Usando el axioma de asignación:
            \begin{equation*}
                \{a \geq 0\} \equiv \{a \geq 0 \land a^2 \geq a^2\} \equiv \{x\geq 0 \land x^2\geq a^2\}^x_{a}\ x=a\ \{x\geq 0 \land x^2\geq a^2\}
            \end{equation*}
            Como $\{a>0\}\rightarrow\{a\geq 0\}$, usamos la segunda regla de la consecuencia y tenemos el primer triple demostrado.
        \item Usando el axioma de asignación:
            \begin{equation*}
                \{a \leq 0\} \equiv \{-a \geq 0 \land a^2 \geq a^2\} \equiv \{x\geq 0 \land x^2\geq a^2\}^x_{-a}\ x=-a\ \{x\geq 0 \land x^2\geq a^2\}
            \end{equation*}
    \end{enumerate}
    Y acabamos de probar el triple que nos pedía el ejercicio.
\end{ejercicio}

\begin{ejercicio}
    El siguiente fragmento de código tiene $\{P\}\equiv \left\{sum = \frac{j(j-1)}{2}\right\}$ como precondición y poscondición. Demostrar que es verdadero:
    \begin{equation*}
        \{P\}\ sum=sum+j;\ j=j+1;\ \{P\}
    \end{equation*}

    Queremos demostrar el triple:
    \begin{equation*}
        \left\{sum=\frac{j(j-1)}{2} \right\}\ sum=sum+j;\ j=j+1;\ \left\{sum=\frac{j(j-1)}{2}\right\}
    \end{equation*}
    Para ello, será suficiente con demostrar los triples
    \begin{gather*}
        \left\{sum = \frac{j(j-1)}{2} \right\}\ sum = sum + j;\ \left\{sum = \frac{(j+1)j}{2} \right\} \\
        \left\{sum = \frac{(j+1)j}{2} \right\}\ j=j+1;\ \left\{sum = \frac{j(j-1)}{2}\right\}
    \end{gather*}
    y aplicar la regla de composición.

    \begin{enumerate}
        \item Para demostrar el primer triple, usamos el axioma de asignación:
            \begin{equation*}
                \left\{sum = \frac{(j+1)j}{2} \right\}^{sum}_{sum+j}\ sum = sum + j;\ \left\{sum = \frac{(j+1)j}{2} \right\} 
            \end{equation*}
            Usando la definición de Sustitución Textual, tenemos que:
            \begin{gather*}
                \left\{sum = \frac{(j+1)j}{2} \right\}^{sum}_{sum+j} \equiv \left\{sum+j = \frac{(j+1)j}{2} \right\} \equiv \\ 
                \equiv \left\{sum = \frac{(j+1)j}{2}-j \right\} \equiv \left\{sum = \frac{j(j-1)}{2} \right\}
            \end{gather*}
        \item Para el segundo, usamos también el axioma de asignación:
            \begin{equation*}
                \left\{sum = \frac{j(j-1)}{2}\right\}^j_{j+1} j=j+1;\ \left\{sum = \frac{j(j-1)}{2}\right\}
            \end{equation*}
            Usando de nuevo la definición de Sustitución Textual, tenemos que:
            \begin{gather*}
                \left\{sum = \frac{j(j-1)}{2}\right\}^j_{j+1} \equiv \left\{sum = \frac{(j+1)(j+1-1)}{2}\right\} \equiv \left\{sum = \frac{(j+1)j}{2}\right\}
            \end{gather*}
    \end{enumerate}
    Por lo que el triple del enunciado es cierto.

\end{ejercicio}

\begin{ejercicio}
    Demostrar que
    \begin{equation*}
        \{i\ast j + 2 \ast j + 3\ast i = 0\}\ j=j+3;\ i=i+2;\ \{i\ast j=6\}
    \end{equation*}

    Vamos buscando aplicar la regla de la composición. Para ello, y como desconocemos el estado intermedio por el que debemos pasar, usamos directamente la Sustitución Textual al final, para así obtener la precondición del segundo triple.
    \begin{equation*}
        \{i\ast j=6\}_{i+2}^i \equiv \{(i+2)\ast j=6\} \equiv \{i\ast j+2\ast j=6\} \equiv \{j\ast(i+2)=6\}
    \end{equation*}

    Usando esa precondición, el segundo libre se demuestra directamente con el axioma de asignación. Demostramos ahora el primer triple:
    \begin{gather*}
        \{i\ast j + 2 \ast j + 3\ast i = 0\}\ j=j+3;\ \{j\ast(i+2)=6\}
    \end{gather*}

    Usando la sustitución textual, tenemos que:
    \begin{align*}
        \{j\ast(i+2)=6\}_{j+3}^j &\equiv \{(j+3)\ast(i+2)=6\} \equiv
        \{j\ast(i+2)+3\ast(i+2)=6\} \equiv \\ &\equiv \{i\ast j + 2 \ast j + 3\ast i = 0\}
    \end{align*}

    Por lo que, tras usar la regla de la composición, vemos que el triple del enunciado es cierto.
\end{ejercicio}

\begin{ejercicio}
    ¿Por qué en la regla del \verb|while B|, la condición \verb|B| debe ser verdadera al comienzo del bucle?

    % // TODO: No entiendo a qué se refiere
\end{ejercicio}

\begin{ejercicio}
    Considerar una función con dos argumentos que se usa en un programa. Explicar por qué el uso de alias puede ser un problema en este caso.

    % // TODO: No entiendo a qué se refiere
\end{ejercicio}

\begin{ejercicio}
    Demostrar la corrección parcial del siguiente fragmento de programa:
    \begin{minted}{pascal}
        sum:= 0; j:= 1;
        while j <> c do begin {<> es !=}
            sum:= sum+j;
            j:= j+1;
        end
        {sum = c*(c-1)/2}
    \end{minted}

    Para ello, tenemos que hacer uso de la regla de la iteración: 
    \begin{equation*}
        \dfrac{\{I\land B\}S\{I\}}{\{I\}\ \texttt{while B do}\ S\ \texttt{end do}\ \{I\land \lnot B\}}
    \end{equation*}

    Identificando términos, sean:
    \begin{align*}
        I &\equiv sum = \frac{j(j-1)}{2} \\
        B &\equiv j \neq c \\
        S &= sum = sum+j;\ j = j+1
    \end{align*}
    Luego tendremos que probar que se cumple el triple
    \begin{equation*}
        \left\{sum=\frac{j(j-1)}{2} \land j\neq c\right\}\ sum=sum+j;j=j+1;\ \left\{sum=\frac{j(j-1)}{2}\right\}
    \end{equation*}
    Para ello, será suficiente con demostrar los triples
    \begin{gather*}
        \left\{sum = \frac{j(j-1)}{2} \land j\neq c\right\}\ sum = sum + j;\ \left\{sum = \frac{(j+1)j}{2} \land j \neq c\right\} \\
        \left\{sum = \frac{(j+1)j}{2} \land j\neq c\right\}\ j=j+1;\ \left\{sum = \frac{j(j-1)}{2}\right\}
    \end{gather*}
    y aplicar la regla de composición.

    \begin{enumerate}
        \item Para demostrar el primer triple, usamos el axioma de asignación:
            \begin{equation*}
                \left\{sum = \frac{(j+1)j}{2} \land j \neq c\right\}^{sum}_{sum+j}\ sum = sum + j;\ \left\{sum = \frac{(j+1)j}{2} \land j \neq c\right\} 
            \end{equation*}
            \begin{gather*}
                \left\{sum = \frac{(j+1)j}{2} \land j \neq c\right\}^{sum}_{sum+j} \equiv \left\{sum+j = \frac{(j+1)j}{2} \land j \neq c\right\} \equiv \\ 
                \equiv \left\{sum = \frac{(j+1)j}{2}-j \land j \neq c\right\} \equiv \left\{sum = \frac{j(j-1)}{2} \land j \neq c\right\}
            \end{gather*}
        \item Para el segundo, usamos también el axioma de asignación:
            \begin{equation*}
                \left\{sum = \frac{j(j-1)}{2}\right\}^j_{j+1} j=j+1;\ \left\{sum = \frac{j(j-1)}{2}\right\}
            \end{equation*}
            \begin{gather*}
                \left\{sum = \frac{j(j-1)}{2}\right\}^j_{j+1} \equiv \left\{sum = \frac{(j+1)(j+1-1)}{2}\right\} \equiv \left\{sum = \frac{(j+1)j}{2}\right\}
            \end{gather*}
            Además, tenemos que se tiene:
            \begin{equation*}
                \left\{sum = \frac{(j+1)j}{2} \land j \neq c\right\} \rightarrow \left\{sum = \frac{(j+1)j}{2}\right\}
            \end{equation*}
            Por tanto, usando la segunda regla de la consecuencia, tenemos que el segundo triple es cierto.
    \end{enumerate}

    Por tanto, mediante la regla de la iteración, tenemos que:
    \begin{equation*}
        \left\{sum=\frac{j(j-1)}{2}\right\}\ \texttt{while j <> c do begin sum:= sum+j; j:= j+1 end}\ \left\{sum=\frac{c(c-1)}{2}\right\}
    \end{equation*}
    Como inicialmente $j=1$ y $sum=0$, tenemos que este triple coincide con el enunciado del ejercicio.
\end{ejercicio}

\begin{ejercicio}
    Demostrar la corrección del siguiente triple:
    \begin{equation*}
        \{a[i]\geq 0\}\ a[i]=a[i]+a[j];\ \{a[i]\geq a[j]\}
    \end{equation*}

    Para ello, basta aplicar el axioma de la asignación:
    \begin{equation*}
        \{a[i]\geq a[j]\}_{a[i]+a[j]}^{a[i]}\ a[i]=a[i]+a[j];\ \{a[i]\geq a[j]\}
    \end{equation*}
    \begin{equation*}
        \{a[i]\geq a[j]\}_{a[i]+a[j]}^{a[i]}\ \equiv \{a[i]+a[j]\geq a[j]\} \equiv \{a[i] \geq 0\}
    \end{equation*}
\end{ejercicio}

\begin{ejercicio}
    Verificar el siguiente segmento de programa:
    \begin{gather*}
        \{n\geq 0\}\\ 
        i = 1;\\
        \texttt{while\ } i\leq n \texttt{\ do\ begin}\\
        a[i] = b[i]; \\
        i=i+1;\\
        \texttt{end} \\
        \left\{\bigwedge_{i=1}^n (a[i]=b[i])\right\}
    \end{gather*}

    % // TODO: Arturo se ha quedado por este ejercicio
\end{ejercicio}

\begin{ejercicio}
    El siguiente fragmento de programa calcula $\displaystyle\sum_{i=1}^{n} i!$. Demostrar que es correcto.
    \begin{minted}{pascal}
        i = 1; sum = 0; f = 1;
        while i <> n+1 do begin
          sum = sum + f;
          i = i + 1;
          f = f * 1;
        end
    \end{minted}
\end{ejercicio}

\begin{ejercicio}
    Hallar la precondición $\{P\}$ que hace que el siguiente triple sea correcto:
    \begin{equation*}
        \{P\}\ a[i]=2\ast b;\ \{j\leq i ~\land~ k < i ~\land~ a[i] + a[j-1]+a[k]>b\}
    \end{equation*}
    Para ello, basta aplicar el axioma de asignación:
    \begin{gather*}
        \{j\leq i~\land~ k < i ~\land~ a[i] + a[j-1]+a[k]>b\}_{2\cdot b}^{a[i]} a[i]=2\ast b;\\ \{j\leq i ~\land~ k < i ~\land~ a[i] + a[j-1]+a[k]>b\}
    \end{gather*}

    Usando la definición de Sustitución Textual, tenemos que:
    \begin{gather*}
        \{j\leq i~\land~ k < i ~\land~ a[i] + a[j-1]+a[k]>b\}_{2\cdot b}^{a[i]} \equiv \\ \equiv \{j\leq i ~\land~ k < i ~\land~ 2\cdot b + a[j-1]+a[k]>b\}
        \ \equiv\\ \equiv \{j\leq i ~\land~ k < i ~\land~ b + a[j-1]+a[k]>0\}
    \end{gather*}
    Luego estamos buscando la precondición:
    \begin{equation*}
        \{P\} \equiv \{j\leq i ~\land~ k < i ~\land ~b + a[j-1]+a[k]>0\}
    \end{equation*}
\end{ejercicio}

\begin{ejercicio}
    Demostrar que para $n>0$ el siguiente fragmento de programa termina.
    \begin{minted}{pascal}
        i = 1; f = 1;
        while i <> n do begin
          i = i + 1;
          f = f * r;
        end
    \end{minted}
    % // TODO: JJ. Cómo demostrar que termina?
\end{ejercicio}

\begin{ejercicio}
    Hallar la precondición de la terna:
    \begin{equation*}
        \{P\}\ a[i]=b;\ \{a[j] = 2\ast a[i]\}
    \end{equation*}
    Para ello, simplemente aplicamos el axioma de asignación:
    \begin{equation*}
        \{a[j]=2\ast b\} \equiv \{a[j]=2\ast a[i]\}_{b}^{a[i]}\ a[i]=b;\ \{a[j] = 2\ast a[i]\}
    \end{equation*}
    % // TODO: JJ. Tenemos que considerar i=j?
\end{ejercicio}

\begin{ejercicio}
    Para cada uno de los siguientes fragmentos de código, obtener la poscondición apropiada:
    \begin{enumerate}
        \item $\{i<10\}\ i=2\ast i+1;$.\\
            La poscondición es $\{i<21\}$:
            \begin{equation*}
                \{i<10\}\ i=2\ast i+1; \{i<21\}
            \end{equation*}
            que puede demostrarse aplicando el axioma de asignación.
        \item $\{i>0\}\ i=i-1;$.\\
            La poscondición es $\{i>-1\}$:
            \begin{equation*}
                \{i>0\}\ i=i-1; \{i>-1\}
            \end{equation*}
            que puede demostrarse aplicando el axioma de asignación.
        \item $\{i>j\}\ i=i+1;\ j=j+1;$.\\
            La poscondición es $\{i>j\}$:
            \begin{gather*}
                \{i>j\}\ i=i+1;\ \{i>j+1\} \\
                \{i>j+1\}\ j=j+1;\ \{i>j\}
            \end{gather*}
            Ambos pueden demostrarse aplicando el axioma de asignación y finalmente tenemos que:
            \begin{equation*}
                \{i>j\}\ i=i+1;\ j=j+1;\ \{i>j\}
            \end{equation*}
            aplicando la regla de composición.
        \item $\{V\}\ i=3;\ j=2\ast i$.\\
            La poscondición es $\{i=3 \land j=6\}$:
            \begin{gather*}
                \{V\}\ i=3;\ \{i=3\} \\
                \{i=3\}\ j=2\ast i;\ \{i=3 \land j = 6\}
            \end{gather*}
    \end{enumerate}
\end{ejercicio}

\begin{ejercicio}
    Para cada uno de los siguientes fragmentos de código, obtener las precondiciones apropiadas.
    \begin{enumerate}
        \item $i=3\ast k;\ \{i>6\}$.\\
            Aplicando el axioma de asignación:
            \begin{equation*}
                \{k>2\} \equiv \{3\cdot k>6\} \equiv \{i>6\}^i_{3\cdot k}\ i=3\ast k;\ \{i>6\}
            \end{equation*}
            obtenemos que la precondición es $\{k>2\}$.
        \item $a=b\ast c;\ \{a=1\}$.\\
            Aplicando el axioma de asignación:
            \begin{equation*}
                \{b = c^{-1}\} \equiv \{b\cdot c = 1\} \equiv \{a=1\}^a_{b\cdot c}\ a=b\ast c;\ \{a=1\}
            \end{equation*}
        La precondición es $\{b=c^{-1}\}$.
        \item $b=c-2;\ a=a/b;$\footnote{Esto es lo que pone en la relación, falta la poscondición para poder hacerse.}
    \end{enumerate}
\end{ejercicio}

\begin{ejercicio}
    Verificar el siguiente código. Indicar todas las reglas usadas.
    \begin{equation*}
        \{y>0\}\ xa=x+y;\ xb=x-y;
    \end{equation*}
    % // TODO: JJ. No hay poscondición??
\end{ejercicio}

\begin{ejercicio}
    Verificar el siguiente código, indicando todas las reglas usadas.
    \begin{equation*}
        \{V\}\ if\ x < 0\ then\ x = -x\ \{x\geq 0\}
    \end{equation*}
    Para comenzar, probamos que $\{x<0\}\ x=-x;\ \{x\geq 0\}$ usando el axioma de asignación:
    \begin{equation*}
        \{x<0\} \equiv \{-x\geq 0\} \equiv \{x\geq 0\}^x_{-x}\ x=-x;\ \{x\geq 0\}
    \end{equation*}
    Posteriormente, como sabemos que $\{V \land (x<0)\}\equiv \{x<0\}$ y que ${\{x\geq 0\}\ null\ \{x\geq 0\}}$ por el axioma de la sentencia nula, podemos aplicar la regla del \texttt{if}:
    \begin{equation*}
        \dfrac{\{V \land (x<0)\}\ x=-x;\ \{x\geq 0\}, \{V \land x\geq 0\}\ null\ \{x\geq 0\}}{\{V\}\ \texttt{if\ }x<0 \texttt{\ then\ }x=-x \texttt{\ else\ null\ } \{x\geq 0\}}
    \end{equation*}
\end{ejercicio}

\begin{ejercicio}\label{ej_rel1_44}
    Verificar el siguiente segmento de programa:
    \begin{gather*}
        max = a[1]; i = 1; \\
        \texttt{while\ } i <> n+1 \texttt{\ do\ begin} \\
        \texttt{if\ } a[i] \geq max \texttt{\ then\ } max = a[i];\\
        i=i+1; \\
        \texttt{end} \\
        \left\{\bigwedge_{i=1}^n (max \geq a[i])\right\}
    \end{gather*}
\end{ejercicio}

\begin{ejercicio}
    Demostrar la corrección parcial del siguiente código:
    \begin{gather*}
        max = a[1]; i = 1; \\
        \texttt{while\ } i < n \texttt{\ do\ begin} \\
        i = i +1; \\
        \texttt{if\ } a[i] \geq max \texttt{\ then\ } max = a[i]; \\
        \texttt{end} \\
        \left\{\bigwedge_{i=1}^n (max \geq a[i]), \bigvee_{j=1}^n (max = a[i]) \right\}
    \end{gather*}
\end{ejercicio}

\begin{ejercicio}
    Demostrar la corrección parcial del siguiente código:
    \begin{gather*}
        i = 0; j= n; \\
        \texttt{while\ } i < n \texttt{\ do\ begin} \\
        i = i + 1; \\
        j = j - 1; \\
        a[i] = b[j] \\
        \texttt{end} \\
        \left\{\bigwedge_{i=1}^n(a[i] = b[n-i])\right\}
    \end{gather*}
\end{ejercicio}

\begin{ejercicio}
    Demostrar la corrección parcial del siguiente código:
    \begin{minted}{pascal}
        i = 0;
        s = 0;
        while i <= n do begin
          s = s + a[i];
          a[i] = s;
          i = i + 1;
        end
    \end{minted}
\end{ejercicio}

\begin{ejercicio}
    Dados $n\geq 0$, $i\leq n$, demostrar que el siguiente segmento de progama evalúa
    \begin{equation*}
        \dfrac{n!}{i!(n-i)!}
    \end{equation*}
    \begin{minted}{pascal}
        k = 0; fact = 1;
        while k <> n do begin
          k = k + 1;
          fact = fact * k;
          if k <= i then afact = fact;
          if k <= n-i then bfact = fact;
        end
        bcof = fact/(afact*bfact);
    \end{minted}
\end{ejercicio}

\begin{ejercicio}
    Demostrar la terminación del fragmento de programa dado en el problema~\ref{ej_rel1_44} ¿Qué condición se debe imponer para realizar la demostración?
\end{ejercicio}

    %\section{Cambios de Varible}

\begin{ejercicio}
    Estudie las soluciones de la ecuación
    \begin{equation*}
        x' = \dfrac{t-5}{x^2}
    \end{equation*}
    dando en cada caso su intervalo maximal de definición.\\

    Tenemos que se trata de una ecuación de variables separadas de la forma dada por $x' = p(t)q(x)$, con:
    \Func{p}{I}{\bb{R}}{t}{t-5}
    \Func{q}{J}{\bb{R}}{x}{\dfrac{1}{x^2}}
    donde consideramos $I=\bb{R}$ y, para que el dominio sea conexo, podemos considerar $J=\bb{R}^+$ o $J=\bb{R}^-$.\\

    Usamos por tanto el método de variables separadas. En primer lugar, comprobamos que $q$ no tiene raíces en $J$:
    \begin{equation*}
        q(x)=0 \Longleftrightarrow \dfrac{1}{x^2}=0 \Longleftrightarrow 1=0
    \end{equation*}
    Una vez comprobado esto, procedemos a resolver la ecuación usando el método de variables separadas:
    \begin{align*}
        \dfrac{dx}{dt} = \dfrac{t-5}{x^2} &\Longrightarrow x^2dx = (t-5)dt \Longrightarrow \int x^2dx = \int (t-5)dt \Longrightarrow \\ &\Longrightarrow \dfrac{x^3}{3} = \dfrac{t^2}{2} - 5t + C' \qquad C'\in \bb{R}
    \end{align*}

    Despejando $x$ obtenemos la solución de la ecuación diferencial:
    \begin{equation*}
        x(t) = \sqrt[3]{\dfrac{3}{2}t^2 - 15t + C} \qquad C\in \bb{R}
    \end{equation*}

    Busquemos ahora su intervalo maximal de definición (llamémoslo $\wh{I}\subset I$). Necesitamos que $x(t)\in J$ para todo $t\in \wh{I}$ y que $x$ sea derivable en $\wh{I}$. Distinguimos casos:
    \begin{itemize}
        \item \ul{$J=\bb{R}^+$}: En este caso, necesitamos que $x(t)>0$ para todo $t\in \wh{I}$. Para ello, basta con que el radicando sea positivo:
        \begin{equation*}
            \dfrac{3}{2}t^2 - 15t + C > 0
        \end{equation*}

        Veamos en qué puntos se anula el radicando:
        \begin{align*}
            \dfrac{3}{2}t^2 - 15t + C &= 0 \Longrightarrow t = \dfrac{15\pm\sqrt{225-6C}}{3} = 5\pm\sqrt{25-\dfrac{2C}{3}}
        \end{align*}

        Distinguimos en función de $C$:
        \begin{equation*}
            25 - \dfrac{2C}{3} = 0 \Longrightarrow C = \dfrac{75}{2}
        \end{equation*}
        \begin{itemize}
            \item \ul{$C>\nicefrac{75}{2}$}: En este caso, el último radicando es negativo, luego no se anula el radicando de $x$, y este es siempre positivo. Por tanto, $x(t)>0$ para todo $t\in I$; es decir, $x(t)\in J$ para todo $t\in I$. Además, $x$ es derivable en $I$, luego el intervalo maximal de definición es $I$, $\wh{I}=I$.
            
            \item \ul{$C=\nicefrac{75}{2}$}: En este caso, el último radicando se anula en $t=5$. Por tanto, $x(t)>0$ para $t\in I\setminus \{5\}$. Por tanto, como el intervalo de definición de la solución debe ser conexo, consideramos las dos siguientes opciones:
            \begin{equation*}
                I_1 = \left]-\infty, 5\right[ \qquad I_2 = \left]5, +\infty\right[
            \end{equation*}

            En ambos casos, como $x(t)\in J$ para todo $t\in I_1$ y todo $t\in I_2$, y $x$ es derivable en $I_1$ y $I_2$, el intervalo maximal de definición es $\wh{I}=I_1$ o $\wh{I}=I_2$.

            \item \ul{$C<\nicefrac{75}{2}$}: En este caso, el último radicando es positivo, luego se anula en dos puntos, $t_1$ y $t_2$ dados por:
            \begin{equation*}
                t_1 = 5-\sqrt{25-\dfrac{2C}{3}} \qquad t_2 = 5+\sqrt{25-\dfrac{2C}{3}}
            \end{equation*}

            Por tanto, $x(t)>0$ para $t\in I\setminus [t_1,t_2]$. Por tanto, como el intervalo de definición de la solución debe ser conexo, consideramos las dos siguientes opciones:
            \begin{equation*}
                I_1 = \left]-\infty, t_1\right[ \qquad I_2 = \left]t_2, +\infty\right[
            \end{equation*}

            En todos los casos, como $x(t)\in J$ para todo $t\in I_1$ y todo $t\in I_2$, y $x$ es derivable en $I_1$ y $I_2$, el intervalo maximal de definición es $\wh{I}=I_1$ o $\wh{I}=I_2$.
        \end{itemize}
        
        \item \ul{$J=\bb{R}^-$}: En este caso, necesitamos que $x(t)<0$ para todo $t\in \wh{I}$. Para ello, basta con que el radicando sea negativo:
        \begin{equation*}
            \dfrac{3}{2}t^2 - 15t + C < 0
        \end{equation*}

        Distinguimos en función de $C$:
        \begin{itemize}
            \item \ul{$C>\nicefrac{75}{2}$}: En este caso, el último radicando es negativo, luego no se anula el radicando de $x$. Además, $x(t)>0$ para todo $t\in I$, por lo que no hay solución en este caso.
            
            \item \ul{$C=\nicefrac{75}{2}$}: En este caso, el último radicando se anula en $t=5$. Por tanto, $x(t)>0$ para $t\in I\setminus \{5\}$. Además, como el intervalo de definición de la solución debe ser abierto y conexo, no hay solución en este caso.

            \item \ul{$C<\nicefrac{75}{2}$}: En este caso, el último radicando es positivo, luego se anula en dos puntos, $t_1$ y $t_2$ dados por:
            \begin{equation*}
                t_1 = 5-\sqrt{25-\dfrac{2C}{3}} \qquad t_2 = 5+\sqrt{25-\dfrac{2C}{3}}
            \end{equation*}

            Por tanto, $x(t)<0$ para $t\in [t_1,t_2]$. Como en el abierto es derivable, el intervalo maximal de definición es $\wh{I}=\left]t_1,t_2\right[$.
        \end{itemize}
        % // TODO: Revisar JJ. Joshoccas lo tiene distinto
    \end{itemize}


\end{ejercicio}

\begin{ejercicio}
    En Dinámica de Poblaciones, dos modelos muy conocidos son la ecuación de Verhulst o logística
    \begin{equation*}
        P' = P(\alpha - \beta P)
    \end{equation*}
    y la ecuación de Gompertz
    \begin{equation*}
        P' = P(\alpha - \beta \ln P)
    \end{equation*}
    siendo $P(t)$ la población a tiempo $t$ de una determinada especie y $\alpha, \beta$ parámetros positivos. Calcule en cada caso la solución con condición inicial $P(0) = 100$.\\

    Resolvamos en primer lugar la ecuación de Verhulst. Se trata de una ecuación de variables separadas de la forma $P' = p(t)q(P)$, con:
    \Func{p}{I}{\bb{R}}{t}{1}
    \Func{q}{J}{\bb{R}}{P}{P(\alpha - \beta P)}
    donde consideramos $I=J=\bb{R}$. Comprobamos las raíces de $q$ en $J$:
    \begin{equation*}
        q(P) = 0 \Longleftrightarrow P(\alpha - \beta P) = 0 \Longleftrightarrow P=0, \dfrac{\alpha}{\beta}
    \end{equation*}

    Por tanto, dos soluciones de la ecuación son, para todo $t\in I$:
    \begin{equation*}
        P(t) = 0 \qquad P(t) = \dfrac{\alpha}{\beta}
    \end{equation*}

    Procedemos ahora a resolver la ecuación de variables separadas. Para ello, trabajaremos con $J_1=\bb{R}^-$, $J_2=\left]0, \nicefrac{\alpha}{\beta}\right[$ y $J_3=\left]\nicefrac{\alpha}{\beta}, +\infty\right[$,
    ya que necesitamos que $q(P)\neq 0$ para todo $P$ en el la segunda componente del dominio.
    \begin{itemize}
        \item \ul{$J_1=\bb{R}^-$}:
        
        Como en este caso no cumple que $P(0)=100\in J_1$, no nos interesa este dominio.

        \item \ul{$J_2=\left]0, \nicefrac{\alpha}{\beta}\right[$}:
        
        Veamos qué hemos de de imponer para considerar este dominio; es decir, para que $P(0)=100\in J_2$:
        \begin{equation*}
            100\in J_2 \Longleftrightarrow 100<\dfrac{\alpha}{\beta} \Longleftrightarrow 100\beta<\alpha
        \end{equation*}
        
        En este caso, resolvemos la ecuación de variables separadas con dominio $I\times J_2$:
        \begin{align*}
            P' = P(\alpha - \beta P) &\Longrightarrow \dfrac{dP}{dt} = P(\alpha - \beta P) \Longrightarrow \dfrac{dP}{P(\alpha - \beta P)} = dt \Longrightarrow \\ &\Longrightarrow \int \dfrac{dP}{P(\alpha - \beta P)} = \int dt
        \end{align*}

        Para resolver la primera integral, aplicamos el método de descomposición en fracciones simples:
        \begin{align*}
            \dfrac{1}{P(\alpha - \beta P)} &= \dfrac{A}{P} + \dfrac{B}{\alpha - \beta P} = \dfrac{A(\alpha - \beta P) + BP}{P(\alpha - \beta P)}
        \end{align*}
        \begin{itemize}
            \item \ul{Para $P=0$}: $1=A \cdot \alpha \Longrightarrow A=\nicefrac{1}{\alpha}$.
            \item \ul{Para $P=\nicefrac{\alpha}{\beta}$}: $1=B \cdot \nicefrac{\alpha}{\beta} \Longrightarrow B=\nicefrac{\beta}{\alpha}$.
        \end{itemize}

        Por tanto, tenemos que:
        \begin{align*}
            \Longrightarrow \int \dfrac{dP}{P(\alpha - \beta P)} = \int dt
            &\Longrightarrow \frac{1}{\alpha}\int \dfrac{dP}{P} + \frac{\beta}{\alpha}\int \dfrac{dP}{\alpha - \beta P} = \int dt \Longrightarrow \\ &\Longrightarrow \dfrac{1}{\alpha}\ln(P) - \dfrac{1}{\alpha}\ln(\alpha - \beta P) = t + C \qquad C\in \bb{R}
        \end{align*}
        donde en la última implicación hemos usado que $P\in J_2$.

        Operando con la solución obtenida, llegamos a:
        \begin{align*}
            \ln\left(\dfrac{P}{\alpha - \beta P}\right) &= \alpha(t + C) \Longrightarrow \dfrac{P}{\alpha - \beta P} = e^{\alpha (t + C)}
            \Longrightarrow\\ &\Longrightarrow P(1+\beta e^{\alpha (t + C)}) = \alpha e^{\alpha (t + C)} \Longrightarrow P = \dfrac{\alpha e^{\alpha (t + C)}}{1+\beta e^{\alpha (t + C)}}
        \end{align*}

        Por tanto, para $P\in J_2$, la familia de soluciones es:
        \begin{equation*}
            P(t) = \dfrac{\alpha e^{\alpha (t + C)}}{1+\beta e^{\alpha (t + C)}} \qquad C\in \bb{R}
        \end{equation*}

        Estableciendo la condición inicial $P(0)=100$, obtenemos:
        \begin{align*}
            P(0) &= 100 = \dfrac{\alpha e^{\alpha C}}{1+\beta e^{\alpha C}} \Longrightarrow 100(1+\beta e^{\alpha C}) = \alpha e^{\alpha C} \Longrightarrow \\ &\Longrightarrow 100 + 100\beta e^{\alpha C} = \alpha e^{\alpha C} \Longrightarrow \\ &\Longrightarrow 100 = \alpha e^{\alpha C} - 100\beta e^{\alpha C} \Longrightarrow \\ &\Longrightarrow 100 = e^{\alpha C}(\alpha - 100\beta) \Longrightarrow \\ &\Longrightarrow e^{\alpha C} = \dfrac{100}{\alpha - 100\beta} \Longrightarrow \\ &\Longrightarrow C =\dfrac{1}{\alpha} \ln\left(\dfrac{100}{\alpha - 100\beta}\right)
        \end{align*}
        
        Por tanto, la solución con condición inicial $P(0)=100$ es:
        \begin{equation*}
            P(t) = \dfrac{\alpha e^{\alpha (t + C)}}{1+\beta e^{\alpha (t + C)}}, \quad t\in I, \qquad C=\dfrac{1}{\alpha}\ln\left(\dfrac{100}{\alpha - 100\beta}\right)
        \end{equation*}

        \item \ul{$J_3=\left]\nicefrac{\alpha}{\beta}, +\infty\right[$}:
        
        Veamos qué hemos de de imponer para considerar este dominio; es decir, para que $P(0)=100\in J_3$:
        \begin{equation*}
            100\in J_3 \Longleftrightarrow 100>\dfrac{\alpha}{\beta} \Longleftrightarrow 100\beta>\alpha
        \end{equation*}

        En este caso, resolvemos la ecuación de variables separadas con dominio $I\times J_3$.
        Por los cáculos realizados en el caso anterior, tenemos que:
        \begin{equation*}
            \dfrac{1}{\alpha} \ln(P) - \dfrac{1}{\alpha} \ln(\beta P - \alpha) = t + C \qquad C\in \bb{R}
        \end{equation*}

        Operando con la solución obtenida, llegamos a que la familia de soluciones para $P\in J_3$ es:
        \begin{equation*}
            P(t) = \dfrac{-\alpha e^{\alpha (t + C)}}{1-\beta e^{\alpha (t + C)}}
            = \dfrac{\alpha e^{\alpha (t + C)}}{-1+\beta e^{\alpha (t + C)}} \qquad C\in \bb{R}
        \end{equation*}

        Estableciendo la condición inicial $P(0)=100$, obtenemos:
        \begin{align*}
            P(0) &= 100 = \dfrac{\alpha e^{\alpha C}}{-1+\beta e^{\alpha C}} \Longrightarrow 100(-1+\beta e^{\alpha C}) = \alpha e^{\alpha C} \Longrightarrow \\ &\Longrightarrow -100 + 100\beta e^{\alpha C} = \alpha e^{\alpha C} \Longrightarrow \\ &\Longrightarrow -100 = \alpha e^{\alpha C} - 100\beta e^{\alpha C} \Longrightarrow \\ &\Longrightarrow -100 = e^{\alpha C}(\alpha - 100\beta) \Longrightarrow \\ &\Longrightarrow e^{\alpha C} = \dfrac{-100}{\alpha - 100\beta} \Longrightarrow \\ &\Longrightarrow C =\dfrac{1}{\alpha} \ln\left(\dfrac{100}{100\beta-\alpha}\right)
        \end{align*}

        Por tanto, la solución con condición inicial $P(0)=100$ es:
        \begin{equation*}
            P(t) = \dfrac{\alpha e^{\alpha (t + C)}}{-1+\beta e^{\alpha (t + C)}}, \quad t\in I, \qquad C=\dfrac{1}{\alpha}\ln\left(\dfrac{100}{100\beta-\alpha}\right)
        \end{equation*}
    \end{itemize}

    Por tanto, y a modo de resumen, las soluciones de la ecuación de Verhulst con condición inicial $P(0)=100$ son, en función de los parámetros $\alpha, \beta$:
    \begin{itemize}
        \item \ul{$100=\nicefrac{\alpha}{\beta}$}: En este caso, se trata de la solución constante, luego:
        \begin{equation*}
            P(t) = 100 = \frac{\alpha}{\beta} \qquad t\in I
        \end{equation*}

        \item \ul{$100<\nicefrac{\alpha}{\beta}$}: En este caso, la solución está en $J_2$, luego:
        \begin{equation*}
            P(t) = \dfrac{\alpha e^{\alpha (t + C)}}{1+\beta e^{\alpha (t + C)}}, \quad t\in I, \qquad C=\dfrac{1}{\alpha}\ln\left(\dfrac{100}{\alpha - 100\beta}\right)
        \end{equation*}

        \item \ul{$100>\nicefrac{\alpha}{\beta}$}: En este caso, la solución está en $J_3$, luego:
        \begin{equation*}
            P(t) = \dfrac{\alpha e^{\alpha (t + C)}}{-1+\beta e^{\alpha (t + C)}}, \quad t\in I, \qquad C=\dfrac{1}{\alpha}\ln\left(\dfrac{100}{100\beta-\alpha}\right)
        \end{equation*}
    \end{itemize}~\\

    Resolvamos ahora la ecuación de Gompertz. Se trata de una ecuación de variables separadas de la forma $P' = p(t)q(P)$, con:
    \Func{p}{I}{\bb{R}}{t}{1}
    \Func{q}{J}{\bb{R}}{P}{P(\alpha - \beta \ln P)}
    donde consideramos $I=\bb{R}$, $J=\bb{R}^+$. Comprobamos las raíces de $q$ en $J$:
    \begin{equation*}
        q(P) = 0 \Longleftrightarrow P(\alpha - \beta \ln P) = 0 \Longleftrightarrow P=0, e^{\nicefrac{\alpha}{\beta}}
    \end{equation*}

    La solución $P=0$ no es válida, puesto que no pertenece al dominio $J$. Por tanto, la única solución constante es:
    \begin{equation*}
        P(t) = e^{\nicefrac{\alpha}{\beta}} \qquad t\in I
    \end{equation*}

    Procedemos ahora a resolver la ecuación de variables separadas. Para ello, trabajaremos con $J_1=\left]0,e^{\nicefrac{\alpha}{\beta}}\right[$ y $J_2=\left]e^{\nicefrac{\alpha}{\beta}}, +\infty\right[$, ya que necesitamos que $q(P)\neq 0$ para todo $P$ en el la segunda componente del dominio.
    \begin{itemize}
        \item \ul{$J_1=\left]0,e^{\nicefrac{\alpha}{\beta}}\right[$}:
        
        Veamos qué hemos de de imponer para considerar este dominio; es decir, para que $P(0)=100\in J_1$:
        \begin{equation*}
            100\in J_1 \Longleftrightarrow 100<e^{\nicefrac{\alpha}{\beta}} \Longleftrightarrow \ln 100 < \dfrac{\alpha}{\beta}
        \end{equation*}

        En este caso, resolvemos la ecuación de variables separadas con dominio $I\times J_1$:
        \begin{align*}
            P' = P(\alpha - \beta \ln P) &\Longrightarrow \dfrac{dP}{dt} = P(\alpha - \beta \ln P) \Longrightarrow \dfrac{dP}{P(\alpha - \beta \ln P)} = dt \Longrightarrow \\ &\Longrightarrow \int \dfrac{dP}{P(\alpha - \beta \ln P)} = \int dt
        \end{align*}

        Para resolver la integral del logaritmo, aplicamos el cambio de variable $P=e^u$, luego $dP=e^u du$:
        \begin{align*}
            \int \dfrac{dP}{P(\alpha - \beta \ln P)} &= \int \dfrac{e^u du}{e^u(\alpha - \beta u)} = \int \dfrac{du}{\alpha - \beta u}
            = -\dfrac{1}{\beta} \ln(\alpha - \beta u) + C' =\\
            &= -\dfrac{1}{\beta} \ln(\alpha - \beta \ln P) + C' \qquad C'\in \bb{R}
        \end{align*}
        donde hemos hecho uso de que:
        \begin{equation*}
            \alpha-\beta u>0\Longleftrightarrow
            u<\dfrac{\alpha}{\beta}\Longleftrightarrow
            \ln P<\dfrac{\alpha}{\beta}\Longleftrightarrow
            0<P<e^{\nicefrac{\alpha}{\beta}} \Longleftrightarrow P\in J_1
        \end{equation*}

        Operando, llegamos a que:
        \begin{equation*}
            -\dfrac{1}{\beta} \ln(\alpha - \beta \ln P) = t + C \Longrightarrow \alpha - \beta \ln P = e^{-\beta(t + C)} \Longrightarrow \ln P = \dfrac{\alpha - e^{-\beta(t + C)}}{\beta}
        \end{equation*}

        Por tanto, la solución uniparamétrica de la ecuación de Gompertz en $J_1$ es:
        \begin{equation*}
            P(t) = \exp\left(\dfrac{\alpha - e^{-\beta(t + C)}}{\beta}\right) \qquad C\in \bb{R}
        \end{equation*}

        Estableciendo la condición inicial $P(0)=100$, obtenemos:
        \begin{align*}
            P(0) &= 100 = \exp\left(\dfrac{\alpha - e^{-\beta C}}{\beta}\right) \Longrightarrow \ln(100) = \dfrac{\alpha - e^{-\beta C}}{\beta} \Longrightarrow \\ &\Longrightarrow e^{-\beta C} = \alpha-\beta\ln(100) \Longrightarrow C=-\dfrac{1}{\beta}\ln(\alpha-\beta\ln(100))
        \end{align*}

        Por tanto, la solución con condición inicial $P(0)=100$ es:
        \begin{equation*}
            P(t) = \exp\left(\dfrac{\alpha - e^{-\beta(t +C)}}{\beta}\right) \qquad t\in I, \qquad C=-\dfrac{1}{\beta}\ln(\alpha-\beta\ln(100))
        \end{equation*}

        \item \ul{$J_2=\left]e^{\nicefrac{\alpha}{\beta}}, +\infty\right[$}:
        
        Veamos qué hemos de de imponer para considerar este dominio; es decir, para que $P(0)=100\in J_2$:
        \begin{equation*}
            100\in J_2 \Longleftrightarrow 100>e^{\nicefrac{\alpha}{\beta}} \Longleftrightarrow \ln 100 > \dfrac{\alpha}{\beta}
        \end{equation*}

        En este caso, resolvemos la ecuación de variables separadas con dominio $I\times J_2$.
        Por los cáculos realizados en el caso anterior, tenemos que:
        \begin{equation*}
            -\dfrac{1}{\beta} \ln(\beta \ln P-\alpha) = t + C \qquad C\in \bb{R}
        \end{equation*}
        donde hemos hecho uso de que:
        \begin{equation*}
            \alpha-\beta u<0\Longleftrightarrow
            u>\dfrac{\alpha}{\beta}\Longleftrightarrow
            \ln P>\dfrac{\alpha}{\beta}\Longleftrightarrow
            P>e^{\nicefrac{\alpha}{\beta}} \Longleftrightarrow P\in J_2
        \end{equation*}

        Operando, llegamos a que la solución uniparamétrica de la ecuación de Gompertz en $J_2$ es:
        \begin{equation*}
            P(t) = \exp\left(\dfrac{-\alpha + e^{-\beta(t + C)}}{\beta}\right) \qquad C\in \bb{R}
        \end{equation*}

        Estableciendo la condición inicial $P(0)=100$, y repitiendo los cálculos del apartado anterior, llegamos a:
        \begin{align*}
            P(0) &= 100 \Longrightarrow C=-\dfrac{1}{\beta}\ln(\beta\ln(100)-\alpha)
        \end{align*}

        Por tanto, la solución con condición inicial $P(0)=100$ es:
        \begin{equation*}
            P(t) = \exp\left(\dfrac{-\alpha + e^{-\beta(t + C)}}{\beta}\right) \qquad t\in I, \qquad C=-\dfrac{1}{\beta}\ln(\beta\ln(100)-\alpha)
        \end{equation*}
    \end{itemize}

    Por tanto, y a modo de resumen, las soluciones de la ecuación de Gompertz con condición inicial $P(0)=100$ son, en función de los parámetros $\alpha, \beta$:
    \begin{itemize}
        \item \ul{$100=e^{\nicefrac{\alpha}{\beta}}$}: En este caso, se trata de la solución constante, luego:
        \begin{equation*}
            P(t) = 100 = e^{\nicefrac{\alpha}{\beta}} \qquad t\in I
        \end{equation*}

        \item \ul{$100<e^{\nicefrac{\alpha}{\beta}}$}: En este caso, la solución está en $J_1$, luego:
        \begin{equation*}
            P(t) = \exp\left(\dfrac{\alpha - e^{-\beta(t + C)}}{\beta}\right) \qquad t\in I, \qquad C=-\dfrac{1}{\beta}\ln(\alpha-\beta\ln(100))
        \end{equation*}

        \item \ul{$100>e^{\nicefrac{\alpha}{\beta}}$}: En este caso, la solución está en $J_2$, luego:
        \begin{equation*}
            P(t) = \exp\left(\dfrac{-\alpha + e^{-\beta(t + C)}}{\beta}\right) \qquad t\in I, \qquad C=-\dfrac{1}{\beta}\ln(\beta\ln(100)-\alpha)
        \end{equation*}
    \end{itemize}

\end{ejercicio}

\begin{ejercicio}
    Nos planteamos resolver la ecuación
    \begin{equation*}
        x' = \cos(t - x)
    \end{equation*}
    Compruebe que el cambio $y = t - x$ nos lleva a una ecuación de variables separadas. Resuelva e invierta el cambio para llegar a una expresión explícita de $x(t)$. Repase el procedimiento por si se ha perdido alguna solución por el camino.\\

    Al no indicarnos nada sobre la primera variable, suponemos que esta no varía, luego el cambio de variable a aplicar es:
    \begin{equation*}
        \left\{
            \begin{aligned}
                s &= t,\\
                y &= t - x.
            \end{aligned}
        \right.
    \end{equation*}

    Al no especificar dominio de la ecuación, suponemos que está definida en $\bb{R}^2$. Consideramos entonces las siguientes funciones:
    \Func{\varphi=(\varphi_1,\varphi_2)}{\bb{R}^2}{\bb{R}^2}{(t,x)}{(s,y)=(t, t-x)}
    \Func{\psi=(\psi_1,\psi_2)}{\bb{R}^2}{\bb{R}^2}{(s,y)}{(t,x)=(s, s-y)}

    Para emplear el cambio de variable, en primer lugar hemos de comprobar que $\varphi$ es un difeomorfismo.
    Para ello, hemos demostrar que $\varphi$ es biyectiva y que $\varphi,\varphi^{-1}\in C^1(\bb{R}^2)$.
    Demostraremos en primer lugar que $\varphi\circ \psi = Id_{\bb{R}^2}= \psi\circ \varphi$:
    \begin{align*}
        \varphi\circ \psi(s,y) &= \varphi(s, s-y) = (s, s-(s-y)) = (s, y) = Id_{\bb{R}^2}(s,y) \qquad \forall (s,y)\in \bb{R}^2\\
        \psi\circ \varphi(t,x) &= \psi(t, t-x) = (t, t-(t-x)) = (t, x) = Id_{\bb{R}^2}(t,x) \qquad \forall (t,x)\in \bb{R}^2
    \end{align*}
    Por tanto, hemos demostrado que $\varphi$ es biyectiva y que $\varphi^{-1}=\psi$. Además, como ambas componentes de $\varphi,\psi$ son de clase $1$, tenemos que $\varphi,\varphi^{-1}\in C^1(\bb{R}^2)$, luego $\varphi$ es un difeomorfismo.

    A continuación, hemos de comprobar que el cambio es admisible. Aunque lo sabemos puesto que no varía la primera variable, lo demostramos:
    \begin{align*}
        \dfrac{\partial \varphi_1}{\partial t} + \dfrac{\partial \varphi_1}{\partial x} f(t,x) &= 1 + 0 \cdot \cos(t-x) = 1\neq 0
    \end{align*}

    Por tanto, tenemos que el cambio de variable es admisible. Procedemos a aplicarlo a la ecuación dada:
    \begin{align*}
        \dfrac{dy}{ds} = \dfrac{dy}{dt}\cdot \dfrac{dt}{ds} = \dfrac{dy}{dt} = 1-\cos(t-x) = 1-\cos y
    \end{align*}

    Usando la notación usual, la nueva ecuación diferencial, con dominio $\bb{R}^2$, es:
    \begin{equation*}
        y' = 1-\cos y
    \end{equation*}

    Esta es de la forma $y' = p(s)q(y)$, con:
    \Func{p}{\bb{R}}{\bb{R}}{s}{1}
    \Func{q}{\bb{R}}{\bb{R}}{y}{1-\cos y}

    Buscamos en primer lugar los valores en los que se anula $q$:
    \begin{equation*}
        q(y) = 0 \Longleftrightarrow 1-\cos y = 0 \Longleftrightarrow \cos y = 1 \Longleftrightarrow y = 2\pi k, \quad k\in \bb{Z}
    \end{equation*}

    Fijado ahora $k\in \bb{Z}$, restringimos ahora el dominio de $q$ a $J=\left]2\pi k, 2\pi(k+1)\right[$. En $J$ tenemos que $q$ no se anula, luego procedemos a resolver la ecuación de variables separadas:
    \begin{align*}
        \dfrac{dy}{ds} = 1-\cos y &\Longrightarrow \dfrac{dy}{1-\cos y} = ds \Longrightarrow \int \dfrac{dy}{1-\cos y} = \int ds
    \end{align*}

    %https://www.youtube.com/watch?v=kxHFFkfIoaU
    % // TODO: Cómo se resuelve la integral de coseno?

\end{ejercicio}

\begin{ejercicio}
    Experimentalmente, se sabe que la resistencia al aire de un cuerpo en caída libre es proporcional al cuadrado de la velocidad del mismo. Por tanto, si $v(t)$ es la velocidad a tiempo $t$, la ecuación de Newton nos dice que
    \begin{equation*}
        v' + \dfrac{k}{m}v^2 = g,
    \end{equation*}
    donde $m$ es la masa del cuerpo, $g$ es la constante de gravitación universal y $k > 0$ depende de la geometría (aerodinámica) del cuerpo. Si se supone que $v(0) = 0$, calcule la solución explícita y describa el comportamiento a largo plazo.\\

    Se trata de una ecuación de variables separadas de la forma $v' = p(t)q(v)$, con:
    \Func{p}{\bb{R}}{\bb{R}}{t}{1}
    \Func{q}{\bb{R}}{\bb{R}}{v}{-\dfrac{k}{m}v^2 + g}

    Comprobamos las raíces de $q$ en $\bb{R}$:
    \begin{equation*}
        q(v) = 0 \Longleftrightarrow -\dfrac{k}{m}v^2 + g = 0 \Longleftrightarrow v^2 = \dfrac{gm}{k} \Longleftrightarrow \left\{
            \begin{aligned}
                v_1 &= -\sqrt{\dfrac{gm}{k}},\\
                v_2 &= \sqrt{\dfrac{gm}{k}}.
            \end{aligned}
        \right.
    \end{equation*}

    Como buscamos la solución que cumple $v(0)=0$, consideramos el dominio dado por $J=\left]v_1,v_2\right[$. En este dominio, $q$ no se anula, luego procedemos a resolver la ecuación de variables separadas:
    \begin{align*}
        \dfrac{dv}{dt} = -\dfrac{k}{m}v^2 + g &\Longrightarrow \dfrac{dv}{-\dfrac{k}{m}v^2 + g} = dt \Longrightarrow \int \dfrac{dv}{-\dfrac{k}{m}v^2 + g} = \int dt
        \Longrightarrow \\ &\Longrightarrow m\int \dfrac{dv}{-kv^2 + gm} = t + C \qquad C\in \bb{R}
    \end{align*}

    Para resolver la integral, descomponemos en fracciones simples:
    \begin{align*}
        \dfrac{1}{-kv^2 + gm} &= \dfrac{A}{v-v_1} + \dfrac{B}{v-v_2} = \dfrac{A(v-v_2) + B(v-v_1)}{(v-v_1)(v-v_2)}
    \end{align*}
    \begin{itemize}
        \item \ul{Para $v=v_1$}: $1=A(v_1-v_2) \Longrightarrow A=\dfrac{1}{v_1-v_2}=-\dfrac{1}{2}\sqrt{\dfrac{k}{mg}}$.
        \item \ul{Para $v=v_2$}: $1=B(v_2-v_1) \Longrightarrow B=\dfrac{1}{v_2-v_1}=\dfrac{1}{2}\sqrt{\dfrac{k}{mg}}$.
    \end{itemize}

    Por tanto, tenemos que:
    \begin{align*}
        m\int \dfrac{dv}{-kv^2 + gm} &= m\int \left(\dfrac{A}{v-v_1} + \dfrac{B}{v-v_2}\right)dv
        =\\&=m\left(A\ln(v-v_1) + B\ln(v_2-v)\right) + C'
        =\\&= \frac{m}{2}\sqrt{\dfrac{k}{mg}}\left(\ln(v_2-v) - \ln(v-v_1)\right) + C'
        =\\&= \frac{m}{2}\sqrt{\dfrac{k}{mg}}\ln\left(\dfrac{v_2-v}{v-v_1}\right) + C'
        = \frac{m}{2}\sqrt{\dfrac{k}{mg}}\ln\left(\dfrac{\sqrt{\dfrac{gm}{k}}-v}{\sqrt{\dfrac{gm}{k}}+v}\right) + C'
    \end{align*}

    Por tanto, la famlia de soluciones en $J$ es:
    \begin{equation*}
        \dfrac{m}{2}\sqrt{\dfrac{k}{mg}}\ln\left(\dfrac{\sqrt{\dfrac{gm}{k}}-v}{\sqrt{\dfrac{gm}{k}}+v}\right) = t + C \qquad C\in \bb{R}
    \end{equation*}

    Operando, llegamos a:
    \begin{align*}
        \dfrac{\sqrt{\dfrac{gm}{k}}-v}{\sqrt{\dfrac{gm}{k}}+v} &= \exp\left((t+C)\cdot \frac{2}{m}\sqrt{\dfrac{mg}{k}}\right)\\
        -v\left[1+\exp\left((t+C)\cdot \frac{2}{m}\sqrt{\dfrac{mg}{k}}\right)\right] &= \sqrt{\dfrac{gm}{k}}\exp\left((t+C)\cdot \frac{2}{m}\sqrt{\dfrac{mg}{k}}\right) - \sqrt{\dfrac{gm}{k}}\\
        v&=-\dfrac{\sqrt{\dfrac{gm}{k}}\exp\left((t+C)\cdot \dfrac{2}{m}\sqrt{\dfrac{mg}{k}}\right) - \sqrt{\dfrac{gm}{k}}}{1+\exp\left((t+C)\cdot \dfrac{2}{m}\sqrt{\dfrac{mg}{k}}\right)} \\
        v&=\dfrac{\sqrt{\dfrac{gm}{k}}\left[1-\exp\left((t+C)\cdot \dfrac{2}{m}\sqrt{\dfrac{mg}{k}}\right)\right]}{1+\exp\left((t+C)\cdot \dfrac{2}{m}\sqrt{\dfrac{mg}{k}}\right)}
    \end{align*}

    Por tanto, la solución con condición inicial $v(0)=0$ es:
    \begin{align*}
        v(0)&=0=\dfrac{\sqrt{\dfrac{gm}{k}}\left[1-\exp\left(C\cdot \dfrac{2}{m}\sqrt{\dfrac{mg}{k}}\right)\right]}{1+\exp\left(C\cdot \dfrac{2}{m}\sqrt{\dfrac{mg}{k}}\right)} \Longrightarrow 1 = \exp\left(C\cdot \dfrac{2}{m}\sqrt{\dfrac{mg}{k}}\right)
        \Longrightarrow \\&\Longrightarrow C\cdot \dfrac{2}{m}\sqrt{\dfrac{mg}{k}} = 0 \Longrightarrow C=0
    \end{align*}

    Por tanto, la solución con condición inicial $v(0)=0$ es:
    \begin{equation*}
        v(t)=\dfrac{\sqrt{\dfrac{gm}{k}}\left[1-\exp\left(t\cdot \dfrac{2}{m}\sqrt{\dfrac{mg}{k}}\right)\right]}{1+\exp\left(t\cdot \dfrac{2}{m}\sqrt{\dfrac{mg}{k}}\right)} \qquad t\in \bb{R}
    \end{equation*}

    Estudiemos ahora el comportamiento a largo plazo de la solución. Para ello, consideramos el límite cuando $t\to +\infty$:
    \begin{align*}
        \lim_{t\to +\infty} v(t) &= -\sqrt{\dfrac{gm}{k}}
    \end{align*}

    % // TODO: Arturo, Revisar cálculos. Velocidad negativa?
\end{ejercicio}

\begin{ejercicio}
    Calcule la solución de la ecuación
    \begin{equation*}
        y' = \dfrac{x + y - 3}{x - y - 1}
    \end{equation*}
    que verifica $y(0) = 1$.\\

    En este caso, se trata de una ecuación reducible a homogénea.
    Estudiemos qué cambio de variable hemos de emplear:
    \begin{equation*}
        \left|\begin{array}{cc}
            1 & 1\\
            1 & -1
        \end{array}\right| = -1-1=-2\neq 0
    \end{equation*}
    Por tanto, hemos de emplear una traslación según el vector $(x_\ast,y_\ast)\in \bb{R}^2$.
    En primer lugar, el dominio de la ecuación diferencial ha de ser una de las componentes conexas de:
    \begin{equation*}
        \left\{(x,y)\in \bb{R}^2 \mid x-y-1\neq 0\right\}
    \end{equation*}

    Para que el punto $(0,y(0))=(0,1)$ esté en el dominio, el dominio será:
    \begin{equation*}
        D=\left\{(x,y)\in \bb{R}^2 \mid x-y-1< 0\right\}
    \end{equation*}

    Consideramos entonces el cambio de variable:
    \Func{\varphi=(\varphi_1,\varphi_2)}{D}{D_1}{(x,y)}{(u,v)=(x-x_\ast,y-y_\ast)}
    donde el codominio, $D_1$, es:
    \begin{align*}
        D_1=\varphi(D)&=\left\{\varphi(x,y)\in \bb{R}^2 \mid x-y-1< 0\right\} = \left\{(u,v)\in \bb{R}^2 \mid (u+x_\ast, v+y_\ast)\in D\right\} =\\
        &= \left\{(u,v)\in \bb{R}^2 \mid u+x_\ast - v-y_\ast-1<0\right\}
    \end{align*}
    La inversa de $\varphi$, puesto que podemos despejar cada una de las componentes de forma única, es:
    \Func{\varphi^{-1}}{D_1}{D}{(u,v)}{(x,y)=(u+x_\ast, v+y_\ast)}

    Por tanto, $\varphi, \varphi^{-1}$ son biyectivas. Como además son continuas por ser ambas componentes polinómicas, tenemos que $\varphi$ es un difeomorfismo. Además, es admisible, ya que:
    \begin{equation*}
        \dfrac{\partial \varphi_1}{\partial x} + \dfrac{\partial \varphi_1}{\partial y} y' = 1 + 0\cdot y'=1 \neq 0
    \end{equation*}

    Por tanto, procedemos a aplicar el cambio de variable a la ecuación dada:
    \begin{align*}
        \dfrac{dv}{du} &= \dfrac{dv}{dx}\cdot \dfrac{dx}{du} = \dfrac{dy}{dx} = \dfrac{x + y - 3}{x - y - 1} = \dfrac{u+x_\ast + v+y_\ast - 3}{u+x_\ast - v-y_\ast - 1}
    \end{align*}

    Buscamos ahora resolver el siguiente sistema:
    \begin{equation*}
        \left\{
            \begin{aligned}
                x_\ast +y_\ast &= 3,\\
                x_\ast -y_\ast &= 1.
            \end{aligned}
        \right\}\Longrightarrow
        \left\{
            \begin{aligned}
                x_\ast &= 2,\\
                y_\ast &= 1.
            \end{aligned}
        \right.
    \end{equation*}

    Por tanto, tras haber aplicado el cambio de variable según la traslación de $(2,1)$, la ecuación diferencial se convierte en:
    \begin{equation*}
        v' = \dfrac{u + v}{u - v}
    \end{equation*}
    Como buscamos que $y(0)=1$, tenemos que $v(-2)=0$.
    Esta es una ecuación homogénea, con dominio:
    \begin{align*}
        D_1'&=\left\{(u,v)\in D_1\mid u<0\right\}=\\
        &= \left\{(u,v)\in \bb{R}^2 \mid u<0,~u-v<0\right\}
    \end{align*}
    
    El cambio de variable a aplicar es:
    \Func{\varphi'=(\varphi_1',\varphi_2')}{D_1'}{D_2}{(u,v)}{(t,s)=(u, \nicefrac{v}{u})}
    donde el codominio, $D_2$, es:
    \begin{align*}
        D_2=\varphi'(D_1')&=\left\{\varphi'(u,v)\in \bb{R}^2 \mid (u,v)\in D_1'\right\} = \left\{(t,s)\in \bb{R}^2 \mid (t, s\cdot t)\in D_1'\right\} =\\
        &= \left\{(t,s)\in \bb{R}^2 \mid t<0,~t-s\cdot t<0\right\}
    \end{align*}

    La inversa de $\varphi'$, puesto que podemos despejar cada una de las componentes de forma única, es:
    \Func{(\varphi')^{-1}}{D_2}{D_1'}{(t,s)}{(u,v)=(t, s\cdot t)}

    Por tanto, $\varphi', (\varphi')^{-1}$ son biyectivas. Como además son continuas por ser ambas componentes polinómicas, tenemos que $\varphi'$ es un difeomorfismo. Además, es admisible, ya que la primera componente no varía. La demostración de que es admisible es:
    \begin{equation*}
        \dfrac{\partial \varphi_1'}{\partial u} + \dfrac{\partial \varphi_1'}{\partial v} v' = 1 + 0\cdot v'=1 \neq 0
    \end{equation*}

    Por tanto, procedemos a aplicar el cambio de variable a la ecuación dada:
    \begin{align*}
        \dfrac{ds}{dt} &= \dfrac{ds}{du}\cdot \cancelto{1}{\dfrac{du}{dt}} = -\dfrac{v}{u^2} + \dfrac{v'}{u} = -\dfrac{1}{u}\cdot \dfrac{v}{u} + \dfrac{1}{u}\cdot \dfrac{1+\nicefrac{v}{u}}{1-\nicefrac{v}{u}} = -\dfrac{s}{t} + \dfrac{1}{t}\cdot \dfrac{1+s}{1-s} 
    \end{align*}

    Por tanto, la ecuación homogénea se convierte en:
    \begin{equation*}
        s' = \dfrac{1}{t}\left(\dfrac{1+s}{1-s} - s\right) \qquad \text{con dominio } D_2=\bb{R}^- \times \left]-\infty,1\right[
    \end{equation*}

    Esta es una ecuación de variables separadas, luego procedemos a resolverla. Buscamos las raíces de la función con variable $s$:
    \begin{equation*}
        \dfrac{1+s}{1-s} - s = 0 \Longleftrightarrow 1+s = s(1-s) = s-s^2 \Longleftrightarrow s^2+1=0
    \end{equation*}
    Vemos que no tiene soluciones constantes, luego aplicamos el método de resolución:
    \begin{align*}
        \dfrac{ds}{\dfrac{1+s}{1-s} - s} &= \dfrac{dt}{t} \Longrightarrow \int \dfrac{1-s}{1+s^2} ~ds = \int \dfrac{dt}{t} \Longrightarrow \int \dfrac{1}{1+s^2} ~ds - \int \dfrac{s}{1+s^2} ~ds = \ln(-t)+C'
        \Longrightarrow \\ &\Longrightarrow \arctan(s) - \dfrac{1}{2}\ln(1+s^2) = \ln(-t)+C
    \end{align*}
    Por la teoría vista de cambio de variable, tenemos que esto define una función implícita $s(t)$ en $\bb{R}^-$.
    Deshacemos el segundo cambio de variable:
    \begin{equation*}
        \arctan\left(\dfrac{v}{u}\right) - \dfrac{1}{2}\ln\left(1+\left(\dfrac{v}{u}\right)^2\right) = \ln(-u)+C
    \end{equation*}
    Esto define una función implícita $v(u)$ en $\bb{R}^-$.
    Por último, deshacemos el primer cambio de variable:
    \begin{equation*}
        \arctan\left(\dfrac{y-1}{x-2}\right) - \dfrac{1}{2}\ln\left(1+\left(\dfrac{y-1}{x-2}\right)^2\right) = \ln(-x+2)+C
    \end{equation*}
    Esta ecuación define una función implícita $y(x)$ en el dominio $\left]-\infty, 1\right[$. Estableciendo la condición inicial $y(0)=1$, obtenemos:
    \begin{align*}
        \arctan\left(\dfrac{1-1}{0-2}\right) - \dfrac{1}{2}\ln\left(1+\left(\dfrac{1-1}{0-2}\right)^2\right) &= \ln(-0+2)+C \Longrightarrow \\ &\Longrightarrow \arctan\left(0\right) - \dfrac{1}{2}\ln\left(1+0\right) = \ln(2)+C \Longrightarrow \\ &\Longrightarrow C=-\ln(2)
    \end{align*}

    Por tanto, la solución de la ecuación dada que verifica $y(0)=1$ es la función implícita $y(x)$ definida por la ecuación siguiente en el dominio $\left]-\infty, 1\right[$:
    \begin{equation*}
        \arctan\left(\dfrac{y-1}{x-2}\right) - \dfrac{1}{2}\ln\left(1+\left(\dfrac{y-1}{x-2}\right)^2\right) = \ln(-x+2)-\ln(2)
    \end{equation*}
\end{ejercicio}

\begin{ejercicio}
    Resuelva los siguientes problemas lineales
    \begin{enumerate}
        \item $x' + 3x = e^{-3t}$, $x(1) = 5$
        \item $x' - \dfrac{x}{t} = \dfrac{1}{1+t^2}$, $x(2) = 0$
        \item $x' = \cosh t \cdot x + \sinh t$, $x(0) = 1$
    \end{enumerate}
\end{ejercicio}

\begin{ejercicio}
    Sean $a, b : \bb{R} \to \bb{R}$ funciones continuas con $a(t) \geq c > 0$ para todo $t$ y
    \begin{equation*}
        \lim_{t \to +\infty} b(t) = 0.
    \end{equation*}
    Demuestre que todas las soluciones de la ecuación $x' = -a(t)x + b(t)$ tienden a cero cuando $t \to +\infty$. (Indicación: regla de L'Hôpital en la fórmula de variación de constantes)
\end{ejercicio}

\begin{ejercicio}
    La ecuación de Bernoulli tiene la forma
    \begin{equation*}
        x' = a(t)x + b(t)x^n,
    \end{equation*}
    donde $a, b : I \to \bb{R}$ son funciones continuas y $n \in \bb{R}$. Compruebe que el cambio de variable $y = x^\alpha$ lleva la ecuación de Bernoulli a una ecuación del mismo tipo, y ajuste el valor de $\alpha$ para que la ecuación obtenida sea lineal ($n = 0$). Usando el cambio anterior, resuelva los problemas de valores iniciales
    \begin{equation*}
        x' = x + t\sqrt{x}, \quad x(0) = 1.
    \end{equation*}
\end{ejercicio}

\begin{ejercicio}
    Se considera la ecuación de Ricatti
    \begin{equation*}
        y' = -\dfrac{1}{x^2} - \dfrac{y}{x} + y^2.
    \end{equation*}
    Encuentre una solución particular de la forma $y(x) = x^\alpha$ con $\alpha \in \bb{R}$. Usando esta solución particular, calcule la solución que cumple $y(1) = 2$ y estudie su intervalo maximal de definición.
\end{ejercicio}

\begin{ejercicio}
    Encuentre una curva $y = y(x)$ que pase por el punto $(1, 2)$ y cumpla la siguiente propiedad: la distancia de cada punto de la curva al origen coincide con la segunda coordenada del punto de corte de la recta tangente y el eje de ordenadas. (C. Sturm, Cours d’Analyse 1859, Vol 2, pag 41).
\end{ejercicio}

\begin{ejercicio}
    Identifique la clase de ecuaciones invariantes por el grupo de transformaciones $s = \lambda t$, $y = \lambda^2 x$, con $\lambda > 0$.
\end{ejercicio}

\begin{ejercicio}
    Resuelva los problemas 42 y 45 (pag. 79) del libro de Nagle-Saff-Sneider.
\end{ejercicio}
\end{document}
